\chapter{Relative Picard sheaf --- \texttt{Scheme.Modules.tensorObj} substrate (A.1.c.SubT)}
\label{chap:Picard_TensorObjSubstrate}
% archon:covers AlgebraicJacobian/Picard/TensorObjSubstrate.lean
% archon:covers AlgebraicJacobian/Picard/TensorObjSubstrate/StalkTensor.lean
% archon:covers AlgebraicJacobian/Picard/TensorObjSubstrate/Vestigial.lean
% archon:covers AlgebraicJacobian/Picard/TensorObjSubstrate/DualInverse.lean
% archon:covers AlgebraicJacobian/Picard/TensorObjSubstrate/DualInverse/SliceTransport.lean
% archon:covers AlgebraicJacobian/Picard/TensorObjSubstrate/PresheafInternalHom.lean
% archon:covers AlgebraicJacobian/Picard/TensorObjInverse.lean

\section{Motivation: the abelian-group-valued relative Picard sheaf}
\label{sec:tensorobj_motivation}

Fix a base field \(k\) and a smooth proper geometrically integral curve \(C\) over
\(k\). In Kleiman's framework (\cref{chap:Picard_RelPicFunctor},
\texttt{rel\_pic\_sharp}), the relative Picard functor
\[
  \Pic_{C/k} \;\;\colon\;\; (\Sch/k)^{op} \;\longrightarrow\; \AddCommGrpCat,
  \qquad
  T \;\longmapsto\; \Pic(C \times_k T) \, / \, \pi_T^* \Pic(T),
\]
is, by definition, abelian-group-valued. The abelian-group structure on each
fiber \(\Pic(C \times_k T)\) comes from tensor product of line bundles
(Stacks tag 01CR / Hartshorne~II.\S 6), and the structure descends to the
quotient because \(\pi_T^*\) is a group homomorphism.

% SOURCE: [Kleiman], "The Picard scheme", \S 2 (FGA Explained Ch.~9 \S 9.2),
% Definitions df:aPf + df:Pfs (the absolute and relative Picard functors are
% abelian-group-valued; the relative one is defined as a quotient of abelian
% groups)
% (read from references/kleiman-picard-src/kleiman-picard.tex, L1274-L1318).
% SOURCE QUOTE:
% "\begin{dfn}\label{df:aPf}
%  The {\it absolute Picard functor} $\Pic_X$ is the functor from  the
%  category of (locally Noetherian) $S$-schemes $T$ to the category of
% abelian groups defined by the formula
% 	$$\Pic_X(T):=\Pic(X_T).$$
% \end{dfn}
%  [\dots]
% \begin{dfn}\label{df:Pfs}
%   The {\it relative Picard functor} $\Pic_{X/S}$ is defined by
% 	$$\Pic_{X/S}(T):=\Pic(X_T)\big/\Pic(T).$$
%  Denote its associated sheaves in the Zariski, \'etale, and fppf
% topologies  by
% 	$$\Pic_{(X/S)\zar},\ \Pic_{(X/S)\et},\
% 	\Pic_{(X/S)\fppf}.$$
%  \end{dfn}"
\textit{Source: [Kleiman], ``The Picard scheme'', \S 2, Defs.~df:aPf + df:Pfs.
The absolute and relative Picard functors take values in the category of
abelian groups; the relative functor is the quotient of one such group by
another, and the structure morphisms of the resulting presheaf are group
homomorphisms.}

Concretely, the abelian-group law on \(\Pic(C \times_k T)\) is the assignment
\(\bigl[\,\mathcal{L}\,\bigr] + \bigl[\,\mathcal{L}'\,\bigr] :=
\bigl[\,\mathcal{L} \otimes_{\mathcal{O}_{C \times_k T}} \mathcal{L}'\,\bigr]\),
with neutral element the class of the structure sheaf
\(\mathcal{O}_{C \times_k T}\) and inverse
\(- \bigl[\,\mathcal{L}\,\bigr] := \bigl[\,\mathcal{L}^{-1}\,\bigr]\), where
\(\mathcal{L}^{-1} = \mathcal{H}om_{\mathcal{O}_{C \times_k T}}(\mathcal{L},
\mathcal{O}_{C \times_k T})\) is the dual invertible sheaf. Pullback of
\(\mathcal{O}\)-modules along \(\pi_T : C \times_k T \to T\) preserves tensor
products and the structure sheaf, so the image
\(\pi_T^* \Pic(T) \subset \Pic(C \times_k T)\) is a subgroup, and the quotient
\(\Pic(C \times_k T) / \pi_T^* \Pic(T)\) inherits a canonical abelian-group
structure. This is the content of
\cref{lem:rel_pic_sharp_groupoid} in chapter \cref{chap:Picard_RelPicFunctor}.

The mathematics is straightforward; the obstacle to formalisation is purely
infrastructural, and --- crucially --- much smaller than a full
monoidal-category construction. The group law lives on \emph{isomorphism
classes} of line bundles: each group axiom is a \emph{proposition} asserting
the existence of an isomorphism (\(\mathrm{Nonempty}(\,\cdots \cong \cdots)\)),
and a proposition is insensitive to which isomorphism witnesses it. In
particular the coherence data of a monoidal category --- the pentagon,
triangle and hexagon identities, and a fortiori any closed structure --- is
\emph{never consumed} by a group law on iso-classes. What the law genuinely
requires on the carrier is therefore only:
\begin{enumerate}
  \item a binary tensor-product operation
    \(\otimes\, : \Scheme.\mathtt{Modules}\,X \times \Scheme.\mathtt{Modules}\,X
                  \to \Scheme.\mathtt{Modules}\,X\), well-defined on
    isomorphism classes;
  \item the structure sheaf \(\mathcal{O}_X\) as a designated unit object
    for \(\otimes\);
  \item the \(\otimes\)-invertibility predicate
    (\(\exists N,\ M \otimes_X N \cong \mathcal{O}_X\)) together with the three
    monoidal coherence isomorphisms of \(\otimes_X\): associativity
    \((M \otimes N) \otimes P \cong M \otimes (N \otimes P)\), unit
    \(\mathcal{O}_X \otimes M \cong M\) (and on the right), and commutativity
    \(M \otimes N \cong N \otimes M\). The inverse is carried by the
    predicate, not by a separate lemma.
\end{enumerate}
At Mathlib's pinned commit one might be tempted to obtain (3) by instantiating a
\emph{full} symmetric monoidal category on \(\Scheme.\mathtt{Modules}\,X\) and
reading the three isomorphisms off it. Two Mathlib facts make such a route
conceivable. First, the category \(\mathtt{PresheafOfModules}\,R\) of presheaves
of modules over a presheaf of commutative rings \(R\) is already a (symmetric)
monoidal category in Mathlib (\(\mathtt{PresheafOfModules.monoidalCategory}\)),
with the sectionwise relative tensor \((M \otimes_R N)(U) = M(U) \otimes_{R(U)}
N(U)\); this applies verbatim to the \emph{varying} structure sheaf
\(R = \mathcal{O}_X\). Second, Mathlib's localization-of-monoidal-categories API
(\(\mathtt{CategoryTheory.Localization.Monoidal.LocalizedMonoidal}\)) takes a
monoidal category \(C\), a localization functor \(L : C \to D\) at a morphism
property \(W\), and an instance \(W.\mathtt{IsMonoidal}\), and produces a
\emph{full} monoidal-category structure on the localization \(D\); sheafification
of presheaves of modules is the localization \(L = \mathtt{sheafification}\) at the
locally-bijective class \(W = J.W\), so a full monoidal structure on
\(\Scheme.\mathtt{Modules}\,X\) would follow once \(J.W.\mathtt{IsMonoidal}\) were
supplied.

This full-monoidal-category route is, however, \emph{not} the goal, and is not
pursued. The goal is the commutative \emph{group} on locally-trivial iso-classes
(à la \(\mathtt{CommRing.Pic}\)), and that group law consumes only propositions, so
it never needs a coherent monoidal category --- neither the
pentagon/triangle/hexagon identities nor the \(J.W.\mathtt{IsMonoidal}\) instance
that the \(\mathtt{LocalizedMonoidal}\) API requires. What it needs is exactly the
\emph{existence} of the three coherence isomorphisms of (3), and these are produced
directly: the associator by gluing canonical local isomorphisms
(\cref{lem:tensorobj_assoc_iso}, via the single linchpin
\cref{lem:tensorobj_restrict_iso}), the unitors by the cheap \(\mathtt{mapIso}\)
pattern (\cref{lem:tensorobj_unit_iso}), and the braiding likewise
(\cref{lem:tensorobj_comm_iso}). With those three existence facts the group law
assembles \emph{by hand} the commutative group of \(\otimes\)-invertible
\(\otimes\)-iso-classes --- the Picard group
(\cref{lem:tensorobj_isoclass_commgroup}) --- each group-axiom field fed a single
existence-of-isomorphism. This is the scheme-theoretic analogue of Mathlib's
fixed-ring \(\mathtt{CommRing.Pic}\,R := \mathrm{Skeleton}(\mathtt{Module.Invertible}\,R)\),
whose \(\mathtt{CommGroup}\) is read off the \(\mathrm{Skeleton}\) of the coherent
monoidal category \(\mathtt{MonoidalCategory}(\mathtt{Module.Invertible}\,R)\) that
Mathlib supplies over a fixed ring; the present construction cannot take that
\(\mathrm{Skeleton}\) route, since the coherent monoidal category on
\(\Scheme.\mathtt{Modules}\,X\) for the varying structure sheaf is exactly what it
avoids. The
\(\otimes\)-invertibility predicate (\cref{def:scheme_modules_isinvertible},
Mathlib's \(\mathtt{Module.Invertible}\) idiom) carries each inverse existentially.
The \emph{full} \(\mathtt{LocalizedMonoidal}\) / \(J.W.\mathtt{IsMonoidal}\) coherent
monoidal category is therefore not built for the group law and is off the critical
path; what the group law does still need from this circle of ideas is the
\emph{existence} of the associator (\cref{lem:tensorobj_assoc_iso}), realized as a
sheafification transport whose single open left-whisker obligation is closed by the
stalk--tensor commutation (d.2). The stalk apparatus (d.1/d.2) is accordingly
\emph{on} the critical path, while the coherent monoidal category itself is not.

\section{Mathlib API survey}
\label{sec:tensorobj_api_survey}

The relevant pieces of the Mathlib API at the pinned commit
(\texttt{b80f227}) were originally assembled into a single \emph{realization route}
(route~(e)): obtaining a full monoidal structure on the substrate by instantiating
Mathlib's abstract localization-of-monoidal-categories API rather than
hand-assembling the associator. That route is \emph{superseded}. The group law on
iso-classes never consumes a coherent monoidal category, and the associator is
obtained by direct gluing through \cref{lem:tensorobj_restrict_iso}, so the
\(J.W.\mathtt{IsMonoidal}\) instance on which route~(e) hinges is never required.
The survey below is retained as a Mathlib-API reference and to record why the route
is vestigial; the four Mathlib facts are all verified against the on-disk Mathlib
source.

\paragraph{(i) The presheaf-of-modules monoidal category.} Mathlib's
\texttt{Mathlib.Algebra.Category.ModuleCat.Presheaf.Monoidal} constructs a
monoidal-category instance
\(\mathtt{PresheafOfModules.monoidalCategory}\) on
\(\mathtt{PresheafOfModules}\,(R \circ \mathtt{forget}_2\,\mathtt{CommRingCat}\,
\mathtt{RingCat})\), with the sectionwise relative tensor
\[
  (M \otimes_R N)(U) \;\;=\;\; M(U) \otimes_{R(U)} N(U)
\]
% SOURCE: [Mathlib], `PresheafOfModules.monoidalCategoryStruct` /
% `PresheafOfModules.monoidalCategory`, in
% `Mathlib/Algebra/Category/ModuleCat/Presheaf/Monoidal.lean` (L104, L125)
% (read from .lake/packages/mathlib/Mathlib/Algebra/Category/ModuleCat/Presheaf/Monoidal.lean).
% SOURCE QUOTE:
% "noncomputable instance monoidalCategoryStruct :
%     MonoidalCategoryStruct (PresheafOfModules.{u} (R ⋙ forget₂ _ _)) where
%   tensorObj := tensorObj
%   whiskerLeft _ _ _ g := tensorHom (𝟙 _) g
%   whiskerRight f _ := tensorHom f (𝟙 _)
%   ...
%   tensorUnit := unit _
%   associator M₁ M₂ M₃ := isoMk (fun _ ↦ α_ _ _ _) ...
% noncomputable instance monoidalCategory :
%     MonoidalCategory (PresheafOfModules.{u} (R ⋙ forget₂ _ _)) where ..."
\textit{Source: [Mathlib], \texttt{PresheafOfModules.monoidalCategoryStruct} /
\texttt{PresheafOfModules.monoidalCategory}, in
\texttt{Mathlib/Algebra/Category/ModuleCat/Presheaf/Monoidal.lean}, L104--L125.}
together with associator, left/right unitors and braiding making it symmetric
monoidal. The base presheaf of rings \(R \circ \mathtt{forget}_2\,
\mathtt{CommRingCat}\,\mathtt{RingCat}\) is exactly the project's
\(\mathtt{Sheaf.val}\,X.\mathtt{ringCatSheaf}\) carrier (equal by \(\mathtt{rfl}\)),
so this monoidal structure applies verbatim to the \emph{varying} structure sheaf
\(\mathcal{O}_X\) --- no project work. This is the presheaf-level analogue of
Stacks tag 03DM (relative tensor product of \(\mathcal{O}_X\)-modules).

\paragraph{(ii) The localization-of-monoidal-categories API.} Mathlib's
\texttt{Mathlib.CategoryTheory.Localization.Monoidal.Basic} provides the class
\(\mathtt{MorphismProperty.IsMonoidal}\) and the construction
\(\mathtt{Localization.Monoidal.LocalizedMonoidal}\). The class records exactly
the whiskering-stability data:
% SOURCE: [Mathlib], `CategoryTheory.MorphismProperty.IsMonoidal` and
% `CategoryTheory.Localization.Monoidal.LocalizedMonoidal`, in
% `Mathlib/CategoryTheory/Localization/Monoidal/Basic.lean` (L44, L86, L440)
% (read from .lake/packages/mathlib/Mathlib/CategoryTheory/Localization/Monoidal/Basic.lean).
% SOURCE QUOTE:
% "class IsMonoidal : Prop extends W.IsMultiplicative where
%   whiskerLeft (X : C) {Y₁ Y₂ : C} (g : Y₁ ⟶ Y₂) (hg : W g) : W (X ◁ g)
%   whiskerRight {X₁ X₂ : C} (f : X₁ ⟶ X₂) (hf : W f) (Y : C) : W (f ▷ Y)
%  ...
% def LocalizedMonoidal (L : C ⥤ D) (W : MorphismProperty C)
%     [W.IsMonoidal] [L.IsLocalization W] {unit : D} (_ : L.obj (𝟙_ C) ≅ unit) := D
%  ...
% noncomputable instance :
%     MonoidalCategory (LocalizedMonoidal L W ε) where ..."
\textit{Source: [Mathlib], \texttt{MorphismProperty.IsMonoidal} and
\texttt{Localization.Monoidal.LocalizedMonoidal}, in
\texttt{Mathlib/CategoryTheory/Localization/Monoidal/Basic.lean}, L44--L440.}
Given \([\MonoidalCategory\,C]\), \([W.\mathtt{IsMonoidal}]\), a localization
functor \(L : C \to D\) with \([L.\mathtt{IsLocalization}\,W]\) and an isomorphism
\(\varepsilon : L(\mathbf{1}_C) \cong \mathtt{unit}\), the localized category
\(\mathtt{LocalizedMonoidal}\,L\,W\,\varepsilon\) is a \(\MonoidalCategory\) with
\emph{all} coherence (associator, unitors, pentagon, triangle, naturalities)
derived. There is \emph{no} \(\mathtt{MonoidalClosed}\) anywhere in the
hypothesis chain.

\paragraph{(iii) The proof template for $W$-monoidality.} Mathlib's
\texttt{Mathlib.CategoryTheory.Sites.Point.IsMonoidalW} proves
\((J.W\,(A := A)).\mathtt{IsMonoidal}\) \emph{stalkwise} for presheaves valued in a
\emph{fixed} monoidal concrete category \(A\) (i.e.\ \(C^{op} \to A\)), via the
conservative-family characterisation \(\mathtt{hP.W\_iff}\),
\(\mathtt{Functor.Monoidal.map\_tensor}\) and typeclass inference:
% SOURCE: [Mathlib], `CategoryTheory.ObjectProperty.IsConservativeFamilyOfPoints.isMonoidal_W`
% and the `(J.W).IsMonoidal` instance, in
% `Mathlib/CategoryTheory/Sites/Point/IsMonoidalW.lean` (L48, L57)
% (read from .lake/packages/mathlib/Mathlib/CategoryTheory/Sites/Point/IsMonoidalW.lean).
% SOURCE QUOTE:
% "lemma ObjectProperty.IsConservativeFamilyOfPoints.isMonoidal_W
%     [J.HasSheafCompose (forget A)] :
%     (J.W (A := A)).IsMonoidal :=
%   .mk' _ (fun f g hf hg ↦ by
%     simp only [hP.W_iff (A := A)] at hf hg ⊢
%     intro Φ
%     rw [Functor.Monoidal.map_tensor]
%     infer_instance)"
\textit{Source: [Mathlib],
\texttt{ObjectProperty.IsConservativeFamilyOfPoints.isMonoidal\_W}, in
\texttt{Mathlib/CategoryTheory/Sites/Point/IsMonoidalW.lean}, L48--L57.}
This is the \emph{template} for the proof technique used in route~(e), but it does
\emph{not} apply directly to \(\mathtt{PresheafOfModules}\,R\): modules over a
\emph{varying} ring are not of the form \(C^{op} \to A\) for a fixed monoidal
\(A\).

\paragraph{The gap and route~(e), superseded.} The substrate
\(\Scheme.\mathtt{Modules}.\mathtt{tensorObj}\) of this chapter is the
sheafification of \(\mathtt{PresheafOfModules.Monoidal.tensorObj}\), and
sheafification of presheaves of modules is exactly the localization
\(L = \mathtt{sheafification}\) at the locally-bijective class \(W = J.W\). Route~(e)
proposed to obtain a \emph{full} monoidal structure on
\(\Scheme.\mathtt{Modules}\,X = \mathtt{SheafOfModules}\,X.\mathtt{ringCatSheaf}\)
by applying \(\mathtt{LocalizedMonoidal}\) (ii) to the already-monoidal category (i),
which would require the instance \(J.W.\mathtt{IsMonoidal}\) (its two
whiskering-stability fields \cref{lem:whisker_of_W}, with load-bearing
local-injectivity half \cref{lem:islocallyinjective_whisker_of_W}). This is
\emph{not} how the group law is realized. The group law consumes only the
\emph{existence} of the three coherence isomorphisms, and the associator is built
directly by gluing canonical local isomorphisms through
\cref{lem:tensorobj_restrict_iso} (\cref{lem:tensorobj_assoc_iso}); it never invokes
the \emph{full} \(J.W.\mathtt{IsMonoidal}\) instance or a coherent monoidal category.
The \emph{full monoidal-category} route~(e) is therefore off the critical path. The
associator does, however, consume a \emph{single} left-whisker existence lemma
\cref{lem:islocallyinjective_whiskerleft_via_stalk}, closed by the d.2 stalk--tensor
commutation; so the stalk apparatus (d.1/d.2) itself is on the critical path, even
though the full coherent monoidal category built from it is not.

The stalk-infrastructure analysis that route~(e) turned on is recorded here for
reference only, since that route is superseded. The \emph{stalk module structure}
is already present in Mathlib: the file
\texttt{Mathlib.Algebra.Category.ModuleCat.Stalk} supplies, for a topological
space \(X\), a presheaf of commutative rings \(R\) on \(X\) and a point
\(x \in X\), the \(R.\mathtt{stalk}\,x\)-module structure on the stalk
\(\mathtt{stalk}\,M.\mathtt{presheaf}\,x\) of a presheaf of modules \(M\),
together with the germ--scalar compatibility \(\mathtt{germ\_smul}\); a linearity
packaging of the induced stalk map of a morphism was built project-side
(\cref{lem:stalk_linear_map}, the \emph{linearity half} of ingredient (d.1)).
Route~(e) additionally required a (d.1)-\emph{bridge} between the site-level
\(J.W\) and the topological stalkwise-isomorphism criterion on
\(\mathtt{Opens}\,X\), and (d.2) the commutation \((A \otimes^{p} B)_x \cong A_x
\otimes_{R_x} B_x\) of the stalk with the relative module-presheaf tensor; there is
moreover no monoidal \(\mathtt{SheafOfModules}\) in Mathlib. The full \(J.W.\mathtt{IsMonoidal}\) coherent monoidal category that route~(e) would
build is not pursued; but the (d.1)-linearity packaging and the (d.2) stalk--tensor
commutation are \emph{on} the critical path, since the associator's single open
left-whisker obligation \cref{lem:islocallyinjective_whiskerleft_via_stalk} is closed
precisely by (d.2) (\cref{sec:tensorobj_stalk_tensor}). What is superseded is only the
\emph{pair}-of-whiskering-fields \(J.W.\mathtt{IsMonoidal}\) packaging and the
\(\mathtt{LocalizedMonoidal}\) monoidal category, not the stalk infrastructure.

The project's policy on this gap is \emph{project-side}: the substrate is
built inside \texttt{AlgebraicJacobian/Picard/} (the file
\(\texttt{AlgebraicJacobian/Picard/TensorObjSubstrate.lean}\) is its home)
and the consumer instance is closed against it. A parallel Mathlib upstream
PR with the same content is welcome but not on the critical path for the
project; the consumer is closed against the project-side substrate
independently of any upstream timing.

\section{The substrate \(\Scheme.\mathtt{Modules}.\mathtt{tensorObj}\)}
\label{sec:tensorobj_substrate}

\begin{definition}


\leanok
[The substrate operation $\otimes$ on $\Scheme.\mathtt{Modules}\,X$]
  \label{def:scheme_modules_tensorobj}
  \lean{AlgebraicGeometry.Scheme.Modules.tensorObj}
  Let \(X\) be a scheme. For
  \(M, N \in \Scheme.\mathtt{Modules}\,X\), define
  \[
    M \otimes_X N \;\;\in\;\; \Scheme.\mathtt{Modules}\,X
  \]
  as follows. Choose an affine open cover \(\{U_i = \Spec A_i\}_{i \in I}\)
  of \(X\). On each \(U_i\), the restriction
  \(M|_{U_i}\) corresponds (via the equivalence between
  \(\mathcal{O}_{\Spec A_i}\)-modules and \(A_i\)-modules) to an
  \(A_i\)-module \(M_i\), and similarly
  \(N|_{U_i}\) corresponds to an \(A_i\)-module \(N_i\). Define
  \((M \otimes_X N)|_{U_i}\) as the sheafification of the presheaf
  \(U_i \mapsto M_i \otimes_{A_i} N_i\). On overlaps
  \(U_i \cap U_j = \Spec A_{ij}\) the resulting modules agree via the
  canonical compatibility
  \((M_i \otimes_{A_i} N_i) \otimes_{A_i} A_{ij}
     \;\cong\;
     (M_j \otimes_{A_j} N_j) \otimes_{A_j} A_{ij}\),
  so the data \(\{(M \otimes_X N)|_{U_i}\}\) glue to a single object
  \(M \otimes_X N\) in \(\Scheme.\mathtt{Modules}\,X\). The result is
  canonically independent of the affine cover chosen.

  Equivalently: the underlying presheaf of \(M \otimes_X N\) is the
  presheaf-of-modules tensor product
  \(\mathtt{PresheafOfModules.Monoidal.tensorObj}\) of the underlying
  presheaves of \(M\) and \(N\), composed with the sheafification
  functor on the small Zariski site of \(X\). The substrate operation
  \(\otimes_X\) is, by construction, a \(\Scheme.\mathtt{Modules}\,X\)
  object reflecting the relative tensor product of
  \(\mathcal{O}_X\)-modules described in Stacks tag 03DM.
\end{definition}

\begin{lemma}


\leanok
[Functoriality of $\otimes_X$]
  \label{lem:scheme_modules_tensorobj_functoriality}
  \lean{AlgebraicGeometry.Scheme.Modules.tensorObj_functoriality}
  \uses{def:scheme_modules_tensorobj}
  Let \(X\) be a scheme. The assignment
  \((M, N) \mapsto M \otimes_X N\) of
  \cref{def:scheme_modules_tensorobj} extends to a bifunctor
  \[
    \otimes_X \;\;\colon\;\; \Scheme.\mathtt{Modules}\,X \;\times\;
                          \Scheme.\mathtt{Modules}\,X \;\;\longrightarrow\;\;
                          \Scheme.\mathtt{Modules}\,X
  \]
  natural in both arguments: a pair of morphisms
  \(f : M \to M'\) and \(g : N \to N'\) determines a morphism
  \(f \otimes g : M \otimes_X N \to M' \otimes_X N'\) compatible with
  identities and composition. This bifunctorial action on morphisms is
  the only datum the present lemma provides. The further coherence data
  of a monoidal structure --- the unitors
  \[
    \lambda_M \;\colon\; \mathcal{O}_X \otimes_X M \;\xrightarrow{\sim}\; M,
    \qquad
    \rho_M \;\colon\; M \otimes_X \mathcal{O}_X \;\xrightarrow{\sim}\; M,
  \]
  the associator
  \(\alpha_{M, N, P} : (M \otimes_X N) \otimes_X P \xrightarrow{\sim}
   M \otimes_X (N \otimes_X P)\), and the braiding
  \(\beta_{M, N} : M \otimes_X N \xrightarrow{\sim} N \otimes_X M\) ---
  are \emph{not} outputs of this lemma; they are off the critical path,
  recorded in \cref{rem:scheme_modules_monoidal_off_path}.
\end{lemma}

\begin{proof}
  \uses{def:scheme_modules_tensorobj}
  \leanok
  Pick an affine open cover \(\{U_i = \Spec A_i\}\) of \(X\). On each
  affine, the assignment is the standard tensor product of modules over a
  commutative ring \(A_i\), which is bifunctorial. The bifunctorial action
  on morphisms glues across the cover because the underlying gluing data of
  \cref{def:scheme_modules_tensorobj} respect the canonical morphisms on
  overlaps; equivalently, the morphism action is inherited from
  \(\mathtt{PresheafOfModules.Monoidal.tensorObj}\) (Mathlib module
  \texttt{Mathlib.CategoryTheory.Monoidal.PresheafOfModules}) under
  sheafification. The unitor/associator/braiding coherence data are
  \emph{not} constructed here; they are off the critical path
  (\cref{rem:scheme_modules_monoidal_off_path}).
\end{proof}

\begin{remark}
[The full monoidal category on all of $\Scheme.\mathtt{Modules}\,X$ is off-path; the group law consumes only propositions]
  \label{rem:scheme_modules_monoidal_off_path}
  \uses{def:scheme_modules_tensorobj, lem:scheme_modules_tensorobj_functoriality, lem:whisker_of_W}
  One could ask for the bifunctor \(\otimes_X\) of
  \cref{lem:scheme_modules_tensorobj_functoriality}, together with
  \(\mathcal{O}_X\), \(\alpha\), \(\lambda\), \(\rho\) and \(\beta\), to
  assemble into a full symmetric \(\MonoidalCategory\) structure on
  \emph{all} of \(\Scheme.\mathtt{Modules}\,X\), discharging the pentagon,
  triangle and hexagon coherence identities. This is explicitly \emph{off-path}
  and \emph{not pursued}. One could in principle obtain it as the output of
  \(\mathtt{Localization.Monoidal.LocalizedMonoidal}\) applied to the
  already-monoidal presheaf category \(\mathtt{PresheafOfModules}\,\mathcal{O}_X\)
  and the sheafification localizer \(J.W\), supplying the instance
  \(J.W.\mathtt{IsMonoidal}\) (\cref{lem:whisker_of_W}); but the group law never
  consumes a coherent monoidal category, so this construction is superseded and not
  carried out. The substrate provides only the \emph{existence} of the three
  coherence isomorphisms that the group law actually consumes.

  Were that superseded route pursued, the instance \(J.W.\mathtt{IsMonoidal}\)
  would need \emph{no} closed structure: its
  \(\mathtt{whiskerLeft}\) field \(W\,g \Rightarrow W(F \triangleleft g)\) for an
  \emph{arbitrary} module \(F\) is the flatness-free stalkwise statement
  \cref{lem:islocallyinjective_whisker_of_W} (a \(J.W\)-morphism \(g\) is a
  stalkwise isomorphism, so \((F \triangleleft g)_x = \mathrm{id}_{F_x} \otimes
  g_x\) is again an isomorphism for any \(F\)), the same technique Mathlib uses for
  fixed-base presheaves in
  \(\mathtt{CategoryTheory.Sites.Point.IsMonoidalW}\). It would use neither
  \(\mathtt{MonoidalClosed}(\mathtt{PresheafOfModules}\,R)\) nor the
  sectionwise-flatness hypothesis of \texttt{flat\_whisker\_localizer} (which is
  false over a non-affine open and is likewise off the critical path). The \emph{full}
  \(J.W.\mathtt{IsMonoidal}\) instance and the \(\mathtt{LocalizedMonoidal}\) monoidal
  category are not on the critical path: the group law is realized without them. The
  single left-whisker existence lemma and its underlying d.1/d.2 stalk apparatus, by
  contrast, \emph{are} on the critical path, since the associator's sheafification
  transport consumes exactly that one whisker obligation
  (\cref{lem:islocallyinjective_whiskerleft_via_stalk}, closed via d.2).

  Independently of how the monoidal structure is built, the relative Picard
  group law of \cref{lem:rel_pic_sharp_groupoid} is a structure on
  \emph{isomorphism classes} of line bundles, and every group axiom there is a
  \emph{proposition} of the form \(\mathrm{Nonempty}(\cdots \cong \cdots)\).
  Such a proposition is insensitive to which isomorphism witnesses it, so the
  coherence data of the monoidal category --- pentagon, triangle, hexagon, and
  a fortiori any monoidal-closed structure --- is never \emph{consumed} by the
  group law. The group law needs only the existence of the three coherence
  isomorphisms \cref{lem:tensorobj_assoc_iso,lem:tensorobj_unit_iso,%
  lem:tensorobj_comm_iso}, supplied directly: the associator by gluing canonical
  local isomorphisms through the restriction-compatibility isomorphism
  \cref{lem:tensorobj_restrict_iso} (\cref{lem:tensorobj_assoc_iso}), and the
  unitors and braiding by the cheap \(\mathtt{mapIso}\) pattern
  (\cref{lem:tensorobj_unit_iso,lem:tensorobj_comm_iso}). The carrier --- iso-classes
  of the substrate tensor \(\otimes_X\) --- is the scheme-theoretic analogue of
  Mathlib's Picard group \(\mathtt{CommRing.Pic}\,R := \mathrm{Skeleton}(
  \mathtt{Module.Invertible}\,R)\) (\texttt{Mathlib.RingTheory.PicardGroup}), but it
  is built \emph{differently}: Mathlib reads its \(\mathtt{CommGroup}\) off the
  \(\mathrm{Skeleton}\) of the coherent monoidal category
  \(\mathtt{MonoidalCategory}(\mathtt{Module.Invertible}\,R)\) over a fixed ring,
  whereas here --- with no coherent monoidal category on
  \(\Scheme.\mathtt{Modules}\,X\) available for the varying structure sheaf --- the
  group is assembled by hand from the existence-of-isomorphism facts.
\end{remark}

\section{The lift through \(\texttt{LineBundle.OnProduct}\)}
\label{sec:tensorobj_onproduct_lift}

The consumer in chapter \cref{chap:Picard_RelPicFunctor} works not with
\(\Scheme.\mathtt{Modules}(C \times_k T)\) directly but with the subtype
\(\texttt{LineBundle.OnProduct}\,\pi_C\,\pi_T\) of locally-trivial
modules on the fiber product \(C \times_k T\), introduced in
\cref{chap:Picard_LineBundlePullback}. This section records the
\emph{existence-of-isomorphism facts on line bundles} (associativity, unit,
commutativity and inverse) that the relative Picard group law consumes. These facts
are supplied \emph{directly}: associativity by the direct-gluing associator
\cref{lem:tensorobj_assoc_iso} (through the linchpin
\cref{lem:tensorobj_restrict_iso}), unit and commutativity by the cheap
\(\mathtt{mapIso}\) pattern (\cref{lem:tensorobj_unit_iso,lem:tensorobj_comm_iso}),
and the inverse existentially from \(\otimes\)-invertibility. They are \emph{not}
read off any \(\MonoidalCategory\) structure: the route that would derive them from
\texttt{jw\_ismonoidal} (route~(e)) is superseded and off-path. The section also
records the closure of \(\otimes\) on the locally-trivial
\(\texttt{LineBundle.OnProduct}\) subtype (\cref{lem:tensorobj_lift_onproduct}) and
the interaction with the pullback functor \(\pi_T^*\). The open-immersion
restriction-compatibility isomorphism \cref{lem:tensorobj_restrict_iso} is on the
critical path precisely because it is consumed by the direct-gluing associator
\cref{lem:tensorobj_assoc_iso}; it is \emph{not} tied to the superseded
whiskering-stability lemma \cref{lem:islocallyinjective_whisker_of_W}, which belongs
to the off-path route~(e).

\begin{lemma}

\leanok
[Sectionwise lax-monoidal structure on \(\mathtt{restrictScalars}\,\varphi\)]
  \label{lem:restrictscalars_laxmonoidal}
  \lean{PresheafOfModules.restrictScalarsLaxMonoidal}
  \uses{def:scheme_modules_tensorobj}
  % NOTE: this CommRingCat-level lax-monoidal instance
  % (`PresheafOfModules.restrictScalarsLaxMonoidal`, with helpers
  % `restrictScalarsLaxε`, `restrictScalarsLaxμ`) is a self-contained
  % supplement. It is NOT used by `lem:tensorobj_restrict_iso`, whose proof
  % (open-immersion sectionwise base change) identifies restriction along an
  % open immersion with base change along a structure-sheaf isomorphism and
  % needs no lax-monoidal pushforward. Retained for potential reuse; off the
  % critical path for the tensor group law.
  Let \(\varphi : R \to S\) be a morphism of presheaves of \emph{commutative}
  rings on a site (so that both \(R\) and \(S\) factor through
  \(\mathtt{CommRingCat}\) and the category of presheaves of modules over each
  is monoidal). Then the presheaf-of-modules restriction-of-scalars functor
  \(\mathtt{PresheafOfModules.restrictScalars}\,\varphi\) is lax monoidal.
  Consequently, composing with the already-monoidal Mathlib functor
  \(\mathtt{pushforward}_0\mathtt{OfCommRingCat}\) via
  \(\mathtt{Functor.LaxMonoidal.comp}\), the full pushforward
  \[
    \mathtt{pushforward}\,\varphi
      \;\;:=\;\;
    \mathtt{pushforward}_0\,F\,R \;\circ\; \mathtt{restrictScalars}\,\varphi
  \]
  is lax monoidal. This lemma is a self-contained
  \(\mathtt{CommRingCat}\)-level lax-monoidal supplement: it stands on its own
  and is retained for potential reuse, but it is \emph{off the critical path}
  for the tensor group law. In particular it is \emph{not} an ingredient of
  \cref{lem:tensorobj_restrict_iso}, whose proof identifies restriction along an
  open immersion with base change along a structure-sheaf isomorphism and
  consequently needs no lax-monoidal structure on the pushforward.
\end{lemma}

\begin{proof}

\leanok
  This is a \emph{sectionwise} lift of the existing fact that, for a ring map
  \(f : R \to S\), the module restriction-of-scalars functor
  \(\mathtt{ModuleCat.restrictScalars}\,f\) is lax monoidal
  (\texttt{Mathlib.Algebra.Category.ModuleCat.Monoidal.Adjunction}, where it is
  obtained as the \(\mathtt{rightAdjointLaxMonoidal}\) of the
  extension/restriction adjunction \(\mathtt{extendRestrictScalarsAdj}\) ---
  i.e.\ Mathlib already uses the same mate idiom one level down). The presheaf
  monoidal structure is computed sectionwise,
  \((M_1 \otimes M_2).\mathrm{obj}\,X = M_1.\mathrm{obj}\,X \otimes M_2.\mathrm{obj}\,X\),
  and the presheaf functor \(\mathtt{restrictScalars}\,\varphi\) acts
  sectionwise as \(\mathtt{ModuleCat.restrictScalars}\,(\varphi.\mathrm{app}\,X)\).
  Hence the lax structure maps \(\mu\) (laxity) and \(\varepsilon\) (unit)
  assemble from the per-section \(\mathtt{ModuleCat}\) ones, the required
  naturality squares being the naturality of the sectionwise maps against the
  presheaf restriction morphisms. Packaging these as a
  \(\mathtt{Functor.CoreLaxMonoidal}\) yields the lax-monoidal structure; no
  new diagram chase is needed beyond transporting the sectionwise coherence.
  The composite \(\mathtt{pushforward}\,\varphi\) is then lax monoidal by
  \(\mathtt{Functor.LaxMonoidal.comp}\) together with the monoidality of
  \(\mathtt{pushforward}_0\mathtt{OfCommRingCat}\) (Mathlib), modulo a definitional
  transport across the unfolding of \(\mathtt{pushforward}\,\varphi\).

  The commutativity hypothesis is essential: the presheaf-of-modules monoidal
  category requires \(\mathtt{CommRingCat}\)-valued presheaves of rings, so the
  instance must be stated over \(\mathtt{CommRingCat}\)-factored \(R, S\). In
  the project this holds: the relevant \(\varphi\) is
  \((\mathtt{Scheme.Hom.toRingCatSheafHom}\,f).\mathrm{hom}\), which is
  structure-sheaf-valued and hence commutative-ring-valued.
\end{proof}

\begin{lemma}


\leanok
[Tensor product commutes with restriction along an open immersion]
  \label{lem:tensorobj_restrict_iso}
  \lean{AlgebraicGeometry.Scheme.Modules.tensorObj_restrict_iso}
  \uses{def:scheme_modules_tensorobj, lem:restrictscalars_ringiso_tensorequiv}
  Let \(X\) be a scheme, let \(M, N \in \Scheme.\mathtt{Modules}\,X\) be
  \emph{arbitrary} \(\mathcal{O}_X\)-modules, and let
  \(f : U \hookrightarrow X\) be an open immersion, with \((-)|_f\) the
  restriction of \(\mathcal{O}_X\)-modules along \(f\). The canonical comparison
  morphism
  \[
    (M \otimes_X N)\big|_f
      \;\xrightarrow{\;\;}\;
    M\big|_f \;\otimes_U\; N\big|_f
  \]
  is an isomorphism. No local-freeness, flatness, or line-bundle hypothesis on
  \(M, N\) is required: the statement holds for all \(\mathcal{O}_X\)-modules.
  The reason is structural to open immersions --- restriction along \(f\) is
  base change along the structure-sheaf \emph{isomorphism} induced by \(f\),
  not along a general ring map --- and base change along a ring isomorphism is
  trivially invertible and commutes with tensor products.
\end{lemma}

\begin{proof}
  \uses{def:scheme_modules_tensorobj}
  \leanok
  Write \(\varphi\) for the morphism of structure-sheaf-valued presheaves of
  rings induced by the open immersion \(f\) (concretely
  \(\varphi = (\mathtt{Scheme.Hom.toRingCatSheafHom}\,f).\mathrm{hom}\)), so that
  restriction along \(f\) is computed by the inverse-image (pullback) functor on
  presheaves of modules. The proof exploits the defining feature of an open
  immersion: the comparison of structure-sheaf rings over the image is an
  \emph{isomorphism}, not merely a ring map. For an open \(V \subseteq U\) the
  section ring \(\Gamma\bigl(X, f(V)\bigr)\) of \(\mathcal{O}_X\) over the image
  of \(V\) is identified with \(\Gamma(U, V)\) by the canonical \emph{ring
  isomorphism}
  \(f.\mathtt{appIso}\,V : \Gamma\bigl(X, f(V)\bigr) \xrightarrow{\sim}
  \Gamma(U, V)\)
  (the \(\mathtt{CommRingCat}\) isomorphism
  \(\mathtt{AlgebraicGeometry.Scheme.Hom.appIso}\)). Restriction along \(f\) is
  therefore extension of scalars along
  this ring isomorphism, and the proof needs no flatness anywhere.

  \emph{Step 1 (reduce restriction to the abstract pullback).} The restriction
  \((-)|_f\) of \(\mathcal{O}_X\)-modules along the open immersion \(f\) is
  naturally isomorphic to the abstract presheaf-of-modules pullback
  \((\mathtt{PresheafOfModules.pullback}\,\varphi).\mathrm{obj}(-)\), via the
  present comparison \(\mathtt{Scheme.Modules.restrictFunctorIsoPullback}\). It
  therefore suffices to produce the comparison and prove it an isomorphism for
  the abstract pullback functor.

  \emph{Step 2 (move the pullback inside the sheafification).} The substrate
  \(\otimes_X\) is sheafification of the presheaf-of-modules tensor product
  (\cref{def:scheme_modules_tensorobj}). The pullback commutes with
  sheafification: there is a natural isomorphism
  \(\mathtt{sheafification} \circ \mathtt{pullback}\,\varphi
   \;\cong\;
   \mathtt{PresheafOfModules.pullback}\,\varphi \circ \mathtt{sheafification}\)
  (the present \(\mathtt{SheafOfModules.sheafificationCompPullback}\)). Hence the
  whole question descends to the presheaf level, where it is the base-change
  comparison
  \[
    (\mathtt{pullback}\,\varphi).\mathrm{obj}(A \otimes B)
      \;\longrightarrow\;
    (\mathtt{pullback}\,\varphi).\mathrm{obj}\,A
      \;\otimes\;
    (\mathtt{pullback}\,\varphi).\mathrm{obj}\,B
  \]
  for the underlying presheaves of modules \(A, B\) of \(M\) and \(N\).

  % SOURCE: [Stacks Project], Modules on Ringed Spaces, Lemma
  % \ref{lemma-exactness-pushforward-pullback} (the pair $(f^*, f_*)$ of
  % pullback/pushforward of $\mathcal{O}$-modules along a morphism of ringed
  % spaces is an adjoint pair)
  % (read from references/stacks-modules.tex, L252-273).
  % SOURCE QUOTE:
  % "\begin{lemma}
  % \label{lemma-exactness-pushforward-pullback}
  % Let $f : (X, \mathcal{O}_X) \to (Y, \mathcal{O}_Y)$
  % be a morphism of ringed spaces.
  % \begin{enumerate}
  % \item The functor
  % $f_* : \textit{Mod}(\mathcal{O}_X) \to \textit{Mod}(\mathcal{O}_Y)$
  % is left exact. In fact it commutes with all limits.
  % \item The functor
  % $f^* : \textit{Mod}(\mathcal{O}_Y) \to \textit{Mod}(\mathcal{O}_X)$
  % is right exact. In fact it commutes with all colimits.
  %  [\dots]
  % \end{enumerate}
  % \end{lemma}
  %  [\dots]
  % Parts (1) and (2) hold because $(f^*, f_*)$ is an adjoint pair
  % of functors, see
  % Sheaves, Lemma \ref{sheaves-lemma-adjoint-pullback-pushforward-modules}"
  \textit{Source: [Stacks Project], Modules on Ringed Spaces, Lemma
  \texttt{lemma-exactness-pushforward-pullback}: pullback and pushforward of
  $\mathcal{O}$-modules along a morphism form an adjoint pair $(f^*, f_*)$, the
  abstract universal property on which the comparison below rests.}

  \emph{Step 3 (the genuine residual: H1 alone).} Completing this comparison at the
  presheaf level requires a sectionwise handle on the abstract functor
  \(\mathtt{PresheafOfModules.pullback}\,\varphi\), which is defined as the left
  adjoint of pushforward (the presheaf-of-modules avatar of the
  \(f^* \dashv f_*\) adjunction of the cited Stacks lemma) and admits no
  sectionwise formula in Mathlib at the pinned commit. The closure passes through
  two presheaf-level ingredients, of which exactly one now remains open.

  The first ingredient --- a strong-monoidal upgrade of
  \(\mathtt{restrictScalars}\) along the structure-sheaf ring \emph{isomorphism}
  \(f.\mathtt{appIso}\), so that base change along the isomorphism commutes with
  \(\otimes\) --- is \emph{closed in this chapter} as
  \cref{lem:restrictscalars_ringiso_tensorequiv}
  (\(\mathtt{restrictScalarsRingIsoTensorEquiv}\)): for a ring isomorphism
  \(e : R \simeq S\) the canonical map \(a \otimes_R b \mapsto a \otimes_S b\) is an
  \(R\)-linear equivalence \(\mathtt{restrictScalars}_e(M) \otimes_R
  \mathtt{restrictScalars}_e(N) \xrightarrow{\sim} \mathtt{restrictScalars}_e(M
  \otimes_S N)\).

  The \emph{sole remaining residual} is therefore the presheaf-level identification
  \[
    \mathtt{pushforward}\,\beta \;\cong\; \mathtt{pullback}\,\varphi
  \]
  via \(\mathtt{Adjunction.leftAdjointUniq}\). One exhibits a
  presheaf-level adjunction
  \(\mathtt{pushforward}\,\beta \dashv \mathtt{pushforward}\,\varphi\) and concludes
  that \(\mathtt{pushforward}\,\beta\), being a left adjoint of
  \(\mathtt{pushforward}\,\varphi\), is canonically isomorphic to the abstract left
  adjoint \(\mathtt{pullback}\,\varphi\). Here \(\beta\) is the structure map of the
  restriction functor, so that \((M|_f).\mathrm{val} =
  (\mathtt{pushforward}\,\beta).\mathrm{obj}\,M.\mathrm{val}\) definitionally;
  composing the resulting comparison with
  \cref{lem:restrictscalars_ringiso_tensorequiv} yields the asserted isomorphism.

  This presheaf-level adjunction is the presheaf analogue of the sheaf-level
  \(\mathtt{SheafOfModules.pushforwardPushforwardAdj}\), and must be built from
  presheaf-level \(\mathtt{pushforwardNatTrans}\) and
  \(\mathtt{pushforwardCongr}\). These two building blocks are \emph{Mathlib-absent}
  at the pinned commit: only the \emph{sheaf}-level versions exist
  (\texttt{Mathlib/.../SheafOfModules/Sheaf/PushforwardContinuous.lean}), whereas
  \(\mathtt{pushforwardId}\) and \(\mathtt{pushforwardComp}\) do already exist at the
  presheaf level. This residual is being constructed \emph{directly}, mirroring the
  existing sheaf-level template app-for-app at the presheaf level (a Mathlib-gradient
  build, not a deferral to an upstream PR).

  This lemma is the single substrate linchpin of the relative Picard group law for
  \emph{arbitrary} \(\mathcal{O}_X\)-modules \(M, N\): the associator
  \cref{lem:tensorobj_assoc_iso} consumes the comparison
  \((M \otimes_X N)|_f \cong M|_f \otimes_U N|_f\) on its sole consumer
  \cref{lem:tensorobj_preserves_locally_trivial}, where the restriction is commuted
  past the \emph{sheafified} tensor \(M \otimes_X N\) \emph{before} any local
  trivialisation of \(M, N\) is introduced --- so the comparison is required for the
  global, untrivialised modules, and no free-cover specialisation can stand in for it.
  Consequently the H1 residual \(\mathtt{pushforward}\,\beta \cong
  \mathtt{pullback}\,\varphi\) is genuinely on the critical path, and closing it is
  the \emph{single open substrate obligation} of this chapter. The superseded
  route~(e) \((J.W).\mathtt{IsMonoidal} \to \mathtt{Sheaf.monoidalCategory}\)
  (\texttt{jw\_ismonoidal}) remains named only as a last-resort fallback should the
  presheaf-level adjunction itself bottom out; it is not the plan.
\end{proof}

\begin{lemma}

\leanok
[Base change along a ring isomorphism commutes with $\otimes$]
  \label{lem:restrictscalars_ringiso_tensorequiv}
  \lean{restrictScalarsRingIsoTensorEquiv}
  Let \(e : R \simeq S\) be an isomorphism of commutative rings, and let \(M, N\) be
  \(S\)-modules. Restriction of scalars along \(e\) commutes with the tensor
  product: the canonical map
  \[
    \mathtt{restrictScalars}_e(M) \otimes_R \mathtt{restrictScalars}_e(N)
      \;\xrightarrow{\ \sim\ }\;
    \mathtt{restrictScalars}_e\bigl(M \otimes_S N\bigr),
    \qquad a \otimes_R b \;\longmapsto\; a \otimes_S b,
  \]
  is an \(R\)-linear equivalence. Equivalently, \(\mathtt{restrictScalars}_e\) is
  \emph{strong} monoidal, not merely lax, when \(e\) is a ring \emph{isomorphism}:
  base change along an isomorphism \(R \simeq S\) identifies the two ground rings, so
  the lax structure map of \(\mathtt{ModuleCat.restrictScalars}\) is invertible. The
  forward map is \(R\)-linear, but its inverse \(a \otimes_S b \mapsto a \otimes_R b\)
  is only \emph{additive} (it is not \(S\)-linear), so the inverse is constructed as
  a homomorphism of additive groups and then verified to be the two-sided inverse of
  the forward \(R\)-linear map. This closes the strong-monoidal ingredient of
  \cref{lem:tensorobj_restrict_iso}, leaving the presheaf-level pushforward
  adjunction as the sole open residual there.
\end{lemma}

\begin{lemma}


\leanok
[Tensor product of locally-trivial modules is locally trivial]
  \label{lem:tensorobj_preserves_locally_trivial}
  \lean{AlgebraicGeometry.Scheme.Modules.tensorObj_isLocallyTrivial}
  \uses{def:scheme_modules_tensorobj}
  Let \(X\) be a scheme, and let
  \(L, L' \in \Scheme.\mathtt{Modules}\,X\) be locally-trivial modules of
  rank one (i.e.\ line bundles). The tensor product \(L \otimes_X L'\) of
  \cref{def:scheme_modules_tensorobj} is again locally trivial of rank
  one. Consequently the bifunctor \(\otimes_X\) restricts to a bifunctor
  on the full subcategory \(\texttt{LineBundle}\,X \subset \Scheme.\mathtt{Modules}\,X\)
  of locally-trivial rank-one modules.
\end{lemma}

\begin{proof}



\leanok
  Pick a common affine open cover \(\{U_i\}\) of \(X\) on which both
  \(L|_{U_i}\) and \(L'|_{U_i}\) are free of rank one. On each \(U_i\)
  there are trivialisations \(L|_{U_i} \cong \mathcal{O}_{U_i}\) and
  \(L'|_{U_i} \cong \mathcal{O}_{U_i}\). The tensor product of two free
  rank-one modules over a commutative ring is again free of rank one;
  hence \((L \otimes_X L')|_{U_i} \cong \mathcal{O}_{U_i}\). So
  \(L \otimes_X L'\) is locally trivial of rank one, i.e.\ an object of
  \(\texttt{LineBundle}\,X\).
\end{proof}

\emph{Closed standalone result --- off the critical path.} The following lemma
\emph{is} formalized and sorry-free; it is not, however, consumed by the relative
Picard group law. The associator \cref{lem:tensorobj_assoc_iso} is realized
\emph{unconditionally} through the varying-ring stalk--tensor commutation
\cref{lem:islocallyinjective_whiskerleft_via_stalk}
(\cref{sec:tensorobj_stalk_tensor}), which closes the single left-whiskering
obligation for \emph{all} modules with no flatness and no local-triviality
hypothesis. The sectionwise-flatness whiskering statement below is therefore
neither the live whiskering-stability route nor on the critical path: it is a valid,
self-contained result about flat whiskering of sheafification-localizer morphisms,
retained for potential reuse. (Unlike the genuinely abandoned
\cref{lem:whisker_of_W}, it is a closed lemma, not a stub to be skipped.)

\emph{Superseded route --- off path, not to be formalized.} The following lemma
belongs to the abandoned coherent-monoidal-category / sheafification-whiskering
approach to the associator. The group law on iso-classes consumes only the
\emph{existence} of the coherence isomorphisms (\cref{lem:tensorobj_isoclass_commgroup}),
and the associator is realized as the sheafification transport of the presheaf
associator (\cref{lem:tensorobj_assoc_iso}), whose single open left-whisker obligation
is closed by the d.2 route of \cref{lem:islocallyinjective_whiskerleft_via_stalk}. The
\emph{full} \((J.W).\mathtt{IsMonoidal}\) / \(\mathtt{LocalizedMonoidal}\)
monoidal-category packaging is not required. This block is retained for the historical
record only and must \emph{not} be formalized.

\begin{lemma}
% LOAD-BEARING (NOT pending deletion): the Lean lemma is invoked by the associator `tensorObj_assoc_iso`, by `pullbackUnitIso`, and by `isIso_sheafifyDelta_unitPair_of_isIso_sheafifyEta`. It is wired into the graph as a dependency of \cref{lem:tensorobj_assoc_iso}.

\leanok
[Sheafification inverts a localizer morphism of presheaves of modules]
  \label{lem:isiso_sheafification_map_of_W}
  \lean{PresheafOfModules.isIso_sheafification_map_of_W}
  Let \(R_0\) be a presheaf of commutative rings on a site carrying a
  Grothendieck topology \(J\), let \(R_{sh}\) be its sheafification, and let
  \(\alpha : R_0 \to R_{sh}.\mathrm{obj}\) be the (locally bijective) comparison
  map of presheaves of rings. Let \(f : A \to B\) be a morphism of presheaves of
  \(R_0\)-modules whose underlying morphism of presheaves of abelian groups
  \((\mathtt{toPresheaf}\,R_0).\mathrm{map}\,f\) lies in the sheafification
  localizer class \(J.W\) --- equivalently, is locally bijective for \(J\). Then
  the presheaf-of-modules sheafification functor
  \(\mathtt{PresheafOfModules.sheafification}\,\alpha\) sends \(f\) to an
  isomorphism:
  \[
    (\mathtt{sheafification}\,\alpha).\mathrm{map}\,f
    \quad\text{is an isomorphism.}
  \]
  This is a closed go/no-go bridge --- it converts ``\,a morphism lies in
  \(J.W\)\,'' into ``\,its sheafification is invertible\,''. Under route~(e) the
  associator is the API-derived one (\texttt{jw\_ismonoidal}), so this bridge is
  retained as a supplement rather than as the associator's load-bearing step.
\end{lemma}

\begin{proof}

\leanok
  This is a one-morphism reading of the structural identity that the
  sheafification of presheaves of \(R_0\)-modules \emph{is} the localization of
  \(\mathtt{PresheafOfModules}\,R_0\) at the class
  \(J.W.\mathtt{inverseImage}\,(\mathtt{toPresheaf}\,R_0)\) of morphisms whose
  underlying map of presheaves of abelian groups is locally bijective. Concretely,
  the underlying \(\mathtt{AddCommGrp}\)-valued sheafification inverts exactly the
  locally bijective morphisms, and Mathlib's
  \(\mathtt{PresheafOfModules.inverseImage\_W\_toPresheaf\_eq\_inverseImage\_%
  isomorphisms}\) identifies the preimage of \(J.W\) under
  \(\mathtt{toPresheaf}\,R_0\) with the preimage of the class of isomorphisms under
  the module-level sheafification functor. Reading that equality of morphism
  classes at the single morphism \(f\): the hypothesis
  \(J.W\,\bigl((\mathtt{toPresheaf}\,R_0).\mathrm{map}\,f\bigr)\) places \(f\) in
  that preimage, hence \((\mathtt{sheafification}\,\alpha).\mathrm{map}\,f\) is an
  isomorphism. No new construction is required beyond instantiating the existing
  Mathlib identity at one morphism; this lemma is a closed thin wrapper, retained
  as a supplement under route~(e).
\end{proof}

\section{Route~(e): the \(J.W\)-monoidality instance and the localized monoidal structure}
\label{sec:tensorobj_route_e}

\emph{Full-monoidal-category construction --- off path, not to be formalized.} The
construction below realizes a \emph{full} symmetric monoidal category on all of
\(\Scheme.\mathtt{Modules}\,X\) by instantiating Mathlib's
localization-of-monoidal-categories API. It was the original plan for supplying the
associator. The relative Picard group law, however, is a structure on
\emph{isomorphism classes} of invertible sheaves and consumes only the
\emph{existence} of the associator, unitor and braiding isomorphisms
(\cref{lem:tensorobj_isoclass_commgroup}, built by hand as the scheme-theoretic
analogue of Mathlib's \(\mathtt{CommRing.Pic}\), not off a coherent monoidal
category's \(\mathtt{Skeleton}\)). The associator is obtained as the sheafification
transport of the presheaf associator (\cref{lem:tensorobj_assoc_iso}), whose
\emph{single} open left-whisker obligation is closed by the d.2 stalk--tensor
commutation (\cref{lem:islocallyinjective_whiskerleft_via_stalk,%
lem:stalk_tensor_commutation}); it uses \emph{no} \emph{full}
\((J.W).\mathtt{IsMonoidal}\) instance and \emph{no} \(\mathtt{LocalizedMonoidal}\)
monoidal category. Accordingly the \emph{full monoidal-category} target of this
section --- the \emph{pair} of whiskering-stability fields (\cref{lem:whisker_of_W})
and the \(\mathtt{LocalizedMonoidal}\) instance (\texttt{jw\_ismonoidal}) --- is off
path and must not be re-attempted; it is kept for the historical record. The d.1/d.2
stalk apparatus it shares with the live single whisker obligation is \emph{not}
vestigial, however: it is on the critical path through
\cref{lem:islocallyinjective_whiskerleft_via_stalk}.

The monoidal structure on the substrate is realized by \emph{instantiating}
Mathlib's localization-of-monoidal-categories API
(\cref{sec:tensorobj_api_survey}): apply
\(\mathtt{Localization.Monoidal.LocalizedMonoidal}\) to the already-monoidal
presheaf category \(\mathtt{PresheafOfModules}\,\mathcal{O}_X\) and the
sheafification localizer \(J.W\). The single new input the API requires is the
instance \(J.W.\mathtt{IsMonoidal}\), whose two fields are the
whiskering-stability lemmas below.

\paragraph{What the group law consumes.} The group law is assembled directly on
isomorphism classes of \emph{locally trivial} (invertible) sheaves
(\cref{lem:tensorobj_isoclass_commgroup}), mirroring Mathlib's ring-level Picard
group \(\mathtt{Module.Invertible}\) / \(\mathtt{CommRing.Pic}\)
(\texttt{Mathlib/RingTheory/PicardGroup.lean}). It needs only the \emph{existence} of
the associator \cref{lem:tensorobj_assoc_iso}, the unitors \cref{lem:tensorobj_unit_iso}
and the braiding \cref{lem:tensorobj_comm_iso}, together with \(\otimes\)-invertibility
supplying inverses; it uses \emph{no} \(J.W.\mathtt{IsMonoidal}\) instance, \emph{no}
\(\mathtt{LocalizedMonoidal}\) monoidal category, and \emph{no} \(\mathtt{Skeleton}\).
It does, however, consume the associator \cref{lem:tensorobj_assoc_iso}, which is
realized as the sheafification transport of the presheaf associator and whose
\emph{single} open left-whisker obligation \cref{lem:islocallyinjective_whiskerleft_via_stalk}
is closed by the stalk--tensor commutation (d.2, \cref{lem:stalk_tensor_commutation}).
Hence the stalk ingredients (d.1)/(d.2) are \emph{on} the critical path of the group
law --- through the associator's single whisker obligation --- even though the full
\(J.W.\mathtt{IsMonoidal} \to \mathtt{LocalizedMonoidal}\) monoidal category on all of
\(\mathtt{SheafOfModules}\) (the subject of the present section) is \emph{not} built:
that full monoidal generality, with its \emph{pair} of whiskering-stability fields
(\cref{lem:whisker_of_W}) and all derived coherence, is what the group law does not
consume. The single open obligation of the present section's full-monoidal target is
thus the local-injectivity half packaged for an arbitrary whisker object
(\cref{lem:islocallyinjective_whisker_of_W}), itself a consequence of the same d.2
commutation.

The following lemma is the \(R_x\)-linear stalk-map packaging (the ``d.1''
ingredient). It is consumed by the stalk--tensor commutation of
\cref{sec:tensorobj_stalk_tensor} (\cref{lem:stalk_tensor_commutation,
lem:islocallyinjective_whiskerleft_via_stalk}) and is therefore a live ingredient
of the unconditional associator route.

\begin{lemma}

\leanok
% Live ingredient: consumed by lem:stalk_tensor_commutation and
% lem:islocallyinjective_whiskerleft_via_stalk (sec:tensorobj_stalk_tensor).

[The $R.\mathtt{stalk}\,x$-linear stalk map of a morphism of presheaves of modules]
  \label{lem:stalk_linear_map}
  \lean{PresheafOfModules.stalkLinearMap}
  \uses{def:scheme_modules_tensorobj}
  Let \(X\) be a topological space, \(R\) a presheaf of commutative rings on
  \(X\), and \(x \in X\) a point. For a morphism \(g : M \to N\) of presheaves of
  \(R\)-modules, the induced map on stalks
  \[
    g_x \;\colon\; M_x \;\longrightarrow\; N_x
  \]
  is \(R.\mathtt{stalk}\,x\)-linear (\(\mathtt{stalkLinearMap}\)), and is
  characterised on germs by
  \(g_x(\mathrm{germ}_x\,s) = \mathrm{germ}_x(g_U\,s)\)
  (\(\mathtt{stalkLinearMap\_germ}\)). When the underlying map of
  \(\mathtt{Ab}\)-valued stalks is an isomorphism, \(g_x\) is bijective
  (\(\mathtt{stalkLinearMap\_bijective\_of\_isIso}\)); in that case it bundles to
  an \(R.\mathtt{stalk}\,x\)-linear equivalence
  \(M_x \xrightarrow{\sim} N_x\) (\(\mathtt{stalkLinearEquivOfIsIso}\)).

  These four declarations are the \emph{(d.1)-partial implementation} --- the
  \(R_x\)-linearity packaging of stalk maps --- consumed by the
  \(\mathrm{id} \otimes g_x\) step in the proofs of
  \cref{lem:stalk_tensor_commutation} and
  \cref{lem:islocallyinjective_whiskerleft_via_stalk}. They build on the
  stalk module structure of \texttt{Mathlib.Algebra.Category.ModuleCat.Stalk};
  the linearity packaging itself is project-original.
\end{lemma}

\begin{proof}
  \uses{def:scheme_modules_tensorobj}
  \leanok
  The stalk \(M_x = \mathtt{stalk}\,M.\mathtt{presheaf}\,x\) carries its
  \(R.\mathtt{stalk}\,x\)-module structure from
  \texttt{Mathlib.Algebra.Category.ModuleCat.Stalk}. The underlying additive map
  \(g_x\) on stalks is the filtered colimit of the section maps \(g_U\); since
  each section map is \(R(U)\)-linear and the scalar action on the stalk is the
  colimit of the section actions, \(g_x\) respects the
  \(R.\mathtt{stalk}\,x\)-action, which gives the linear map together with its
  germ characterisation. Bijectivity is read off the underlying
  \(\mathtt{Ab}\)-stalk isomorphism (a bijective linear map is a linear
  equivalence), and bundling the inverse yields the linear equivalence.
\end{proof}

\emph{Off-path duplicate --- the full-generality packaging, not to be formalized.}
The following lemma is the load-bearing whiskering-stability field of the \emph{full}
\((J.W).\mathtt{IsMonoidal}\) / \(\mathtt{LocalizedMonoidal}\) monoidal-category route
to the associator. That \emph{full monoidal category} is not built: the group law on
iso-classes consumes only the \emph{existence} of the associator
(\cref{lem:tensorobj_isoclass_commgroup}). The associator instead arises as the
sheafification transport whose single open left-whisker obligation is closed by the
live lemma \cref{lem:islocallyinjective_whiskerleft_via_stalk}. The stalk ingredients
(d.1)/(d.2) themselves --- the \(R_x\)-linear stalk map \cref{lem:stalk_linear_map}
and the stalk--tensor commutation \cref{lem:stalk_tensor_commutation} --- are
\emph{on} the critical path through that live lemma; what is off path is only the
\emph{full-generality} \(J.W.\mathtt{IsMonoidal}\) packaging of the following
duplicate, retained for the historical record.

\begin{lemma}
% SUPERSEDED route (route-(d)/(e) whiskering/stalk machinery); pending deletion once the assoc re-route (morphism-level descent) lands. Declaration still present and still backing the current `tensorObj_assoc_iso` proof.

\leanok
[Local injectivity of a whiskered localizer morphism (stalkwise, flatness-free)]
  \label{lem:islocallyinjective_whisker_of_W}
  % \lean: this superseded, off-critical-path obligation is the SAME statement
  % (left-whiskering of a J.W-morphism by an arbitrary F is locally injective) as
  % the LIVE lemma lem:islocallyinjective_whiskerleft_via_stalk, and is realised by
  % the same sorry-free decl. Both blueprint nodes legitimately share the pin
  % (leandag flags only Lean-side name collisions, not shared blueprint pins).
  \lean{PresheafOfModules.isLocallyInjective_whiskerLeft_of_W}
  \uses{def:scheme_modules_tensorobj, lem:stalk_linear_map, lem:tensorobj_restrict_iso, lem:tensorobj_unit_iso}
  Let \(R\) be a presheaf of commutative rings on a site carrying a Grothendieck
  topology \(J\), and let \(J.W\) be the sheafification localizer (the locally
  bijective morphisms). Concretely the site carries the
  \(\mathtt{WEqualsLocallyBijective}\) property for \(\mathtt{Ab}\), so \(J.W\) is
  exactly the class of \(J\)-locally bijective morphisms; this is recorded as a
  typeclass hypothesis \([J.\mathtt{WEqualsLocallyBijective}\,\mathtt{Ab}]\) on
  the statement. Let \(F\) be an \emph{arbitrary} presheaf of
  \(R\)-modules, and let \(g : M \to N\) be a morphism of presheaves of
  \(R\)-modules whose underlying morphism of presheaves of abelian groups
  \((\mathtt{toPresheaf}\,R).\mathrm{map}\,g\) lies in \(J.W\). Then the
  left-whiskered morphism
  \[
    F \triangleleft g \;\;=\;\; \mathrm{id}_F \otimes g
  \]
  is \emph{\(J\)-locally injective}. No flatness and no local-triviality
  hypothesis on \(F\) is required.

  This is the \(\mathtt{whiskerLeft}\)-field content (the injectivity half) of the
  \emph{full} instance \(J.W.\mathtt{IsMonoidal}\) (\texttt{jw\_ismonoidal}), the
  load-bearing field of the off-path \(\mathtt{LocalizedMonoidal}\) monoidal-category
  path (arbitrary \(F\)). Its content for an arbitrary \(F\) is exactly the live lemma
  \cref{lem:islocallyinjective_whiskerleft_via_stalk}, which the associator's
  sheafification transport consumes and which is closed \emph{unconditionally} by the
  d.2 stalk--tensor commutation (\cref{lem:stalk_tensor_commutation}); the d.1/d.2
  ingredients are therefore on the critical path. The two proofs recorded below --- a
  sheaf-level argument specialised to locally trivial \(F\), and the general stalkwise
  argument via (d.1)/(d.2) --- are retained for the historical record; the live group
  law closes this obligation in its general (d.2) form through
  \cref{lem:islocallyinjective_whiskerleft_via_stalk}.
\end{lemma}

\begin{proof}
  \uses{def:scheme_modules_tensorobj, lem:stalk_linear_map, lem:tensorobj_restrict_iso, lem:tensorobj_unit_iso}

\leanok
  \emph{PRIMARY route (\(F\) locally trivial --- the consumer's case; sheaf-level,
  no stalks, no (d.2)).} Suppose \(F\) is locally trivial
  (\(\mathtt{LineBundle.IsLocallyTrivial}\,F\)), which is exactly the hypothesis its
  sole consumer --- the associator \cref{lem:tensorobj_assoc_iso} --- supplies.
  \(J\)-local injectivity of a morphism is \emph{local on a covering family}: it
  suffices to produce a cover of \(X\) on each member of which \(F \triangleleft g\)
  restricts to a locally injective morphism. Take a trivialising open cover
  \(\{V\}\) of \(X\) for \(F\), so that \(F|_V \cong \mathcal{O}_V\) on each \(V\).
  On such a \(V\):
  \begin{itemize}
    \item the substrate--restriction compatibility \cref{lem:tensorobj_restrict_iso}
      gives a canonical isomorphism
      \[
        (F \triangleleft g)\big|_V \;\cong\; (F|_V) \triangleleft (g|_V)
      \]
      --- restriction along the open immersion \(V \hookrightarrow X\) commutes with
      \(\otimes\), and the whiskering \(F \triangleleft g = \mathrm{id}_F \otimes g\)
      is \(\otimes\) with an identity, so it restricts to \(\mathrm{id}_{F|_V}
      \otimes (g|_V) = (F|_V) \triangleleft (g|_V)\);
    \item the trivialisation \(F|_V \cong \mathcal{O}_V\) together with the left
      unitor \cref{lem:tensorobj_unit_iso} (\(\mathcal{O}_V \otimes_V (-) \cong
      (-)\)) identifies \((F|_V) \triangleleft (g|_V) \cong g|_V\).
  \end{itemize}
  Composing, \((F \triangleleft g)|_V\) is isomorphic to \(g|_V\). Now \(g \in J.W\),
  and membership in \(J.W\) (local bijectivity) is stable under restriction along an
  open immersion, so \(g|_V \in J.W\); in particular \(g|_V\) is \(J\)-locally
  injective. An isomorphism preserves local injectivity, hence \((F \triangleleft
  g)|_V\) is locally injective on \(V\). As the \(\{V\}\) cover \(X\) and local
  injectivity is local on the cover, \(F \triangleleft g\) is \(J\)-locally
  injective. \emph{No stalks and no (d.2) stalk--tensor commutation enter}; the
  single comparison lemma \cref{lem:tensorobj_restrict_iso} (a presheaf-level
  comparison, decomposed into its two open-immersion base-change halves) carries the
  argument, and thereby moves \emph{onto} the critical path.

  This is \emph{not} the section-level route that tensors a bare section map \(g(V)\)
  --- only injective --- and stalls on a \(\mathrm{Tor}_1\)/flatness obligation. Here
  the argument is sheaf-level and uses that \(g\) is locally \emph{bijective} (a
  \(J.W\) morphism), reduced through the unitor to \(g|_V\) (again a \(J.W\)
  morphism); a mere injection is never tensored, so no flatness is needed.

  \emph{FALLBACK route (arbitrary \(F\); stalkwise, required only for the
  full-generality route~(e) \(\mathtt{LocalizedMonoidal}\) path).} For an
  \emph{arbitrary} (not necessarily locally trivial) \(F\), local injectivity is
  proved stalkwise, by the argument Mathlib uses for fixed-base presheaves in
  \(\mathtt{CategoryTheory.Sites.Point.IsMonoidalW}\)
  (\cref{sec:tensorobj_api_survey}), ported to \(\mathtt{PresheafOfModules}\). A
  morphism in \(J.W\) is, on the topological site of \(X\), a \emph{stalkwise
  isomorphism}: at each point \(x\) the induced map on stalks \(g_x\) is an
  isomorphism. The presheaf-of-modules tensor is sectionwise, so it commutes with
  the stalk (a filtered colimit), giving
  \[
    (F \triangleleft g)_x
      \;=\; \mathrm{id}_{F_x} \otimes_{R_x} g_x
      \;\colon\; F_x \otimes_{R_x} M_x \;\longrightarrow\; F_x \otimes_{R_x} N_x .
  \]
  Tensoring an isomorphism \(g_x\) by the identity \(\mathrm{id}_{F_x}\) yields an
  isomorphism --- no flatness is needed, since \(g_x\) is invertible, not merely
  injective. Hence \(F \triangleleft g\) is a stalkwise isomorphism, in particular
  locally bijective, in particular locally injective, for \emph{any} \(F\).

  The stalkwise computation rests on three ingredients with distinct status at the
  pinned commit:
  \begin{enumerate}
    \item[(d.1-done)] the \(R.\mathtt{stalk}\,x\)-linear stalk-map packaging
      (\(\mathtt{stalkLinearMap}\) and its companions, \cref{lem:stalk_linear_map}):
      a morphism of presheaves of modules induces an \(R.\mathtt{stalk}\,x\)-linear
      map on stalks, characterised on germs, that becomes a linear equivalence once
      the underlying \(\mathtt{Ab}\)-stalk map is an isomorphism. This is built
      project-side, axiom-clean, and furnishes the \(\mathrm{id} \otimes g_x\) step
      above.
    \item[(d.1-bridge)] (remaining) the site-\(J.W\) \(\Longleftrightarrow\)
      stalkwise-iso characterisation on \(\mathtt{Opens}\,X\): that a morphism lies
      in \(J.W\) iff its stalk maps are isomorphisms at every point. The concrete
      Mathlib assembly route uses the \(\mathtt{WEqualsLocallyBijective}\) structure
      together with the topological stalk criteria --- local injectivity
      \(\Longleftrightarrow\) injectivity on stalks and local surjectivity
      \(\Longleftrightarrow\) surjectivity on stalks --- bridging the site-level
      \(\mathtt{Presheaf.IsLocallyInjective}\) / \(\mathtt{IsLocallySurjective}\)
      to the \(\mathtt{TopCat.Presheaf}\) stalk versions. This half is still
      Mathlib-absent at the \(\mathtt{PresheafOfModules}\) level.
    \item[(d.2)] (remaining) the commutation \((A \otimes^{p} B)_x \cong A_x
      \otimes_{R_x} B_x\) of the stalk (a filtered colimit) with the relative
      module-presheaf tensor, under which \((F \triangleleft g)_x\) is identified
      with \(\mathrm{lTensor}\,F_x\,(g_x) = \mathrm{id}_{F_x} \otimes g_x\) (using
      that \(\mathtt{tensorLeft}\) / \(\mathtt{tensorRight}\) preserve filtered
      colimits over a module category). This is genuinely Mathlib-absent over a
      \emph{varying} ring and is the largest remaining piece.
  \end{enumerate}
  Mathlib supplies the analogues for presheaves valued in a fixed monoidal concrete
  category, but not for modules over a varying ring; porting (d.1-bridge) and (d.2)
  to \(\mathtt{PresheafOfModules}\) is the remaining work. Once (d.2) lands, the
  conclusion is immediate and flatness-free: by
  \(\mathtt{stalkLinearMap\_bijective\_of\_isIso}\) the stalk map \(g_x\) is a
  linear equivalence, and the linear-equivalence \(\mathrm{lTensor}\) by \(F_x\)
  carries it to an isomorphism, finishing the proof.

  This stalkwise fallback is required \emph{only} if the full-generality
  \(J.W.\mathtt{IsMonoidal}\) / \(\mathtt{LocalizedMonoidal}\) route~(e) --- the
  monoidal category on \emph{all} of \(\Scheme.\mathtt{Modules}\,X\), for an
  arbitrary whisker object \(F\) --- is later wanted. For the relative Picard group
  law the whisker objects are locally trivial, and the PRIMARY route above suffices,
  trading the Mathlib-absent (d.2) for the single comparison
  \cref{lem:tensorobj_restrict_iso}.
\end{proof}

\emph{Superseded route --- off path, not to be formalized.} The following lemma
supplies the two whiskering-stability fields of the abandoned
\((J.W).\mathtt{IsMonoidal}\) instance (\texttt{jw\_ismonoidal}). The group law
needs no such instance: the associator is built directly by gluing canonical local
isomorphisms (\cref{lem:tensorobj_assoc_iso}, via \cref{lem:tensorobj_restrict_iso}).
This block is retained for the historical record only and must \emph{not} be
formalized.

\begin{lemma}

\leanok
% SUPERSEDED route (route-(d)/(e) whiskering/stalk machinery); pending deletion once the assoc re-route (morphism-level descent) lands. Declaration still present and still backing the current `tensorObj_assoc_iso` proof.
[Flatness-free whiskering preserves the sheafification localizer]
  \label{lem:whisker_of_W}
  \lean{PresheafOfModules.W_whiskerLeft_of_W}
  \uses{def:scheme_modules_tensorobj, lem:islocallyinjective_whisker_of_W}
  With notation as in \cref{lem:islocallyinjective_whisker_of_W}: for an
  \emph{arbitrary} presheaf of \(R\)-modules \(F\) and a morphism \(g\) whose
  underlying additive-presheaf map lies in \(J.W\), both whiskerings again lie in
  \(J.W\):
  \[
    g \in J.W \;\Longrightarrow\; F \triangleleft g \in J.W,
    \qquad
    g \in J.W \;\Longrightarrow\; g \triangleright F \in J.W .
  \]
  These two statements are precisely the \(\mathtt{whiskerLeft}\) and
  \(\mathtt{whiskerRight}\) data required by the class
  \(\mathtt{MorphismProperty.IsMonoidal}\) (\cref{sec:tensorobj_api_survey}) for
  the morphism property \(J.W\); together they furnish the instance
  \(J.W.\mathtt{IsMonoidal}\) of \texttt{jw\_ismonoidal}. No flatness and no
  local triviality is used --- they are the flatness-free replacements for the
  flat variants of \texttt{flat\_whisker\_localizer}.
\end{lemma}

\begin{proof}
  \uses{lem:islocallyinjective_whisker_of_W}
  \leanok
  By the locally-bijective characterisation of \(J.W\), it suffices to show
  \(F \triangleleft g\) is both locally injective and locally surjective.
  \emph{Local surjectivity} is free and needs no hypothesis on \(F\): the presheaf
  tensor is right exact, so \(F \otimes^{p} -\) carries the locally-zero cokernel
  of \(g\) to the cokernel of \(F \triangleleft g\), which is therefore locally
  zero (Mathlib \(\mathtt{isLocallySurjective\_whiskerLeft}\)). \emph{Local
  injectivity} is \cref{lem:islocallyinjective_whisker_of_W}. The two together
  give \(F \triangleleft g \in J.W\). The right-whiskered statement
  \(g \triangleright F\) follows by conjugating \(F \triangleleft g\) with the
  braiding of the symmetric presheaf monoidal structure, which is an isomorphism
  and hence preserves membership in \(J.W\).
\end{proof}

\emph{Superseded route --- off path, not to be formalized.} The following lemma is
the route~(e) target: a full \(\MonoidalCategory\) structure on all of
\(\Scheme.\mathtt{Modules}\,X\) from \(\mathtt{LocalizedMonoidal}\). It is not
needed. The group law consumes only the existence of the associator, unitor and
braiding isomorphisms (\cref{lem:tensorobj_isoclass_commgroup}, assembled by hand
rather than off a coherent monoidal category's \(\mathtt{Skeleton}\)), and the
associator is now built directly by
gluing (\cref{lem:tensorobj_assoc_iso}, via \cref{lem:tensorobj_restrict_iso}). This
block is retained for the historical record only and must \emph{not} be formalized.

\begin{lemma}


\leanok
[Associator for $\otimes_X$ (unconditional)]
  \label{lem:tensorobj_assoc_iso}
  \lean{AlgebraicGeometry.Scheme.Modules.tensorObj_assoc_iso}
  \uses{def:scheme_modules_tensorobj, lem:islocallyinjective_whiskerleft_via_stalk, lem:isiso_sheafification_map_of_W}
  Let \(X\) be a scheme and let
  \(M, N, P \in \Scheme.\mathtt{Modules}\,X\). There is an isomorphism
  \[
    (M \otimes_X N) \otimes_X P
      \;\xrightarrow{\ \sim\ }\;
    M \otimes_X (N \otimes_X P)
  \]
  in \(\Scheme.\mathtt{Modules}\,X\), natural in \(M, N, P\). The group law on
  \(\otimes\)-iso-classes (\cref{lem:tensorobj_isoclass_commgroup}) consumes only
  the \emph{existence} of this isomorphism, the proposition
  \(\mathrm{Nonempty}\bigl((M \otimes_X N) \otimes_X P \cong M \otimes_X (N
  \otimes_X P)\bigr)\); no pentagon, triangle, or naturality coherence of it is
  ever required.

  The Lean pin is unconditional: it holds for arbitrary
  \(\mathcal{O}_X\)-modules \(M, N, P\), with no local-triviality or
  invertibility hypothesis. (Earlier drafts carried vestigial
  \(\mathtt{IsLocallyTrivial}\) hypotheses to match the carrier; these were
  dropped once the construction below was seen to use none of them.)
\end{lemma}

\begin{proof}
  \uses{def:scheme_modules_tensorobj, lem:islocallyinjective_whiskerleft_via_stalk}
  \leanok
  \textbf{Realization (sheafification transport, unconditional via d.2).} The
  associator is the sheafification transport of the presheaf-of-modules associator
  \(\mathtt{PresheafOfModules.monoidalCategoryStruct.associator}\), which exists for
  \emph{all} presheaves of modules over the varying ring. The transport requires the
  single left-whisker \(F \triangleleft \mathtt{toSheafify}\) to lie in the
  sheafification localizer \(J.W\) for the relevant \(F\); that single open
  left-whiskering obligation \(\mathtt{isLocallyInjective\_whiskerLeft\_of\_W}\) is
  closed \emph{unconditionally} --- for an arbitrary whiskering object, with no
  flatness and no local triviality --- by the d.2 stalk--tensor commutation
  \cref{lem:islocallyinjective_whiskerleft_via_stalk} (\cref{sec:tensorobj_stalk_tensor}):
  a \(J.W\)-morphism \(g\) is a stalkwise isomorphism, and by the d.2 commutation
  \(\bigl(F \triangleleft g\bigr)_x \cong \mathrm{id}_{F_x} \otimes g_x\) is again an
  isomorphism, so \(F \triangleleft g \in J.W\). Hence the associator is itself
  unconditional and axiom-clean once d.2 (\cref{lem:stalk_tensor_commutation}) is
  formalized; it does \emph{not} require the full \(J.W.\mathtt{IsMonoidal}\) instance
  or the \(\mathtt{LocalizedMonoidal}\) monoidal category --- only the one left-whisker
  lemma --- and no sectionwise flatness enters. The local-gluing description below
  records the equivalent concrete local realization through the
  restriction-compatibility comparison \cref{lem:tensorobj_restrict_iso}.

  Concretely, the same isomorphism is described by gluing canonical local isomorphisms
  through the restriction-compatibility comparison \cref{lem:tensorobj_restrict_iso}.
  Fix a
  common open cover \(\{U\}\) of \(X\). On each member \(U\) form the canonical
  composite
  \[
    \bigl((M \otimes_X N) \otimes_X P\bigr)\big|_U
      \;\xrightarrow{\ \sim\ }\;
    \bigl(M|_U \otimes_U N|_U\bigr) \otimes_U P|_U
      \;\xrightarrow{\ \sim\ }\;
    M|_U \otimes_U \bigl(N|_U \otimes_U P|_U\bigr)
      \;\xrightarrow{\ \sim\ }\;
    \bigl(M \otimes_X (N \otimes_X P)\bigr)\big|_U ,
  \]
  in which:
  \begin{itemize}
    \item the first arrow is two applications of the restriction-compatibility
      isomorphism \cref{lem:tensorobj_restrict_iso} along the open immersion
      \(U \hookrightarrow X\): first \((M \otimes_X N)|_U \cong M|_U \otimes_U
      N|_U\) tensored with \(P|_U\), then \(\bigl((M \otimes_X N) \otimes_X
      P\bigr)|_U \cong (M \otimes_X N)|_U \otimes_U P|_U\);
    \item the middle arrow is the \emph{canonical presheaf-level associator}
      \(\alpha\) of the (symmetric) monoidal category
      \(\mathtt{PresheafOfModules}\,\mathcal{O}_U\)
      (\(\mathtt{PresheafOfModules.monoidalCategoryStruct.associator}\),
      \cref{sec:tensorobj_api_survey}(i)), read on the restricted modules;
    \item the third arrow is again two applications of
      \cref{lem:tensorobj_restrict_iso}, in the opposite direction.
  \end{itemize}
  Each of the three arrows is a \emph{canonical} isomorphism: the
  restriction-compatibility comparison \cref{lem:tensorobj_restrict_iso} is induced
  by the open immersion with no choices, hence natural in its module arguments, and
  the presheaf associator \(\alpha\) is a fixed natural transformation. Therefore,
  on a double overlap \(U \cap U'\), restricting the \(U\)-composite and the
  \(U'\)-composite to \(U \cap U'\) yields the same morphism: each of the three
  arrows commutes with the further restriction to \(U \cap U'\) by the naturality of
  \cref{lem:tensorobj_restrict_iso} and of the presheaf associator, so the two local
  composites agree on the overlap.

  Because the internal hom \(\mathcal{H}om(-,-)\) of sheaves of
  \(\mathcal{O}_X\)-modules is itself a sheaf on \(X\), a family of morphisms defined
  on the members of an open cover and agreeing on the overlaps glues to a single
  global morphism. Applying this both to the local composites above and to their
  local inverses produces a global morphism
  \((M \otimes_X N) \otimes_X P \to M \otimes_X (N \otimes_X P)\) whose restriction
  to each \(U\) is the local composite, together with a two-sided inverse glued from
  the local inverses; hence the glued morphism is a global isomorphism --- the
  associator. The single non-formal input of this re-route is the
  restriction-compatibility comparison \cref{lem:tensorobj_restrict_iso}; the presheaf
  associator (\cref{sec:tensorobj_api_survey}(i)) and the sheaf-of-modules
  internal-hom gluing are already available.

  Note that, unlike the pointwise-existence statement
  \cref{lem:tensorobj_preserves_locally_trivial} --- a \emph{proposition} proved
  cover-pointwise with \emph{no} overlap obligation --- the associator is
  \emph{global data}, a single isomorphism of sheaves on all of \(X\). Assembling it
  from local pieces is therefore genuinely stronger than that pattern: the
  overlap-naturality compatibility of the local composites must be discharged before
  gluing, and is discharged above. The locally-trivial hypotheses of the statement
  are not consumed; the comparison \cref{lem:tensorobj_restrict_iso} is applied to the
  global, untrivialised restrictions \((M \otimes_X N)|_U\), so no free-cover
  specialisation can replace it. What this proof does \emph{not} require is the
  \emph{full} \((J.W).\mathtt{IsMonoidal}\) instance or the
  \(\mathtt{LocalizedMonoidal}\) monoidal category
  (\cref{lem:whisker_of_W}): only the single left-whisker obligation
  \cref{lem:islocallyinjective_whiskerleft_via_stalk}, closed via the d.2 stalk--tensor
  commutation \cref{lem:stalk_tensor_commutation}, is consumed by the sheafification
  transport. The group law on \(\otimes\)-iso-classes
  consumes only the \emph{existence} of the resulting isomorphism, never its pentagon,
  triangle, or naturality coherence.
\end{proof}

\begin{lemma}
[Left and right unitors for $\otimes_X$]
  \label{lem:tensorobj_unit_iso}
  \lean{AlgebraicGeometry.Scheme.Modules.tensorObj_left_unitor, AlgebraicGeometry.Scheme.Modules.tensorObj_right_unitor}
  \uses{def:scheme_modules_tensorobj}
  Let \(X\) be a scheme and \(M \in \Scheme.\mathtt{Modules}\,X\) an
  \emph{arbitrary} \(\mathcal{O}_X\)-module. With
  \(\mathcal{O}_X = \mathtt{SheafOfModules.unit}\,X.\mathtt{ringCatSheaf}\) the
  structure-sheaf unit object, there are natural isomorphisms
  \[
    \mathcal{O}_X \otimes_X M \;\xrightarrow{\ \sim\ }\; M,
    \qquad
    M \otimes_X \mathcal{O}_X \;\xrightarrow{\ \sim\ }\; M .
  \]
\end{lemma}

\begin{proof}
  \uses{def:scheme_modules_tensorobj}
  The presheaf monoidal structure on \(\mathtt{PresheafOfModules}\,R\) supplies
  the left and right unitors \(\lambda, \rho\) at the presheaf level. Their
  sheafifications \(\mathtt{sheafification.mapIso}(\lambda)\),
  \(\mathtt{sheafification.mapIso}(\rho)\), composed with the sheafification
  counit identifying the sheafified structure-sheaf presheaf with
  \(\mathcal{O}_X\) (the step already used in the structure-sheaf-unit
  comparison \(\mathtt{tensorObj\_unit\_iso}\) in the substrate file), give the
  two isomorphisms. This route is the cheap \(\mathtt{mapIso}\) pattern and uses no
  abstract pullback.
\end{proof}

\begin{lemma}


\leanok
[Braiding for $\otimes_X$]
  \label{lem:tensorobj_comm_iso}
  \lean{AlgebraicGeometry.Scheme.Modules.tensorObj_braiding}
  \uses{def:scheme_modules_tensorobj}
  Let \(X\) be a scheme and \(M, N \in \Scheme.\mathtt{Modules}\,X\)
  \emph{arbitrary} \(\mathcal{O}_X\)-modules. There is a natural isomorphism
  \[
    M \otimes_X N \;\xrightarrow{\ \sim\ }\; N \otimes_X M .
  \]
\end{lemma}

\begin{proof}
  \uses{def:scheme_modules_tensorobj}
  \leanok
  The presheaf-of-modules monoidal category is symmetric, so it carries a
  braiding \(\beta\) at the presheaf level
  (\texttt{Mathlib.CategoryTheory.Monoidal.PresheafOfModules}). Its
  sheafification \(\mathtt{sheafification.mapIso}(\beta)\) is the asserted
  isomorphism, again by the cheap \(\mathtt{mapIso}\) pattern.
\end{proof}

\begin{lemma}




\leanok
[Inverse of an invertible module]
  \label{lem:tensorobj_inverse_invertible}
  \lean{AlgebraicGeometry.Scheme.Modules.exists_tensorObj_inverse}
  \uses{def:scheme_modules_tensorobj, lem:tensorobj_preserves_locally_trivial, lem:tensorobj_restrict_iso}
  % NOTE: ON the critical path under the PRIMARY locally-trivial group law (it
  % supplies inverses to lem:tensorobj_isoclass_commgroup, in whose uses it sits).
  % Only under the route-(e) fallback is it dispensable: there the group law uses
  % the tensor-invertibility predicate IsInvertible (def:scheme_modules_isinvertible),
  % which carries the inverse existentially, so no separate ``a tensor inverse
  % exists'' lemma is required.
  % SOURCE: [Stacks Project], Tag 01CR (Modules on Ringed Spaces, \S Invertible
  % modules), Lemma~lemma-constructions-invertible (item~2: the dual of an
  % invertible is invertible and the contraction map is an isomorphism) +
  % Definition~definition-pic (the Picard group as the abelian group of
  % iso-classes of invertibles)
  % (read from references/stacks-modules.tex, L4200--L4213 and L4350--L4357).
  % SOURCE QUOTE:
  % "\begin{lemma}
  % \label{lemma-constructions-invertible}
  % Let $(X, \mathcal{O}_X)$ be a ringed space.
  % ...
  % \item If $\mathcal{L}$ is an invertible $\mathcal{O}_X$-module, then so is
  % $\SheafHom_{\mathcal{O}_X}(\mathcal{L}, \mathcal{O}_X)$ and the evaluation map
  % $\mathcal{L} \otimes_{\mathcal{O}_X}
  % \SheafHom_{\mathcal{O}_X}(\mathcal{L}, \mathcal{O}_X) \to \mathcal{O}_X$
  % is an isomorphism.
  % ...
  % \begin{definition}
  % \label{definition-pic}
  % Let $(X, \mathcal{O}_X)$ be a ringed space.
  % The {\it Picard group} $\Pic(X)$ of $X$ is the
  % abelian group whose elements are isomorphism classes of
  % invertible $\mathcal{O}_X$-modules, with addition
  % corresponding to tensor product.
  % \end{definition}"
  \textit{Source: [Stacks Project], Tag~01CR, Lemma~\texttt{lemma-constructions-%
  invertible} (item~2, the dual of an invertible is invertible and the contraction
  is an isomorphism) and Definition~\texttt{definition-pic} (the Picard group is
  the abelian group of iso-classes of invertible \(\mathcal{O}_X\)-modules under
  tensor product).}
  Let \(X\) be a scheme and let \(L \in \texttt{LineBundle}\,X\) be a line
  bundle. The dual
  \[
    L^{-1} \;\;:=\;\; \mathcal{H}om_{\mathcal{O}_X}(L, \mathcal{O}_X)
  \]
  is again a line bundle on \(X\), and there is a canonical natural
  isomorphism
  \(\varepsilon_L : L \otimes_X L^{-1} \xrightarrow{\sim} \mathcal{O}_X\).
  Consequently every line bundle has a two-sided tensor inverse under
  \(\otimes_X\), and the monoid \(\bigl(\texttt{LineBundle}\,X / \sim, \otimes_X,
  \mathcal{O}_X\bigr)\) of isomorphism classes is an abelian group ---
  the Picard group \(\Pic(X)\) (Stacks tag 01CR).
\end{lemma}

\begin{proof}
  \uses{lem:tensorobj_preserves_locally_trivial, lem:tensorobj_restrict_iso, lem:tensorobj_unit_iso, lem:internal_hom_eval, lem:internal_hom_isSheaf, lem:dual_isLocallyTrivial, lem:dual_restrict_iso, lem:sheafofmodules_hom_of_local_compat, lem:hom_of_local_compat_restrict, lem:isiso_of_isiso_restrict, lem:trivialisation_restrict_compat, lem:tensorobj_unit_self_duality_collapse}
  The dual object and its evaluation are now supplied by the internal-hom
  infrastructure of \cref{sec:tensorobj_dual_infra}. Take
  \(L^{-1} := \mathcal{H}om_{\mathcal{O}_X}(L, \mathcal{O}_X)\) to be the
  sheaf-level dual \cref{lem:internal_hom_isSheaf} --- a genuine object of
  \(\Scheme.\mathtt{Modules}\,X\), assembled from the presheaf internal hom
  \cref{def:presheaf_internal_hom} by sheafification --- and let
  \(\varepsilon_L : L \otimes_X L^{-1} \to \mathcal{O}_X\) be the evaluation
  (contraction) morphism \cref{lem:internal_hom_eval}. With both in hand the
  argument is three steps.

  \emph{Step 1 (\(L^{-1}\) is a line bundle).} Local triviality of \(L\)
  (\(\mathtt{LineBundle.IsLocallyTrivial}\,L\)) furnishes an open cover \(\{U\}\) of
  \(X\) with trivialisations \(L|_U \cong \mathcal{O}_U\). The internal hom commutes
  with restriction along the open immersion \(U \hookrightarrow X\), so on each \(U\)
  \[
    L^{-1}\big|_U \;\cong\; \mathcal{H}om_{\mathcal{O}_U}(L|_U, \mathcal{O}_U)
                 \;\cong\; \mathcal{H}om_{\mathcal{O}_U}(\mathcal{O}_U, \mathcal{O}_U)
                 \;\cong\; \mathcal{O}_U ,
  \]
  the last identification being the evaluation-at-\(1\) isomorphism for the dual of a
  free rank-one module. Hence \(L^{-1}\) is locally trivial of rank one, i.e.\ an
  object of \(\texttt{LineBundle}\,X\); this is exactly
  \cref{lem:dual_isLocallyTrivial}.

  \emph{Step 2 (the contraction is a local isomorphism).} The evaluation/contraction
  morphism
  \[
    \varepsilon_L \;\colon\; L \otimes_X L^{-1}
      \;=\; L \otimes_X \mathcal{H}om_{\mathcal{O}_X}(L, \mathcal{O}_X)
      \;\longrightarrow\; \mathcal{O}_X,
    \qquad s \otimes \phi \longmapsto \phi(s),
  \]
  is a morphism in \(\Scheme.\mathtt{Modules}\,X\). Restricting to a trivialising
  open \(U\), the restriction-compatibility isomorphism \cref{lem:tensorobj_restrict_iso}
  commutes \(\otimes_X\) past \((-)|_U\):
  \[
    (L \otimes_X L^{-1})\big|_U \;\cong\; L|_U \otimes_U L^{-1}|_U
      \;\cong\; \mathcal{O}_U \otimes_U \mathcal{O}_U ,
  \]
  using the trivialisations of Step~1. Under these identifications \(\varepsilon_L|_U\)
  becomes the contraction \(\mathcal{O}_U \otimes_U \mathcal{O}_U \to \mathcal{O}_U\) of
  the free rank-one module with its dual, which is precisely the left unitor
  \cref{lem:tensorobj_unit_iso} (\(\mathcal{O}_U \otimes_U (-) \cong (-)\)) and hence an
  isomorphism. Thus \(\varepsilon_L|_U\) is an isomorphism on each member of the cover.

  \emph{Step 3 (glue to a global isomorphism).} The local trivialisations are induced
  by the open immersions with no choices, so the local contraction isomorphisms
  \(\varepsilon_L|_U\) agree on the overlaps \(U \cap U'\) by the naturality of
  \cref{lem:tensorobj_restrict_iso} and of the unitor under further restriction. As
  ``being an isomorphism'' is local on \(X\) --- a morphism of sheaves of
  \(\mathcal{O}_X\)-modules whose restriction to every member of an open cover is an
  isomorphism is itself an isomorphism --- the globally-defined \(\varepsilon_L\) is an
  isomorphism \(L \otimes_X L^{-1} \xrightarrow{\sim} \mathcal{O}_X\). Equivalently, the
  local inverses glue (they agree on overlaps for the same reason) to a two-sided
  inverse of \(\varepsilon_L\).

  Consequently every line bundle \(L\) has the two-sided tensor inverse \(L^{-1}\)
  under \(\otimes_X\), and the monoid of \(\otimes\)-iso-classes of line bundles is
  an abelian group --- the Picard group \(\Pic(X)\) (Stacks tag 01CR). The
  local-triviality of \(L^{-1}\) (Step~1) is \cref{lem:dual_isLocallyTrivial}, the
  dual object and its evaluation are \cref{lem:internal_hom_isSheaf} and
  \cref{lem:internal_hom_eval}, and the local-iso\(\Rightarrow\)global-iso passage
  (Step~3) is the already-closed \cref{lem:tensorobj_restrict_iso}.

  \emph{Formal dependency.} The proof of
  \(\mathtt{exists\_tensorObj\_inverse}\) is \emph{not} a one-line consequence: it is
  the gluing argument \cref{rem:dual_discharges_inverse}, which assembles the global
  contraction \(\varepsilon_L\) from the chart-local left-unitor contractions. Its load-
  bearing inputs are: the gluing of compatible local morphisms
  \cref{lem:sheafofmodules_hom_of_local_compat} together with its restriction-value
  connector \cref{lem:hom_of_local_compat_restrict}; the local-iso\(\Rightarrow\)global-iso
  passage \cref{lem:isiso_of_isiso_restrict}; the overlap cocycle equality, whose two
  hardest constituents are the trivialisation restriction-compatibility chain
  \cref{lem:trivialisation_restrict_compat} and the unit self-duality collapse
  \cref{lem:tensorobj_unit_self_duality_collapse}. The dual restriction commutation
  \cref{lem:dual_restrict_iso} is one input among several. The principal open obligation is
  the restriction-naturality of the trivialisation chain
  (\cref{lem:trivialisation_restrict_compat}), itself a chain of per-constituent naturality
  squares.
\end{proof}

\begin{lemma}




\leanok
[Closure of \(\otimes\) on the locally-trivial carrier
  \(\texttt{LineBundle.OnProduct}\)]
  \label{lem:tensorobj_lift_onproduct}
  \lean{AlgebraicGeometry.Scheme.Modules.tensorObjOnProduct}
  \uses{def:scheme_modules_tensorobj, lem:tensorobj_preserves_locally_trivial}
  Let \(C \to S\) and \(T \to S\) be \(S\)-schemes. The carrier is the
  \emph{locally-trivial} subtype
  \[
    \texttt{LineBundle.OnProduct}\,\pi_C\,\pi_T
    \;\;=\;\;
    \bigl\{\,M \in \Scheme.\mathtt{Modules}(C \times_S T)
            \,\bigm|\, \mathtt{IsLocallyTrivial}\,M \,\bigr\}
  \]
  introduced in \cref{chap:Picard_LineBundlePullback}. The bifunctor
  \(\otimes_{C \times_S T}\) of
  \cref{lem:scheme_modules_tensorobj_functoriality} restricts to this subtype:
  if \(\mathtt{IsLocallyTrivial}\,M\) and \(\mathtt{IsLocallyTrivial}\,M'\) then
  \(\mathtt{IsLocallyTrivial}(M \otimes_{C \times_S T} M')\). Hence \(\otimes\)
  closes on \(\texttt{LineBundle.OnProduct}\,\pi_C\,\pi_T\), supplying the binary
  operation the relative group law multiplies with.
\end{lemma}

\begin{proof}
  \uses{lem:tensorobj_preserves_locally_trivial}
  \leanok
  Closure is exactly
  \cref{lem:tensorobj_preserves_locally_trivial}: the tensor product of two
  locally-trivial rank-one modules is again locally trivial of rank one
  (\(\mathtt{tensorObj\_isLocallyTrivial}\)). The Lean construction
  \(\mathtt{tensorObjOnProduct}\) applies this directly to a pair of objects of
  the \(\mathtt{IsLocallyTrivial}\) subtype \(\texttt{LineBundle.OnProduct}\),
  returning the tensor with its locally-trivial witness; no
  \(\otimes\)-invertibility predicate and no coherence isomorphism enters this
  closure statement. (The \(\otimes\)-invertibility carrier
  \(\mathtt{IsInvertible}\) of \cref{def:scheme_modules_isinvertible} is a
  separate predicate, connected to \(\mathtt{IsLocallyTrivial}\) only by an
  off-critical-path equivalence; the units-group packaging is
  \cref{lem:tensorobj_isoclass_commgroup}, not this lemma.)
\end{proof}

\begin{lemma}
[\(\pi_T^*\) is a tensor-functor]
  \label{lem:pullback_compatible_with_tensorobj}
  \lean{AlgebraicGeometry.LineBundle.OnProduct.pullback_tensorObj_iso}
  % NOTE: \lean{} pin target absent from Lean source as of iter-271 (verified by grep).
  % Forward reference to a not-yet-formalized declaration — NOT a stale rename; there is
  % no existing Lean decl under any name to repin to. The `unmatched_lean` status is the
  % normal "unformalized target" state, not a pin defect.
  \uses{lem:tensorobj_lift_onproduct}
  % NOTE: pullback preserves the structure sheaf and tensor products (and hence
  % tensor-invertibility) by the standard pullback-tensor compatibility of
  % \(\mathcal{O}\)-modules (Stacks \texttt{lemma-tensor-product-pullback} /
  % \texttt{lemma-pullback-invertible}), independently of the open-immersion
  % restriction-compatibility iso.
  With notation as in \cref{lem:tensorobj_lift_onproduct}, the pullback
  functor
  \[
    \pi_T^{*} \;\;\colon\;\; \texttt{LineBundle}\,T \;\longrightarrow\;
                          \texttt{LineBundle.OnProduct}\,\pi_C\,\pi_T
  \]
  of \cref{def:pullback_along_projection} is a (strong) monoidal
  functor: there is a canonical natural isomorphism
  \(\pi_T^*(\mathcal{N} \otimes_T \mathcal{N'})
     \xrightarrow{\sim}
   \pi_T^* \mathcal{N} \otimes_{C \times_S T} \pi_T^* \mathcal{N'}\)
  for \(\mathcal{N}, \mathcal{N'} \in \texttt{LineBundle}\,T\), and a
  canonical isomorphism
  \(\pi_T^* \mathcal{O}_T \xrightarrow{\sim} \mathcal{O}_{C \times_S T}\),
  both inherited from the analogous data on the underlying
  \(\Scheme.\mathtt{Modules}\) functors via the
  pullback-tensor-compatibility of \(\mathcal{O}\)-modules.
\end{lemma}

\begin{proof}
  \uses{lem:tensorobj_lift_onproduct}
  On any affine open of \(T\) over which \(\mathcal{N}\) and
  \(\mathcal{N'}\) trivialise simultaneously, the assertion reduces to
  the standard fact that the extension of scalars
  \(B \otimes_A (M \otimes_A M')
   \;\cong\; (B \otimes_A M) \otimes_B (B \otimes_A M')\)
  for a ring map \(A \to B\) and \(A\)-modules \(M, M'\). This affine
  identity glues across an affine trivialising cover of \(T\): on each
  affine open where \(\mathcal{N}\) and \(\mathcal{N}'\) are free of rank one
  over \(A\), both sides evaluate to the free rank-one \(B\)-module, so
  the comparison is an isomorphism. No flatness of \(A \to B\) is required
  at this level, as the modules involved are locally free. The compatibility
  with the unit, the associator, and the unitors is verified pointwise on the
  same cover.
\end{proof}

\section{The varying-ring stalk--tensor commutation (d.2) and the unconditional
left-whiskering}
\label{sec:tensorobj_stalk_tensor}

This section supplies the one genuinely Mathlib-absent ingredient on which the
\emph{unconditional} associator \cref{lem:tensorobj_assoc_iso} rests: the commutation
of the stalk functor with the relative-ring tensor product of presheaves of modules
(``d.2''). Together with the already-built \(R_x\)-linear stalk comparison
(\cref{lem:stalk_linear_map}, ``d.1'') it closes the single open left-whiskering
obligation for an \emph{arbitrary} whiskering object --- with no flatness and no
local-triviality hypothesis --- which is exactly what makes the sheafification
transport of the presheaf associator work for all modules. The construction is to be
formalized in the dedicated Lean file
\(\mathtt{AlgebraicJacobian/Picard/TensorObjSubstrate/StalkTensor.lean}\).

\begin{lemma}

\leanok
[Stalk of a tensor product (d.2): varying-ring stalk--tensor commutation]
  \label{lem:stalk_tensor_commutation}
  \lean{PresheafOfModules.stalkTensorIso}
  \uses{def:scheme_modules_tensorobj, lem:stalk_linear_map}
  % SOURCE: [Stacks Project], Tag (Modules on Ringed Spaces, \S Tensor product),
  % Lemma lemma-stalk-tensor-product (stalk of a tensor product of modules is the
  % tensor product of the stalks over the stalk ring)
  % (read from references/stacks-modules.tex, L2332-L2344).
  % SOURCE QUOTE:
  % "\begin{lemma}
  % \label{lemma-stalk-tensor-product}
  % Let $(X, \mathcal{O}_X)$ be a ringed space.
  % Let $\mathcal{F}$, $\mathcal{G}$ be $\mathcal{O}_X$-modules.
  % Let $x \in X$. There is a canonical isomorphism
  % of $\mathcal{O}_{X, x}$-modules
  % $$
  % (\mathcal{F} \otimes_{\mathcal{O}_X} \mathcal{G})_x
  % =
  % \mathcal{F}_x \otimes_{\mathcal{O}_{X, x}} \mathcal{G}_x
  % $$
  % functorial in $\mathcal{F}$ and $\mathcal{G}$.
  % \end{lemma}"
  \textit{Source: [Stacks Project], Modules on Ringed Spaces,
  Lemma~\texttt{lemma-stalk-tensor-product}.}
  Let \(R : (\mathtt{Opens}\,X)^{op} \to \mathtt{CommRingCat}\) be the structure
  presheaf of a scheme \(X\), and let \(A, B\) be presheaves of
  \(R\)-modules (objects of \(\mathtt{PresheafOfModules}\,(R \circ
  \mathtt{forget_2}\,\mathtt{CommRingCat}\,\mathtt{RingCat})\)). For each point
  \(x \in X\) there is a canonical isomorphism of \((R \circ
  \mathtt{forget_2}).\mathtt{stalk}\,x\)-modules
  \[
    (A \otimes^{p} B).\mathtt{stalk}\,x
      \;\xrightarrow{\ \sim\ }\;
    A.\mathtt{stalk}\,x \;\otimes_{R.\mathtt{stalk}\,x}\; B.\mathtt{stalk}\,x,
  \]
  natural in \(A\) and \(B\), where \(\otimes^{p}\) is the presheaf-of-modules tensor
  product \(\mathtt{PresheafOfModules.Monoidal.tensorObj}\).
\end{lemma}

\begin{proof}
  \uses{def:scheme_modules_tensorobj, lem:stalk_linear_map}
  \leanok
  The stalk at \(x\) is the filtered colimit over the neighbourhoods \(U \ni x\) of
  the sections functor, and sectionwise the tensor product is computed pointwise:
  \((A \otimes^{p} B)(U) = A(U) \otimes_{R(U)} B(U)\)
  (\(\mathtt{PresheafOfModules.Monoidal.tensorObj\_obj}\)). The neighbourhood system
  \(\{U \ni x\}\) is cofiltered under inclusion, so the colimit is filtered. The
  isomorphism is assembled in five named stages. Stages (i)--(ii) construct the
  forward additive comparison map and its germ characterisation; stages (iii)--(v)
  upgrade it to an \(R_x\)-linear map, construct the reverse map, and check mutual
  inversion.

  \emph{(i) Per-neighbourhood balanced bilinear map and its section-level descent}
  (\(\mathtt{stalkTensorBilin}\), \(\mathtt{stalkTensorDescU}\)). For a fixed
  neighbourhood \(U \ni x\) the assignment
  \[
    (a, b) \;\longmapsto\; \mathrm{germ}_x\,a \,\otimes\, \mathrm{germ}_x\,b,
    \qquad a \in A(U),\ b \in B(U),
  \]
  is \(R(U)\)-balanced bilinear into \(A_x \otimes_{R_x} B_x\): it is additive in each
  variable, and for \(r \in R(U)\) the two ways of moving \(r\) across the tensor agree
  after passing to germs, since \(\mathrm{germ}_x(r \cdot a) = (\mathrm{germ}_x r)
  \cdot \mathrm{germ}_x a\) and the stalk tensor is balanced over
  \(R_x = R.\mathtt{stalk}\,x\) (the germ--scalar compatibility of
  \cref{lem:stalk_linear_map}). By the universal property of the relative tensor
  product it descends to an additive map on the section module
  \[
    \mathtt{stalkTensorDescU}_U \;\colon\;
      (A \otimes^{p} B)(U) = A(U) \otimes_{R(U)} B(U)
        \;\longrightarrow\; A_x \otimes_{R_x} B_x,
    \qquad a \otimes b \longmapsto \mathrm{germ}_x a \otimes \mathrm{germ}_x b.
  \]

  \emph{(ii) Colimit descent to the forward comparison map}
  (\(\mathtt{stalkTensorDesc}\), \(\mathtt{stalkTensorDesc\_germ\_tmul}\)). The
  maps \(\mathtt{stalkTensorDescU}_U\) are compatible with the restriction maps of the
  presheaf \(A \otimes^{p} B\) as \(U\) shrinks toward \(x\) (a smaller neighbourhood
  computes the same germs), so they form a cocone over the filtered diagram
  \(\{(A \otimes^{p} B)(U)\}_{U \ni x}\). The induced map out of the colimit is the
  forward additive comparison map
  \[
    \mathtt{stalkTensorDesc} \;\colon\;
      (A \otimes^{p} B).\mathtt{stalk}\,x \;\longrightarrow\; A_x \otimes_{R_x} B_x,
  \]
  characterised on germs of simple tensors by
  \(\mathtt{stalkTensorDesc}\bigl(\mathrm{germ}_x(a \otimes b)\bigr) =
  \mathrm{germ}_x a \otimes \mathrm{germ}_x b\)
  (\(\mathtt{stalkTensorDesc\_germ\_tmul}\)). This is the forward half of the asserted
  isomorphism (\texttt{stalk\_tensor\_desc\_forward}).

  \emph{(iii) \(R_x\)-linearity packaging} (\(\mathtt{stalkTensorLinearMap}\)). The
  additive map of (ii) is upgraded to an \(R.\mathtt{stalk}\,x\)-linear map: on a germ
  \(\mathrm{germ}_x r\) of a section \(r \in R(U)\) one checks
  \(\mathtt{stalkTensorDesc}\bigl((\mathrm{germ}_x r)\cdot \xi\bigr) =
  (\mathrm{germ}_x r)\cdot \mathtt{stalkTensorDesc}(\xi)\) on germs of simple tensors
  and extends by additivity. The one subtlety is a \emph{carrier-duality obstacle}:
  the section-level tensor \(A(U) \otimes_{R(U)} B(U)\) is a module over the
  \(\mathtt{RingCat}\) carrier \((R \circ \mathtt{forget_2})(U)\), whereas the natural
  scalar \(r : R(U)\) lives in the \(\mathtt{CommRingCat}\) carrier \(R(U)\); the two
  carriers are canonically identified but not definitionally equal, so the
  scalar-pulling lemma \(\mathtt{TensorProduct.smul\_tmul'}\) applies only after
  inserting a canonical \(\mathtt{RingEquiv}\) bridge between the
  \(\mathtt{CommRingCat}\) and \(\mathtt{RingCat}\) carriers. Once that bridge is
  in place the linearity computation mirrors the germ-representative pattern used
  for the d.1 stalk map \(\mathtt{stalkLinearMap}\) (\cref{lem:stalk_linear_map}).

  \emph{(iv) The reverse map.} The map
  \[
    A_x \otimes_{R_x} B_x \;\longrightarrow\; (A \otimes^{p} B).\mathtt{stalk}\,x
  \]
  is built by the universal property of the tensor product over \(R_x\) from the
  \(R_x\)-bilinear map of the two stalks
  \[
    (\mathrm{germ}_x a,\ \mathrm{germ}_x b) \;\longmapsto\;
      \mathrm{germ}_x\bigl(a|_{U \cap V} \otimes b|_{U \cap V}\bigr),
  \]
  where \(a \in A(U)\) and \(b \in B(V)\) are chosen representatives; this is a
  \emph{nested} colimit descent, descending first the \(A\)-variable out of the
  filtered colimit \(\{A(U)\}_{U \ni x}\) and then the \(B\)-variable out of
  \(\{B(V)\}_{V \ni x}\), passing to the common refinement \(U \cap V\) to form the
  simple tensor and its germ. Well-definedness on germs (independence of the chosen
  representatives and of the neighbourhoods \(U, V\)) is exactly filteredness of the
  neighbourhood system.

  It remains to check that the reverse map is \(R_x\)-bilinear, and in particular that
  its underlying bilinear map is \emph{balanced} over \(R_x\), i.e. that pulling a scalar
  out of the left variable agrees with pulling it out of the right. The key point is to
  establish this balancing \emph{at the stalk level}, with every scalar living in the
  stalk ring \(R_x = R.\mathtt{stalk}\,x\) throughout. Concretely, for a germ
  \(\mathrm{germ}_x r\) of a section \(r \in R(W)\) one verifies the identity
  \[
    \mathtt{revBihom}\bigl((\mathrm{germ}_x r)\cdot \mathrm{germ}_x a,\ \mathrm{germ}_x b\bigr)
      \;=\; (\mathrm{germ}_x r)\cdot \mathtt{revBihom}\bigl(\mathrm{germ}_x a,\ \mathrm{germ}_x b\bigr),
  \]
  and symmetrically on the right variable, where the scalar \(\mathrm{germ}_x r \in R_x\)
  never leaves the stalk ring. This holds because of the germ--scalar compatibility
  \(\mathrm{germ}_x(r \cdot s) = (\mathrm{germ}_x r)\cdot(\mathrm{germ}_x s)\) of the germ
  map \(R(W) \to R_x\) (the Lean lemma \(\mathtt{germ\_smul}\)) — the same germ--scalar
  compatibility already used in stages (i) and (iii). Multiplying a representative
  \(r \cdot a\) and then taking germs produces \((\mathrm{germ}_x r)\cdot \mathrm{germ}_x a\),
  so the action of \(r\) is transparently the \(R_x\)-action on the stalk and may be
  moved freely across the tensor of stalks, where balancedness over \(R_x\) is built into
  \(A_x \otimes_{R_x} B_x\). Because the whole computation takes place over the single
  ring \(R_x\), no carrier identification intervenes and the balancing follows directly.

  One must \emph{not} instead reduce this balancing to a section-level identity. If the
  scalar were pushed through the presheaf restriction to a common neighbourhood \(W\) and
  the balancing read off inside the section module \(A(W) \otimes_{R(W)} B(W)\), then the
  scalar action appearing there is the one produced by the presheaf-of-modules structure
  of \(A\) via restriction of scalars: it lives over the \(\mathtt{RingCat}\) carrier
  \((R \circ \mathtt{forget_2})(W)\), \emph{not} over the \(\mathtt{CommRingCat}\) carrier
  \(R(W)\) that annotates the section tensor. This is precisely the
  \(\mathtt{CommRingCat}\)-vs-\(\mathtt{RingCat}\) carrier-duality obstacle already met in
  stage (iii), here compounded by the extra restriction-of-scalars wrapper, so the naive
  ``move the scalar across the tensor'' step (\(\mathtt{TensorProduct.smul\_tmul'}\)) does
  not apply at the section level without first inserting the carrier bridge. Working at the
  stalk level as above sidesteps this wrapper entirely.

  Where a \(\mathtt{CommRingCat}\)-vs-\(\mathtt{RingCat}\) carrier identification is
  genuinely unavoidable, it is the \emph{same} canonical bridge already invoked for the
  stage-(iii) \(R_x\)-linearity packaging (the \(\mathtt{stalkTensorDescU\_smul}\)/%
  \(\mathtt{stalkTensorLinearMap}\) step, \texttt{stalk\_tensor\_linear\_map}): the two
  carriers are canonically identified but not definitionally equal, and the same
  \(\mathtt{RingEquiv}\) bridge resolves the obstruction here. Thus the stage-(iv)
  balancing reuses the technique already resolved in stage (iii) rather than re-deriving
  it.

  \emph{(v) Mutual inversion on germ generators} (\(\mathtt{stalkTensorIso}\)). The two
  composites are checked to be identities on generators. The composite
  \(A_x \otimes_{R_x} B_x \to (A \otimes^{p} B).\mathtt{stalk}\,x \to A_x
  \otimes_{R_x} B_x\) is the identity on a simple tensor \(\mathrm{germ}_x a \otimes
  \mathrm{germ}_x b\) by the germ characterisation \(\mathtt{stalkTensorDesc\_germ\_tmul}\)
  of (ii), and extends to all of \(A_x \otimes_{R_x} B_x\) by
  \(\mathtt{TensorProduct.induction\_on}\). The opposite composite
  \((A \otimes^{p} B).\mathtt{stalk}\,x \to A_x \otimes_{R_x} B_x \to (A \otimes^{p}
  B).\mathtt{stalk}\,x\) is the identity on a germ \(\mathrm{germ}_x(a \otimes b)\) of
  a simple tensor by the same characterisation; since germs generate the stalk and the
  germ maps are jointly epimorphic (\(\mathtt{stalk\_hom\_ext}\)), together with
  \(\mathtt{TensorProduct.induction\_on}\) applied to each section, the composite is the
  identity. Bundling the \(R_x\)-linear forward map of (iii) with the reverse map of
  (iv) and these two inversion identities yields the asserted isomorphism
  \(\mathtt{stalkTensorIso}\).

  Naturality in \(A, B\) is inherited from naturality of the sectionwise tensor and of
  the germ maps. This is the presheaf-of-modules, varying-ring analogue of the Stacks
  lemma above; the only subtlety beyond the ringed-space statement is that the ground
  ring varies with \(U\), handled by the ``tensor of filtered colimits over the colimit
  ring'' identification together with the carrier-duality bridge of (iii). Conceptually
  the tensor product of modules commutes with filtered colimits in each variable and
  the ground ring \(R(U)\) itself colimits to the stalk ring \(R.\mathtt{stalk}\,x\),
  which is the reason the forward map (ii) and the reverse map (iv) are mutually
  inverse.
\end{proof}

The proof of the live whiskering lemma below factors through two named
sub-lemmas: the \emph{d.1-bridge} relating membership in the sheafification
localizer \(J.W\) to stalkwise isomorphism, and the \emph{second-factor
naturality} of the d.2 commutation isomorphism \cref{lem:stalk_tensor_commutation}.
We state both before the lemma that consumes them.

\begin{lemma}

\leanok
[Unconditional left-whiskering of the localizer, via d.2]
  \label{lem:islocallyinjective_whiskerleft_via_stalk}
  \lean{PresheafOfModules.isLocallyInjective_whiskerLeft_of_W}
  \uses{def:scheme_modules_tensorobj, lem:stalk_tensor_commutation, lem:stalk_linear_map}
  Let \(R\), \(J\) be as above with \(J.\mathtt{WEqualsLocallyBijective}\,\mathtt{Ab}\),
  let \(F\) be an \emph{arbitrary} presheaf of \(R\)-modules, and let \(g : M \to N\)
  be a morphism whose underlying morphism lies in \(J.W\). Then the left-whiskered
  morphism \(F \triangleleft g = \mathrm{id}_F \otimes g\) is \(J\)-locally injective
  (indeed lies in \(J.W\)). No flatness and no local-triviality hypothesis on \(F\) is
  required. This discharges the open left-whiskering obligation
  \cref{lem:islocallyinjective_whisker_of_W} that the unconditional associator
  \cref{lem:tensorobj_assoc_iso} consumes.
\end{lemma}

% SOURCE QUOTE PROOF: (Archon-original assembly from the Stacks stalk-tensor lemma
% above and the stalkwise criterion for membership in the sheafification localizer;
% no single external source proof.)
\begin{proof}
  \leanok
  \uses{lem:stalk_tensor_commutation, lem:stalk_linear_map}
  The proof runs in three movements: the d.1-bridge turns membership in \(J.W\) into a
  stalkwise statement, the second-factor naturality of the d.2 isomorphism transports
  the whiskered stalk into a tensor of the stalk maps, and a flatness-free
  ``tensoring by an equivalence'' argument concludes.

  \emph{(a) The d.1-bridge.} Specialise to the topological site \(\mathtt{Opens}\,X\),
  where stalks exist. By the hypothesis \(g \in J.W\) and the d.1-bridge
  \texttt{W\_implies\_stalkwise\_iso}, the \(\mathtt{Ab}\)-stalk map \(g_x : M_x \to N_x\)
  is an isomorphism at every point \(x \in X\). By \cref{lem:stalk_linear_map} this
  stalk map carries an \(R_x\)-linear structure, and the bijectivity just obtained
  upgrades it to an \(R_x\)-linear equivalence
  \(\mathtt{stalkLinearEquivOfIsIso} : M_x \xrightarrow{\sim} N_x\).

  \emph{(b) Identifying the whiskered stalk.} By the second-factor naturality of the
  stalk--tensor comparison \texttt{stalk\_tensor\_commutation\_naturality\_right}, the
  isomorphism \(\mathtt{stalkTensorIso}\) of \cref{lem:stalk_tensor_commutation}
  identifies the stalk \((F \triangleleft g)_x\) of the whiskered morphism with the
  map
  \[
    \mathrm{id}_{F_x} \otimes g_x \;=\; \mathtt{LinearMap.lTensor}\,F_x\,g_x
      \;\colon\; F_x \otimes_{R_x} M_x \longrightarrow F_x \otimes_{R_x} N_x .
  \]

  \emph{(c) Conclusion.} Tensoring the \(R_x\)-linear equivalence
  \(g_x : M_x \xrightarrow{\sim} N_x\) of (a) on the left by the identity of \(F_x\)
  yields an \(R_x\)-linear equivalence
  \(\mathtt{LinearMap.lTensor}\,F_x\,g_x = \mathrm{id}_{F_x} \otimes g_x\)
  (this is \(\mathtt{LinearEquiv.lTensor}\); it needs \emph{no} flatness of \(F_x\),
  because lifting an isomorphism through \(F_x \otimes_{R_x} (-)\) only uses the
  functoriality of the tensor product, and the inverse is supplied by the inverse
  equivalence). By (b), \((F \triangleleft g)_x\) is therefore an isomorphism of
  \(\mathtt{Ab}\)-stalks for every \(x\). Applying the converse half of the d.1-bridge
  \texttt{W\_implies\_stalkwise\_iso} to the morphism \(F \triangleleft g\) gives
  \(F \triangleleft g \in J.W\), and in particular it is \(J\)-locally injective
  (\(\mathtt{GrothendieckTopology.W.isLocallyInjective}\)). No flatness, no local
  triviality, and the whiskering object \(F\) is arbitrary.

  The statement holds over any Grothendieck site \(J\) satisfying
  \(J.\mathtt{WEqualsLocallyBijective}\,\mathtt{Ab}\), and is applied in particular
  at the topological site \(\mathtt{Opens}\,X\) of a scheme, where the stalk
  arguments of (a)--(c) are available and through which the unconditional associator
  \cref{lem:tensorobj_assoc_iso} is realized.
\end{proof}

\subsection{Stalk--tensor comparison helper declarations}
\label{subsec:tensorobj_stalktensor_helpers}

The following are the internal Lean helper declarations assembling the forward
comparison map \texttt{stalk\_tensor\_desc\_forward}, its \(R_x\)-linear packaging
\texttt{stalk\_tensor\_linear\_map}, the reverse map, and the mutual-inversion germ
identities feeding the stalk--tensor isomorphism \cref{lem:stalk_tensor_commutation}.
Each is proved directly in Lean and carries the dependency edges of the construction.

\section{\(\otimes\)-invertibility and the commutative-monoid group-law engine}
\label{sec:tensorobj_invertibility}

The relative Picard group law is assembled from \emph{\(\otimes\)-invertibility}
and the monoidal coherence isomorphisms of
\cref{lem:tensorobj_assoc_iso,lem:tensorobj_unit_iso,lem:tensorobj_comm_iso},
following Mathlib's Picard-group idiom (\(\mathtt{Module.Invertible}\),
\(\mathtt{CommRing.Pic} = \mathrm{Units}(\mathrm{Skeleton}(\mathtt{ModuleCat}\,R))\),
\(\mathtt{instCommGroupPic}\); \texttt{Mathlib.RingTheory.PicardGroup}) as the
design template. Under the route~(e) fallback the inverse of an invertible object
is carried by the predicate itself, so neither the open-immersion
restriction-compatibility isomorphism \cref{lem:tensorobj_restrict_iso} nor a
separate tensor-inverse lemma is needed there. Under the PRIMARY locally-trivial
route, by contrast, both \cref{lem:tensorobj_restrict_iso} (via the PRIMARY proof
of \cref{lem:islocallyinjective_whisker_of_W}) and the tensor-inverse lemma
\cref{lem:tensorobj_inverse_invertible} (supplying inverses to
\cref{lem:tensorobj_isoclass_commgroup}) are on the critical path.

\begin{definition}\leanok
[\(\otimes\)-invertibility of an \(\mathcal{O}_X\)-module]
  \label{def:scheme_modules_isinvertible}
  \lean{AlgebraicGeometry.Scheme.Modules.IsInvertible}
  \uses{def:scheme_modules_tensorobj}
  % SOURCE: [Stacks Project], Tag 01CR (Modules on Ringed Spaces, \S Invertible
  % modules), Definition 01CS (definition-invertible) + Lemma 0B8K
  % (lemma-invertible, the "$\exists\,\mathcal{N}$" characterisation)
  % (read from references/stacks-modules.tex, L4046-L4079).
  % SOURCE QUOTE:
  % "\begin{definition}
  % \label{definition-invertible}
  % Let $(X, \mathcal{O}_X)$ be a ringed space. An
  % {\it invertible $\mathcal{O}_X$-module} is a sheaf
  % of $\mathcal{O}_X$-modules $\mathcal{L}$ such that
  % the functor
  % $$
  % \textit{Mod}(\mathcal{O}_X) \longrightarrow \textit{Mod}(\mathcal{O}_X),\quad
  % \mathcal{F} \longmapsto \mathcal{L} \otimes_{\mathcal{O}_X} \mathcal{F}
  % $$
  % is an equivalence of categories. We say that $\mathcal{L}$ is
  % {\it trivial} if it is isomorphic as an $\mathcal{O}_X$-module
  % to $\mathcal{O}_X$.
  % \end{definition}
  % [\dots]
  % \begin{lemma}
  % \label{lemma-invertible}
  % Let $(X, \mathcal{O}_X)$ be a ringed space. Let $\mathcal{L}$
  % be an $\mathcal{O}_X$-module. Equivalent are
  % \begin{enumerate}
  % \item $\mathcal{L}$ is invertible, and
  % \item there exists an $\mathcal{O}_X$-module $\mathcal{N}$
  % such that
  % $\mathcal{L} \otimes_{\mathcal{O}_X} \mathcal{N} \cong \mathcal{O}_X$.
  % \end{enumerate}
  % In this case $\mathcal{L}$ is locally a direct summand of a finite free
  % $\mathcal{O}_X$-module and the module $\mathcal{N}$ in (2) is isomorphic to
  % $\SheafHom_{\mathcal{O}_X}(\mathcal{L}, \mathcal{O}_X)$.
  % \end{lemma}"
  \textit{Source: [Stacks Project], Tag 01CR (Modules on Ringed Spaces,
  ``Invertible modules''), Definition~01CS and Lemma~0B8K. An invertible
  \(\mathcal{O}_X\)-module is one for which \(\mathcal{L} \otimes_{\mathcal{O}_X}
  (-)\) is an equivalence; equivalently there exists \(\mathcal{N}\) with
  \(\mathcal{L} \otimes_{\mathcal{O}_X} \mathcal{N} \cong \mathcal{O}_X\), and
  then \(\mathcal{N} \cong \mathcal{H}om_{\mathcal{O}_X}(\mathcal{L},
  \mathcal{O}_X)\).}
  Let \(X\) be a scheme and \(M \in \Scheme.\mathtt{Modules}\,X\). We say \(M\)
  is \emph{\(\otimes\)-invertible} when it admits a tensor inverse: there is an
  object \(N \in \Scheme.\mathtt{Modules}\,X\) together with an isomorphism
  \[
    M \otimes_X N \;\cong\;
      \mathtt{SheafOfModules.unit}\,X.\mathtt{ringCatSheaf}
      \;=\; \mathcal{O}_X ,
  \]
  i.e.\ \(\mathtt{IsInvertible}(M) :\equiv \exists\, N,\
  \mathrm{Nonempty}\bigl(M \otimes_X N \cong
  \mathtt{SheafOfModules.unit}\,X.\mathtt{ringCatSheaf}\bigr)\), with
  \(\otimes_X\) the substrate of \cref{def:scheme_modules_tensorobj}. The
  predicate is a \emph{proposition} carrying its inverse witness existentially;
  consequently the dual/inverse needed by the group law is definitional and no
  separate ``a tensor inverse exists'' lemma is consumed downstream. This is the
  scheme-level analogue of Mathlib's \(\mathtt{Module.Invertible}\), transported
  one categorical level up to \(\Scheme.\mathtt{Modules}\,X\) with the substrate
  \(\otimes_X\).
\end{definition}

\begin{lemma}[The commutative monoid of \(\otimes\)-iso-classes and its units]
  \label{lem:tensorobj_isoclass_commgroup}
  % NOTE: the monolithic `tensorObjIsoclassCommMonoid` of the original design was
  % refactored into the carrier `PicGroup` (def:pic_carrier) and its group law
  % `picCommGroup` (thm:pic_commgroup); pin corrected to those two decls (review iter-244).
  \lean{AlgebraicGeometry.Scheme.Modules.PicGroup, AlgebraicGeometry.Scheme.Modules.picCommGroup}
  \uses{def:scheme_modules_tensorobj, def:scheme_modules_isinvertible, lem:scheme_modules_tensorobj_functoriality, lem:tensorobj_lift_onproduct, lem:tensorobj_inverse_invertible, lem:tensorobj_assoc_iso, lem:tensorobj_unit_iso, lem:tensorobj_comm_iso, def:tensorobjisoofiso, def:tensorobj_unit_self_iso, def:tensorobj_middlefour, def:pic_setoid, def:pic_mul, def:pic_inv}
  % SOURCE: [Stacks Project], Tag 01CR (Modules on Ringed Spaces, \S Invertible
  % modules), Definition 01CX (definition-pic) + Lemma lemma-constructions-
  % invertible (tensor of invertibles is invertible; dual is the inverse)
  % (read from references/stacks-modules.tex, L4200-L4213 and L4350-L4357).
  % SOURCE QUOTE:
  % "\begin{lemma}
  % \label{lemma-constructions-invertible}
  % Let $(X, \mathcal{O}_X)$ be a ringed space.
  % \begin{enumerate}
  % \item If $\mathcal{L}$, $\mathcal{N}$ are invertible
  % $\mathcal{O}_X$-modules, then so is
  % $\mathcal{L} \otimes_{\mathcal{O}_X} \mathcal{N}$.
  % \item If $\mathcal{L}$ is an invertible $\mathcal{O}_X$-module, then so is
  % $\SheafHom_{\mathcal{O}_X}(\mathcal{L}, \mathcal{O}_X)$ and the evaluation map
  % $\mathcal{L} \otimes_{\mathcal{O}_X}
  % \SheafHom_{\mathcal{O}_X}(\mathcal{L}, \mathcal{O}_X) \to \mathcal{O}_X$
  % is an isomorphism.
  % \end{enumerate}
  % \end{lemma}
  % [\dots]
  % \begin{definition}
  % \label{definition-pic}
  % Let $(X, \mathcal{O}_X)$ be a ringed space.
  % The {\it Picard group} $\Pic(X)$ of $X$ is the
  % abelian group whose elements are isomorphism classes of
  % invertible $\mathcal{O}_X$-modules, with addition
  % corresponding to tensor product.
  % \end{definition}"
  % SOURCE: [Mathlib], `CommRing.Pic` definition, `CommGroup` instance,
  % `mul_eq_tensor` and `inv_eq_dual`, in
  % `Mathlib/RingTheory/PicardGroup.lean` (L405--L412, L492, L495)
  % (read from .lake/packages/mathlib/Mathlib/RingTheory/PicardGroup.lean).
  % SOURCE QUOTE:
  % "/-- The Picard group of a commutative semiring R consists of the invertible R-modules,
  % up to isomorphism. -/
  % def CommRing.Pic (R : Type u) [CommSemiring R] : Type u :=
  %   Shrink (Skeleton <| SemimoduleCat.{u} R)ˣ
  % ...
  % noncomputable instance : CommGroup (Pic R) := fast_instance% (equivShrink _).symm.commGroup
  % ...
  % theorem inv_eq_dual (M : Pic R) : M⁻¹ = Pic.mk R (Dual R M) := ...
  % theorem mul_eq_tensor (M N : Pic R) : M * N = Pic.mk R (M ⊗[R] N) := ..."
  \textit{Source: [Stacks Project], Tag 01CR (Modules on Ringed Spaces,
  ``Invertible modules''), Definition~01CX and Lemma~\texttt{lemma-constructions-%
  invertible}. The isomorphism classes of invertible \(\mathcal{O}_X\)-modules
  form an abelian group \(\Pic(X)\) under \(\otimes\), with unit
  \(\mathcal{O}_X\) and inverse the dual; the tensor product of invertibles is
  invertible and the dual is a tensor inverse.}
  Let \(X\) be a scheme. The isomorphism classes of \(\otimes\)-invertible
  objects of \(\Scheme.\mathtt{Modules}\,X\)
  (\cref{def:scheme_modules_isinvertible}) form a commutative monoid under the
  substrate tensor product \(\otimes_X\) in which every element is a unit ---
  equivalently, an abelian group --- the Picard group \(\Pic(X)\) (Stacks tag
  01CR).

  \textbf{Primary construction (directly on locally-trivial iso-classes, à la
  \(\mathtt{Module.Invertible}\)).} The commutative group is built directly on
  isomorphism classes of \emph{locally-trivial} (equivalently \(\otimes\)-invertible)
  sheaves, mirroring Mathlib's ring-level Picard group
  \(\mathtt{Module.Invertible}\) / \(\mathtt{CommRing.Pic}\)
  (\texttt{Mathlib/RingTheory/PicardGroup.lean}). The group operation is
  \([L]\cdot[M] := [L \otimes_X M]\), well-defined on iso-classes by the
  bifunctoriality of \(\otimes_X\)
  (\cref{lem:scheme_modules_tensorobj_functoriality}) and again locally trivial by
  the closure lemma \cref{lem:tensorobj_lift_onproduct}; the identity is
  \([\mathcal{O}_X]\); and the inverse of \([L]\) is the class of the dual
  \([L^{-1}]\) supplied by \cref{lem:tensorobj_inverse_invertible} (with
  \(L \otimes_X L^{-1} \cong \mathcal{O}_X\)).

  The four group-axiom fields are defined \emph{by hand}, each discharged by a
  \emph{single} existence-of-isomorphism read as an equality of iso-classes --- passage
  to iso-classes turns an isomorphism of sheaves into an equality of classes:
  \[
    \mathtt{one\_mul} \;\leftarrow\; \langle \lambda_X \rangle,
    \quad
    \mathtt{mul\_one} \;\leftarrow\; \langle \rho_X \rangle,
    \quad
    \mathtt{mul\_assoc} \;\leftarrow\; \langle \alpha_{X,Y,Z} \rangle,
    \quad
    \mathtt{mul\_comm} \;\leftarrow\; \langle \beta_{X,Y} \rangle,
  \]
  where the unit laws are the unitors \cref{lem:tensorobj_unit_iso}, associativity is
  the associator \cref{lem:tensorobj_assoc_iso}, and commutativity is the braiding
  \cref{lem:tensorobj_comm_iso}, each consumed only as the proposition
  \(\mathrm{Nonempty}(\cdots \cong \cdots)\); the pentagon, triangle and hexagon
  coherence carried by a \(\MonoidalCategory\) typeclass is \emph{never} invoked in
  any of the four axiom proofs.

  This by-hand assembly does \emph{not} go through
  \(\mathtt{CategoryTheory.monoidOfSkeletalMonoidal}\) or \(\mathtt{Skeleton}\)
  (\texttt{Mathlib/CategoryTheory/Monoidal/Skeleton.lean}): those constructions
  \emph{require} a coherent \([\mathtt{MonoidalCategory}]\) (and, for the commutative
  law, \([\mathtt{BraidedCategory}]\)) on the carrier category --- precisely the
  coherent monoidal category on \(\Scheme.\mathtt{Modules}\,X\) over the \emph{varying}
  structure sheaf that this construction explicitly avoids. Consequently the group law
  requires \emph{no} \(J.W.\mathtt{IsMonoidal}\), \emph{no}
  \(\mathtt{LocalizedMonoidal}\) and \emph{no} \(\mathtt{Skeleton}\) instance on
  \(\Scheme.\mathtt{Modules}\,X\). The genuine precedent is Mathlib's fixed-ring
  Picard group \(\mathtt{CommRing.Pic}\,R := \mathrm{Skeleton}(
  \mathtt{Module.Invertible}\,R)\) (\texttt{Mathlib/RingTheory/PicardGroup.lean}):
  there the \(\mathtt{CommGroup}\) \emph{is} read off the \(\mathrm{Skeleton}\) of the
  coherent monoidal category \(\mathtt{MonoidalCategory}(\mathtt{Module.Invertible}\,R)\)
  that Mathlib supplies over a fixed ring. The present construction \emph{cannot} take
  that \(\mathrm{Skeleton}\) route --- no coherent monoidal category on
  \(\Scheme.\mathtt{Modules}\,X\) for the varying structure sheaf is available --- so
  its \(\mathtt{CommGroup}\) is built directly from the existence-of-isomorphism
  propositions above, the construction ultimately resting, as
  \(\mathtt{CommRing.Pic}\) does, on iso-classes of invertible modules. The pinned
  carrier \(\mathtt{tensorObjIsoclassCommMonoid}\) realizes this commutative-group
  structure; the scheme-level construction is genuinely absent from Mathlib at the
  pinned commit.

  This is the \emph{lower-risk} path: every ingredient is either already proved
  (the unitors \cref{lem:tensorobj_unit_iso}, the braiding
  \cref{lem:tensorobj_comm_iso}) or locally-trivial-scoped and assembled (the
  associator \cref{lem:tensorobj_assoc_iso}). The only remaining obligations on this
  critical path are the restriction-compatibility comparison
  \cref{lem:tensorobj_restrict_iso} and the
  dual/inverse \cref{lem:tensorobj_inverse_invertible}. The associator
  \cref{lem:tensorobj_assoc_iso} is supplied by the unconditional route of
  \cref{sec:tensorobj_stalk_tensor} and transitively invokes the d.2 stalk--tensor
  commutation \cref{lem:stalk_tensor_commutation}.

  \textbf{Fallback (full route~(e)).} Alternatively, once
  \texttt{jw\_ismonoidal} lands, \(\Scheme.\mathtt{Modules}\,X =
  \mathtt{SheafOfModules}\,X.\mathtt{ringCatSheaf}\) is a full (symmetric)
  \(\MonoidalCategory\) via \(\mathtt{Localization.Monoidal.LocalizedMonoidal}\),
  and the Picard group is the units of its skeleton,
  \(\Pic(X) = \mathrm{Units}(\mathrm{Skeleton}(\Scheme.\mathtt{Modules}\,X))\),
  exactly as Mathlib builds \(\mathtt{CommRing.Pic}\,R =
  \mathrm{Units}(\mathrm{Skeleton}(\mathtt{ModuleCat}\,R))\) and transports
  \(\mathtt{instCommGroupPic}\). This fallback yields the group on \emph{all}
  invertible modules and is retained for the full monoidal generality, but it is
  not required for the relative Picard group law, which the primary construction
  already supplies.
\end{lemma}

\begin{proof}
  \uses{def:scheme_modules_tensorobj, def:scheme_modules_isinvertible, lem:scheme_modules_tensorobj_functoriality, lem:tensorobj_lift_onproduct, lem:tensorobj_inverse_invertible, lem:tensorobj_assoc_iso, lem:tensorobj_unit_iso, lem:tensorobj_comm_iso, def:tensorobjisoofiso, def:tensorobj_unit_self_iso, def:tensorobj_middlefour, def:pic_setoid, def:pic_mul, def:pic_inv}
  Work directly on isomorphism classes of locally-trivial sheaves, as in Mathlib's
  \(\mathtt{Module.Invertible}\) / \(\mathtt{CommRing.Pic}\)
  (\texttt{Mathlib/RingTheory/PicardGroup.lean}). The binary operation
  \([L]\cdot[M] := [L \otimes_X M]\) is well-defined on iso-classes by the
  bifunctoriality of \(\otimes_X\)
  (\cref{lem:scheme_modules_tensorobj_functoriality}) --- an isomorphism in either
  argument tensors to an isomorphism --- and its value is again locally trivial by
  the closure lemma \cref{lem:tensorobj_lift_onproduct}, so the operation closes on
  the carrier. The identity element is \([\mathcal{O}_X]\). The four group-axiom
  fields are each defined \emph{by hand} from a \emph{single} existence-of-isomorphism,
  read as an equality of iso-classes,
  \[
    \mathtt{mul\_assoc} \leftarrow \langle\alpha\rangle,
    \quad
    \mathtt{one\_mul} \leftarrow \langle\lambda\rangle,
    \quad
    \mathtt{mul\_one} \leftarrow \langle\rho\rangle,
    \quad
    \mathtt{mul\_comm} \leftarrow \langle\beta\rangle:
  \]
  associativity is \cref{lem:tensorobj_assoc_iso}, the unit laws are the unitors
  \cref{lem:tensorobj_unit_iso}, and commutativity is the braiding
  \cref{lem:tensorobj_comm_iso}. No pentagon, triangle or hexagon identity is
  invoked, since each law asserts only \(\mathrm{Nonempty}(\cdots \cong \cdots)\).
  This is the same one-isomorphism-per-law shape that Mathlib's
  \(\mathtt{monoidOfSkeletalMonoidal}\)
  (\texttt{Mathlib/CategoryTheory/Monoidal/Skeleton.lean}) uses on a skeletal
  monoidal category via its skeletality map \(\mathtt{hC}\), but we do \emph{not}
  instantiate that construction: it requires a coherent \([\mathtt{MonoidalCategory}]\)
  (and \([\mathtt{BraidedCategory}]\)) on \(\Scheme.\mathtt{Modules}\,X\), which is
  exactly the structure for the varying structure sheaf that is unavailable.
  Finally every element is invertible: the dual \(L^{-1}\) of
  \cref{lem:tensorobj_inverse_invertible} satisfies \(L \otimes_X L^{-1} \cong
  \mathcal{O}_X\), so \([L^{-1}]\) is a two-sided inverse of \([L]\); equivalently,
  \([L]\) is a unit precisely when \(L\) is \(\otimes\)-invertible
  (\cref{def:scheme_modules_isinvertible}), the inverse witness carried
  existentially by \(\mathtt{IsInvertible}\). This assembles the abelian group ---
  the Picard group of \(X\) (Stacks tag 01CR) --- with \emph{no}
  \(J.W.\mathtt{IsMonoidal}\), \(\mathtt{LocalizedMonoidal}\) or \(\mathtt{Skeleton}\),
  exactly as \(\mathtt{CommRing.Pic}\) carries its \(\mathtt{CommGroup}\) on
  iso-classes of invertible modules (there off a fixed-ring \(\mathtt{Skeleton}\),
  here by hand).

  Alternatively (fallback), once \texttt{jw\_ismonoidal} provides the full
  \(\MonoidalCategory\) structure on \(\Scheme.\mathtt{Modules}\,X\) via
  \(\mathtt{LocalizedMonoidal}\), the group is the units of the skeleton,
  \(\mathrm{Units}(\mathrm{Skeleton}(\Scheme.\mathtt{Modules}\,X))\), exactly the
  Mathlib \(\mathtt{Units}(\mathrm{Skeleton}(\cdots))\) template used for
  \(\mathtt{CommRing.Pic}\); this route is not needed for the group law.
\end{proof}

\section{Pullback-monoidality: invertibility under a general scheme morphism}
\label{sec:tensorobj_pullback_monoidality}

The relative Picard consumer (\cref{chap:Picard_RelPicFunctor}) carries its line
bundles on the \emph{local-triviality} predicate \(\mathtt{IsLocallyTrivial}\)
(\cref{def:IsLocallyTrivial}): its carrier \(\mathtt{OnProduct}\) is the
subobject \(\{\,L \in \Scheme.\mathtt{Modules}\,(C \times_S T) \mid
\mathtt{IsLocallyTrivial}\,L\,\}\) of locally-trivial \(\mathcal{O}\)-modules, and
its functorial structure maps are \emph{pullback} maps of \(\mathcal{O}\)-modules:
the projection pullback \(\pi_T^*\) along \(C \times_S T \to T\) and the
base-change pullback along \(C \times_S T' \to C \times_S T\). For these maps to
descend to the carrier and to be \emph{additive} (an \(\mathtt{AddMonoidHom}\) of
isomorphism classes) one needs the comparison isomorphism
\[
  f^*(L \otimes_X L')
    \;\cong\;
  f^*L \otimes_Y f^*L'
\]
for \emph{locally-trivial} \(L, L'\) and a \emph{general} morphism of schemes
\(f : Y \to X\) --- the projection and the base-change morphism are neither open
immersions nor flat. This restricted comparison isomorphism, on locally-trivial
pairs, is the load-bearing result of the section
(\cref{lem:pullback_tensor_iso_loctriv}); the invertibility corollary
\texttt{isinvertible\_pullback} is then a three-line consequence. The comparison
is needed \emph{only} on locally-trivial pairs --- the carrier never presents a
general \(\mathcal{O}\)-module --- so no comparison for arbitrary modules is
required.

This is a genuinely different fact from the open-immersion compatibility
\cref{lem:tensorobj_restrict_iso}, and not a corollary of it. That lemma handles
restriction along an open immersion, and it does so by routing through the
\emph{lax right} adjoint \(\mathtt{pushforward}/\mathtt{restrictScalars}\), which is
strong monoidal there \emph{only} because an open immersion's structure-sheaf ring
map is a local \emph{isomorphism} (\(\mathtt{appIso}\,f\)). A general scheme
morphism has no such isomorphism, so the right-adjoint route is unavailable.

\paragraph{The reduction to two facts about pullback.} The corollary
\texttt{isinvertible\_pullback} is a three-line consequence of two facts about
pullback, exactly as in the Stacks proof of \texttt{lemma-pullback-invertible}. The
first is that pullback commutes with the tensor product on the \emph{locally-trivial}
pairs the consumer presents: the canonical comparison
\(f^*(L \otimes_X L') \cong f^*L \otimes_Y f^*L'\) is an isomorphism whenever
\(L, L'\) are locally trivial, for every \(f\)
(\cref{lem:pullback_tensor_iso_loctriv}). The second is that pullback preserves the
structure sheaf, \(f^*\mathcal{O}_X \cong \mathcal{O}_Y\)
(\cref{lem:pullback_unit_iso}, already in hand). Given a locally-trivial inverse
\(N\) with \(L \otimes_X N \cong \mathcal{O}_X\), pulling back yields
\[
  f^*L \otimes_Y f^*N
    \;\cong\; f^*(L \otimes_X N)
    \;\cong\; f^*\mathcal{O}_X
    \;\cong\; \mathcal{O}_Y,
\]
so \(f^*N\) is a tensor inverse of \(f^*L\). The locally-trivial comparison
\cref{lem:pullback_tensor_iso_loctriv} is the load-bearing input.

\paragraph{Why the apparent obstruction does not apply.} It is tempting to argue that
the comparison can never be an isomorphism because an (op)lax monoidal functor need
not preserve invertible objects --- global sections \(\Gamma\) is lax monoidal, yet
\(\Gamma(\mathbb{P}^1, \mathcal{O}(1)) = 0\) is not invertible. But that
counterexample concerns the \emph{lax} tensorator of \emph{pushforward}: \(\Gamma\)
is a right adjoint (global sections), and its comparison
\(\Gamma(\mathcal{F}) \otimes \Gamma(\mathcal{G}) \to \Gamma(\mathcal{F} \otimes
\mathcal{G})\) genuinely fails to be an isomorphism. The comparison we need is the
\emph{oplax} \(\delta\) of \emph{pullback}, a \emph{left} adjoint, running the other
way, \(f^*(M \otimes N) \to f^*M \otimes f^*N\). Pullback \emph{provably preserves
local triviality} (\cref{lem:IsLocallyTrivial_pullback}): if \(M\) is trivialised
by \(\mathcal{O}_U\) on a cover \(\{U_i\}\), then \(f^*M\) is trivialised by
\(\mathcal{O}_{f^{-1}U_i}\) on \(\{f^{-1}U_i\}\) --- using only
\(f^*\mathcal{O} \cong \mathcal{O}\) and the compatibility of pullback with
restriction along open immersions. Consequently, on locally-trivial objects, both
sides of \(\delta\) are locally trivial, and \(\delta\) restricts to a comparison
between the structure sheaf and itself on each chart, where it is forced to be an
isomorphism. The bare oplax structure supplies \(\delta\) as a \emph{morphism} for
free (\cref{lem:pullback_tensor_map}); the entire remaining content is the
chart-local check that, on a trivialising cover, \(\delta\) is an isomorphism.

\paragraph{The construction route.} The isomorphism is obtained by \emph{upgrading
the already-built oplax comparison map} \(\delta_{\mathrm{sheaf}} =
\mathtt{pullbackTensorMap}\) (\cref{lem:pullback_tensor_map}) to an isomorphism
through a chart-chase --- the same shape as the proven local-triviality results
\cref{lem:IsLocallyTrivial_pullback} and
\cref{lem:tensorobj_preserves_locally_trivial}, and \emph{not} via any concrete
inverse-image model or filtered-colimit construction. The argument decomposes into
four steps, carried out in the sub-lemmas cited by
\cref{lem:pullback_tensor_iso_loctriv}:
\begin{itemize}
  \item[(D1$'$)] \emph{Naturality of \(\delta\) in \((M,N)\).} The sheaf-level
    comparison \(\delta_{\mathrm{sheaf}}\) is natural in both arguments, assembled from
    the naturality of the abstract oplax \(\delta\) (Mathlib's
    \(\mathtt{Functor.OplaxMonoidal.}\delta\mathtt{\_natural}\)) and of the
    sheafification reconciliations used to build it
    (\cref{lem:pullback_tensor_map_natural}). This naturality is what lets a chart
    replace \(M|_{U_i}, N|_{U_i}\) by the trivial bundle \(\mathcal{O}\) without
    changing whether \(\delta\) is an isomorphism there.
  \item[(D2$'$)] \emph{\(\delta\) on the unit pair \((\mathcal{O},\mathcal{O})\) is an
    isomorphism.} On the trivial bundle the comparison
    \(\delta_{(\mathcal{O},\mathcal{O})} : f^*(\mathcal{O} \otimes \mathcal{O}) \to
    f^*\mathcal{O} \otimes f^*\mathcal{O}\) is reduced, by the \(\delta\)-wrapping
    \(\mathtt{isIso\_sheafifyDelta\_unitPair\_of\_isIso\_sheafifyEta}\) (the presheaf
    oplax left-unitality factored through sheafification, using
    \(\mathtt{W\_of\_isIso\_sheafification}\) and the flatness-free right-whiskering
    \(\mathtt{W\_whiskerRight\_of\_W}\)), to the single \emph{unit square}
    \(a_Y.\mathrm{map}(\eta\,F) \circ \mathtt{sheafifyUnitIso} = \mathtt{pullbackValIso}
    \circ \mathtt{pullbackObjUnitToUnit}\) (\texttt{eta\_bridge\_unit\_square}). That
    square, transposed across the sheaf pullback--pushforward adjunction, bottoms out in
    the structure-sheaf comparison \(\mathtt{pullbackUnitIso}\)
    (\cref{lem:pullback_unit_iso}, an isomorphism for \emph{all} \(f\)); hence
    \(\delta_{(\mathcal{O},\mathcal{O})}\) is an isomorphism
    (\cref{lem:pullback_tensor_iso_unit}).
  \item[(D3$'$)] \emph{Composition coherence of \(\delta\)} (the genuinely new
    ingredient). For composable scheme morphisms \(h, f\) the comparison \(\delta\) for
    the composite \(h \circ f\) factors through the comparisons for \(f\) and for \(h\),
    conjugated by the pullback pseudofunctoriality isomorphism
    \(\mathtt{pullbackComp}\,h\,f\) (\cref{lem:pullback_tensor_map_basechange}).
    Specialising to \(h := j_i'\) at an open \(U_i\) of \(X\) with preimage
    \(f^{-1}U_i\) --- with \(g_i : f^{-1}U_i \to U_i\) the restriction of \(f\) and
    \(j_i : U_i \hookrightarrow X\), \(j_i' : f^{-1}U_i \hookrightarrow Y\) the open
    immersions, so \(f \circ j_i' = j_i \circ g_i\) --- this intertwines the
    restriction-to-\(f^{-1}U_i\) of \(\delta\) for \(f\) with the \(\delta\) for
    \(g_i\). Because \(\mathtt{pullbackTensorMap}\) is a hand-built four-fold composite
    rather than an adjunction transpose, the unit-coherence mate calculus of
    \cref{lem:pullbackObjUnitToUnit_comp} does \emph{not} transfer; the coherence is
    instead established as a paste of four composition-coherence squares, decomposing
    \(\delta\) of the composite pullback via Mathlib's
    \(\mathtt{Functor.OplaxMonoidal.comp\_}\delta\) and the conjugation
    \(\mathtt{conjugateEquiv\_pullbackComp\_inv}\) of the pseudofunctor isomorphisms.
  \item[(D4$'$)] \emph{Chart-chase assembly.} Choose a common affine cover \(\{U_i\}\)
    of \(X\) trivialising both \(M\) and \(N\) (the common refinement, exactly as in
    \cref{lem:tensorobj_preserves_locally_trivial}). On each \(f^{-1}U_i\), the
    base-change coherence (D3$'$) presents the restriction of \(\delta\) for \(f\) as the
    \(\delta\) for \(g_i\) \emph{flanked} by the open-immersion comparisons
    \(\delta^{j_i'}, \delta^{j_i}\); cancelling those flanking factors --- legitimate only
    because each is invertible by the open-immersion comparison
    \cref{lem:pullback_tensor_map_isiso_open_immersion} (K1) --- aligns the restriction with
    the \(\delta\) for \(g_i\); naturality (D1$'$) replaces the trivialised
    \(M|_{U_i}, N|_{U_i}\) by \(\mathcal{O}\); and (D2$'$) makes the resulting
    comparison on the unit pair an isomorphism. Thus \(\delta\) restricts to an
    isomorphism on every member of the cover \(\{f^{-1}U_i\}\) of \(Y\), and the
    restriction-detects-isomorphism criterion
    \cref{lem:isiso_of_isiso_restrict} promotes it to a global isomorphism. The
    structure mirrors the chase of \cref{lem:IsLocallyTrivial_pullback}.
\end{itemize}
Two sub-steps carry genuinely new content: (D3$'$) the composition coherence of
\(\delta\), and --- as the realised chart-chase makes explicit --- (K1) the
open-immersion comparison \cref{lem:pullback_tensor_map_isiso_open_immersion}, whose
invertibility is what licenses cancelling the flanking factors \(\delta^{j_i'}, \delta^{j_i}\)
in (D4$'$) and reducing to the \(\delta\) for \(g_i\). The remaining ingredients are lighter:
(D1$'$) is Mathlib naturality, (D2$'$) is the unit coherence backed by
\(\mathtt{pullbackUnitIso}\), and the cover assembly of (D4$'$) reuses the already-proven
machinery (\cref{lem:tensorobj_preserves_locally_trivial},
\cref{lem:isiso_of_isiso_restrict}). No concrete inverse-image model, no
filtered-colimit/\(\otimes\) interchange, and no general comparison for arbitrary
modules is constructed.

\begin{lemma}


\leanok
[The presheaf pushforward is lax monoidal]
  \label{lem:presheaf_pushforward_laxmonoidal}
  \lean{AlgebraicGeometry.Scheme.Modules.presheafPushforwardLaxMonoidal}
  Let \(\varphi\) be a morphism of \(\mathtt{CommRingCat}\)-valued ring presheaves,
  presenting a pushforward of presheaves of modules. Then the presheaf pushforward
  \(\mathtt{PresheafOfModules.pushforward}\,\varphi\) is \emph{lax monoidal}: it
  carries the comparison
  \[
    \mu :
      (\mathtt{pushforward}\,\varphi).\mathrm{obj}\,A
        \otimes (\mathtt{pushforward}\,\varphi).\mathrm{obj}\,B
      \;\longrightarrow\;
      (\mathtt{pushforward}\,\varphi).\mathrm{obj}\,(A \otimes B),
  \]
  obtained as the composite of the strong-monoidal pushforward of presheaves over a
  fixed ring presheaf (\(\mathtt{pushforward}_0\mathtt{OfCommRingCat}\)) with the lax
  monoidal restriction of scalars \(\mathtt{restrictScalars}\,\varphi\) along the
  ring-presheaf map. This is a statement at the \emph{presheaf} level of
  \(\mathcal{O}\)-modules; it does not require a monoidal structure on
  \(\mathtt{SheafOfModules}\).
\end{lemma}

\begin{proof}



\leanok
  Restriction of scalars along a homomorphism of ring presheaves is lax monoidal --- the
  comparison sends \((a \otimes b)\) to itself under the identification of the
  restricted tensor product with the source tensor product --- and the pushforward of
  presheaves of modules over a \emph{fixed} ring presheaf is strong monoidal. The
  presheaf pushforward along \(\varphi\) factors as the composite of these two functors,
  so its lax monoidal structure is the composite of the strong-monoidal structure on the
  first factor with the lax structure on the second; the comparison \(\mu\) is the
  whiskered composite of the two comparisons.
\end{proof}

\begin{lemma}


\leanok
[The presheaf pullback is oplax monoidal; the comparison map $\delta$]
  \label{lem:presheaf_pullback_oplaxmonoidal}
  \lean{AlgebraicGeometry.Scheme.Modules.presheafPullbackOplaxMonoidal}
  \uses{lem:presheaf_pushforward_laxmonoidal}
  In the same setting, the presheaf pullback \(\mathtt{PresheafOfModules.pullback}\,\varphi\)
  --- the left adjoint of the lax-monoidal presheaf pushforward of
  \cref{lem:presheaf_pushforward_laxmonoidal} --- is \emph{oplax monoidal}. In
  particular it carries the canonical comparison \emph{morphism}
  \[
    \delta :
      (\mathtt{pullback}\,\varphi).\mathrm{obj}\,(A \otimes B)
      \;\longrightarrow\;
      (\mathtt{pullback}\,\varphi).\mathrm{obj}\,A
        \otimes (\mathtt{pullback}\,\varphi).\mathrm{obj}\,B,
  \]
  produced by the doctrinal adjunction transfer
  \(\mathtt{Adjunction.leftAdjointOplaxMonoidal}\): the oplax structure on a left
  adjoint is the adjunction mate of the lax structure on its right adjoint. The
  comparison \(\delta\) is at the \emph{presheaf} level, and it is only a
  \emph{morphism}: upgrading \(\delta\) to an isomorphism is the remaining content, and
  is exactly what the chart-chase of \cref{lem:pullback_tensor_iso_loctriv}
  supplies for a locally-trivial pair.
\end{lemma}

\begin{proof}
  \uses{lem:presheaf_pushforward_laxmonoidal}
  \leanok
  The presheaf pullback is by construction the left adjoint of the presheaf pushforward,
  which is lax monoidal by \cref{lem:presheaf_pushforward_laxmonoidal}. Mathlib's
  \(\mathtt{Adjunction.leftAdjointOplaxMonoidal}\) turns the lax-monoidal structure on a
  right adjoint into an oplax-monoidal structure on the left adjoint, with the oplax
  comparison \(\delta\) defined as the mate of the lax comparison \(\mu\) under the
  adjunction. Applying this to the pullback--pushforward adjunction yields the asserted
  oplax structure and the comparison morphism \(\delta\).
\end{proof}

% NOTE: ABANDONED general route, OFF the critical path. This lemma asserts the
% comparison iso for ARBITRARY modules M,N. The consumer (the relative Picard carrier
% OnProduct = { L // IsLocallyTrivial L }) never presents a general module: it needs
% the comparison only on LOCALLY-TRIVIAL pairs, which is supplied by the chart-chase
% lem:pullback_tensor_iso_loctriv WITHOUT any concrete inverse-image / filtered-colimit
% build. The general statement here is retained for the record but is no longer
% load-bearing; its proof (Steps D1-D4 below, via lem:pullback_lan_decomposition and
% lem:pullback0_tensor_iso) is the Mathlib-scale build the project has pivoted off.
% lem:isinvertible_pullback no longer \uses this lemma.

% NOTE: OFF the critical path. General presheaf-pullback infrastructure, retained
% (and proven), but no longer load-bearing: it served only the abandoned general
% comparison lem:pullback_tensor_iso. The committed loc-triv route
% (lem:pullback_tensor_iso_loctriv) does not use it.

% NOTE: ABANDONED general route, OFF the critical path. This is the Mathlib-scale
% filtered-colimit/tensor interchange (the irreducible piece D3 of the abandoned general
% comparison lem:pullback_tensor_iso). The consumer never needs it: the loc-triv
% comparison lem:pullback_tensor_iso_loctriv is obtained by a chart-chase upgrading the
% already-built oplax map, with no concrete inverse-image colimit model. Retained for the
% record only.

\begin{lemma}


\leanok
[The sheaf-level pullback--tensor comparison map $\delta_{\mathrm{sheaf}}$]
  \label{lem:pullback_tensor_map}
  \lean{AlgebraicGeometry.Scheme.Modules.pullbackTensorMap}
  \uses{lem:presheaf_pullback_oplaxmonoidal, def:scheme_modules_tensorobj, def:sheafify_tensor_unit_iso, def:pullback_val_iso}
  Let \(f : Y \to X\) be a morphism of schemes and let
  \(M, N \in \Scheme.\mathtt{Modules}\,X\) be \emph{arbitrary} \(\mathcal{O}_X\)-modules.
  There is a canonical comparison \emph{morphism}
  \[
    \delta_{\mathrm{sheaf}} :
      (\mathtt{Scheme.Modules.pullback}\,f).\mathrm{obj}\,(M \otimes_X N)
      \;\longrightarrow\;
      (\mathtt{Scheme.Modules.pullback}\,f).\mathrm{obj}\,M
        \otimes_Y (\mathtt{Scheme.Modules.pullback}\,f).\mathrm{obj}\,N
  \]
  of \(\mathcal{O}_Y\)-modules, natural in \(M\) and \(N\). It is a \emph{map only}: in
  general it is not asserted to be an isomorphism.
\end{lemma}

\begin{proof}
  \uses{lem:presheaf_pullback_oplaxmonoidal, def:scheme_modules_tensorobj, def:sheafify_tensor_unit_iso, def:pullback_val_iso}
  \leanok
  The presheaf pullback is oplax monoidal (\cref{lem:presheaf_pullback_oplaxmonoidal}),
  so at the presheaf level it carries the comparison morphism
  \[
    \delta : (\mathtt{pullback}\,\varphi).\mathrm{obj}\,(A \otimes B)
      \;\longrightarrow\;
      (\mathtt{pullback}\,\varphi).\mathrm{obj}\,A
        \otimes (\mathtt{pullback}\,\varphi).\mathrm{obj}\,B,
  \]
  where \(\varphi\) is the ring-presheaf map presenting \(f\). One transports \(\delta\)
  through sheafification to the sheaf level, using the same
  \(\mathtt{sheafificationCompPullback}\) device that supplies Steps 1--2 of
  \cref{lem:tensorobj_restrict_iso}: the abstract pullback of \(\mathcal{O}\)-modules
  agrees, up to canonical natural isomorphism, with the sheafification of the presheaf
  pullback, and likewise the sheafified tensor product \(\otimes_Y\) is the
  sheafification of the presheaf tensor. The right-hand side reconciliation of the unit
  and tensor factors with their sheafified forms is supplied by the
  sheafification-monoidality brick \(\mathtt{sheafifyTensorUnitIso}\). Whiskering and
  composing these canonical isomorphisms with the sheafification of \(\delta\) produces
  the asserted sheaf-level morphism \(\delta_{\mathrm{sheaf}}\), natural in \(M, N\)
  because each ingredient is. No claim of invertibility is made: \(\delta_{\mathrm{sheaf}}\)
  is the carrier on which the invertible-case isomorphism argument of
  \texttt{isinvertible\_pullback} acts.
\end{proof}

\begin{lemma}


\leanok
[Reduction of $\mathtt{pullbackTensorMap}$ iso-ness to the sheafified presheaf $\delta$]
  \label{lem:isiso_pullbacktensormap_of_sheafifydelta}
  \lean{AlgebraicGeometry.Scheme.Modules.isIso_pullbackTensorMap_of_isIso_sheafifyDelta}
  \uses{lem:pullback_tensor_map, lem:isiso_sheafify_tensorhom_pullbackvaliso}
  Let \(f : Y \to X\) be a morphism of schemes and let
  \(M, N \in \Scheme.\mathtt{Modules}\,X\). Write
  \[
    \varphi' :
      (X.\mathtt{presheaf} \circ \mathtt{forget}_2)
      \;\longrightarrow\;
      (\mathtt{Opens.map}\,f.\mathtt{base})^{\mathrm{op}}
        \circ (Y.\mathtt{presheaf} \circ \mathtt{forget}_2)
  \]
  for the induced presheaf-of-modules map presenting \(f\), and let
  \(a_Y = \mathtt{PresheafOfModules.sheafification}\,(\mathbf{1}_{Y.\mathtt{ringCatSheaf}})\)
  denote sheafification on \(Y\). If the sheafified presheaf-level oplax comparison
  \[
    a_Y.\mathrm{map}\bigl(
      \delta\,(\mathtt{PresheafOfModules.pullback}\,\varphi')\,M.\mathtt{val}\,N.\mathtt{val}
    \bigr)
  \]
  is an isomorphism, then the sheaf-level comparison
  \(\mathtt{pullbackTensorMap}\,f\,M\,N : f^*(M \otimes_X N) \to f^*M \otimes_Y f^*N\)
  of \cref{lem:pullback_tensor_map} is an isomorphism.
\end{lemma}

\begin{proof}
  \uses{lem:pullback_tensor_map, lem:isiso_sheafify_tensorhom_pullbackvaliso}
  \leanok
  Unfolding \(\mathtt{pullbackTensorMap}\) (\cref{lem:pullback_tensor_map}) exhibits it as a
  four-fold composite
  \[
    \mathtt{pullbackTensorMap}\,f\,M\,N
      \;=\;
    p_1 \;\mathbin{;}\; a_Y.\mathrm{map}\,\delta \;\mathbin{;}\; p_3 \;\mathbin{;}\; p_4,
  \]
  where:
  \(p_1 = (\mathtt{sheafificationCompPullback}).\mathrm{hom}\), reconciling the abstract
  \(\mathcal{O}\)-module pullback with the sheafification of the presheaf pullback;
  \(a_Y.\mathrm{map}\,\delta\) is the sheafification of the presheaf-level oplax comparison
  \(\delta\); \(p_3 = (\mathtt{sheafifyTensorUnitIso}).\mathrm{hom}\), the
  sheafification-monoidality brick reconciling the sheafified tensor with the sheaf tensor;
  and \(p_4 = a_Y.\mathrm{map}\) of the tensor \(\otimes_{\mathrm{m}}\) of the two
  \(\mathtt{pullbackValIso}\) comparisons (one for \(M\), one for \(N\)) identifying the
  sheafified pullback value with the pullback.

  The outer three factors are isomorphisms unconditionally: \(p_1\) and \(p_3\) are the
  \(.\mathrm{hom}\) components of isomorphisms, and \(p_4\) is the image under the additive
  functor \(a_Y.\mathrm{map}\) of a tensor product of the two \(\mathtt{pullbackValIso}\)
  \emph{isomorphisms}, hence itself an isomorphism. The hypothesis supplies the one
  conditional middle factor \(a_Y.\mathrm{map}\,\delta\). A composite of isomorphisms is an
  isomorphism, so \(\mathtt{pullbackTensorMap}\,f\,M\,N\) is an isomorphism.
\end{proof}

\begin{lemma}


\leanok
[Composition coherence of the pushforward unit comparison]
  \label{lem:unitToPushforwardObjUnit_comp}
  \lean{AlgebraicGeometry.Scheme.Modules.unitToPushforwardObjUnit_comp}
  Let \(h : Z \to Y\) and \(f : Y \to X\) be composable morphisms of schemes. Write
  \(\mathrm{upu}\) for the pushforward (right-adjoint) unit comparison
  \(\mathtt{SheafOfModules.unitToPushforwardObjUnit}\), the canonical map
  \(\mathcal{O} \to f_*(\mathcal{O})\) comparing the structure sheaf of the target with
  the pushforward of the structure sheaf of the source. Then \(\mathrm{upu}\) satisfies
  the composition coherence
  \[
    \mathrm{upu}(h \mathbin{;} f)
      \;=\;
    \mathrm{upu}\,f \;\mathbin{;}\;
      (\mathtt{pushforward}\,f).\mathrm{map}\bigl(\mathrm{upu}\,h\bigr) \;\mathbin{;}\;
      (\mathtt{pushforwardComp}\,h\,f).\mathrm{hom},
  \]
  where \(\mathtt{pushforwardComp}\,h\,f : (h \mathbin{;} f)_* \cong f_* \circ h_*\) is the
  pseudofunctoriality isomorphism of the pushforward. Sectionwise both sides are nothing
  but the functoriality of the structure-sheaf ring maps, so the identity holds on the
  nose after the comparison of sections. This is the right-adjoint half of the unit
  composition coherence, and it is the input consumed by the mate transport that produces
  the pullback-side coherence \cref{lem:pullbackObjUnitToUnit_comp}.
\end{lemma}

\begin{proof}



\leanok
  Both sides are morphisms of \(\mathcal{O}\)-modules; it suffices to compare them on
  sections over each open. Each comparison \(\mathrm{upu}\) acts on sections as the
  structure-sheaf ring homomorphism induced by the relevant morphism of schemes, and the
  pseudofunctor coherence \(\mathtt{pushforwardComp}\) acts as the identification of
  \((h \mathbin{;} f)\)-pushforward sections with iterated \(f_*\,h_*\)-pushforward
  sections. The composition coherence is therefore the functoriality identity
  \((h \mathbin{;} f)^\sharp = h^\sharp \circ f^\sharp\) of the structure-sheaf ring maps,
  which holds definitionally; the two sides agree section by section.
\end{proof}

\begin{lemma}


\leanok
[Composition coherence of the pullback unit comparison]
  \label{lem:pullbackObjUnitToUnit_comp}
  \lean{AlgebraicGeometry.Scheme.Modules.pullbackObjUnitToUnit_comp}
  \uses{lem:unitToPushforwardObjUnit_comp}
  Let \(h : Z \to Y\) and \(f : Y \to X\) be composable morphisms of schemes. Write
  \(\mathrm{pbu}\) for the pullback (left-adjoint) unit comparison
  \(\mathtt{SheafOfModules.pullbackObjUnitToUnit}\), the canonical map
  \(f^*\mathcal{O} \to \mathcal{O}\). Then \(\mathrm{pbu}\) satisfies the composition
  coherence
  \[
    \mathrm{pbu}(h \mathbin{;} f)
      \;=\;
    (\mathtt{pullbackComp}\,h\,f).\mathrm{inv} \;\mathbin{;}\;
      (\mathtt{pullback}\,h).\mathrm{map}\bigl(\mathrm{pbu}\,f\bigr) \;\mathbin{;}\;
      \mathrm{pbu}\,h,
  \]
  where \(\mathtt{pullbackComp}\,h\,f : (h \mathbin{;} f)^* \cong h^* \circ f^*\) is the
  pseudofunctoriality isomorphism of the pullback. This is the genuinely-new ingredient
  for \cref{lem:pullback_unit_iso}: the left-adjoint pullback is defined abstractly, with
  no sectionwise value, so this identity is \emph{not} a sectionwise statement and cannot
  be checked on sections.
\end{lemma}

\begin{proof}
  \uses{lem:unitToPushforwardObjUnit_comp}
  \leanok
  The pullback \(f^*\) is the left adjoint of the pushforward \(f_*\), and likewise for
  \(h\) and \(h \mathbin{;} f\); the comparison \(\mathrm{pbu}\,f\) is the adjunction
  \emph{mate} of the pushforward comparison \(\mathrm{upu}\,f\)
  (\cref{lem:unitToPushforwardObjUnit_comp}) across the pullback--pushforward adjunction,
  and similarly for \(h\) and \(h \mathbin{;} f\). Since the mate correspondence is a
  bijection, the asserted identity follows once both sides are transposed under the
  adjunction \((h \mathbin{;} f)^* \dashv (h \mathbin{;} f)_*\).

  Under this transpose the left-hand side \(\mathrm{pbu}(h \mathbin{;} f)\) becomes the
  pushforward comparison \(\mathrm{upu}(h \mathbin{;} f)\). The right-hand side, a
  composite involving \((\mathtt{pullbackComp}\,h\,f).\mathrm{inv}\) and the mates of
  \(\mathrm{pbu}\,f\) and \(\mathrm{pbu}\,h\), transposes --- using the conjugation of the
  pseudofunctoriality isomorphisms (\(\mathtt{pullbackComp}\,h\,f\) on the left adjoints
  corresponds to \(\mathtt{pushforwardComp}\,h\,f\) on the right adjoints) and the unit of
  the composite adjunction \((f^* \dashv f_*) \circ (h^* \dashv h_*)\) --- exactly into the
  right-hand side of the pushforward coherence
  \cref{lem:unitToPushforwardObjUnit_comp}, namely
  \(\mathrm{upu}\,f \mathbin{;} f_*(\mathrm{upu}\,h) \mathbin{;}
  (\mathtt{pushforwardComp}\,h\,f).\mathrm{hom}\). Thus the two transposed sides coincide
  by \cref{lem:unitToPushforwardObjUnit_comp}, and the original identity follows by the
  injectivity of the transpose.
\end{proof}

\begin{lemma}


\leanok
[Pullback preserves the structure sheaf]
  \label{lem:pullback_unit_iso}
  \lean{AlgebraicGeometry.Scheme.Modules.pullbackUnitIso}
  \uses{def:scheme_modules_tensorobj, def:pullback_obj_unit_to_unit_iso}
  Let \(f : Y \to X\) be a morphism of schemes. The pullback of the structure sheaf
  is the structure sheaf: there is a canonical isomorphism
  \[
    f^*\bigl(\mathtt{SheafOfModules.unit}\,X.\mathtt{ringCatSheaf}\bigr)
      \;\xrightarrow{\;\sim\;}\;
    \mathtt{SheafOfModules.unit}\,Y.\mathtt{ringCatSheaf},
    \qquad\text{i.e.}\qquad
    f^* \mathcal{O}_X \;\cong\; \mathcal{O}_Y,
  \]
  where \(\mathtt{SheafOfModules.unit}\,X.\mathtt{ringCatSheaf} = \mathcal{O}_X\) is
  the monoidal unit of \(\Scheme.\mathtt{Modules}\,X\).
\end{lemma}

\begin{proof}
  \uses{def:scheme_modules_tensorobj, def:pullback_obj_unit_to_unit_iso}
  \leanok
  This is a one-step consequence of a Mathlib instance, with no chart-chase. The
  canonical comparison map is the genuine Mathlib morphism
  \[
    \mathtt{SheafOfModules.pullbackObjUnitToUnit}\,f :
      f^*\mathcal{O}_X \;\longrightarrow\; \mathcal{O}_Y,
  \]
  the unit comparison \(\eta\) of \(f^*\); abbreviate it \(\mathrm{pbu}\,f\). Mathlib
  proves this map is an isomorphism whenever the comparison functor
  \(\mathtt{Opens.map}\,f.\mathtt{base}\) on opens is \(\mathtt{Final}\): the instance
  \(\mathtt{SheafOfModules.instIsIsoPullbackObjUnitToUnitOfFinal}\).

  The point is that the finality hypothesis holds \emph{unconditionally}, for every
  morphism of schemes \(f\). The preimage functor \(\mathtt{Opens.map}\,f.\mathtt{base}
  : \mathtt{Opens}\,X \to \mathtt{Opens}\,Y\) is a frame homomorphism: it preserves
  arbitrary unions and finite intersections, hence preserves all finite limits. A
  functor between small categories that preserves finite limits is representably flat,
  so \(\mathtt{final\_of\_representablyFlat}\) gives \((\mathtt{Opens.map}\,f.\mathtt{base}).\mathtt{Final}\)
  with no hypothesis on \(f\) whatsoever. (There is no need to pass to affine charts:
  finality is not a special feature of open immersions but a general property of
  preimage-on-opens.)

  Consequently the instance \(\mathtt{instIsIsoPullbackObjUnitToUnitOfFinal}\) fires
  directly, and \(\mathrm{pbu}\,f\) is an isomorphism for every \(f\). The asserted
  isomorphism \(f^*\mathcal{O}_X \cong \mathcal{O}_Y\) is the bundled iso
  \(\mathtt{pullbackUnitIso}\,f\) of this comparison.
\end{proof}

\begin{lemma}


\leanok
[Naturality of the sheaf-level pullback--tensor comparison $\delta_{\mathrm{sheaf}}$ (D1$'$)]
  \label{lem:pullback_tensor_map_natural}
  \lean{AlgebraicGeometry.Scheme.Modules.pullbackTensorMap_natural}
  \uses{lem:pullback_tensor_map, lem:presheaf_pullback_oplaxmonoidal, lem:sheafify_tensor_unit_iso_natural, lem:pullback_val_iso_natural}
  Let \(f : Y \to X\) be a morphism of schemes. The sheaf-level comparison map
  \(\delta_{\mathrm{sheaf}} = \mathtt{pullbackTensorMap}\) of
  \cref{lem:pullback_tensor_map} is \emph{natural} in both module arguments: for
  morphisms \(a : M \to M'\) and \(b : N \to N'\) in \(\Scheme.\mathtt{Modules}\,X\)
  the square
  \[
    \begin{array}{ccc}
      f^*(M \otimes_X N)
        & \xrightarrow{\;\delta_{\mathrm{sheaf}}\;}
        & f^*M \otimes_Y f^*N \\[2pt]
      \big\downarrow{\scriptstyle f^*(a \otimes b)}
        & & \big\downarrow{\scriptstyle f^*a \,\otimes\, f^*b} \\[4pt]
      f^*(M' \otimes_X N')
        & \xrightarrow{\;\delta_{\mathrm{sheaf}}\;}
        & f^*M' \otimes_Y f^*N'
    \end{array}
  \]
  commutes.
\end{lemma}

\begin{proof}
  \uses{lem:pullback_tensor_map, lem:presheaf_pullback_oplaxmonoidal, lem:sheafify_tensor_unit_iso_natural, lem:pullback_val_iso_natural}
  \leanok
  The map \(\delta_{\mathrm{sheaf}}\) is built (\cref{lem:pullback_tensor_map}) by
  transporting the abstract oplax comparison \(\delta\) of the presheaf pullback
  (\cref{lem:presheaf_pullback_oplaxmonoidal}) through the sheafification
  reconciliations. The asserted square is the pasting of four naturality squares, one for
  each ingredient of this composite; we record them in order, the only delicate one being the
  naturality of the oplax comparison itself.

  \emph{Square 2: naturality of the oplax comparison \(\delta\).} The comparison \(\delta\) is
  the oplax-monoidal comparison of the presheaf-level pullback functor
  \(F = \mathtt{PresheafOfModules.pullback}\,\varphi'\), whose defining ring map
  \(\varphi' = (\mathtt{toRingCatSheafHom}\,f).\mathrm{hom}\) has, as the source of its domain,
  the ring presheaf of \(X\). For the monoidal structure on \(\mathtt{PresheafOfModules}\) to
  be available on this domain the ring presheaf must be presented in its \emph{canonical}
  form as the composite
  \(\mathcal{O}_X \circ \mathtt{forget}_2\,\mathtt{CommRingCat}\,\mathtt{RingCat}\) --- the
  presentation on which that monoidal structure is defined. When the comparison map's
  construction is unfolded, however, the domain ring re-presents in a form definitionally
  equal, but not explicitly identical, to the canonical one, and on that form the
  monoidal structure is not directly available. The naturality of \(\delta\),
  \[
    \delta \circ F(a \otimes b) \;=\; (Fa \otimes Fb) \circ \delta,
  \]
  which holds for any oplax monoidal functor, is therefore applied by re-establishing the
  canonical domain-ring spelling at the functor argument itself: one ascribes the defining
  ring map \(\varphi'\) explicitly to the canonical composite-functor source, so that the
  domain monoidal structure resolves, and the naturality identity then matches the goal
  definitionally. This is a structural device --- it changes only the presentation of
  \(F\)'s defining ring map. No restatement of \(\mathtt{pullbackTensorMap}\) or of its helper
  isomorphisms is needed, and the unit-case lemma \(\mathtt{pullbackTensorMap\_unit\_isIso}\)
  (D2\(^\prime\)) is unaffected. After this step the comparison \(\delta\) commutes past
  \(F(a \otimes b)\), producing the factor \(Fa \otimes Fb\).

  \emph{Squares 1, 3, 4 and bifunctoriality.} The remaining ingredients of
  \(\delta_{\mathrm{sheaf}}\) are honest natural transformations and contribute already-closed
  naturality squares: the sheafification, the tensor/unit reconciliation
  \(\mathtt{sheafifyTensorUnitIso}\) (Square 3, \cref{lem:sheafify_tensor_unit_iso_natural}),
  and the \(\mathtt{pullbackValIso}\) reconciliation (Square 4,
  \cref{lem:pullback_val_iso_natural}). Finally, the tensor of the two component morphisms
  factors as \(f^*a \otimes f^*b\) by the bifunctoriality of \(\otimes_Y\) --- the interchange
  law \((f_1 \otimes f_2) \circ (g_1 \otimes g_2) = (f_1 \circ g_1) \otimes (f_2 \circ g_2)\),
  here \(\mathtt{tensorHom\_comp\_tensorHom}\). Pasting the four squares together with this
  bifunctoriality identity yields the asserted commuting square.
\end{proof}

\begin{lemma}


\leanok
  [Section-level naturality of the tensor/unit reconciliation
   \(\mathtt{sheafifyTensorUnitIso}\)]
  \label{lem:sheafify_tensor_unit_iso_natural}
  \lean{AlgebraicGeometry.Scheme.Modules.sheafifyTensorUnitIso_hom_natural}
  \uses{lem:pullback_tensor_map, def:sheafify_tensor_unit_iso, lem:sheafify_tensor_unit_iso_hom_eq_prime, lem:restrictscalarsid_map}
  The tensor/unit reconciliation isomorphism \(\mathtt{sheafifyTensorUnitIso}\) entering the
  fourth naturality square of \cref{lem:pullback_tensor_map_natural} is natural in both
  module arguments. Equivalently: after descending to sections over each open \(U\) and
  passing to the underlying module components, the naturality square is the
  bilinear, sectionwise naturality of the sheafification unit \(\eta\), verified one pure
  tensor \(p \otimes q\) at a time, using bilinearity.
\end{lemma}

\begin{proof}
  \uses{lem:pullback_tensor_map}
  \leanok
  Write \(\Phi\) for the natural transformation whose naturality square (in the two module
  arguments \(M, N\), along morphisms \(a : M \to M'\) and \(b : N \to N'\)) is to be shown
  to commute; \(\Phi\) is assembled from the reconciliation isomorphism
  \(\mathtt{sheafifyTensorUnitIso}\) (the intermediate construction inside
  \cref{lem:pullback_tensor_map}), so its components are morphisms of presheaves of modules.
  The square is closed at the categorical (tensor-of-morphisms) level, without any descent to
  sections or to elements.

  \emph{Step 1: a single tensor-of-morphisms presentation.} The \(\mathrm{hom}\) leg of
  \(\mathtt{sheafifyTensorUnitIso}\) is rewritten as a \emph{single} tensor of morphisms of the
  form \(a.\mathrm{map}(\eta_P \otimes \eta_Q)\), where
  \(\eta = (\mathtt{sheafificationAdjunction}).\mathrm{unit}\) is the unit
  \(\mathbf{1} \Rightarrow (\text{sheafify}) \circ (\text{inclusion})\) of the sheafification
  adjunction and \(a\) is the sheafification functor. The essential point is that all factors
  are kept inside \emph{one} monoidal structure on \(\mathtt{PresheafOfModules}\) --- the
  canonical presentation as a composite with the forgetful functor \(\mathtt{forget}_2\) ---
  so that the bifunctoriality law of \(\otimes\) can be applied consistently within one
  monoidal structure.

  \emph{Step 2: merge the two functor-image factors.} Both composites around the square are
  images, under the sheafification functor, of a tensor of morphisms. The two functor-image
  factors are merged into a single image of a composite by functoriality,
  \(F(g) \circ F(h) = F(g \circ h)\).

  \emph{Step 3: bifunctoriality and unit naturality.} The resulting equality of two tensors of
  morphisms is closed by the bifunctoriality (interchange) law
  \[
    (f_1 \otimes f_2) \circ (g_1 \otimes g_2) \;=\; (f_1 \circ g_1) \otimes (f_2 \circ g_2),
  \]
  applied within the single chosen monoidal structure, together with the two
  single-component naturality squares of the sheafification unit \(\eta\), one in each module
  argument:
  \[
    p \circ \eta_{M'} = \eta_{M} \circ (f^*a),
    \qquad
    q \circ \eta_{N'} = \eta_{N} \circ (f^*b).
  \]
  Each of these holds because \(\eta\) is a natural transformation. Combining the two
  single-component squares through the interchange law equates the two tensors of morphisms,
  which is exactly the naturality square of \(\mathtt{sheafifyTensorUnitIso}\).
\end{proof}

\begin{lemma}


\leanok
  [Naturality of the \(\mathtt{pullbackValIso}\) comparison]
  \label{lem:pullback_val_iso_natural}
  \lean{AlgebraicGeometry.Scheme.Modules.pullbackValIso_hom_natural}
  \uses{def:pullback_val_iso}
  The \(\mathtt{pullbackValIso}\) comparison --- the isomorphism reconciling the underlying
  presheaf of a pullback sheaf of modules with the pullback of its underlying presheaf ---
  is natural in its module argument. This is one of the (closed) input squares of the
  pasting that proves \cref{lem:pullback_tensor_map_natural}.
\end{lemma}

\begin{lemma}


\leanok
[The comparison on the unit pair $(\mathcal{O},\mathcal{O})$ is an isomorphism (D2$'$)]
  \label{lem:pullback_tensor_iso_unit}
  \lean{AlgebraicGeometry.Scheme.Modules.pullbackTensorMap_unit_isIso}
  \uses{lem:pullback_tensor_map, lem:pullback_unit_iso, lem:isiso_pullbacktensormap_of_sheafifydelta}
  Let \(f : Y \to X\) be a morphism of schemes. On the monoidal unit pair
  \((\mathcal{O}_X, \mathcal{O}_X)\) the sheaf-level comparison
  \[
    \delta_{(\mathcal{O},\mathcal{O})} :
      f^*(\mathcal{O}_X \otimes_X \mathcal{O}_X)
      \;\longrightarrow\;
      f^*\mathcal{O}_X \otimes_Y f^*\mathcal{O}_X
  \]
  of \cref{lem:pullback_tensor_map} is an isomorphism.
\end{lemma}

\begin{proof}
  \uses{lem:pullback_tensor_map, lem:isiso_pullbacktensormap_of_sheafifydelta}
  \leanok
  By \cref{lem:isiso_pullbacktensormap_of_sheafifydelta} applied on the unit pair,
  it suffices to verify that the sheafified presheaf comparison
  \(a_Y.\mathrm{map}\,(\delta\,(\mathtt{pullback}\,\varphi')\,\mathbf{1}\,\mathbf{1})\)
  is an isomorphism; the underlying presheaf of the structure-sheaf unit coincides with
  the presheaf monoidal unit \(\mathbf{1}\).

  On the unit pair the presheaf-level oplax comparison \(\delta\) is governed by Mathlib's
  oplax-monoidal \emph{unit coherence}: \(\mathtt{Functor.OplaxMonoidal.left\_unitality\_hom}\)
  expresses \(\delta\,(\mathtt{pullback}\,\varphi')\,\mathbf{1}\,\mathbf{1}\) as a composite of
  the functorial unitor \(F.\mathrm{map}\,(\lambda_{\mathbf{1}}).\mathrm{hom}\), the target
  unitor \(\lambda_{F.\mathrm{obj}\,\mathbf{1}}\), and the whiskered presheaf unit comparison
  \(\eta\,(\mathtt{pullback}\,\varphi') \rhd F.\mathrm{obj}\,\mathbf{1}\). The unitors are
  isomorphisms, so after sheafification by \(a_Y\) the genuine remaining content is
  \[
    \mathtt{IsIso}\bigl(a_Y.\mathrm{map}\,(\eta\,(\mathtt{pullback}\,\varphi'))\bigr),
  \]
  the sheafified presheaf unit comparison (the \(\delta\)-wrapping reduction). By \texttt{isiso\_sheafifyeta\_of\_unitsquare} it suffices, in turn, to prove the
  \emph{unit square}
  \[
    (\mathtt{pullbackValIso}\,f\,\mathcal{O}_X).\mathrm{inv} \mathbin{;}
      a_Y.\mathrm{map}\,(\eta\,(\mathtt{pullback}\,\varphi')) \mathbin{;}
      \mathtt{sheafifyUnitIso}.\mathrm{hom}
      \;=\; \mathtt{pullbackObjUnitToUnit}\,\varphi,
  \]
  which is \texttt{eta\_bridge\_unit\_square}. Granting that square,
  \texttt{isiso\_sheafifyeta\_of\_unitsquare} yields
  \(\mathtt{IsIso}\bigl(a_Y.\mathrm{map}\,(\eta\,(\mathtt{pullback}\,\varphi'))\bigr)\) --- the
  right-hand side \(\mathtt{pullbackObjUnitToUnit}\,\varphi = \mathtt{pullbackUnitIso}\,f :
  f^*\mathcal{O}_X \xrightarrow{\sim} \mathcal{O}_Y\) being an isomorphism for \emph{every}
  \(f\) (\cref{lem:pullback_unit_iso}) --- and the \(\delta\)-wrapping reduction above then
  delivers that \(\delta_{(\mathcal{O},\mathcal{O})}\) is an isomorphism.
\end{proof}

\subsection{The unit square (D2$'$): a mate-calculus telescope}

The proof of \cref{lem:pullback_tensor_iso_unit} bottoms out in the single morphism equation
labelled the \emph{unit square} below. Its proof is a telescope across three nested adjunction
layers (sheafification, presheaf pullback--pushforward, sheaf pullback--pushforward). This
subsection isolates the load-bearing intermediate identities as named lemmas so each can be
discharged independently, then assembles them.

We fix throughout a morphism of schemes \(f : Y \to X\) and write
\(\varphi = f.\mathtt{toRingCatSheafHom}\), \(\varphi' = \varphi.\mathrm{hom}\). Let \(a_X, a_Y\)
denote sheafification of presheaves of modules on \(X\), \(Y\); \(\mathbf{1}^{\mathrm{p}}\) the
presheaf monoidal unit; \(F = \mathtt{pullback}\,\varphi'\) the presheaf pullback;
\(\eta\,F : F.\mathrm{obj}\,\mathbf{1}^{\mathrm{p}} \to \mathbf{1}^{\mathrm{p}}\) its oplax unit
comparison; and \(\varepsilon(G)\) the lax unit comparison of a lax monoidal functor \(G\). The
isomorphisms \(\mathtt{pullbackValIso}\,f\,\mathcal{O}_X\) and \(\mathtt{sheafifyUnitIso}\) are
the file's comparison morphisms relating the sheaf-level pullback value to the sheafification of
the presheaf-level pullback. Write
\(\mathtt{sheafCounit}_\bullet = (\mathtt{sheafificationAdjunction}\,(\mathbf{1}_{\bullet.\mathtt{ringCatSheaf}.\mathrm{val}})).\mathrm{counit}\).

\begin{lemma}\leanok
[Composite-adjunction homEquiv factorisation]
  \label{lem:comp_homequiv_factor_sheafify_pullback}
  \lean{AlgebraicGeometry.Scheme.Modules.compHomEquivFactor}
  \uses{lem:presheaf_pushforward_laxmonoidal}
  Let \(\mathtt{adj}_1 : L_1 \dashv R_1\) and \(\mathtt{adj}_2 : L_2 \dashv R_2\) be composable
  adjunctions and \(A = \mathtt{adj}_1.\mathrm{comp}\,\mathtt{adj}_2 : L_1 \circ L_2 \dashv
  R_2 \circ R_1\) their composite. Then the hom-set bijection of the composite factors as the
  composite of the two factor bijections: for every \(g\),
  \[
    A.\mathrm{homEquiv}\,g
      \;=\;
    \mathtt{adj}_1.\mathrm{homEquiv}\bigl(\mathtt{adj}_2.\mathrm{homEquiv}\,g\bigr).
  \]
  Consequently, for the composite adjunction \(A\) of the unit square (left adjoint
  \(a_X \circ \mathtt{pullback}_{\mathrm{sheaf}}\,\varphi\)), the head of step~2 of
  \texttt{eta\_bridge\_unit\_square} equals
  \((\mathtt{sheafAdj}_X.\mathrm{homEquiv})^{-1}\bigl(A.\mathrm{homEquiv}\,g\bigr)\) with
  \(g = (\mathtt{sheafificationCompPullback}\,\varphi).\mathrm{hom}.\mathrm{app}\,\mathbf{1}^{\mathrm{p}}\).
\end{lemma}

\begin{proof}
  \uses{lem:presheaf_pushforward_laxmonoidal}
  \leanok
  The composite adjunction \(\mathtt{Adjunction.comp}\) is defined so that its unit is the
  pasting \(\mathtt{adj}_1.\mathrm{unit} \mathbin{;} R_1(\mathtt{adj}_2.\mathrm{unit})_{L_1(-)}\)
  and its counit dually; equivalently its hom-set bijection \(A.\mathrm{homEquiv}\) is the
  composite of the two factor bijections in the order \(\mathtt{adj}_2.\mathrm{homEquiv}\) then
  \(\mathtt{adj}_1.\mathrm{homEquiv}\). This is the standard factorisation
  \(\mathtt{Adjunction.comp\_homEquiv}\) (the naturality of \(\mathtt{homEquiv}\) under
  composition of adjunctions): a morphism \(L_1 L_2\,c \to d\) is transposed first across
  \(\mathtt{adj}_1\) to \(L_2\,c \to R_1\,d\) and then across \(\mathtt{adj}_2\) to
  \(c \to R_2 R_1\,d\). Applying this with \(\mathtt{adj}_1 =
  \mathtt{sheafificationAdjunction}\,(\mathbf{1}_X)\), \(\mathtt{adj}_2 =
  \mathtt{pullbackPushforwardAdjunction}_{\mathrm{sheaf}}\,\varphi\) and transposing back along
  \(\mathtt{adj}_1\) (i.e.\ applying \((\mathtt{sheafAdj}_X.\mathrm{homEquiv})^{-1}\)) gives the
  stated rewrite of the head of step~2.
\end{proof}

\begin{lemma}


\leanok
[The sheafify-then-pullback comparison is the left-adjoint-uniqueness isomorphism]
  \label{lem:sheafification_comp_pullback_eq_leftadjointuniq}
  \lean{AlgebraicGeometry.Scheme.Modules.sheafificationCompPullback_eq_leftAdjointUniq}
  Write \(A\) for the composite adjunction obtained by following the sheafification adjunction
  on \(X\) with the sheaf-level pullback--pushforward adjunction of \(\varphi\), namely
  \(A = (\mathtt{sheafificationAdjunction}\,(\mathbf{1}_X)).\mathrm{comp}\,
  (\mathtt{pullbackPushforwardAdjunction}_{\mathrm{sheaf}}\,\varphi)\), whose left adjoint is
  \(a_X \circ \mathtt{pullback}_{\mathrm{sheaf}}\,\varphi\); and write \(B\) for the composite
  adjunction obtained by following the presheaf-level pullback--pushforward adjunction of
  \(\varphi'\) with the sheafification adjunction on \(Y\), namely
  \(B = (\mathtt{pullbackPushforwardAdjunction}\,\varphi').\mathrm{comp}\,
  (\mathtt{sheafificationAdjunction}\,(\mathbf{1}_Y))\), whose left adjoint is \(F \circ a_Y\).
  These two left adjoints are canonically isomorphic, and the canonical ``sheafify-then-pullback''
  comparison agrees on the nose with the left-adjoint-uniqueness isomorphism:
  \[
    \mathtt{sheafificationCompPullback}\,\varphi
      \;=\;
    \mathtt{Adjunction.leftAdjointUniq}\,A\,B.
  \]
\end{lemma}

\begin{proof}



\leanok
  The two composite adjunctions \(A\) and \(B\) share the same right adjoint:
  pushing forward along \(\varphi\) and then forgetting to presheaves of modules on \(X\) is, on
  the nose, the same as forgetting to presheaves on \(Y\) and then pushing forward along
  \(\varphi'\) --- the pushforward commutes with the forgetful functor to presheaves by
  definition. Hence \(A\) and \(B\) are two adjunctions sharing the same right adjoint, so the
  left-adjoint-uniqueness isomorphism \(\mathtt{Adjunction.leftAdjointUniq}\,A\,B\) between their
  left adjoints is defined. Unfolding the canonical comparison
  \(\mathtt{sheafificationCompPullback}\,\varphi\) and unfolding
  \(\mathtt{Adjunction.leftAdjointUniq}\,A\,B\) produces the same morphism --- both sides are the
  composite isomorphism assembled from the two factor adjunctions in the same order. The equality
  therefore holds by reflexivity.
\end{proof}

\begin{lemma}\leanok
[leftAdjointUniq transport of the composite unit]
  \label{lem:leftadjointuniq_app_unit_eta}
  \lean{AlgebraicGeometry.Scheme.Modules.leftAdjointUniqUnitEta}
  \uses{lem:comp_homequiv_factor_sheafify_pullback, lem:sheafification_comp_pullback_eq_leftadjointuniq}
  Let \(A\), \(B\) be the two composite adjunctions of \texttt{eta\_bridge\_unit\_square} (with
  left adjoints \(a_X \circ \mathtt{pullback}_{\mathrm{sheaf}}\,\varphi\) and \(F \circ a_Y\)
  respectively, which agree up to the canonical iso
  \(\mathtt{sheafificationCompPullback}\,\varphi = \mathtt{Adjunction.leftAdjointUniq}\,A\,B\)),
  and let \(g = (\mathtt{sheafificationCompPullback}\,\varphi).\mathrm{hom}.\mathrm{app}\,\mathbf{1}^{\mathrm{p}}\).
  Then
  \[
    A.\mathrm{homEquiv}\,g
      \;=\;
    B.\mathrm{unit}.\mathrm{app}\,\mathbf{1}^{\mathrm{p}}
      \;=\;
    \mathtt{presheafAdj}.\mathrm{unit}.\mathrm{app}\,\mathbf{1}^{\mathrm{p}}
      \;\mathbin{;}\;
    (\mathtt{pushforward}\,\varphi').\mathrm{map}\bigl(\mathtt{toSheafify}_Y(F.\mathrm{obj}\,\mathbf{1}^{\mathrm{p}})\bigr).
  \]
\end{lemma}

\begin{proof}
  \uses{lem:comp_homequiv_factor_sheafify_pullback, lem:sheafification_comp_pullback_eq_leftadjointuniq}
  \leanok
  The comparison \(\mathtt{sheafificationCompPullback}\,\varphi\) is by definition
  \(\mathtt{Adjunction.leftAdjointUniq}\,A\,B\)
  (\cref{lem:sheafification_comp_pullback_eq_leftadjointuniq}), the canonical isomorphism between two left
  adjoints of the same functor. Mathlib's
  \(\mathtt{Adjunction.homEquiv\_leftAdjointUniq\_hom\_app}\,A\,B\) states that applying
  \(A.\mathrm{homEquiv}\) to the component of this comparison at an object \(c\) returns
  \(B.\mathrm{unit}.\mathrm{app}\,c\); taking \(c = \mathbf{1}^{\mathrm{p}}\) gives the first
  equality. For the second, expand \(B.\mathrm{unit}\) by
  \(\mathtt{Adjunction.comp\_unit\_app}\) on \(B = (\mathtt{pullbackPushforwardAdjunction}\,
  \varphi').\mathrm{comp}\,(\mathtt{sheafificationAdjunction}\,(\mathbf{1}_Y))\): the unit of a
  composite adjunction at \(\mathbf{1}^{\mathrm{p}}\) is the presheaf unit
  \(\mathtt{presheafAdj}.\mathrm{unit}.\mathrm{app}\,\mathbf{1}^{\mathrm{p}}\) followed by the
  pushforward of the sheafification unit \(\mathtt{toSheafify}_Y\) at
  \(F.\mathrm{obj}\,\mathbf{1}^{\mathrm{p}}\), which is the displayed right-hand side.
\end{proof}

\begin{lemma}\leanok
[leftAdjointUniq transport of the composite unit ($P$-general form)]
  \label{lem:leftadjointuniq_app_unit_eta_general}
  \lean{AlgebraicGeometry.Scheme.Modules.leftAdjointUniqUnitEta_app}
  \uses{lem:leftadjointuniq_app_unit_eta, lem:sheafification_comp_pullback_eq_leftadjointuniq}
  This is the \(P\)-general twin of \cref{lem:leftadjointuniq_app_unit_eta}: the latter is the
  specialisation to the monoidal unit \(P = \mathbf{1}^{\mathrm{p}}\), whereas the present form
  holds for an \emph{arbitrary} presheaf \(P\) over \(X\). Fix a morphism of schemes
  \(\varphi\) (with structure-ring map \(\varphi'\)), and let
  \[
    A_{\varphi}
      \;=\;
    (\mathtt{PresheafPullbackPushforwardAdj}\,\varphi').\mathrm{comp}\,\mathtt{sheafAdj}
  \]
  be the two-layer composite adjunction (presheaf pullback--pushforward composed with the
  sheafification adjunction), and let \(B_{\varphi}\) denote its sheaf-level counterpart (the
  composite adjunction whose left adjoint is the sheaf-level pullback). Writing
  \(\mathtt{sheafCompPb} = \mathtt{SheafOfModules.sheafificationCompPullback}\) for the
  canonical comparison, the lemma asserts
  \[
    A_{\varphi}.\mathrm{homEquiv}\,P\,\_\,\bigl((\mathtt{sheafCompPb}\,\varphi).\mathrm{hom}.\mathrm{app}\,P\bigr)
      \;=\;
    B_{\varphi}.\mathrm{unit}.\mathrm{app}\,P.
  \]
  It is the key \textbf{step-1} brick of the tail assembly of
  \cref{lem:pullback_tensor_map_basechange}: applying \(A_{\varphi}.\mathrm{homEquiv}\) to the
  comparison component recovers the composite-adjunction unit at \(P\). Its proof is identical
  to that of \cref{lem:leftadjointuniq_app_unit_eta} with \(\mathbf{1}^{\mathrm{p}}\) replaced
  by \(P\): \(\mathtt{sheafCompPb}\,\varphi\) is by definition
  \(\mathtt{Adjunction.leftAdjointUniq}\,A_{\varphi}\,B_{\varphi}\)
  (\cref{lem:sheafification_comp_pullback_eq_leftadjointuniq}), so Mathlib's
  \(\mathtt{Adjunction.homEquiv\_leftAdjointUniq\_hom\_app}\) at \(P\) returns
  \(B_{\varphi}.\mathrm{unit}.\mathrm{app}\,P\).
\end{lemma}

\begin{lemma}


\leanok
[Sheafification kills the \(\mathtt{pullbackComp}\) hom--inv pair (D3$'$ residual brick)]
  \label{lem:sheafify_pullbackcomp_hom_inv_cancel}
  \lean{AlgebraicGeometry.Scheme.Modules.sheafifyMap_pullbackComp_hom_inv_id}
  Let \(h : Z \to Y\) and \(f : Y \to X\) be composable, write
  \(\mathtt{PrPbComp} = \mathtt{PresheafOfModules.pullbackComp}\,\varphi'_f\,\varphi'_h\) for the
  presheaf pullback pseudofunctoriality isomorphism
  \(\mathtt{pullback}\,\varphi'_f \circ \mathtt{pullback}\,\varphi'_h \cong
  \mathtt{pullback}\,\varphi'_{h\circ f}\), and let \(a_Z\) be sheafification of presheaves of
  modules on \(Z\). Then for every presheaf \(T\) over \(X\) the sheafified hom--inv pair cancels:
  \[
    a_Z.\mathrm{map}\bigl(\mathtt{PrPbComp}.\mathrm{hom}.\mathrm{app}\,T\bigr) \;\mathbin{;}\;
    a_Z.\mathrm{map}\bigl(\mathtt{PrPbComp}.\mathrm{inv}.\mathrm{app}\,T\bigr) \;=\;
    \mathrm{id}_{a_Z.\mathrm{map}\,(\cdots)}.
  \]
\end{lemma}
\begin{proof}


\leanok
  The two factors are \(a_Z.\mathrm{map}\) of the components at \(T\) of the hom--inv pair of the
  isomorphism \(\mathtt{PrPbComp}\). By \(\mathtt{Iso.hom\_inv\_id\_app}\) the underlying composite
  \(\mathtt{PrPbComp}.\mathrm{hom}.\mathrm{app}\,T \mathbin{;} \mathtt{PrPbComp}.\mathrm{inv}.\mathrm{app}\,T\)
  is the identity; applying \(a_Z.\mathrm{map}\) (\(\mathtt{congrArg}\), then \(\mathtt{Functor.map\_comp}\)
  to split and \(\mathtt{Functor.map\_id}\) to collapse the identity) yields the claim. The
  cancellation is built as this \emph{standalone} equation and then spliced into the D3\(^\prime\)
  paste by \(\mathtt{erw}\); a plain \(\mathtt{rw}\) does not fire across the parenthesised group
  boundary because of the \(\mathtt{Sheaf.val}\,Z\) carrier-spelling mismatch.
\end{proof}

\begin{lemma}
[Sq3 --- composition coherence of \(\mathtt{sheafifyTensorUnitIso}\) (D3$'$)]
  \label{lem:sheafify_tensor_unit_iso_comp}
  % NOTE: no \lean{} pin — exposition sub-lemma, not a standalone build target.
  % Its mathematical content is realised in-place inside the proof of
  % \cref{lem:pullback_tensor_map_basechange}; kept (with \label) only
  % so the final-assembly prose \cref's resolve.
  \uses{def:sheafify_tensor_unit_iso, lem:sheafify_tensor_unit_iso_hom_eq_prime, lem:toringcatsheafhom_comp_hom_reconcile}
  With notation as in \cref{lem:pullback_tensor_map_basechange}, write \(F^{\mathrm p}_\bullet =
  \mathtt{pullback}\,\varphi'_\bullet\) for the presheaf-level pullback and \(a_\bullet\) for
  sheafification. The tensor/unit reconciliation \(\mathtt{sheafifyTensorUnitIso}\) of the
  \emph{composite} \(h \circ f\), evaluated on \(F^{\mathrm p}_{h\circ f}M.\mathrm{val}\) and
  \(F^{\mathrm p}_{h\circ f}N.\mathrm{val}\), agrees --- after transport through the presheaf
  pullback pseudofunctoriality identification
  \(\mathtt{PrPbComp} = \mathtt{PresheafOfModules.pullbackComp}\,\varphi'_f\,\varphi'_h\)
  (\cref{lem:toringcatsheafhom_comp_hom_reconcile} licenses the composite ring map) --- with the
  \(f\)-instance carried by \(h^{*}\) followed by the \(h\)-instance on the pulled-back arguments.
  Concretely, the \(\mathtt{hom}\) legs of the two presentations are equal as morphisms in
  \(\Scheme.\mathtt{Modules}\,Z\).
\end{lemma}
\begin{proof}
  \uses{def:sheafify_tensor_unit_iso, lem:sheafify_tensor_unit_iso_hom_eq_prime, lem:toringcatsheafhom_comp_hom_reconcile}
  By \cref{lem:sheafify_tensor_unit_iso_hom_eq_prime} each \(\mathtt{sheafifyTensorUnitIso}.\mathrm{hom}\)
  is a single \(a.\mathrm{map}\) of a \(\mathtt{tensorHom}\) of sheafification-unit components
  \(\eta_P \otimes \eta_Q\), kept inside one monoidal structure on \(\mathtt{PresheafOfModules}\).
  The composite-vs-interleaved discrepancy is therefore the image, under \(a_Z.\mathrm{map}\), of the
  naturality of the sheafification unit \(\eta\) against the presheaf identification \(\mathtt{PrPbComp}\),
  one tensor factor at a time; bifunctoriality of \(\otimes\) (the interchange law) recombines the two
  single-factor squares into the asserted equality.
\end{proof}

\begin{lemma}


\leanok
[Sq4 --- per-leg presheaf composition coherence of \(\mathtt{pullbackValIso}\) (D3$'$)]
  \label{lem:pullback_val_iso_comp}
  \lean{AlgebraicGeometry.Scheme.Modules.pullbackValIso_comp_leg}
  % NOTE: the Lean target `pullbackValIso_comp_leg` is the presheaf-level PER-LEG identity
  % displayed below (a single application; the M- and N-legs are structurally identical),
  % consumed by `pullbackTensorMap_restrict`. Its two reassembly steps are realised as in-proof
  % `have`s, recorded as \cref{lem:pullback_val_iso_comp_scpb} and
  % \cref{lem:pullback_val_iso_comp_counit}.
  \uses{def:pullback_val_iso, lem:sheafificationcomppullback_comp, lem:pullback_val_iso_comp_scpb, lem:pullback_val_iso_comp_counit}
  Let \(h : Z \to Y\), \(f : Y \to X\) be composable and \(W \in \Scheme.\mathtt{Modules}\,X\).
  Write \(F^{\mathrm p}_\bullet = \mathtt{pullback}\,\varphi'_\bullet\) for the presheaf-level
  pullback, \(\eta^\bullet\) for the sheafification unit on \(\bullet\), \(\mathtt{forget}\) for
  the forgetful functor \(\mathtt{SheafOfModules} \to \mathtt{PresheafOfModules}\),
  \(\mathtt{pVI} = \mathtt{pullbackValIso}\) (\cref{def:pullback_val_iso}),
  \(\mathtt{PrPbComp} = \mathtt{PresheafOfModules.pullbackComp}\,\varphi'_f\,\varphi'_h\), and
  \(\mathtt{pullbackComp}\,h\,f\) for the scheme-level pullback pseudofunctoriality isomorphism
  \((h\circ f)^{*} \cong h^{*}\circ f^{*}\), with \(f^{*}W = (\mathtt{pullback}\,f).\mathrm{obj}\,W\).
  Then the following identity of morphisms in
  \(\mathtt{PresheafOfModules}\,Z.\mathtt{ringCatSheaf}.\mathrm{obj}\) holds:
  \[
    \mathtt{PrPbComp}.\mathrm{hom}.\mathrm{app}\,W.\mathrm{val} \;\mathbin{;}\;
      \eta^Z_{\,F^{\mathrm p}_{h\circ f}W.\mathrm{val}} \;\mathbin{;}\;
      \mathtt{forget}\,(\mathtt{pVI}\,(h\circ f)\,W).\mathrm{hom}
  \]
  \[
    \;=\;
    F^{\mathrm p}_h\!\bigl(\eta^Y_{\,F^{\mathrm p}_f W.\mathrm{val}} \;\mathbin{;}\;
        \mathtt{forget}\,(\mathtt{pVI}\,f\,W).\mathrm{hom}\bigr) \;\mathbin{;}\;
      \eta^Z_{\,F^{\mathrm p}_h (f^{*}W).\mathrm{val}} \;\mathbin{;}\;
      \mathtt{forget}\,(\mathtt{pVI}\,h\,(f^{*}W)).\mathrm{hom} \;\mathbin{;}\;
      \mathtt{forget}\,((\mathtt{pullbackComp}\,h\,f).\mathrm{app}\,W).\mathrm{hom}.
  \]
  This is the per-leg identity actually consumed --- once for \(M\), once for \(N\) --- by the
  final \(\mathtt{tensorObjIsoOfIso}\bigl((\mathtt{pullbackComp}\,h\,f)_M,
  (\mathtt{pullbackComp}\,h\,f)_N\bigr)\) step of \cref{lem:pullback_tensor_map_basechange}; it is
  the presheaf-level form of the composition coherence of the connecting comparison
  \(\mathtt{pullbackValIso}\), relating the \((h\circ f)\)-comparison to the \(f\)- and
  \(h\)-comparisons conjugated by \(\mathtt{pullbackComp}\,h\,f\).
\end{lemma}
\begin{proof}
  \uses{def:pullback_val_iso, lem:sheafificationcomppullback_comp, lem:pullback_val_iso_comp_scpb, lem:pullback_val_iso_comp_counit}
  \leanok
  Expand the connecting comparison on each of the three legs by its definition
  (\cref{def:pullback_val_iso}): for every \(g\) and \(V\),
  \[
    (\mathtt{pVI}\,g\,V).\mathrm{hom} \;=\;
      \bigl((\mathtt{sheafificationCompPullback}\,\varphi_g).\mathrm{app}\,V.\mathrm{val}\bigr)^{-1}
        \;\mathbin{;}\; g^{*}(\varepsilon_V),
  \]
  the inverse of the sheafify--pullback comparison \(\mathtt{sheafCompPb}\) at \(V.\mathrm{val}\)
  followed by the pullback of the sheafification counit \(\varepsilon_V\). Applying
  \(\mathtt{forget}\) and distributing it over the composite splits each
  \(\mathtt{forget}\,(\mathtt{pVI}\,g\,V).\mathrm{hom}\) into a \(\mathtt{sheafCompPb}^{-1}\) factor
  and a pulled-back-counit factor; after this expansion both sides of the asserted identity are
  composites whose factors are exactly these two kinds, threaded by the sheafification units
  \(\eta^\bullet\) and the presheaf comparison \(\mathtt{PrPbComp}\).

  The two families of factors reassemble by the two sub-coherences. The
  \(\mathtt{sheafCompPb}^{-1}\) factors recombine by \cref{lem:pullback_val_iso_comp_scpb}, the
  inverse of the Sq1 coherence \cref{lem:sheafificationcomppullback_comp}. The pulled-back-counit
  factors recombine by \cref{lem:pullback_val_iso_comp_counit}, the naturality
  (pseudofunctoriality) of the sheafification counit across \(\mathtt{pullbackComp}\,h\,f\).
  Substituting both reassemblies and reassociating the intervening \(\eta^\bullet\)-units and the
  \(\mathtt{PrPbComp}\) factor turns the left-hand composite into the right-hand one.

  The two reassemblies are not independent: the \(f\)-leg counit \(f^{*}(\varepsilon_W)\) and the
  \(h\)-leg counit \(\varepsilon_{f^{*}W}\) are bridged by the \(f\)-instance of
  \(\mathtt{sheafCompPb}\), so the counit reassembly of \cref{lem:pullback_val_iso_comp_counit}
  interleaves with the \(\mathtt{sheafCompPb}^{-1}\) reassembly of
  \cref{lem:pullback_val_iso_comp_scpb} through that mediating comparison.
\end{proof}

\begin{lemma}

\leanok
[Sq4a --- the \(\mathtt{sheafCompPb}^{-1}\) reassembly (inverse of Sq1)]
  \label{lem:pullback_val_iso_comp_scpb}
  \lean{AlgebraicGeometry.Scheme.Modules.sheafificationCompPullback_comp_inv}
  \uses{lem:sheafificationcomppullback_comp}
  Let \(h : Z \to Y\), \(f : Y \to X\) be composable and \(P\) a presheaf over \(X\); write
  \(a_X, a_Z\) for the sheafifications, \(F^{\mathrm p}_\bullet = \mathtt{pullback}\,\varphi'_\bullet\),
  \(\mathtt{sheafCompPb} = \mathtt{SheafOfModules.sheafificationCompPullback}\), and
  \(\mathtt{PrPbComp} = \mathtt{PresheafOfModules.pullbackComp}\,\varphi'_f\,\varphi'_h\). Then the
  inverse of the composite sheafify--pullback comparison factors as
  \[
    \bigl((\mathtt{sheafCompPb}\,(h\circ f)).\mathrm{app}\,P\bigr)^{-1}
      \;=\;
      a_Z.\mathrm{map}\bigl(\mathtt{PrPbComp}.\mathrm{inv}_P\bigr) \;\mathbin{;}\;
      \bigl((\mathtt{sheafCompPb}\,h).\mathrm{app}\,(F^{\mathrm p}_f P)\bigr)^{-1} \;\mathbin{;}\;
      h^{*}\!\bigl(((\mathtt{sheafCompPb}\,f).\mathrm{app}\,P)^{-1}\bigr) \;\mathbin{;}\;
      (\mathtt{pullbackComp}\,h\,f).\mathrm{hom}_{a_X P}.
  \]
\end{lemma}
\begin{proof}
  \uses{lem:sheafificationcomppullback_comp}
  \leanok
  The Sq1 coherence \cref{lem:sheafificationcomppullback_comp} expresses the \(\mathrm{hom}\) leg
  of the composite comparison as the ordered composite
  \[
    \bigl((\mathtt{sheafCompPb}\,(h\circ f)).\mathrm{app}\,P\bigr).\mathrm{hom}
      \;=\;
      (\mathtt{pullbackComp}\,h\,f).\mathrm{inv}_{a_X P} \;\mathbin{;}\;
      h^{*}\!\bigl(((\mathtt{sheafCompPb}\,f).\mathrm{app}\,P).\mathrm{hom}\bigr) \;\mathbin{;}\;
      \bigl((\mathtt{sheafCompPb}\,h).\mathrm{app}\,(F^{\mathrm p}_f P)\bigr).\mathrm{hom} \;\mathbin{;}\;
      a_Z.\mathrm{map}\bigl(\mathtt{PrPbComp}.\mathrm{hom}_P\bigr).
  \]
  Taking inverses of both sides reverses the order of the four factors and inverts each. The
  inverse of \(a_Z.\mathrm{map}(\mathtt{PrPbComp}.\mathrm{hom}_P)\) is
  \(a_Z.\mathrm{map}(\mathtt{PrPbComp}.\mathrm{inv}_P)\), since functoriality of \(a_Z\) sends the
  inverse of an isomorphism to the inverse of its image; likewise the inverse of
  \(h^{*}(\,\cdots.\mathrm{hom})\) is \(h^{*}(\,\cdots.\mathrm{inv})\) by functoriality of
  \(h^{*}\); and \((\mathtt{pullbackComp}\,h\,f).\mathrm{inv}^{-1}
  = (\mathtt{pullbackComp}\,h\,f).\mathrm{hom}\). Substituting these gives the displayed
  factorisation.
\end{proof}

\begin{lemma}
[Sq4b --- the counit reassembly (counit pseudofunctoriality)]
  \label{lem:pullback_val_iso_comp_counit}
  \uses{def:pullback_val_iso}
  Let \(h : Z \to Y\), \(f : Y \to X\) be composable and \(W \in \Scheme.\mathtt{Modules}\,X\);
  write \(\varepsilon\) for the sheafification counit (\cref{def:pullback_val_iso}) and
  \(\mathtt{pullbackComp}\,h\,f : (h\circ f)^{*} \cong h^{*}\circ f^{*}\) for the scheme-level
  pullback pseudofunctoriality isomorphism. The pulled-back counit of the composite reassembles
  through \(\mathtt{pullbackComp}\,h\,f\): naturality of \(\mathtt{pullbackComp}\,h\,f\) at the
  counit morphism \(\varepsilon_W : a_X(W.\mathrm{val}) \to W\) gives
  \[
    (h\circ f)^{*}(\varepsilon_W) \;\mathbin{;}\; (\mathtt{pullbackComp}\,h\,f).\mathrm{hom}_W
      \;=\;
    (\mathtt{pullbackComp}\,h\,f).\mathrm{hom}_{a_X(W.\mathrm{val})} \;\mathbin{;}\;
      h^{*}\!\bigl(f^{*}(\varepsilon_W)\bigr).
  \]
  This is the counit-pseudofunctoriality identity that, in the assembly of
  \cref{lem:pullback_val_iso_comp}, moves the composite counit into the \(h\)-image of the
  \(f\)-counit, where it meets the \(h\)-leg counit \(\varepsilon_{f^{*}W}\).
\end{lemma}
\begin{proof}
  \uses{def:pullback_val_iso}
  The displayed equality is the naturality square of the natural isomorphism
  \(\mathtt{pullbackComp}\,h\,f : (h\circ f)^{*} \Rightarrow h^{*}\circ f^{*}\) evaluated at the
  morphism \(\varepsilon_W : a_X(W.\mathrm{val}) \to W\): for a natural transformation, the domain
  image followed by the codomain component equals the source component followed by the codomain
  image, and \((h^{*}\circ f^{*})(\varepsilon_W) = h^{*}(f^{*}(\varepsilon_W))\) by functoriality
  of the composite. No further input is needed; the counit \(\varepsilon\) enters only as the
  morphism at which the natural isomorphism is evaluated.
\end{proof}

% NOTE: Composition coherence of the pullback tensorator \(\delta\) for composable
% morphisms \(h : Z \to Y\), \(f : Y \to X\). The open-immersion base-change square
% of \cref{lem:pullback_tensor_iso_loctriv} is the specialisation \(h := j'\).
% \(\mathtt{pullbackTensorMap}\) is not an adjunction transpose (it is a four-fold
% composite), so the mate-calculus argument for the unit coherence does not apply
% directly; the genuine proof pastes four squares directly.
\begin{lemma}


\leanok
[Composition coherence of the pullback tensorator $\delta$ (D3$'$)]
  \label{lem:pullback_tensor_map_basechange}
  \lean{AlgebraicGeometry.Scheme.Modules.pullbackTensorMap_restrict}
  \uses{lem:pullback_tensor_map, lem:pullback_tensor_map_natural, lem:presheaf_pullback_oplaxmonoidal, def:tensorobjisoofiso, lem:toringcatsheafhom_comp_hom_reconcile, lem:pullbackcomp_delta, lem:sheafificationcomppullback_comp, def:pullback_val_iso, def:sheafify_tensor_unit_iso, lem:sheafify_pullbackcomp_hom_inv_cancel}
  Let \(h : Z \to Y\) and \(f : Y \to X\) be composable morphisms of schemes and let
  \(M, N \in \Scheme.\mathtt{Modules}\,X\). Write \(\delta_{\mathrm{sheaf}} =
  \mathtt{pullbackTensorMap}\) for the sheaf-level pullback--tensor comparison of
  \cref{lem:pullback_tensor_map}, \(h^{*}(-) = (\mathtt{pullback}\,h).\mathrm{map}\) for
  the pullback functor, and \(\mathtt{pullbackComp}\,h\,f\) for the pullback
  pseudofunctoriality isomorphism \((h \circ f)^{*} \cong h^{*}\circ f^{*}\). Then the
  comparison for the composite \(h \circ f\) factors through the comparisons for \(f\)
  and for \(h\), conjugated by \(\mathtt{pullbackComp}\,h\,f\):
  \[
    \delta_{\mathrm{sheaf}}^{\,h \circ f}(M, N)
      \;=\;
      (\mathtt{pullbackComp}\,h\,f).\mathrm{inv}_{M \otimes_X N} \;\mathbin{;}\;
      h^{*}\!\bigl(\delta_{\mathrm{sheaf}}^{\,f}(M, N)\bigr) \;\mathbin{;}\;
      \delta_{\mathrm{sheaf}}^{\,h}(f^{*}M,\, f^{*}N) \;\mathbin{;}\;
      \mathtt{tensorObjIsoOfIso}\bigl((\mathtt{pullbackComp}\,h\,f)_M,\,
        (\mathtt{pullbackComp}\,h\,f)_N\bigr).\mathrm{hom},
  \]
  an equality of morphisms \((h \circ f)^{*}(M \otimes_X N) \to (h^{*}f^{*}M)
  \otimes_Z (h^{*}f^{*}N)\) in \(\Scheme.\mathtt{Modules}\,Z\), where
  \(\mathtt{tensorObjIsoOfIso}\) is the functoriality of \(\otimes\) in a pair of
  isomorphisms.

  \emph{Remark (the open-immersion base-change square).} The base-change form consumed
  downstream by \cref{lem:pullback_tensor_iso_loctriv} is the \emph{specialisation} of
  this coherence. Given \(f : Y \to X\), an open \(U \subseteq X\) with preimage
  \(f^{-1}U\), the restriction \(g : f^{-1}U \to U\) and the open immersions
  \(j : U \hookrightarrow X\), \(j' : f^{-1}U \hookrightarrow Y\), the morphism
  \(f^{-1}U \to X\) admits the two equal factorisations \(j' \circ f = g \circ j\)
  (i.e. \(j' \mathbin{;} f = g \mathbin{;} j\)). Applying the displayed coherence to
  each factorisation and equating the two results identifies the restriction
  \((j')^{*}\bigl(\delta^{f}_{\mathrm{sheaf}}(M,N)\bigr)\) with
  \(\delta^{g}_{\mathrm{sheaf}}(j^{*}M, j^{*}N)\), through the pseudofunctoriality
  isomorphisms; in particular each is an isomorphism if and only if the other is. This
  is the form used in the chart-chase, where \(M, N\) are trivialised on \(U\).
\end{lemma}

\begin{proof}
  \uses{lem:pullback_tensor_map, lem:pullback_tensor_map_natural, lem:presheaf_pullback_oplaxmonoidal, lem:pullbackObjUnitToUnit_comp, def:tensorobjisoofiso, lem:toringcatsheafhom_comp_hom_reconcile, lem:pullbackcomp_delta, lem:sheafifymap_deltacomp_split, lem:sheafificationcomppullback_comp, def:pullback_val_iso, def:sheafify_tensor_unit_iso, lem:sheafify_pullbackcomp_hom_inv_cancel, lem:sheafify_tensor_unit_iso_hom_eq_prime, lem:pullback_val_iso_comp}
  \leanok
  Unfolding \(\delta_{\mathrm{sheaf}} = \mathtt{pullbackTensorMap}\)
  (\cref{lem:pullback_tensor_map}) on both sides exposes each as the same four-fold
  composite
  \[
    S_1 \;\mathbin{;}\; a.\mathrm{map}\,\delta \;\mathbin{;}\; S_3 \;\mathbin{;}\; S_4,
  \]
  where \(S_1\) is the sheafification--pullback comparison
  \(\mathtt{sheafificationCompPullback}\), \(a.\mathrm{map}\,\delta\) is the
  sheafification of the presheaf-level oplax comparison \(\delta\) of
  \cref{lem:presheaf_pullback_oplaxmonoidal}, \(S_3\) is the monoidal-unit
  sheafification isomorphism \(\mathtt{sheafifyTensorUnitIso}\), and \(S_4\) is the
  tensor of the connecting isomorphisms \(\mathtt{pullbackValIso}\). The identity to
  prove is therefore a paste of \emph{four commuting squares}, conjugated throughout by
  the pullback pseudofunctoriality isomorphism \(\mathtt{pullbackComp}\,h\,f\). Unlike
  the \emph{naturality} paste of \cref{lem:pullback_tensor_map_natural} (D1$'$), this is
  a \emph{composition}-coherence paste: the left-hand side is the composite-morphism
  (\(h \circ f\)) instance, while the right-hand side interleaves the \(f\)-instance
  (transported by \(h^{*}\)) with the \(h\)-instance (on the pulled-back modules
  \(f^{*}M, f^{*}N\)).

  We record \emph{why} the unit-comparison mirror does not apply. The unit coherence
  \cref{lem:pullbackObjUnitToUnit_comp} succeeds because
  \(\mathtt{pullbackObjUnitToUnit}\) is, by definition, an adjunction transpose, so its
  composition coherence is obtained by transposing across the pullback--pushforward
  adjunction. The tensorator \(\mathtt{pullbackTensorMap}\) is \emph{not} a transpose ---
  it is the four-fold composite above --- and there is no analogous
  \(\mathtt{homEquiv}\)-bridge, so the transpose-and-inject opening of the unit analog
  leaves an un-evaluable transpose of a concrete composite. The proof instead
  pastes the four squares directly.

  \emph{Sq2 (the \(\delta\) core).} The presheaf-level oplax comparison \(\delta\) of the
  \emph{composite} pullback decomposes by the Mathlib oplax-monoidal coherence
  \(\mathtt{Functor.OplaxMonoidal.comp\_}\delta\), once the composite presheaf pullback
  \(\mathtt{pullback}\,\varphi'_{h \circ f}\) is identified with the composite functor
  \(\mathtt{pullback}\,\varphi'_{f} \,\circ\, \mathtt{pullback}\,\varphi'_{h}\) through the
  Mathlib presheaf coherence \(\mathtt{PresheafOfModules.pullbackComp}\). This
  identification rests on the ring-map reconciliation
  \[
    (\mathtt{toRingCatSheafHom}\,(h \circ f)).\mathrm{hom}
      \;=\; \varphi'_{f} \;\mathbin{;}\;
      (\mathtt{Opens.map}\,f.\mathrm{base})^{\mathrm{op}}.\mathtt{whiskerLeft}\,\varphi'_{h},
  \]
  which is \emph{definitional}: although the two sides are presented in functor categories
  that nominally differ by the equality \(\mathtt{Opens.map}\,(h.\mathrm{base} \circ
  f.\mathrm{base}) = \mathtt{Opens.map}\,f.\mathrm{base} \,\circ\, \mathtt{Opens.map}\,
  h.\mathrm{base}\) (\(\mathtt{Opens.map\_comp}\)), this equality and the underlying
  scheme-composition data already hold up to definitional equality at default
  transparency, so the reconciliation is definitional
  (\(\mathtt{toRingCatSheafHom\_comp\_hom\_reconcile}\)). The genuine content of Sq2 is
  therefore not the reconciliation but the monoidality of \(\mathtt{pullbackComp}\) itself,
  isolated as Sq2b below.

  \emph{Sq2b (monoidality of \(\mathtt{pullbackComp}\) --- the genuinely new ingredient).}
  The decomposition in Sq2 reduces the \(\delta\)-core to a single structural fact, absent
  from Mathlib: the connecting isomorphism
  \[
    \mathtt{pullbackComp}\,\varphi'_{f}\,\varphi'_{h} \;:\;
      \mathtt{pullback}\,\varphi'_{f} \,\circ\, \mathtt{pullback}\,\varphi'_{h}
      \;\cong\; \mathtt{pullback}\,
        \bigl(\varphi'_{f} \;\mathbin{;}\;
          (\mathtt{Opens.map}\,f.\mathrm{base})^{\mathrm{op}}.\mathtt{whiskerLeft}\,
          \varphi'_{h}\bigr)
  \]
  is a \emph{monoidal} natural isomorphism between the oplax-monoidal functors
  \(\bigl(\mathtt{pullback}\,\varphi'_{f} \circ \mathtt{pullback}\,\varphi'_{h},\,
  \mathtt{comp\_}\delta\bigr)\) and \(\bigl(\mathtt{pullback}\,(\varphi'_{f}\mathbin{;}
  \cdots),\,\delta\bigr)\). Writing \(\delta\) for the oplax tensorator of the single
  composite pullback and \(\mathtt{comp\_}\delta = (\mathtt{pullback}\,\varphi'_{h})
  .\mathrm{map}\,\delta^{\varphi'_{f}} \mathbin{;} \delta^{\varphi'_{h}}\) for the
  composite tensorator of the two-step pullback, monoidality is the identity
  \[
    \delta_{M,N}
      \;=\; (\mathtt{pullbackComp}\,\varphi'_{f}\,\varphi'_{h}).\mathrm{inv}_{M\otimes N}
      \;\mathbin{;}\; (\mathtt{comp\_}\delta)_{M,N} \;\mathbin{;}\;
      \bigl((\mathtt{pullbackComp}\,\varphi'_{f}\,\varphi'_{h}).\mathrm{hom}_{M}
        \otimes (\mathtt{pullbackComp}\,\varphi'_{f}\,\varphi'_{h}).\mathrm{hom}_{N}\bigr).
  \]
  This is precisely the \(\delta\)-twin of the unit coherence
  \cref{lem:pullbackObjUnitToUnit_comp}, and it is proved by the \emph{same} mate-calculus
  argument with the unit comparison \(\eta\) replaced by the tensorator \(\delta\). The
  pivot is that, for the canonical oplax structure on a left adjoint
  (\(\mathtt{leftAdjointOplaxMonoidal}\), here
  \cref{lem:presheaf_pullback_oplaxmonoidal}), \(\delta\) is \emph{itself} an adjunction
  transpose: \(\delta_{M,N} = \mathtt{homEquiv}^{-1}\bigl((\eta_{M}\otimes\eta_{N})
  \mathbin{;} \mu_{M,N}\bigr)\), where \(\mu\) is the lax tensorator of the
  \(\mathtt{pushforward}\). Consequently the transpose-and-inject opening that closes the
  unit coherence ports verbatim: applying \(\mathtt{homEquiv}.\mathrm{injective}\) for the
  pullback--pushforward adjunction and transporting \(\mathtt{pullbackComp}\) to
  \(\mathtt{pushforwardComp}\) across the adjunctions via
  \(\mathtt{conjugateEquiv\_pullbackComp\_inv}\) (with \(\mathtt{unit\_conjugateEquiv}\)
  and \(\mathtt{Adjunction.comp\_unit\_app}\)) reduces the goal to the lax-\(\mu\)
  composition coherence of \(\mathtt{PresheafOfModules.pushforward}\) across
  \(\mathtt{pushforwardComp}\). This last residual --- isolated as the lemma
  \(\mathtt{pushforwardComp\_lax\_}\mu\) --- is the \emph{right-adjoint twin} of the
  \(\delta\)-coherence: it is precisely the assertion that the pseudofunctoriality
  isomorphism \(\mathtt{pushforwardComp}\) is \emph{monoidal} for the lax tensorators
  \(\mu\), i.e.\ that comparing the lax-\(\mu\) structure of the single composite
  \(\mathtt{pushforward}\) with the composite of the two individual lax-\(\mu\) structures
  agrees across \(\mathtt{pushforwardComp}\). Two points must be stated precisely.

  First, the mate-calculus reduction proceeds as follows: the transpose-and-inject
  opening, the \(\delta = \mathtt{homEquiv}^{-1}\bigl((\eta\otimes\eta)\mathbin{;}\mu\bigr)\)
  pivot, the \(\mathtt{conjugateEquiv\_pullbackComp\_inv}\) transport across the adjunctions
  (with \(\mathtt{unit\_conjugateEquiv}\) and \(\mathtt{Adjunction.comp\_unit\_app}\)), and
  the two-argument \(\mathtt{tensorHom}\)/\(\delta_{\mathrm{natural}}\) bookkeeping --- which
  follows the template of Mathlib's
  \(\mathtt{CategoryTheory.Adjunction.isMonoidal\_comp}\) (the canonical ``a composite of
  monoidal adjunctions is monoidal'') --- together reduce Sq2b to the single residual
  \(\mathtt{pushforwardComp\_lax\_}\mu\). This is the content of
  \cref{lem:pullbackObjUnitToUnit_comp}'s \(\delta\)-mirror.

  Second, this residual \(\mathtt{pushforwardComp\_lax\_}\mu\) is discharged by a short
  \emph{sectionwise pure-tensor collapse}, not by a change-of-rings calculation. On the
  strongly-monoidal pushforward objects the pushforward tensorator \(\mu\) reduces
  \emph{definitionally} to the lighter \(\mathtt{restrictScalars}\) tensorator \(\mu\) (the
  identity \(\mathtt{pushforward\_}\mu\mathtt{\_eq}\), which holds by reflexivity); after
  this reduction, evaluated on a pure tensor \(m \otimes n\) every
  \(\mathtt{restrictScalars}\) tensorator is the identity (Mathlib's
  \(\mathtt{ModuleCat.restrictScalars\_}\mu\mathtt{\_tmul}\)), so both sides of the
  interchange collapse to the same pure tensor \(m \otimes n\); no change-of-rings
  construction is required. This closes the Sq2b obligation.

  Crucially, this sub-step is stated at the \(\mathtt{PresheafOfModules}\) level --- about
  \(\mathtt{pullback}\,\varphi'\) with \(\varphi'\) the ring map
  \((\mathcal{S}_{0}\mathbin{;}\mathtt{forget}_{2}) \to f^{\mathrm{op}}\mathbin{;}
  (\mathcal{R}_{0}\mathbin{;}\mathtt{forget}_{2})\) carried by \(\mathtt{pullbackTensorMap}\)
  --- rather than at the scheme level, where the composite factorisation is immediate from
  \(\mathtt{pullbackComp}\) and the ring-map reconciliation is definitional.

  The observation above --- that \(\mathtt{pullbackTensorMap}\) is not an adjunction
  transpose --- concerns the \emph{full} four-fold composite and does not constrain
  Sq2b: the inner tensorator \(\delta\) genuinely \emph{is} a transpose
  (\(\mathtt{homEquiv}^{-1}\) of \((\eta\otimes\eta)\mathbin{;}\mu\)), so for Sq2b
  specifically the \(\eta\to\delta\) mirror of \cref{lem:pullbackObjUnitToUnit_comp} is the
  correct route.

  \emph{Sq1 (sheafification \(\leftrightarrow\) pullback) --- the sole open ingredient.} The
  composition coherence of \(\mathtt{SheafOfModules.sheafificationCompPullback}\) across
  \(h \circ f\) --- the analog of \(\mathtt{pullbackComp}\) for the
  ``sheafify-then-pullback'' natural isomorphism --- is \emph{absent from Mathlib} and is
  supplied as a named, private project sub-lemma
  \(\mathtt{sheafificationCompPullback\_comp}\). For a presheaf \(P\) over \(X\), writing
  \(a_X, a_Z\) for the sheafifications and \(\mathtt{sheafCompPb} =
  \mathtt{SheafOfModules.sheafificationCompPullback}\), it asserts
  \[
    ((\mathtt{sheafCompPb}\,(h \circ f)).\mathrm{app}\,P).\mathrm{hom}
      \;=\;
      (\mathtt{pullbackComp}\,h\,f).\mathrm{inv}_{a_X P} \;\mathbin{;}\;
      h^{*}\bigl(((\mathtt{sheafCompPb}\,f).\mathrm{app}\,P).\mathrm{hom}\bigr) \;\mathbin{;}\;
      ((\mathtt{sheafCompPb}\,h).\mathrm{app}\,
        (\mathtt{pullback}\,\varphi'_{f}\,P)).\mathrm{hom} \;\mathbin{;}\;
      a_Z.\mathrm{map}\bigl(
        (\mathtt{PresheafOfModules.pullbackComp}\,\varphi'_{f}\,\varphi'_{h}).\mathrm{hom}_P
      \bigr).
  \]
  It is proved by the \(\mathtt{leftAdjointUniq}\) mate calculus: both \(\mathtt{sheafCompPb}\)
  isomorphisms are \(\mathtt{leftAdjointUniq}\) of composite adjunctions, so transposing the
  identity under the composite adjunction for \(h \circ f\) and evaluating the left-hand side
  by the mate identity \(\mathtt{homEquiv\_leftAdjointUniq\_hom\_app}\) reduces it to a
  unit-level identity; the two \(\mathtt{pullbackComp}\) factors are then transported across
  the adjunction units --- the \(\delta\)-free twin of
  \cref{lem:pullbackObjUnitToUnit_comp}.

  \emph{The reduced tail goal.} After peeling the leading pseudofunctoriality factor
  \((\mathtt{pullbackComp}\,h\,f).\mathrm{inv}\), the remaining obligation is a unit
  identity for the two-layer composite adjunction (presheaf pullback--pushforward adjunction
  composed with the sheafification adjunction), isolated as the named sub-lemma
  \(\mathtt{sheafificationCompPullback\_comp\_tail}\). Writing \(B_{\varphi}\) for the
  sheaf-level composite adjunction of \cref{lem:leftadjointuniq_app_unit_eta_general} and
  \(\mathtt{PresheafPullback}_f\) for the presheaf-level pullback along \(f\), the assembly
  proceeds in the following ordered steps, mirroring the already-closed
  \cref{lem:pullbackObjUnitToUnit_comp} (whose proof is the \(\delta\)-free analog of this one,
  exactly one sheafification layer down):
  \begin{enumerate}
    \item[(a)] \emph{Strip the outer identity wrapper.} The tail goal carries an outer
      \(\mathtt{restrictScalars}(\mathbf{1})\) change-of-rings wrapper coming from the
      identity ring map; this restriction-of-scalars-along-identity factor is removed first,
      so that the genuine comparison factors are exposed.
    \item[(b)] \emph{Distribute the right-hand sheaf composite under \(\mathtt{forget}\).} The
      right-hand side is the sheaf-level composite carrying the two sub-comparison factors
      \(R_1\) (the \(f\)-layer) and \(R_5\) (the \(h\)-layer). Pushing the forgetful functor
      \(\mathtt{forget}\) (sheaf \(\to\) presheaf) through the composite separates these two
      factors \emph{without} disturbing the left-hand composite unit \(B_{h\circ f}.\mathrm{unit}\),
      which is left intact for the final reassembly.
    \item[(c)] \emph{Recover the sub-comparison units.} Apply
      \cref{lem:leftadjointuniq_app_unit_eta_general}
      (\(\mathtt{leftAdjointUniqUnitEta\_app}\), the \(P\)-general step-1 brick) to rewrite
      \(R_1\) as the composite-adjunction unit \(B_f.\mathrm{unit}.\mathrm{app}\,P\) and
      \(R_5\) as \(B_h.\mathrm{unit}.\mathrm{app}\,(\mathtt{PresheafPullback}_f\,P)\). This is
      precisely the uniqueness-of-left-adjoints characterisation
      \(\mathtt{homEquiv\_leftAdjointUniq\_hom\_app}\), now used in its \(P\)-general form.
    \item[(d)] \emph{Slide the \(\mathtt{pushforwardComp}\) comparison past the recovered
      units.} The sheaf-level comparison \(\mathtt{pushforwardComp}\,h\,f\) is moved across the
      two recovered units \(B_f.\mathrm{unit}\) and \(B_h.\mathrm{unit}\) by its naturality.
    \item[(e)] \emph{Reassemble the composite unit.} Expanding
      \(B_{h\circ f}.\mathrm{unit}\) by the composite-adjunction unit formula
      \(\mathtt{comp\_unit\_app}\) and applying \(\mathtt{unit\_naturality}\) identifies the
      slid expression of step~(d) with the left-hand composite unit, closing the tail.
  \end{enumerate}

  \emph{The precise binding obligation.} The one place where the assembly is not yet a
  mechanical paste is the bridge between the \emph{presheaf}-level left-hand data (the
  presheaf pushforward \(\mathtt{pushforward}^{\mathrm{pre}}\) and the composite unit
  \(B_{h\circ f}.\mathrm{unit}\)) and the \emph{sheaf}-level right-hand factors \(R_1, R_5\)
  wrapped in \(\mathtt{forget}\). Crossing this boundary requires the compatibility
  \[
    \mathtt{forget} \circ \mathtt{pushforward}^{\mathrm{sheaf}}
      \;=\;
    \mathtt{pushforward}^{\mathrm{pre}} \circ \mathtt{forget},
  \]
  interacting with the recovered \(B_f\)- and \(B_h\)-units of step~(c). This compatibility is
  the natural unit to isolate as its own named sub-lemma \emph{before} the assembly: once it
  is in hand, steps (a)--(e) become a mechanical paste, whereas without it the presheaf- and
  sheaf-level pushforwards cannot be aligned at the recovered units. The remainder of the tail
  is bookkeeping around this single compatibility.

  This is the sheafification-laden analog of \cref{lem:pullbackObjUnitToUnit_comp}; unlike
  that lemma it is not definitional sectionwise, because the adjunctions involved are
  composite rather than single transposes. The outer Sq1 lemma
  \(\mathtt{sheafificationCompPullback\_comp}\) \emph{consumes} this tail: it peels the leading
  \((\mathtt{pullbackComp}\,h\,f).\mathrm{inv}\) factor and then invokes
  \(\mathtt{sheafificationCompPullback\_comp\_tail}\) for the residual unit identity.
  \emph{Warning:} transposing the \emph{whole} tail back under the composite adjunction --- as
  opposed to recovering the two sub-units individually via step~(c) and reassembling --- is
  \emph{circular} (verified): it returns the very identity one set out to prove, so the
  step-(c) recovery of the individual \(R_1, R_5\) units is essential.

  \emph{Sq3 (unit-iso transport).} The monoidal-unit sheafification isomorphism
  \(\mathtt{sheafifyTensorUnitIso}\) is carried through the same \(\mathtt{pullbackComp}\)
  identification used for Sq2. This composition coherence is isolated as the named brick
  \cref{lem:sheafify_tensor_unit_iso_comp} (\(\mathtt{sheafifyTensorUnitIso\_comp}\)): via
  \cref{lem:sheafify_tensor_unit_iso_hom_eq_prime} the \(\mathtt{hom}\) leg is a single
  \(a.\mathrm{map}\) of \(\eta\otimes\eta\), so the coherence reduces to the naturality of the
  sheafification unit \(\eta\) against \(\mathtt{PrPbComp}\), recombined by bifunctoriality.

  \emph{Sq4 (the connecting iso).} The scheme-level composition coherence of
  \(\mathtt{pullbackValIso}\) is isolated as the named brick \cref{lem:pullback_val_iso_comp}
  (\(\mathtt{pullbackValIso\_comp}\)): it relates \(\mathtt{pullbackValIso}\,(h \circ f)\,M\) to
  \(h^{*}\bigl(\mathtt{pullbackValIso}\,f\,M\bigr)\),
  \(\mathtt{pullbackValIso}\,h\,(f^{*}M)\), and the pseudofunctoriality component
  \((\mathtt{pullbackComp}\,h\,f)_M\), and produces the final
  \(\mathtt{tensorObjIsoOfIso}\bigl((\mathtt{pullbackComp}\,h\,f)_M,
  (\mathtt{pullbackComp}\,h\,f)_N\bigr)\) factor. It is \emph{absent from Mathlib}; but
  \(\mathtt{pullbackValIso}\) factors through \(\mathtt{sheafCompPb}\) (as
  \(\mathtt{sheafCompPb}^{-1}\) followed by the pullback of a counit), so by
  \cref{lem:pullback_val_iso_comp} its composition coherence reduces to that of Sq1
  (\cref{lem:sheafificationcomppullback_comp}, now closed) together with the
  pseudofunctoriality of the counit --- i.e.\ Sq4 is a short corollary once Sq1 is in hand. It
  is the naturality across the connecting-object boundary where the abstract pullback's
  \(\mathtt{.val}\) presentation meets the helper-\(\mathtt{.obj}\) presentation. The prover
  treats it as its own brick (its own statement and connecting-object bookkeeping), not as a
  step inlined in the proof of Sq1.

  \emph{State of the paste and the residual chain.} Sq1 (\(\mathtt{sheafificationCompPullback\_comp}\))
  and Sq2b (\(\mathtt{pullbackComp\_}\delta\) with \(\mathtt{pushforwardComp\_lax\_}\mu\)) are
  \emph{closed}, and the Sq2 ring-map reconciliation is definitional. In the Lean paste these are
  consumed by the committed splice of the equation \(\mathtt{h}\delta\) (the Sq2b decomposition of
  \(a_Z.\mathrm{map}\,\delta_{h\circ f}\) into
  \(a_Z.\mathrm{map}(\mathtt{PrPbComp}.\mathrm{inv}) \mathbin{;} a_Z.\mathrm{map}\,\delta_{\mathrm{comp}}
  \mathbin{;} a_Z.\mathrm{map}\,\mathtt{tcomp}\), via \cref{lem:pullbackcomp_delta} under
  \(\mathtt{congrArg}\,a_Z.\mathrm{map}\)) by \(\mathtt{erw}\). After that splice the residual is a
  three-brick chain:
  \begin{enumerate}
    \item[(i)] \emph{The hom--inv cancellation} \cref{lem:sheafify_pullbackcomp_hom_inv_cancel}
      (\(\mathtt{sheafifyMap\_pullbackComp\_hom\_inv\_id}\)): the adjacent pair
      \(a_Z.\mathrm{map}(\mathtt{PrPbComp}.\mathrm{hom}.\mathrm{app}\,(M\otimes N)) \mathbin{;}
      a_Z.\mathrm{map}(\mathtt{PrPbComp}.\mathrm{inv}.\mathrm{app}\,(M\otimes N))\) created by the
      \(\mathtt{h}\delta\)-splice collapses to the identity. The two cancelling
      factors are adjacent in the composite only after a reassociation; the intervening bracket
      joins through a \(\mathtt{SheafOfModules}\) instance that is definitionally but not
      syntactically equal to the one in the ambient goal, so the cancellation is carried out via
      a private generic middle-cancellation lemma over an abstract category, invoked by
      definitional-equality unification. The helper carries no separate blueprint pin.
    \item[(ii)] \emph{The \(\delta_{\mathrm{comp}}\) split} via
      \(\mathtt{Functor.OplaxMonoidal.comp\_}\delta\) (mechanical): \(\delta_{\mathrm{comp}}\)
      becomes \((\mathtt{pullback}\,\varphi'_h).\mathrm{map}(\delta_f\,M.\mathrm{val}\,N.\mathrm{val})
      \mathbin{;} \delta_h\,(F^{\mathrm p}_f M.\mathrm{val})\,(F^{\mathrm p}_f N.\mathrm{val})\).
    \item[(iii)] \emph{Sq3 \cref{lem:sheafify_tensor_unit_iso_comp} and Sq4
      \cref{lem:pullback_val_iso_comp}}, discharged via the interleaved four-square merge below.
  \end{enumerate}
  The adjunction-mate bridge \(\mathtt{conjugateEquiv\_pullbackComp\_inv}\) (relating
  \(\mathtt{pullbackComp}\) for the left adjoints to \(\mathtt{pushforwardComp}\) for the
  right adjoints) is the device used inside Sq2b and Sq4 to transport the pseudofunctoriality
  coherence across the adjunctions.

  \emph{The final assembly (the interleaved four-square merge).} The four squares do \emph{not}
  paste row-by-row, and step~(iii) is \emph{not} the D1$'$ paste applied verbatim. In D1$'$
  the four ingredients are honest natural transformations whose squares close independently;
  here two distinct obstructions --- a misaligned \(S_1^{h}\) argument and the
  \(\mathtt{pullbackValIso}\) bridge between Sq3 and Sq4 --- force the last three factors into
  a single combined chase.

  \emph{The slide of \(S_1^{h}\).} After the \(\delta_{\mathrm{comp}}\) split~(ii), the
  \(h\)-tensorator \(\delta_h\) acts on the \emph{sheaf}-pullback arguments
  \(\bigl((\mathtt{pullback}\,f).\mathrm{obj}\,M\bigr) \otimes
  \bigl((\mathtt{pullback}\,f).\mathrm{obj}\,N\bigr)\), whereas \(S_1^{h}\) --- the
  \(h\)-instance of \(\mathtt{sheafificationCompPullback}\) (Sq1,
  \cref{lem:sheafificationcomppullback_comp}, closed) --- must be made to act instead on the
  \emph{presheaf}-pullback arguments \(F^{\mathrm p}_{f}M.\mathrm{val} \otimes
  F^{\mathrm p}_{f}N.\mathrm{val}\). This re-alignment is a naturality step in its own right,
  \emph{not} the D1$'$ slide of \(\delta\) past \(F(a\otimes b)\): the connecting comparison
  \(\mathtt{sheafificationCompPullback}\,h\) is a natural transformation, and its naturality
  square at the oplax-tensorator morphism
  \[
    \delta_{f}(M,N) \;:\;
      F^{\mathrm p}_{f}(M.\mathrm{val} \otimes N.\mathrm{val})
      \;\longrightarrow\;
      F^{\mathrm p}_{f}M.\mathrm{val} \otimes F^{\mathrm p}_{f}N.\mathrm{val}
  \]
  slides \(S_1^{h}\) from the sheaf-pullback arguments to the presheaf-pullback arguments.
  It is this bespoke square --- the naturality of the \emph{connecting} transformation at
  \(\delta_f\), not the oplax naturality of a comparison --- that effects the slide.

  \emph{The merged Sq3/Sq4 chase.} The two remaining bricks --- Sq3
  \cref{lem:sheafify_tensor_unit_iso_comp} (the composition coherence of
  \(\mathtt{sheafifyTensorUnitIso}\)) and Sq4 \cref{lem:pullback_val_iso_comp} (the
  composition coherence of \(\mathtt{pullbackValIso}\)) --- do \emph{not} discharge
  separately. Their arguments are bridged by the connecting isomorphism
  \(\mathtt{pullbackValIso}\,f\) (\cref{def:pullback_val_iso}): one side presents the data
  through the sheaf-pullback \(\bigl((\mathtt{pullback}\,f).\mathrm{obj}\,M\bigr).\mathrm{val}\),
  the other through the presheaf-pullback \(F^{\mathrm p}_{f}M.\mathrm{val}\), so a paste of
  Sq3 followed by Sq4 never lines up. They must be discharged \emph{together}, as one chase,
  exactly as in the merge of \cref{lem:pullback_tensor_map_natural} (D1$'$): there the two
  functor-image factors are folded into a single image by functoriality
  \(F(g)\mathbin{;}F(h)=F(g\mathbin{;}h)\), reducing the sheaf-level equality to a
  presheaf-level one. We reproduce that structure here: the composite
  \(a.\mathrm{map}\,\delta \mathbin{;} S_3 \mathbin{;} S_4\) on the left, and the matching
  tail on the right, are each folded into a single \(a.\mathrm{map}\) of one presheaf
  morphism, collapsing the sheaf-level equation to a single presheaf-level identity between
  the two folded presheaf morphisms.

  \emph{Closing the presheaf identity.} That presheaf identity is closed by three facts
  already in hand: the naturality of the oplax tensorator \(\delta\)
  (\cref{lem:presheaf_pullback_oplaxmonoidal}); the single-\(\mathtt{tensorHom}\)
  presentation \(\mathtt{sheafifyTensorUnitIso}.\mathrm{hom} = a.\mathrm{map}(\eta\otimes\eta)\)
  of \cref{lem:sheafify_tensor_unit_iso_hom_eq_prime}, which keeps every factor inside one
  monoidal instance; and the \(\mathtt{pullbackValIso}\) factorisation of
  \cref{def:pullback_val_iso} (the inverse of \(\mathtt{sheafificationCompPullback}\) followed
  by the pulled-back sheafification counit), which rewrites the connecting isomorphism into
  the same presheaf-level form on both sides. With the slide supplying the \(\delta_h\)
  alignment and the closed Sq2b decomposition supplying the \(\delta\)-core, the four squares
  assemble into the displayed identity. The base-change-square consequence (the Remark) is
  then the specialisation \(h := j'\) obtained by equating the two factorisations
  \(j' \mathbin{;} f = g \mathbin{;} j\); the iff-clause is immediate, an isomorphism being
  preserved and reflected by the pseudofunctoriality isomorphism it is conjugated by.
\end{proof}

\subsection{The K1 monoidal-mate bridge ($\mu$-side): two pure-tensor collapses}

The open-immersion comparison lemma K1 (\cref{lem:pullback_tensor_map_isiso_open_immersion})
reduces, by the same two-ingredient model as \cref{lem:tensorobj_restrict_iso}, to the
statement that the canonical left-adjoint-uniqueness isomorphism
\(H_1 = \mathtt{leftAdjointUniq} : G_\beta \xrightarrow{\sim} \mathtt{pullback}\,\varphi'\) is a
\emph{monoidal} natural isomorphism. Here \(f : Y \to X\) is the open immersion,
\(\varphi' = f.\mathtt{toRingCatSheafHom}.\mathrm{hom}\) is its presheaf-level structure map,
\(\beta\) is the sectionwise-bijective structure-ring map of the restriction (sectionwise
\(f.\mathtt{appIso}^{-1}\)), and
\(G_\beta = \mathtt{pushforward}_0^{\mathrm{Comm}}\,f.\mathtt{opensFunctor}\,X.\mathrm{presheaf}
\,\circ\, \mathtt{restrictScalars}\,\beta'\) is the resulting \emph{strong}-monoidal restriction
functor (a left adjoint of \(\mathtt{pushforward}\,\varphi'\), so \(H_1\) is defined). That \(H_1\)
is monoidal is two compatibilities --- one on the unit comparison, one on the tensorator --- and
each is an open-immersion pure-tensor collapse of the bijective structure-ring map
\(f.\mathtt{appIso}\). These are the $\mu$/$\delta$-side twins of the already-proved $\eta$/$\varepsilon$-side
template \cref{lem:presheafunit_comp_map_eta}, and the open-immersion analogues of the D3$'$
composition coherence \cref{lem:pushforwardcomp_lax_mu}; they use the same D3$'$ pure-tensor
section calculus. We isolate them as the next two lemmas; granting them, K1's
monoidal-mate obligation is discharged by the standard mate transport across \(H_1\).

\begin{lemma}


\leanok
[Presheaf-side mate identity for the $\eta$-bridge]
  \label{lem:presheafunit_comp_map_eta}
  \lean{AlgebraicGeometry.Scheme.Modules.presheafUnit_comp_map_eta}
  \uses{lem:presheaf_pushforward_laxmonoidal, lem:presheaf_pullback_oplaxmonoidal}
  For a scheme morphism \(f : Y \to X\) with \(\varphi' = f.\mathtt{toRingCatSheafHom}.\mathrm{hom}\),
  the presheaf-of-modules adjunction unit at the monoidal unit, post-composed with the pushforward
  of the oplax unit comparison \(\eta(\mathtt{pullback}\,\varphi')\), recovers the lax unit
  comparison \(\varepsilon(\mathtt{pushforward}\,\varphi')\):
  \[
    (\mathtt{pullbackPushforwardAdjunction}\,\varphi').\mathrm{unit}.\mathrm{app}\,
      \mathbf{1}^{\mathrm p}
    \;\mathbin{;}\;
    (\mathtt{pushforward}\,\varphi').\mathrm{map}\bigl(\eta(\mathtt{pullback}\,\varphi')\bigr)
    \;=\;
    \varepsilon(\mathtt{pushforward}\,\varphi').
  \]
  This is the presheaf-level driver of the D2$'$ $\eta$-bridge, and the template the K1
  $\mu$-side bridge mirrors.
\end{lemma}
\begin{proof}
  \uses{lem:presheaf_pushforward_laxmonoidal, lem:presheaf_pullback_oplaxmonoidal}
  \leanok
  This is the mate identity
  \(\mathtt{Adjunction.unit\_app\_unit\_comp\_map\_}\eta\) for the adjunction
  \(\mathtt{pullbackPushforwardAdjunction}\,\varphi'\), whose
  \(\mathtt{Adjunction.IsMonoidal}\) structure is supplied by
  \(\mathtt{presheafPushforwardLaxMonoidal}\) (\cref{lem:presheaf_pushforward_laxmonoidal}) and
  \(\mathtt{presheafPullbackOplaxMonoidal}\) (\cref{lem:presheaf_pullback_oplaxmonoidal}).
\end{proof}

\begin{lemma}


\leanok
[K1 $\eta$-side collapse: $H_1$ respects the unit comparison]
  \label{lem:pushforward_eta_appiso_collapse}
  \lean{AlgebraicGeometry.Scheme.Modules.pushforward_eta_appIso_collapse}
  \uses{lem:tensorobj_restrict_iso, lem:restrictscalars_laxmonoidal, lem:presheaf_pullback_oplaxmonoidal, lem:pushforwardpushforwardadj_unit_app_app_apply, lem:restrictscalars_oplaxmonoidal_eta_app_one}
  In the K1 setting above, the strong unit comparison \(\eta\,G_\beta\) of the restriction functor
  agrees, under \(H_1\), with the oplax unit comparison \(\eta(\mathtt{pullback}\,\varphi')\):
  \[
    \eta\,G_\beta
    \;=\;
    H_1.\mathrm{hom}.\mathrm{app}\,\mathbf{1}^{\mathrm p}
    \;\mathbin{;}\;
    \eta(\mathtt{pullback}\,\varphi').
  \]
  Equivalently, \(H_1\) carries the monoidal unit of \(G_\beta\) to that of
  \(\mathtt{pullback}\,\varphi'\).
\end{lemma}
\begin{proof}
  \uses{lem:tensorobj_restrict_iso, lem:restrictscalars_laxmonoidal, lem:pushforwardpushforwardadj_unit_app_app_apply, lem:restrictscalars_oplaxmonoidal_eta_app_one}
  \leanok
  This is the $\eta$/$\varepsilon$-side of ``\(H_1\) is a monoidal natural isomorphism'', the
  open-immersion analogue of the D2$'$ template \cref{lem:presheafunit_comp_map_eta}. It is
  verified sectionwise on an open \(U\). On the unit module the strong unit
  comparison \(\eta\,G_\beta\) of \(\mathtt{restrictScalars}\,\beta'\) is, sectionwise, the
  structure-ring identity of the \emph{bijective} \(f.\mathtt{appIso}\)
  (the strong-monoidal upgrade of base change along a ring isomorphism, as in
  \cref{lem:tensorobj_restrict_iso}); the mate \(\eta(\mathtt{pullback}\,\varphi')\) is the same
  base-change value transported across \(H_1\), so the two sides coincide. The collapse of the
  adjunction-unit composite to a sectionwise structure-ring map uses the presheaf-level value
  \cref{lem:pushforwardpushforwardadj_unit_app_app_apply}, after which both sides reduce to the
  ring-unit identity \(1 \mapsto 1\).
\end{proof}

\begin{lemma}


\leanok
[Presheaf-level value of the pushforward--pushforward adjunction unit]
  \label{lem:pushforwardpushforwardadj_unit_app_app_apply}
  \lean{AlgebraicGeometry.Scheme.Modules.pushforwardPushforwardAdj_unit_app_app_apply}
  For the pushforward--pushforward adjunction \(F \dashv G\) assembled from a morphism of
  presheaves of rings (the adjunction \(\mathtt{pushforwardPushforwardAdj}\,adj\,\varphi\,\psi\,H_1\,H_2\)),
  the unit's sectionwise carrier action is precomposition by the pullback of the adjunction
  counit: for a module \(M\), an open \(U\) and a section \(x\),
  \[
    \bigl(((\mathtt{pushforwardPushforwardAdj}\,adj\,\varphi\,\psi\,H_1\,H_2).\mathrm{unit}.\mathrm{app}\,M).\mathrm{app}\,U\bigr).\mathrm{hom}\,x
    \;=\;
    \bigl(M.\mathrm{map}\,(adj.\mathrm{counit}.\mathrm{app}\,U^{\mathrm{op}})^{\mathrm{op}}\bigr).\mathrm{hom}\,x .
  \]
\end{lemma}
\begin{proof}

\leanok
  Definitional. The constituent isomorphisms \(\mathtt{pushforwardId}\) and
  \(\mathtt{pushforwardComp}\) are \(\mathtt{Iso.refl}\), so the unit composite collapses to the
  \(\mathtt{pushforwardNatTrans}\) of the adjunction counit, whose carrier action is exactly the
  displayed precomposition; the equation holds by \(\mathtt{rfl}\) on the element-applied form.
\end{proof}

\begin{lemma}

\leanok
[Presheaf-level value of the pushforward--pushforward adjunction counit]
  \label{lem:pushforwardpushforwardadj_counit_app_app_apply}
  \lean{AlgebraicGeometry.Scheme.Modules.pushforwardPushforwardAdj_counit_app_app_apply}
  The exact dual of \cref{lem:pushforwardpushforwardadj_unit_app_app_apply}. For the same
  pushforward--pushforward adjunction \(F \dashv G\)
  (\(\mathtt{pushforwardPushforwardAdj}\,adj\,\varphi\,\psi\,H_1\,H_2\)), the \emph{counit}'s
  sectionwise carrier action is precomposition by the pullback of the adjunction \emph{unit}: for a
  module \(N\), an open \(U\) and a section \(y\),
  \[
    \bigl(((\mathtt{pushforwardPushforwardAdj}\,adj\,\varphi\,\psi\,H_1\,H_2).\mathrm{counit}.\mathrm{app}\,N).\mathrm{app}\,U\bigr).\mathrm{hom}\,y
    \;=\;
    \bigl(N.\mathrm{map}\,(adj.\mathrm{unit}.\mathrm{app}\,U^{\mathrm{op}})^{\mathrm{op}}\bigr).\mathrm{hom}\,y .
  \]
  In the K1 setting this is the structure-ring action of the \emph{bijective} \(f.\mathtt{appIso}\)
  on each tensor factor, which is what discharges the counit pair leg of the LHS mate telescope.
\end{lemma}
\begin{proof}

\leanok
  Definitional, exactly as the unit dual. The constituent isomorphisms \(\mathtt{pushforwardId}\)
  and \(\mathtt{pushforwardComp}\) are \(\mathtt{Iso.refl}\), so the counit composite collapses to the
  \(\mathtt{pushforwardNatTrans}\) of the adjunction unit, whose carrier action is the displayed
  precomposition; the equation holds by \(\mathtt{rfl}\) on the element-applied form.
\end{proof}

\begin{lemma}


\leanok
[Unit-preservation of the strong-monoidal $\mathtt{restrictScalars}$ oplax unit]
  \label{lem:restrictscalars_oplaxmonoidal_eta_app_one}
  \lean{AlgebraicGeometry.Scheme.Modules.restrictScalars_oplaxMonoidal_η_app_one}
  \uses{lem:restrictscalars_laxmonoidal}
  Let \(\alpha : R \to S\) be a morphism of presheaves of commutative rings whose every sectionwise
  carrier \(\alpha_W\) is \emph{bijective}, so that \(\mathtt{restrictScalars}\,\alpha\) is
  \emph{strong} monoidal. Then its oplax-monoidal unit comparison
  \(\eta(\mathtt{restrictScalars}\,\alpha)\) preserves the section ring unit: for every object \(W\),
  \[
    \bigl(\eta(\mathtt{restrictScalars}\,\alpha)_W\bigr)(1_{S(W)}) = 1_{R(W)} .
  \]
\end{lemma}
\begin{proof}
  \uses{lem:restrictscalars_laxmonoidal}
  \leanok
  Since \(\alpha\) is sectionwise bijective, \(\mathtt{restrictScalars}\,\alpha\) carries a
  \emph{strong} monoidal structure, in which the oplax unit \(\eta\) is the two-sided inverse of the
  lax unit \(\varepsilon\). The lax unit sends \(1 \mapsto 1\): on sections it is the change-of-rings
  unit, namely the ring map \(\alpha_W\) applied to the unit, and \(\alpha_W(1)=1\). Composing
  \(\varepsilon\) with its inverse \(\eta\) yields the identity (\(\varepsilon \mathbin{;} \eta =
  \mathbf 1\)), so \(\eta\) too sends the unit to the unit.
\end{proof}

\begin{lemma}


\leanok
[Sectionwise comparison of the LHS mate tensorator on a pure tensor]
  \label{lem:pushforward_lax_mu_comparison_lhs_tmul}
  \lean{AlgebraicGeometry.Scheme.Modules.pushforward_lax_mu_comparison_lhs_tmul}
  \uses{lem:pushforwardpushforwardadj_unit_app_app_apply, lem:pushforwardpushforwardadj_counit_app_app_apply, lem:restrictscalars_mu_app, lem:restrictscalars_mu_app_tmul, lem:forget2_restrictscalars_mu_hom_tmul, lem:restrictscalars_delta_app_tmul, lem:pushforward_lax_mu_comparison_rhs_tmul}
  In the K1 setting the adjoint-transported tensorator
  \(\mu(\mathtt{rightAdjointLaxMonoidal}\,\mathtt{hadj}')_W\) on the objects \(G_\beta A\),
  \(G_\beta B\) arises as an adjunction \emph{mate}, whereas the directly-defined composition
  tensorator \(\mu(\mathtt{presheafPushforwardLaxMonoidal}\,\varphi')_W\) is given outright. This
  lemma is the per-section comparison of the two: on every open \(W\), evaluated at a pure tensor
  \(m \otimes_t n\), the mate tensorator agrees with the directly-defined tensorator. Concretely,
  expanding the mate to its explicit adjunction form
  \[
    \mathtt{hadj}'.\mathrm{unit}
      \;\mathbin{;}\;
    (\mathtt{pushforward}\,\varphi').\mathrm{map}
      \bigl(\delta\,G_\beta\,A\,B \;\mathbin{;}\;
        (\mathtt{hadj}'.\mathrm{counit} \otimes_{\mathrm m} \mathtt{hadj}'.\mathrm{counit})\bigr),
  \]
  the comparison reduces to computing its value on \(m \otimes_t n\); that value is the base-change
  value \(m \otimes_t n\), matching the directly-defined tensorator.
\end{lemma}
\begin{proof}
  \uses{lem:pushforwardpushforwardadj_unit_app_app_apply, lem:pushforwardpushforwardadj_counit_app_app_apply, lem:restrictscalars_mu_app, lem:restrictscalars_mu_app_tmul, lem:forget2_restrictscalars_mu_hom_tmul, lem:restrictscalars_delta_app_tmul, lem:pushforward_lax_mu_comparison_rhs_tmul}
  \leanok
  The value lemma \cref{lem:pushforwardpushforwardadj_unit_app_app_apply} characterises the
  adjunction unit in the canonical \(\mathtt{pushforwardPushforwardAdjunction}\) form. Since the
  mate is expressed via the adjunction \(\mathtt{hadj}'\), we identify \(\mathtt{hadj}'\) with
  \(\mathtt{pushforwardPushforwardAdjunction}\) and then evaluate the three legs of the mate
  sectionwise at \(W\) on the pure tensor
  \(m \otimes_t n\). The adjunction unit acts by precomposition with the pullback of the counit,
  \(M.\mathrm{map}\,(\mathtt{hadj}'.\mathrm{counit})^{\mathrm{op}}\), by
  \cref{lem:pushforwardpushforwardadj_unit_app_app_apply}. Reindexing the
  \((\mathtt{pushforward}\,\varphi').\mathrm{map}\) leg to the section over \(F^{\mathrm{op}}W\), the
  strong-monoidal oplax tensorator \(\delta\,G_\beta\) of the composite
  \(\mathtt{pushforward}_0^{\mathrm{Comm}} \cdot \mathtt{restrictScalars}\,\beta'\) is, on a pure
  tensor, the identity: its \(\mathtt{restrictScalars}\,\beta'\) factor is the base-change tensorator,
  which is the identity on pure tensors by \cref{lem:restrictscalars_mu_app},
  \cref{lem:restrictscalars_mu_app_tmul} and \cref{lem:forget2_restrictscalars_mu_hom_tmul}, while its
  \(\mathtt{pushforward}_0^{\mathrm{Comm}}\) factor has identity tensorator. Finally the pair of
  counits \(\mathtt{hadj}'.\mathrm{counit} \otimes_{\mathrm m} \mathtt{hadj}'.\mathrm{counit}\) acts as
  the structure-ring identity of the \emph{bijective} \(f.\mathtt{appIso}\) on each factor. The three
  legs compose to the identity, so the value is \(m \otimes_t n\).

  We spell out the three closing steps on the pure tensor explicitly. \emph{Step 3 (the \(\delta\)
  leg).} After reindexing the inner \(\varepsilon\)-leg, the comonoidal \(\delta\) of the composite
  functor \(G_\beta = \mathtt{pushforward}_0^{\mathrm{Comm}} \cdot \mathtt{restrictScalars}\,\beta'\)
  splits via \(\mathtt{Functor.OplaxMonoidal.comp}\): the \(\mathtt{pushforward}_0^{\mathrm{Comm}}\)
  factor's \(\delta\) is the identity on a pure tensor, leaving only the
  \(\mathtt{restrictScalars}\,\beta'\) factor's \(\delta\), which is discharged on the pure tensor by
  \cref{lem:restrictscalars_delta_app_tmul}. \emph{Step 4 (the counit pair).} Split the
  \(\mathtt{tensorHom}\) \(\mathtt{hadj}'.\mathrm{counit} \otimes_{\mathrm m}
  \mathtt{hadj}'.\mathrm{counit}\) over the pure tensor into its two factors and collapse each counit
  leg by \cref{lem:pushforwardpushforwardadj_counit_app_app_apply}. \emph{Closing (triangle
  cancellation).} The unit-restriction leg produced by the inner-unit discharge and the
  counit-restriction leg of Step 4 are mutually inverse on the image of the open immersion --- this is
  the adjunction triangle identity read on opens --- so the two legs cancel. What remains is
  \(m \otimes_t n\), which matches the right-hand value by
  \cref{lem:pushforward_lax_mu_comparison_rhs_tmul}.
\end{proof}

\begin{lemma}


\leanok
[Sectionwise value of the RHS composition tensorator on a pure tensor]
  \label{lem:pushforward_lax_mu_comparison_rhs_tmul}
  \lean{AlgebraicGeometry.Scheme.Modules.pushforward_lax_mu_comparison_rhs_tmul}
  \uses{lem:pushforward_mu_eq, lem:restrictscalars_mu_app, lem:restrictscalars_mu_app_tmul, lem:forget2_restrictscalars_mu_hom_tmul, lem:pushforward_map_restrictscalars_mu_app_tmul}
  In the K1 setting, the directly-defined composition tensorator
  \(\mu(\mathtt{presheafPushforwardLaxMonoidal}\,\varphi')\) on \(G_\beta A\), \(G_\beta B\),
  evaluated on an open \(W\) at a pure tensor \(m \otimes_t n\), is the base-change value
  \(m \otimes_t n\).
\end{lemma}
\begin{proof}
  \uses{lem:pushforward_mu_eq, lem:restrictscalars_mu_app, lem:restrictscalars_mu_app_tmul, lem:forget2_restrictscalars_mu_hom_tmul, lem:pushforward_map_restrictscalars_mu_app_tmul}
  \leanok
  By \cref{lem:pushforward_mu_eq} the pushforward composition tensorator reduces to the
  \(\mathtt{restrictScalars}\,\varphi'\) tensorator on the \(\mathtt{pushforward}_0^{\mathrm{Comm}}\)
  reindexed objects. Exposing the \(\mathtt{ModuleCat}\) tensorator by
  \cref{lem:restrictscalars_mu_app} and collapsing on the pure tensor by
  \cref{lem:restrictscalars_mu_app_tmul} and \cref{lem:forget2_restrictscalars_mu_hom_tmul}, with the
  outer \((\mathtt{pushforward}\,\varphi').\mathrm{map}\) leg collapsed by
  \cref{lem:pushforward_map_restrictscalars_mu_app_tmul}, the value is \(m \otimes_t n\). This reuses
  the same helper family as the composition coherence of the single pushforward.
\end{proof}

\begin{lemma}


\leanok
[Bare tensorator comparison of the two lax structures on $\mathtt{pushforward}\,\varphi'$]
  \label{lem:pushforward_lax_mu_comparison}
  \lean{AlgebraicGeometry.Scheme.Modules.pushforward_lax_mu_comparison}
  \uses{lem:pushforward_lax_mu_comparison_lhs_tmul, lem:pushforward_lax_mu_comparison_rhs_tmul, lem:pushforward_mu_eq}
  In the K1 setting above, the functor \(\mathtt{pushforward}\,\varphi'\) carries \emph{two} lax
  monoidal structures: the adjoint-transported one
  \(\mathtt{rightAdjointLaxMonoidal}\,(\mathtt{hadj}')\), induced from the strong-monoidal left
  adjoint \(G_\beta\) through the adjunction \(\mathtt{hadj}' : G_\beta \dashv
  \mathtt{pushforward}\,\varphi'\), and the directly-defined
  \(\mathtt{presheafPushforwardLaxMonoidal}\) (\cref{lem:presheaf_pushforward_laxmonoidal}). Their
  tensorators agree on the objects \(G_\beta A\), \(G_\beta B\):
  \[
    \mu\bigl(\mathtt{rightAdjointLaxMonoidal}\,\mathtt{hadj}'\bigr)\,(G_\beta A)\,(G_\beta B)
    \;=\;
    \mu\bigl(\mathtt{presheafPushforwardLaxMonoidal}\bigr)\,(G_\beta A)\,(G_\beta B).
  \]
  This is a \emph{bare} tensorator identity between two \emph{different kinds} of lax structures on
  the same functor, with \emph{no} appeal to any monoidal-adjunction structure on \(\mathtt{hadj}'\):
  that structure is precisely the monoidal adjunction being established, which in turn depends on
  this identity.
\end{lemma}
\begin{proof}
  \uses{lem:pushforward_lax_mu_comparison_lhs_tmul, lem:pushforward_lax_mu_comparison_rhs_tmul, lem:pushforward_mu_eq}
  \leanok
  Crucially the two sides are \emph{not} symmetric. The left-hand tensorator
  \(\mu(\mathtt{rightAdjointLaxMonoidal}\,\mathtt{hadj}')\) is an adjunction \emph{mate} of the strong
  oplax tensorator of \(G_\beta\), whereas the right-hand
  \(\mu(\mathtt{presheafPushforwardLaxMonoidal}\,\varphi')\) is a \emph{composition} structure. They
  carry no common normal form, and must be brought to a common value by two distinct routes. No
  appeal is made to a monoidal structure on the adjunction \(\mathtt{hadj}'\); such an appeal would
  be circular, since that structure is precisely the monoidal adjunction instance being established,
  which in turn consumes this identity.

  We first reduce the right-hand side at the level of \emph{morphisms}, replacing
  \(\mu(\mathtt{presheafPushforwardLaxMonoidal}\,\varphi')\) by the
  \(\mathtt{restrictScalars}\,\varphi'\) tensorator via \cref{lem:pushforward_mu_eq} at the
  \emph{morphism level}, before descending to sections --- \cref{lem:pushforward_mu_eq} is an
  identity of presheaf morphisms and must be applied there, prior to section evaluation. We then argue sectionwise: on an open \(W\) the domain
  \(G_\beta.\mathrm{obj}\,(A \otimes B)\) is, sectionwise, an honest \(\mathtt{ModuleCat}\) tensor, so
  pure-tensor extensionality reduces the goal to checking equality on each pure tensor
  \(m \otimes_t n\). By \cref{lem:pushforward_lax_mu_comparison_lhs_tmul} the left-hand mate evaluates
  to \(m \otimes_t n\), and by \cref{lem:pushforward_lax_mu_comparison_rhs_tmul} the right-hand
  composition tensorator evaluates to the same \(m \otimes_t n\). The two sides therefore agree on
  every pure tensor, hence as presheaf morphisms.
\end{proof}

\begin{lemma}


\leanok
[K1 $\mu$/$\delta$-side collapse: $H_1$ respects the tensorator]
  \label{lem:pushforward_mu_appiso_collapse}
  \lean{AlgebraicGeometry.Scheme.Modules.pushforward_mu_appIso_collapse}
  \uses{lem:pushforward_lax_mu_comparison, lem:pushforward_mu_eq, lem:tensorobj_restrict_iso, lem:presheaf_pushforward_laxmonoidal, lem:presheaf_pullback_oplaxmonoidal, lem:restrictscalars_monoidal_of_bijective, lem:functor_monoidal_muiso, lem:deltaconjofmucomparison, lem:adjunction_leftadjointuniq_trans}
  In the K1 setting above, for all presheaves of modules \(A, B\) the strong tensorator
  \(\delta\,G_\beta\,A\,B\) of the restriction functor agrees with the \(H_1\)-conjugate of the
  oplax tensorator \(\delta(\mathtt{pullback}\,\varphi')\,A\,B\):
  \[
    \delta\,G_\beta\,A\,B
    \;=\;
    H_1.\mathrm{hom}.\mathrm{app}\,(A \otimes B)
    \;\mathbin{;}\;
    \delta(\mathtt{pullback}\,\varphi')\,A\,B
    \;\mathbin{;}\;
    \bigl(H_1.\mathrm{inv}.\mathrm{app}\,A \otimes_{\mathrm m} H_1.\mathrm{inv}.\mathrm{app}\,B\bigr).
  \]
  Equivalently, \(H_1\) is monoidal on tensorators. This is the genuine geometric content of K1
  (the $\mu$/$\delta$-side twin of \cref{lem:presheafunit_comp_map_eta}, and the open-immersion
  analogue of \cref{lem:pushforwardcomp_lax_mu}).
\end{lemma}
\begin{proof}
  \uses{lem:pushforward_lax_mu_comparison, lem:pushforward_mu_eq}
  \leanok
  The genuine geometric content is the bare tensorator comparison
  \cref{lem:pushforward_lax_mu_comparison}; this proof is the formal reduction to it, and uses
  \emph{no} mate machinery (\(\mathtt{hadj}'.\mathtt{IsMonoidal}\)), which would be circular. Four
  steps: (1) cancel the isomorphism \(H_1 = \mathtt{leftAdjointUniq}\), restating the goal on the
  \(\mathtt{pullback}\,\varphi'\) domain; (2) transpose across
  \(\mathtt{hadj}' : G_\beta \dashv \mathtt{pushforward}\,\varphi'\) by hom-equivalence
  injectivity, turning the left adjoint's oplax tensorator into the right adjoint's lax tensorator
  \(\mu(\mathtt{rightAdjointLaxMonoidal}\,\mathtt{hadj}')\); (3) replace \(\delta\,G_\beta\) by
  \(\mu_{\mathrm{Iso},\beta}^{-1}\); (4) reduce
  \(\delta(\mathtt{pullback}\,\varphi')\) to \(\mu(\mathtt{presheafPushforwardLaxMonoidal})\) via
  \cref{lem:pushforward_mu_eq}. The goal is then exactly
  \cref{lem:pushforward_lax_mu_comparison}, which closes it.
\end{proof}

\begin{lemma}


\leanok
[Pullback commutes with $\otimes$ along an open immersion (K1)]
  \label{lem:pullback_tensor_map_isiso_open_immersion}
  \lean{AlgebraicGeometry.Scheme.Modules.pullbackTensorMap_isIso_of_isOpenImmersion}
  \uses{lem:pullback_tensor_map, lem:tensorobj_restrict_iso, lem:isiso_pullbacktensormap_of_sheafifydelta}
  Let \(f : Y \to X\) be an \emph{open immersion} of schemes and let
  \(M, N \in \Scheme.\mathtt{Modules}\,X\) be \emph{arbitrary} \(\mathcal{O}_X\)-modules.
  Then the sheaf-level comparison map \(\mathtt{pullbackTensorMap}\,f\,M\,N\) of
  \cref{lem:pullback_tensor_map},
  \[
    f^*(M \otimes_X N)
      \;\xrightarrow{\;\sim\;}\;
    (f^*M) \otimes_Y (f^*N),
  \]
  is an isomorphism. No local-triviality hypothesis on \(M, N\) is required: along an open
  immersion the pullback functor \emph{is} the restriction functor, which is \emph{strong}
  monoidal, so the comparison is invertible for every pair of modules.
  % SOURCE: [Stacks Project], Modules on Ringed Spaces, Lemma
  % \ref{lemma-tensor-product-pullback} (pullback commutes with the tensor product,
  % functorially)
  % (read from references/stacks-modules.tex, L2392-2400).
  % SOURCE QUOTE:
  % "\begin{lemma}
  % \label{lemma-tensor-product-pullback}
  % Let $f : (X, \mathcal{O}_X) \to (Y, \mathcal{O}_Y)$ be
  % a morphism of ringed spaces. Let $\mathcal{F}$, $\mathcal{G}$
  % be $\mathcal{O}_Y$-modules. Then
  % $f^*(\mathcal{F} \otimes_{\mathcal{O}_Y} \mathcal{G})
  % = f^*\mathcal{F} \otimes_{\mathcal{O}_X} f^*\mathcal{G}$
  % functorially in $\mathcal{F}$, $\mathcal{G}$.
  % \end{lemma}"
  \textit{Source: [Stacks Project], Modules on Ringed Spaces, Lemma
  \texttt{lemma-tensor-product-pullback}: pullback commutes with \(\otimes\) functorially.
  Along an open immersion the pullback is restriction, base change along the structure-sheaf
  \emph{isomorphism}, so the comparison is an isomorphism for all modules.}
\end{lemma}

\begin{proof}
  \uses{lem:pullback_tensor_map, lem:tensorobj_restrict_iso, lem:isiso_pullbacktensormap_of_sheafifydelta, lem:presheafunit_comp_map_eta, lem:pushforward_eta_appiso_collapse, lem:pushforward_mu_appiso_collapse, lem:deltaconjofmucomparison, lem:isiso_oplaxdelta_of_conj}
  \leanok
  \textit{Source: [Stacks Project], Modules on Ringed Spaces, Lemma
  \texttt{lemma-tensor-product-pullback}.}

  Unfold \(\mathtt{pullbackTensorMap}\,f\,M\,N\) by its four-fold factorisation
  (\cref{lem:pullback_tensor_map}, made explicit in
  \cref{lem:isiso_pullbacktensormap_of_sheafifydelta}). The three outer factors
  \(p_1, p_3, p_4\) are isomorphisms unconditionally, so the comparison is an isomorphism as
  soon as the single middle factor --- the sheafification \(a_Y.\mathrm{map}\,\delta\) of the
  presheaf-level oplax comparison
  \(\delta(\mathtt{pullback}\,\varphi')(M.\mathrm{val}, N.\mathrm{val})\) --- is an
  isomorphism, where \(\varphi'\) is the presheaf structure map presenting \(f\).
  Sheafification is a functor and hence preserves isomorphisms, so it suffices to show the
  \emph{presheaf}-level comparison \(\delta(\mathtt{pullback}\,\varphi')\) is an isomorphism.

  For an open immersion this presheaf-level statement is exactly the strong-monoidality of
  pullback \(=\) restriction, and we establish it by the same two-ingredient model that
  closes \cref{lem:tensorobj_restrict_iso}.

  \emph{H1 (the presheaf identification).} Let \(\beta\) be the structure map of the
  restriction functor along the open immersion \(f\). One exhibits the presheaf-level
  adjunction \(\mathtt{pushforward}\,\beta \dashv \mathtt{pushforward}\,\varphi'\) and
  concludes, by \(\mathtt{Adjunction.leftAdjointUniq}\), a canonical isomorphism
  \(\mathtt{pushforward}\,\beta \cong \mathtt{pullback}\,\varphi'\) between the two left
  adjoints of \(\mathtt{pushforward}\,\varphi'\). (This is the very same H1 identification
  that closes \cref{lem:tensorobj_restrict_iso}.)

  \emph{H2 (the strong-monoidal tensorator).} Along an open immersion the structure map
  \(\beta\) is sectionwise the bijective structure-ring isomorphism \(f.\mathtt{appIso}\), so
  \(\mathtt{restrictScalars}\,\beta\), and hence \(\mathtt{pushforward}\,\beta\), is
  \emph{strong} monoidal --- exactly the strong-monoidal upgrade of base change along a ring
  isomorphism used in \cref{lem:tensorobj_restrict_iso}. In particular
  \(\mathtt{pushforward}\,\beta\) carries an \emph{invertible} tensorator.

  The two pure-tensor collapses --- the unit compatibility
  \cref{lem:pushforward_eta_appiso_collapse} (with template \cref{lem:presheafunit_comp_map_eta})
  and the tensorator compatibility \cref{lem:pushforward_mu_appiso_collapse} --- confirm that
  \(H_1\) is a monoidal natural isomorphism. The tensorator compatibility in particular supplies
  the hypothesis of the abstract \(\delta\)-conjugation \cref{lem:deltaconjofmucomparison},
  yielding that \(\delta\,G_\beta\) is the \(H_1\)-conjugate of
  \(\delta(\mathtt{pullback}\,\varphi')\). Since \(G_\beta\) is strong monoidal (\textbf{H2}),
  its oplax tensorator \(\delta\,G_\beta\) is the inverse of its strong tensorator
  \(\mu_{\mathrm{Iso}}\,G_\beta\) (\cref{lem:functor_monoidal_muiso}) and is therefore an
  isomorphism; \cref{lem:isiso_oplaxdelta_of_conj} then concludes that
  \(\delta(\mathtt{pullback}\,\varphi')\) is an isomorphism.

  Hence \(\delta(\mathtt{pullback}\,\varphi')\) is an isomorphism, so
  its sheafification \(a_Y.\mathrm{map}\,\delta\) is too, and therefore
  \(\mathtt{pullbackTensorMap}\,f\,M\,N\) is an isomorphism.
\end{proof}

\begin{lemma}


\leanok
[Pullback commutes with $\otimes$ on locally-trivial pairs (D4$'$)]
  \label{lem:pullback_tensor_iso_loctriv}
  \lean{AlgebraicGeometry.Scheme.Modules.pullbackTensorIsoOfLocallyTrivial}
  % NOTE: the D4' realisation routes through the open-immersion comparison lemma K1
  % (lem:pullback_tensor_map_isiso_open_immersion): isolating the restriction factor from
  % the D3' composite requires the flanking comparisons delta^{j_i'}, delta^{j_i} along
  % the open immersions to be invertible = K1, before the reduction to delta^{g_i} on the
  % restricted modules is valid. The per-chart obligation and the structure-sheaf-pair case
  % are covered by lem:chart_isiso and lem:pullback_tensor_map_isiso_base_unit.
  \uses{def:IsLocallyTrivial, lem:pullback_tensor_map, lem:pullback_tensor_map_natural, lem:pullback_tensor_iso_unit, lem:pullback_tensor_map_basechange, lem:pullback_tensor_map_isiso_open_immersion, lem:isiso_of_isiso_restrict, lem:tensorobj_preserves_locally_trivial, lem:IsLocallyTrivial_pullback}
  Let \(f : Y \to X\) be a morphism of schemes and let
  \(M, N \in \Scheme.\mathtt{Modules}\,X\) be \emph{locally trivial}
  (\cref{def:IsLocallyTrivial}). Then the sheaf-level comparison map
  \(\delta_{\mathrm{sheaf}} = \mathtt{pullbackTensorMap}\) of
  \cref{lem:pullback_tensor_map} is an isomorphism:
  \[
    f^*(M \otimes_X N)
      \;\xrightarrow{\;\sim\;}\;
    (f^*M) \otimes_Y (f^*N),
  \]
  natural in \(M\) and \(N\). This is the comparison the relative Picard consumer needs;
  it is established by promoting the already-built map \(\delta_{\mathrm{sheaf}}\) to an
  isomorphism through a chart-chase, with no comparison for arbitrary modules and no
  filtered-colimit machinery.
  % SOURCE: [Stacks Project], Modules on Ringed Spaces, Lemma
  % \ref{lemma-tensor-product-pullback} (pullback commutes with the tensor product,
  % functorially) and Lemma \ref{lemma-pullback-locally-free} (the pullback of a locally
  % free module is locally free --- the geometric content licensing the chart-chase)
  % (read from references/stacks-modules.tex, L2392-2400 and L2112-2118).
  % SOURCE QUOTE:
  % "\begin{lemma}
  % \label{lemma-tensor-product-pullback}
  % Let $f : (X, \mathcal{O}_X) \to (Y, \mathcal{O}_Y)$ be
  % a morphism of ringed spaces. Let $\mathcal{F}$, $\mathcal{G}$
  % be $\mathcal{O}_Y$-modules. Then
  % $f^*(\mathcal{F} \otimes_{\mathcal{O}_Y} \mathcal{G})
  % = f^*\mathcal{F} \otimes_{\mathcal{O}_X} f^*\mathcal{G}$
  % functorially in $\mathcal{F}$, $\mathcal{G}$.
  % \end{lemma}"
  % SOURCE QUOTE:
  % "\begin{lemma}
  % \label{lemma-pullback-locally-free}
  % Let $f : (X, \mathcal{O}_X) \to (Y, \mathcal{O}_Y)$
  % be a morphism of ringed spaces. If $\mathcal{G}$ is
  % a locally free $\mathcal{O}_Y$-module, then
  % $f^*\mathcal{G}$ is a locally free $\mathcal{O}_X$-module.
  % \end{lemma}"
  \textit{Source: [Stacks Project], Modules on Ringed Spaces, Lemmas
  \texttt{lemma-tensor-product-pullback} and \texttt{lemma-pullback-locally-free}:
  pullback commutes with \(\otimes\) functorially, and the pullback of a locally free
  (locally trivial) module is again locally free; on locally-trivial pairs the comparison
  is therefore an isomorphism, checkable on a trivialising cover.}
\end{lemma}

\begin{proof}
  \uses{def:IsLocallyTrivial, lem:pullback_tensor_map, lem:pullback_tensor_map_natural, lem:pullback_tensor_iso_unit, lem:pullback_tensor_map_basechange, lem:pullback_tensor_map_isiso_open_immersion, lem:isiso_of_isiso_restrict, lem:tensorobj_preserves_locally_trivial, lem:IsLocallyTrivial_pullback}
  \textit{Source: [Stacks Project], Lemmas \texttt{lemma-tensor-product-pullback} and
  \texttt{lemma-pullback-locally-free}.}

  By local triviality of \(M\) and \(N\) (\cref{def:IsLocallyTrivial}) choose a common
  affine open cover \(\{U_i\}\) of \(X\) on which both restrict to the trivial bundle,
  \(M|_{U_i} \cong \mathcal{O}_{U_i}\) and \(N|_{U_i} \cong \mathcal{O}_{U_i}\) --- the
  common refinement of the two trivialising covers, exactly as in the proof that the
  tensor product of locally-trivial modules is locally trivial
  (\cref{lem:tensorobj_preserves_locally_trivial}). The preimages
  \(\{f^{-1}U_i\}\) form an open cover of \(Y\). It suffices, by the
  restriction-detects-isomorphism criterion \cref{lem:isiso_of_isiso_restrict}, to show
  that the canonical map \(\delta_{\mathrm{sheaf}}(M,N)\) restricts to an isomorphism on
  each \(f^{-1}U_i\); the canonical global map being shown locally-iso is exactly what
  lets us conclude (a non-canonical patchwork of local isomorphisms would not glue).

  Fix \(i\) and write \(g_i : f^{-1}U_i \to U_i\) for the restriction of \(f\), with open
  immersions \(j_i : U_i \hookrightarrow X\) and \(j_i' : f^{-1}U_i \hookrightarrow Y\),
  so \(f \circ j_i' = j_i \circ g_i\). The base-change coherence
  \cref{lem:pullback_tensor_map_basechange} expresses the restriction
  \((j_i')^*\delta^{f}_{\mathrm{sheaf}}(M,N)\) as the comparison
  \(\delta^{g_i}_{\mathrm{sheaf}}(j_i^*M, j_i^*N)\) for \(g_i\) on the restricted modules,
  \emph{flanked} by the comparisons \(\delta^{j_i'}\) and \(\delta^{j_i}\) along the two open
  immersions. Isolating the restriction factor \(\delta^{g_i}\) from this four-fold composite
  is legitimate only because those flanking factors are \emph{invertible}: by
  \cref{lem:pullback_tensor_map_isiso_open_immersion} (K1) the comparison along an open
  immersion is an isomorphism for arbitrary modules, so \(\delta^{j_i'}\) and \(\delta^{j_i}\)
  may be cancelled. The claim is thereby reduced to:
  \(\delta^{g_i}_{\mathrm{sheaf}}(j_i^*M, j_i^*N)\) is an isomorphism. On \(U_i\) the restricted modules are trivial, \(j_i^*M \cong
  \mathcal{O}_{U_i}\) and \(j_i^*N \cong \mathcal{O}_{U_i}\); by naturality of
  \(\delta_{\mathrm{sheaf}}\) in both arguments (\cref{lem:pullback_tensor_map_natural})
  these isomorphisms transport \(\delta^{g_i}_{\mathrm{sheaf}}(j_i^*M, j_i^*N)\) to the
  comparison on the unit pair \(\delta^{g_i}_{\mathrm{sheaf}}(\mathcal{O}_{U_i},
  \mathcal{O}_{U_i})\). That comparison is an isomorphism by
  \cref{lem:pullback_tensor_iso_unit} --- the unit-pair case, where the structure-sheaf
  comparison \(\mathtt{pullbackUnitIso}\) (an isomorphism for all \(g_i\)) and the unitors
  do the work. Hence \(\delta^{g_i}_{\mathrm{sheaf}}(j_i^*M, j_i^*N)\) is an isomorphism,
  so is its base-change identification \((j_i')^*\delta^{f}_{\mathrm{sheaf}}(M,N)\), and
  thus \(\delta^{f}_{\mathrm{sheaf}}(M,N)\) restricts to an isomorphism on each
  \(f^{-1}U_i\). (Both sides are indeed locally trivial on this cover --- \(f^*M, f^*N,
  f^*(M\otimes_X N)\) are locally trivial by \cref{lem:IsLocallyTrivial_pullback} and
  \cref{lem:tensorobj_preserves_locally_trivial} --- which is what makes the chart-local
  comparison a map between copies of the structure sheaf.) By
  \cref{lem:isiso_of_isiso_restrict}, \(\delta_{\mathrm{sheaf}}(M,N)\) is a global
  isomorphism, natural in \(M, N\) by \cref{lem:pullback_tensor_map_natural}.
\end{proof}

\begin{lemma}


\leanok
[$\delta_{\mathrm{sheaf}}$ is an isomorphism on a structure-sheaf pair (D4$'$ base case)]
  \label{lem:pullback_tensor_map_isiso_base_unit}
  \lean{AlgebraicGeometry.Scheme.Modules.pullbackTensorMap_isIso_of_base_unit}
  \uses{lem:pullback_tensor_map_natural, lem:pullback_tensor_iso_unit}
  Let \(f : Y \to X\) be a morphism of schemes and let
  \(P, Q \in \Scheme.\mathtt{Modules}\,X\) each be isomorphic to the structure sheaf,
  \(P \cong \mathcal{O}_X\) and \(Q \cong \mathcal{O}_X\). Then the comparison
  \(\mathtt{pullbackTensorMap}\,f\,P\,Q\) of \cref{lem:pullback_tensor_map} is an isomorphism.
\end{lemma}

\begin{proof}
  \uses{lem:pullback_tensor_map_natural, lem:pullback_tensor_iso_unit}
  \leanok
  By naturality of \(\delta_{\mathrm{sheaf}}\) in both arguments
  (\cref{lem:pullback_tensor_map_natural}), the chosen isomorphisms
  \(P \cong \mathcal{O}_X\) and \(Q \cong \mathcal{O}_X\) transport
  \(\mathtt{pullbackTensorMap}\,f\,P\,Q\) onto the comparison on the unit pair
  \(\mathtt{pullbackTensorMap}\,f\,\mathcal{O}_X\,\mathcal{O}_X\), which is an isomorphism by
  \cref{lem:pullback_tensor_iso_unit} (D2$'$). Transport of an isomorphism along
  isomorphisms is again an isomorphism, so \(\mathtt{pullbackTensorMap}\,f\,P\,Q\) is one.
\end{proof}

\begin{lemma}


\leanok
[Per-chart isomorphism obligation for the D4$'$ chart-chase]
  \label{lem:chart_isiso}
  \lean{AlgebraicGeometry.Scheme.Modules.chart_isIso}
  \uses{lem:isiso_of_isiso_restrict, lem:pullback_tensor_map_basechange, lem:pullback_tensor_map_isiso_open_immersion, lem:pullback_tensor_map_isiso_base_unit}
  In the situation of \cref{lem:pullback_tensor_iso_loctriv}, fix a member \(f^{-1}U_i\) of
  the trivialising cover, with restriction \(g_i : f^{-1}U_i \to U_i\) of \(f\) and open
  immersions \(j_i : U_i \hookrightarrow X\), \(j_i' : f^{-1}U_i \hookrightarrow Y\). Then the
  restriction \((j_i')^*\delta^{f}_{\mathrm{sheaf}}(M,N)\) of the comparison to that chart is
  an isomorphism.
\end{lemma}

\begin{proof}
  \uses{lem:isiso_of_isiso_restrict, lem:pullback_tensor_map_basechange, lem:pullback_tensor_map_isiso_open_immersion, lem:pullback_tensor_map_isiso_base_unit}
  \leanok
  The base-change coherence \cref{lem:pullback_tensor_map_basechange}, applied at the two
  factorisations \(f \circ j_i' = j_i \circ g_i\), presents
  \((j_i')^*\delta^{f}_{\mathrm{sheaf}}(M,N)\) as the comparison \(\delta^{g_i}\) flanked by
  the two open-immersion comparisons \(\delta^{j_i'}, \delta^{j_i}\); these flanking factors
  are isomorphisms by \cref{lem:pullback_tensor_map_isiso_open_immersion} (K1, applied twice)
  and cancel, while \(\delta^{g_i}\) on the trivialised modules
  \(j_i^*M \cong \mathcal{O}_{U_i}, j_i^*N \cong \mathcal{O}_{U_i}\) is an isomorphism by the
  structure-sheaf-pair case \cref{lem:pullback_tensor_map_isiso_base_unit}. (The
  restriction criterion \cref{lem:isiso_of_isiso_restrict} is what makes this chart-local
  invertibility the relevant per-cover obligation.) Hence
  \((j_i')^*\delta^{f}_{\mathrm{sheaf}}(M,N)\) is an isomorphism.
\end{proof}

\section{The invertibility-carrier Picard group}
\label{sec:tensorobj_pic_carrier}

The group law of \cref{lem:tensorobj_isoclass_commgroup} is carried on the
\emph{locally-trivial} predicate \(\mathtt{IsLocallyTrivial}\), and on that carrier
the existence of a tensor inverse for each class is an extra obligation
(\cref{lem:tensorobj_inverse_invertible}): one must \emph{manufacture} the dual
object \(L^{-1} = \mathcal{H}om_{\mathcal{O}_X}(L, \mathcal{O}_X)\) and prove the
contraction \(L \otimes_X L^{-1} \cong \mathcal{O}_X\), which is the whole content of
the dual-infrastructure build of \cref{sec:tensorobj_dual_infra}.

This section carries the group law on a \emph{different} predicate for which that
obligation evaporates: the \(\otimes\)-invertibility predicate
\(\mathtt{IsInvertible}\) of \cref{def:scheme_modules_isinvertible},
\[
  \mathtt{IsInvertible}(M)
  \;:\equiv\;
  \exists\, N,\ \mathrm{Nonempty}\bigl(M \otimes_X N \cong \mathcal{O}_X\bigr).
\]
On this carrier the inverse of a class \([M]\) is supplied by the \emph{membership
witness itself} --- the very \(N\) whose existence \(\mathtt{IsInvertible}(M)\)
asserts --- so no inverse object need be constructed and no separate
tensor-inverse lemma is consumed. The cost is shifted to a routine
well-definedness fact (\cref{lem:isinvertible_inverse_welldef}: the witness \(N\) is
determined up to isomorphism), proved from the monoidal coherence isomorphisms
alone. The carrier is the scheme-theoretic analogue of Mathlib's
\(\mathtt{CommRing.Pic}\), whose elements are exactly the iso-classes of invertible
modules (Stacks tag 01CX).

\begin{lemma}


\leanok
[Associator for $\otimes_X$ on invertible objects (corollary of the unconditional
associator)]
  \label{lem:tensorobj_assoc_iso_invertible}
  \lean{AlgebraicGeometry.Scheme.Modules.tensorObj_assoc_iso_invertible}
  \uses{def:scheme_modules_tensorobj, def:scheme_modules_isinvertible, lem:tensorobj_assoc_iso}
  Let \(X\) be a scheme and let
  \(M, N, P \in \Scheme.\mathtt{Modules}\,X\) be \(\otimes\)-invertible
  (\cref{def:scheme_modules_isinvertible}). There is an isomorphism
  \[
    (M \otimes_X N) \otimes_X P
      \;\xrightarrow{\ \sim\ }\;
    M \otimes_X (N \otimes_X P)
  \]
  in \(\Scheme.\mathtt{Modules}\,X\). The invertibility-carrier group law
  (\texttt{pic\_commgroup}) consumes only the \emph{existence} of this
  isomorphism, the proposition \(\mathrm{Nonempty}\bigl((M \otimes_X N) \otimes_X P
  \cong M \otimes_X (N \otimes_X P)\bigr)\); no pentagon, triangle, or naturality
  coherence is ever required.
\end{lemma}

\begin{proof}
  \uses{lem:tensorobj_assoc_iso}
  \leanok
  Immediate specialisation of the \emph{unconditional} associator
  \cref{lem:tensorobj_assoc_iso} to the invertible arguments \(M, N, P\); no
  hypothesis on the arguments is used. (The earlier route argued via ``invertible
  \(\Rightarrow\) locally free \(\Rightarrow\) sectionwise flat'' to feed the
  flat-whiskering lemma \texttt{flat\_whisker\_localizer}; that step is \emph{false}
  --- local freeness does not give global sectionwise flatness \(M(U)\) over
  \(\mathcal{O}_X(U)\) for a non-affine open \(U\) --- and is no longer needed, since
  the general associator is itself unconditional, its single open left-whiskering
  obligation \cref{lem:islocallyinjective_whisker_of_W} being closed via the d.2
  stalk--tensor commutation \cref{lem:islocallyinjective_whiskerleft_via_stalk}.)
\end{proof}

\begin{definition}


\leanok
[The invertibility-carrier Picard group \(\Pic X\)]
  \label{def:pic_carrier}
  \lean{AlgebraicGeometry.Scheme.Modules.PicGroup}
  \uses{def:scheme_modules_tensorobj, def:scheme_modules_isinvertible, def:pic_setoid}
  % SOURCE: [Stacks Project], Tag 01CX (Modules on Ringed Spaces, \S Invertible
  % modules), Definition definition-pic (the Picard group as the abelian group of
  % iso-classes of invertible modules under tensor product)
  % (read from references/stacks-modules.tex, L4350-L4357).
  % SOURCE QUOTE:
  % "\begin{definition}
  % \label{definition-pic}
  % Let $(X, \mathcal{O}_X)$ be a ringed space.
  % The {\it Picard group} $\Pic(X)$ of $X$ is the
  % abelian group whose elements are isomorphism classes of
  % invertible $\mathcal{O}_X$-modules, with addition
  % corresponding to tensor product.
  % \end{definition}"
  \textit{Source: [Stacks Project], Tag 01CX, Definition~\texttt{definition-pic}:
  the Picard group \(\Pic(X)\) is the abelian group of isomorphism classes of
  invertible \(\mathcal{O}_X\)-modules, with addition the tensor product.}
  Let \(X\) be a scheme. Consider the subtype of \(\otimes\)-invertible objects
  \[
    \mathtt{Invertible}\,X
    \;:=\;
    \bigl\{\, M \in \Scheme.\mathtt{Modules}\,X
              \,\bigm|\, \mathtt{IsInvertible}\,M \,\bigr\}
  \]
  of \cref{def:scheme_modules_isinvertible}, and equip it with the setoid whose
  relation is ``there exists an isomorphism'',
  \(M \sim M' :\equiv \mathrm{Nonempty}(M \cong M')\) (a proposition-valued
  equivalence relation: reflexive, symmetric and transitive because isomorphisms
  compose and invert). The \emph{Picard group carrier}
  \[
    \Pic X \;:=\; \mathtt{Invertible}\,X \,/\!\sim
  \]
  is the quotient type of isomorphism classes of \(\otimes\)-invertible
  \(\mathcal{O}_X\)-modules. We write \([M] \in \Pic X\) for the class of an
  invertible \(M\). This is the scheme-theoretic analogue of Mathlib's
  \(\mathtt{CommRing.Pic}\,R\), carried on \(\mathtt{IsInvertible}\) rather than on
  \(\mathtt{IsLocallyTrivial}\); on a scheme the two predicates agree (Stacks tag
  0B8M), but only the invertibility form carries the inverse witness existentially.
\end{definition}

\begin{lemma}


\leanok
[\(\otimes\)-invertibility is closed under tensor product]
  \label{lem:isinvertible_tensor}
  \lean{AlgebraicGeometry.Scheme.Modules.IsInvertible.tensorObj}
  \uses{def:scheme_modules_isinvertible, lem:tensorobj_assoc_iso_invertible, lem:tensorobj_comm_iso, lem:tensorobj_unit_iso}
  % SOURCE: [Stacks Project], Tag 01CR (Modules on Ringed Spaces, \S Invertible
  % modules), Lemma lemma-constructions-invertible (item 1: the tensor product of
  % two invertible modules is invertible)
  % (read from references/stacks-modules.tex, L4200-L4213).
  % SOURCE QUOTE:
  % "\begin{lemma}
  % \label{lemma-constructions-invertible}
  % Let $(X, \mathcal{O}_X)$ be a ringed space.
  % \begin{enumerate}
  % \item If $\mathcal{L}$, $\mathcal{N}$ are invertible
  % $\mathcal{O}_X$-modules, then so is
  % $\mathcal{L} \otimes_{\mathcal{O}_X} \mathcal{N}$.
  % \item If $\mathcal{L}$ is an invertible $\mathcal{O}_X$-module, then so is
  % $\SheafHom_{\mathcal{O}_X}(\mathcal{L}, \mathcal{O}_X)$ and the evaluation map
  % $\mathcal{L} \otimes_{\mathcal{O}_X}
  % \SheafHom_{\mathcal{O}_X}(\mathcal{L}, \mathcal{O}_X) \to \mathcal{O}_X$
  % is an isomorphism.
  % \end{enumerate}
  % \end{lemma}"
  \textit{Source: [Stacks Project], Tag 01CR,
  Lemma~\texttt{lemma-constructions-invertible} (item 1): the tensor product of two
  invertible \(\mathcal{O}_X\)-modules is again invertible.}
  Let \(X\) be a scheme and \(M, M' \in \Scheme.\mathtt{Modules}\,X\). If
  \(\mathtt{IsInvertible}\,M\) and \(\mathtt{IsInvertible}\,M'\), then
  \(\mathtt{IsInvertible}(M \otimes_X M')\).
\end{lemma}

\begin{proof}
  \uses{def:scheme_modules_isinvertible, lem:tensorobj_assoc_iso_invertible, lem:tensorobj_comm_iso, lem:tensorobj_unit_iso}
  \leanok
  Let \(N\) be a tensor inverse of \(M\) (so \(M \otimes_X N \cong \mathcal{O}_X\))
  and \(N'\) a tensor inverse of \(M'\) (so \(M' \otimes_X N' \cong \mathcal{O}_X\)),
  furnished by the two invertibility witnesses. We claim \(N \otimes_X N'\) is a
  tensor inverse of \(M \otimes_X M'\). Indeed, reassociating and commuting the four
  factors with the associator \cref{lem:tensorobj_assoc_iso_invertible} and the
  braiding \cref{lem:tensorobj_comm_iso},
  \[
    (M \otimes_X M') \otimes_X (N \otimes_X N')
    \;\cong\;
    (M \otimes_X N) \otimes_X (M' \otimes_X N')
    \;\cong\;
    \mathcal{O}_X \otimes_X \mathcal{O}_X
    \;\cong\;
    \mathcal{O}_X,
  \]
  the middle step substituting the two witness isomorphisms and the last being the
  unitor \cref{lem:tensorobj_unit_iso}. Hence \(N \otimes_X N'\) witnesses
  \(\mathtt{IsInvertible}(M \otimes_X M')\). Every step is consumed only as the
  existence of an isomorphism, so no coherence identity is required.
\end{proof}

\begin{lemma}


\leanok
[The tensor inverse is determined up to isomorphism]
  \label{lem:isinvertible_inverse_welldef}
  \lean{AlgebraicGeometry.Scheme.Modules.IsInvertible.inverse_unique}
  \uses{def:scheme_modules_isinvertible, lem:tensorobj_unit_iso, lem:tensorobj_assoc_iso_invertible, lem:tensorobj_comm_iso}
  % SOURCE: [Stacks Project], Tag 01CR (Modules on Ringed Spaces, \S Invertible
  % modules), Lemma lemma-invertible (the inverse $\mathcal{N}$ is unique up to
  % isomorphism: it is isomorphic to the dual)
  % (read from references/stacks-modules.tex, L4066-L4079).
  % SOURCE QUOTE:
  % "\begin{lemma}
  % \label{lemma-invertible}
  % Let $(X, \mathcal{O}_X)$ be a ringed space. Let $\mathcal{L}$
  % be an $\mathcal{O}_X$-module. Equivalent are
  % \begin{enumerate}
  % \item $\mathcal{L}$ is invertible, and
  % \item there exists an $\mathcal{O}_X$-module $\mathcal{N}$
  % such that
  % $\mathcal{L} \otimes_{\mathcal{O}_X} \mathcal{N} \cong \mathcal{O}_X$.
  % \end{enumerate}
  % In this case $\mathcal{L}$ is locally a direct summand of a finite free
  % $\mathcal{O}_X$-module and the module $\mathcal{N}$ in (2) is isomorphic to
  % $\SheafHom_{\mathcal{O}_X}(\mathcal{L}, \mathcal{O}_X)$.
  % \end{lemma}"
  \textit{Source: [Stacks Project], Tag 01CR, Lemma~\texttt{lemma-invertible}: when
  \(\mathcal{L}\) is invertible the module \(\mathcal{N}\) with
  \(\mathcal{L} \otimes_{\mathcal{O}_X} \mathcal{N} \cong \mathcal{O}_X\) is
  isomorphic to the dual \(\SheafHom_{\mathcal{O}_X}(\mathcal{L}, \mathcal{O}_X)\),
  hence unique up to isomorphism.}
  Let \(X\) be a scheme and \(M, N, N' \in \Scheme.\mathtt{Modules}\,X\). If
  \(M \otimes_X N \cong \mathcal{O}_X\) and \(M \otimes_X N' \cong \mathcal{O}_X\),
  then \(N \cong N'\). Consequently the assignment \([M] \mapsto [N]\), sending an
  invertible class to the class of its tensor inverse, is well-defined on
  \(\Pic X\).
\end{lemma}

\begin{proof}
  \uses{def:scheme_modules_isinvertible, lem:tensorobj_unit_iso, lem:tensorobj_assoc_iso_invertible, lem:tensorobj_comm_iso}
  \leanok
  This is the standard ``inverse-of-inverse'' argument in a commutative monoid, run
  with isomorphisms in place of equalities. Chain the unitor
  \cref{lem:tensorobj_unit_iso}, the two hypotheses, the associator
  \cref{lem:tensorobj_assoc_iso_invertible} and the braiding
  \cref{lem:tensorobj_comm_iso}:
  \[
    N
    \;\cong\; N \otimes_X \mathcal{O}_X
    \;\cong\; N \otimes_X (M \otimes_X N')
    \;\cong\; (N \otimes_X M) \otimes_X N'
    \;\cong\; \mathcal{O}_X \otimes_X N'
    \;\cong\; N',
  \]
  where the second step substitutes \(\mathcal{O}_X \cong M \otimes_X N'\), the
  third is the associator, the fourth uses the braiding
  \(N \otimes_X M \cong M \otimes_X N \cong \mathcal{O}_X\), and the outer two are
  the unitors. Hence \(N \cong N'\). For well-definedness on classes, note that if
  \(M \cong M'\) then a tensor inverse of \(M\) is also a tensor inverse of \(M'\)
  (tensor with the isomorphism), so the inverse class depends only on \([M]\).
\end{proof}

\section{Consumer: the \(\AddCommGroup\) instance for the relative Picard sheaf}
\label{sec:tensorobj_consumer_acg}

The substrate of \cref{sec:tensorobj_substrate,sec:tensorobj_onproduct_lift}
supplies all the data missing from \cref{lem:rel_pic_sharp_groupoid}'s Lean
encoding. The consumer instance below packages the abelian-group law
\(\bigl[\,\mathcal{L}\,\bigr] + \bigl[\,\mathcal{L}'\,\bigr]
 := \bigl[\,\mathcal{L} \otimes_{C \times_S T} \mathcal{L}'\,\bigr]\) on the
relative Picard quotient via the substrate.

\begin{theorem}

[Abelian-group instance on the relative Picard quotient via \(\Scheme.\mathtt{Modules}.\mathtt{tensorObj}\)]
  \label{thm:rel_pic_addcommgroup_via_tensorobj}
  \lean{AlgebraicGeometry.Scheme.PicSharp.addCommGroup_via_tensorObj}
  % NOTE: \lean{} pin target absent from Lean source as of iter-271 (verified by grep).
  % Forward reference to a not-yet-formalized declaration (the real instance currently in
  % Lean is `PicSharp.addCommGroup`; the `_via_tensorObj` upgrade is not yet built) — NOT a
  % stale rename; no existing Lean decl to repin to under this name.
  % NOTE (carrier pivot, deferred): the \uses below still points at the old
  % locally-trivial group law lem:tensorobj_isoclass_commgroup; this is DEFERRED and
  % will be repointed to thm:pic_commgroup (the invertibility-carrier group, inverse
  % free) once that carrier group is formalized. Left unchanged this iter.
  \uses{lem:tensorobj_lift_onproduct, lem:pullback_compatible_with_tensorobj, def:scheme_modules_isinvertible, lem:tensorobj_isoclass_commgroup, lem:tensorobj_assoc_iso, lem:tensorobj_unit_iso, lem:tensorobj_comm_iso, thm:relative_pic_quotient_well_defined, lem:rel_pic_sharp_groupoid}

  % SOURCE: [Kleiman], "The Picard scheme", \S 2 (FGA Explained Ch.~9 \S 9.2),
  % Definitions df:aPf + df:Pfs (the absolute and relative Picard functors are
  % abelian-group-valued; the relative one is defined as a quotient of abelian
  % groups)
  % (read from references/kleiman-picard-src/kleiman-picard.tex, L1274-L1318).
  % SOURCE QUOTE:
  % "\begin{dfn}\label{df:aPf}
  %  The {\it absolute Picard functor} $\Pic_X$ is the functor from  the
  %  category of (locally Noetherian) $S$-schemes $T$ to the category of
  % abelian groups defined by the formula
  % 	$$\Pic_X(T):=\Pic(X_T).$$
  % \end{dfn}
  %  [\dots]
  % \begin{dfn}\label{df:Pfs}
  %   The {\it relative Picard functor} $\Pic_{X/S}$ is defined by
  % 	$$\Pic_{X/S}(T):=\Pic(X_T)\big/\Pic(T).$$
  %  Denote its associated sheaves in the Zariski, \'etale, and fppf
  % topologies  by
  % 	$$\Pic_{(X/S)\zar},\ \Pic_{(X/S)\et},\
  % 	\Pic_{(X/S)\fppf}.$$
  %  \end{dfn}"
  \textit{Source: [Kleiman], ``The Picard scheme'', \S 2,
  Defs.~df:aPf + df:Pfs (the target category is the category of abelian
  groups; the relative Picard functor is a quotient of abelian groups).}
  Let \(C/k\) be a smooth proper geometrically integral curve over a
  field \(k\), let \(\pi_C : C \to \Spec k\) be the structure morphism,
  and let \(T\) be a \(k\)-scheme with structure morphism
  \(\pi_T : T \to \Spec k\). The set
  \[
    \Pic^{\sharp}_{C/k}(T) \;\;:=\;\; \Pic(C \times_k T) \,/\, \pi_T^*\Pic(T)
  \]
  of \cref{thm:relative_pic_quotient_well_defined} carries a canonical
  abelian-group structure with addition
  \(\bigl[\,\mathcal{L}\,\bigr] + \bigl[\,\mathcal{L}'\,\bigr]
   := \bigl[\,\mathcal{L} \otimes_{C \times_k T} \mathcal{L}'\,\bigr]\),
  neutral element \(\bigl[\,\mathcal{O}_{C \times_k T}\,\bigr]\), and
  inverse
  \(-\bigl[\,\mathcal{L}\,\bigr] := \bigl[\,\mathcal{L}^{-1}\,\bigr]\),
  where \(\otimes_{C \times_k T}\) and \((-)^{-1}\) are the operations
  supplied by the substrate of \cref{def:scheme_modules_tensorobj}
  restricted to the \(\texttt{LineBundle.OnProduct}\,\pi_C\,\pi_T\) carrier
  via \cref{lem:tensorobj_lift_onproduct}. With this structure the quotient
  map
  \(\Pic(C \times_k T) \twoheadrightarrow \Pic^{\sharp}_{C/k}(T)\) is a
  surjective group homomorphism with kernel exactly \(\pi_T^*\Pic(T)\),
  in agreement with \cref{lem:rel_pic_sharp_groupoid}.
\end{theorem}

\begin{proof}
  \uses{lem:tensorobj_lift_onproduct, lem:pullback_compatible_with_tensorobj, def:scheme_modules_isinvertible, lem:tensorobj_isoclass_commgroup, thm:relative_pic_quotient_well_defined, lem:rel_pic_sharp_groupoid}
  Write \(\Pic(C \times_k T)\) for the set of isomorphism classes of objects
  of \(\texttt{LineBundle.OnProduct}\,\pi_C\,\pi_T\) (line bundles on
  \(C \times_k T\)), with the operation
  \([L] + [L'] := [L \otimes_{C \times_k T} L']\). This operation is
  well-defined on isomorphism classes: an isomorphism \(L \cong L_0\) tensors
  to an isomorphism \(L \otimes L' \cong L_0 \otimes L'\) by functoriality of
  \(\otimes\) (\cref{lem:scheme_modules_tensorobj_functoriality}), and
  symmetrically in the second argument. The abelian-group structure on
  \(\Pic(C \times_k T)\) is then exactly the \emph{units group of the
  commutative monoid of \(\otimes\)-iso-classes}
  (\cref{lem:tensorobj_isoclass_commgroup}): multiplication
  \([L]\cdot[L'] = [L \otimes_{C\times_k T} L']\) is well-defined and
  commutative-monoidal by the coherence isomorphisms
  \cref{lem:tensorobj_assoc_iso,lem:tensorobj_unit_iso,lem:tensorobj_comm_iso},
  and an iso-class is a unit precisely when its carrier is
  \(\otimes\)-invertible (\cref{def:scheme_modules_isinvertible}), the inverse
  being carried by the predicate. Each axiom is a \emph{proposition} --- the
  existence of an isomorphism --- so no monoidal-coherence datum (pentagon,
  triangle, hexagon) is invoked; the \(\AddCommGroup\) on
  \(\Pic(C \times_k T)\) is the additive form of this units group. This mirrors
  Mathlib's own Picard group: \(\mathtt{CommRing.Pic}\,R\) carries its
  \(\mathtt{CommGroup}\) (\(\mathtt{instCommGroupPic}\),
  \texttt{Mathlib.RingTheory.PicardGroup}) transported from
  \(\mathrm{Units}\bigl(\mathrm{Skeleton}(\mathtt{ModuleCat}\,R)\bigr)\); here
  the same \(\mathrm{Units}(\cdots)\) shape is transported to the scheme-level
  \(\otimes_X\). (The scheme-level Picard group is genuinely absent from
  Mathlib --- \texttt{Mathlib.RingTheory.PicardGroup} records ``connect to
  invertible sheaves'' as an open task --- so the project supplies it lightly.)
  The invertibility-predicate carrier is the project's \(\mathtt{IsInvertible}\)
  (\cref{def:scheme_modules_isinvertible}), the scheme-level analogue of
  Mathlib's \(\mathtt{Module.Invertible}\); the geometric predicate
  \(\mathtt{IsLocallyTrivial}\) is retained separately and is connected to
  \(\mathtt{IsInvertible}\) only by an off-critical-path equivalence (only the
  geometric \(\mathtt{IsInvertible} \Leftarrow \mathtt{IsLocallyTrivial}\)-style
  direction is nontrivial, and it is not needed here).

  By \cref{lem:pullback_compatible_with_tensorobj}, the pullback functor
  \(\pi_T^{*} : \texttt{LineBundle}\,T \to
   \texttt{LineBundle.OnProduct}\,\pi_C\,\pi_T\) carries tensor products to
  tensor products and the unit to the unit, so the induced map on isomorphism
  classes \(\pi_T^{*} : \Pic(T) \to \Pic(C \times_k T)\) is a group
  homomorphism, and its image is a subgroup.

  By \cref{thm:relative_pic_quotient_well_defined}, the underlying set of
  \(\Pic^{\sharp}_{C/k}(T)\) is the set of cosets of this subgroup. The
  quotient of an abelian group by a subgroup is itself an abelian group, via
  the standard \(\mathtt{QuotientAddGroup}\) construction, and the quotient
  map is a surjective group homomorphism whose kernel is exactly the subgroup
  (Stacks tag 01CR for the Picard-group special case; the general
  quotient-of-abelian-groups statement is in any algebra textbook). Putting
  the pieces together gives the \(\AddCommGroup\) instance on
  \(\Pic^{\sharp}_{C/k}(T)\) of the statement, and identifies it with the
  structure of \cref{lem:rel_pic_sharp_groupoid}.
\end{proof}

\section{Out of scope}
\label{sec:tensorobj_out_of_scope}

The following are downstream of this chapter and live in dedicated
sibling chapters; this chapter does \emph{not} touch them:
\begin{itemize}
  \item \textbf{The sheafified relative Picard functor and its abelian-group target}
    (\texttt{rel\_pic\_etale\_sheafification}, \texttt{rel\_pic\_etale\_sheaf\_group\_structure}) --- once
    \cref{thm:rel_pic_addcommgroup_via_tensorobj} is established, the
    \(\AddCommGroup\)-valued presheaf, its \'etale sheafification, and
    the sheafification preserving the abelian-group target close
    against it (chapter \cref{chap:Picard_RelPicFunctor}); none of those
    declarations is the responsibility of this chapter.
  \item \textbf{Representability of \(\Pic^\sharp_{(C/k)\et}\) by a scheme}
    (A.2.c, the FGA Picard-representability chapter (parent scope)) --- consumes the
    full \(\AddCommGroup\)-valued sheafified presheaf;
    representability is downstream of \texttt{rel\_pic\_etale\_sheaf\_group\_structure}.
  \item \textbf{Identity component \(\Pic^0_{C/k}\)} (A.3,
    the Picard identity-component chapter (parent scope)) --- once representability is
    in hand, the identity component is cut out by Hilbert polynomials;
    that is independent of the tensor-product substrate of this chapter.
  \item \textbf{Albanese universal property} (A.4,
    the Albanese universal-property chapter (parent scope)) --- consumes the full Route~A stack
    downstream of A.3.
  \item \textbf{Mathlib upstream PR of \(\Scheme.\mathtt{Modules}.\mathtt{tensorObj}\)} ---
    the substrate of this chapter is project-side, intentionally; a
    parallel Mathlib upstream PR is welcome but not on the critical
    path. Once the substrate ships in Mathlib, the
    \texttt{TensorObjSubstrate.lean} file can be slimmed to a single
    import + namespace re-export, but no downstream Lean file changes.
\end{itemize}

\section{Internal-consistency check}
\label{sec:tensorobj_consistency_check}

The Lean-pinned declaration blocks of this chapter form a connected
\texttt{\textbackslash uses\{\dots\}} chain whose leaves all resolve
either inside this chapter, inside chapter
\cref{chap:Picard_LineBundlePullback} (A.1.b on disk), or inside
chapter \cref{chap:Picard_RelPicFunctor} (A.1.c on disk):
\begin{itemize}
  \item \cref{def:scheme_modules_tensorobj} introduces the substrate
    operation and depends on no prior block.
  \item \cref{lem:scheme_modules_tensorobj_functoriality} uses
    \cref{def:scheme_modules_tensorobj}.
  \item \cref{rem:scheme_modules_monoidal_off_path} uses
    \cref{def:scheme_modules_tensorobj,
          lem:scheme_modules_tensorobj_functoriality, lem:whisker_of_W}; it
    records that, under route~(e), the full monoidal category is \emph{cheap}
    (the output of \(\mathtt{LocalizedMonoidal}\)) and that the group law still
    consumes only existence-of-isomorphism propositions.
  \item \cref{lem:restrictscalars_laxmonoidal} uses
    \cref{def:scheme_modules_tensorobj}; it is a self-contained
    \(\mathtt{CommRingCat}\)-level supplement, off the critical path and
    \emph{not} an ingredient of \cref{lem:tensorobj_restrict_iso} --- nothing
    downstream depends on it.
  \item \cref{def:scheme_modules_isinvertible} (the \(\otimes\)-invertibility
    predicate, carrying the inverse) uses \cref{def:scheme_modules_tensorobj}.
  \item The route~(e) whisker lemmas
    \cref{lem:islocallyinjective_whisker_of_W} (local injectivity of
    \(F \triangleleft g\) for \(g \in J.W\): the PRIMARY route assumes \(F\) locally
    trivial and proves it sheaf-level via \cref{lem:tensorobj_restrict_iso} and the
    unitor, with no stalks; the FALLBACK route, for arbitrary \(F\), is the
    stalkwise argument with residual ingredients d.1/d.2 --- the same (d.2)
    stalk--tensor commutation \cref{lem:stalk_tensor_commutation} that the \emph{live}
    associator obligation \cref{lem:islocallyinjective_whiskerleft_via_stalk} consumes,
    hence on the critical path) and
    \cref{lem:whisker_of_W} (its closed left/right packaging) supply the two
    fields of the instance \texttt{jw\_ismonoidal}. The latter applies Mathlib's
    \(\mathtt{LocalizedMonoidal}\) to the already-monoidal
    \(\mathtt{PresheafOfModules}\,\mathcal{O}_X\) and the sheafification localizer
    to give the full \(\MonoidalCategory\) structure on
    \(\Scheme.\mathtt{Modules}\,X\); it is stated-but-unformalized, pending the
    sole open proof obligation.
  \item The coherence isomorphisms \cref{lem:tensorobj_unit_iso}
    (left/right unitors) and \cref{lem:tensorobj_comm_iso} (braiding) are the cheap
    \(\mathtt{sheafification.mapIso}\) of the presheaf-level coherence data. The
    associator \cref{lem:tensorobj_assoc_iso} is the sheafification transport of the
    presheaf-level associator, whose single open left-whiskering obligation
    \cref{lem:islocallyinjective_whiskerleft_via_stalk} is closed
    \emph{unconditionally} --- for an arbitrary whiskering object, with no flatness and
    no local triviality --- by the d.2 stalk--tensor commutation
    \cref{lem:stalk_tensor_commutation} (\cref{sec:tensorobj_stalk_tensor}).
    None uses \(\mathtt{MonoidalClosed}\) or sectionwise flatness; the associator's
    pinned \(\mathtt{IsLocallyTrivial}\) hypotheses are vestigial (the construction is
    unconditional).
  \item \cref{lem:isiso_sheafification_map_of_W} (sheafification sends a
    morphism whose underlying map is in \(J.W\) to an isomorphism) depends on no
    prior block; it is a closed go/no-go bridge, retained as a supplement (the
    route~(e) associator is the API-derived one, so this bridge is no longer the
    associator's load-bearing step).
  \item \texttt{flat\_whisker\_localizer} (left/right whiskering of a
    sheafification-localizer morphism by a flat object stays in the localizer
    class) uses \cref{def:scheme_modules_tensorobj}; it is a valid closed
    standalone result, superseded on the critical path by the d.2 route of
    \cref{sec:tensorobj_stalk_tensor}
    (\cref{lem:islocallyinjective_whiskerleft_via_stalk}).
  \item \cref{lem:tensorobj_isoclass_commgroup} (the commutative group on
    locally-trivial \(\otimes\)-iso-classes) uses
    \cref{def:scheme_modules_tensorobj,
          def:scheme_modules_isinvertible,
          lem:scheme_modules_tensorobj_functoriality,
          lem:tensorobj_lift_onproduct,
          lem:tensorobj_inverse_invertible,
          lem:tensorobj_assoc_iso,
          lem:tensorobj_unit_iso,
          lem:tensorobj_comm_iso}. This is the group-law engine: the PRIMARY
    construction is the commutative group on locally-trivial iso-classes à la
    \(\mathtt{Module.Invertible}\) / \(\mathtt{CommRing.Pic}\), needing only
    \cref{lem:tensorobj_restrict_iso} and \cref{lem:tensorobj_inverse_invertible}
    on top of the assembled locally-trivial associator/unitors/braiding; the
    \(\mathrm{Units}(\mathrm{Skeleton}(\Scheme.\mathtt{Modules}\,X))\) presentation
    via \texttt{jw\_ismonoidal} is retained as the route~(e) fallback.
  \item Under the PRIMARY locally-trivial group law, \cref{lem:tensorobj_restrict_iso}
    (consumed by the PRIMARY proof of \cref{lem:islocallyinjective_whisker_of_W}) and
    \cref{lem:tensorobj_inverse_invertible} (in the \(\uses\) of
    \cref{lem:tensorobj_isoclass_commgroup}, supplying inverses) are ON the critical
    path. The remaining supplements
    \cref{lem:restrictscalars_laxmonoidal,lem:tensorobj_preserves_locally_trivial}
    use \cref{def:scheme_modules_tensorobj} and are auxiliary. The group law rests on
    \cref{lem:tensorobj_isoclass_commgroup} and \cref{def:scheme_modules_isinvertible}.
  \item \cref{lem:tensorobj_lift_onproduct} restricts \(\otimes\) to the
    locally-trivial carrier \(\texttt{LineBundle.OnProduct}\) on \(C \times_S T\),
    using \cref{def:scheme_modules_tensorobj,lem:tensorobj_preserves_locally_trivial}
    (the Lean \(\mathtt{tensorObjOnProduct}\) applies
    \(\mathtt{tensorObj\_isLocallyTrivial}\) directly).
  \item \cref{lem:pullback_compatible_with_tensorobj} uses
    \cref{lem:tensorobj_lift_onproduct} and implicitly
    \cref{def:pullback_along_projection} from
    \cref{chap:Picard_LineBundlePullback}; it is \emph{not} blocked on
    \cref{lem:tensorobj_restrict_iso}.
  \item \cref{thm:rel_pic_addcommgroup_via_tensorobj} uses
    \cref{lem:tensorobj_lift_onproduct,
          lem:pullback_compatible_with_tensorobj,
          def:scheme_modules_isinvertible,
          lem:tensorobj_isoclass_commgroup,
          thm:relative_pic_quotient_well_defined,
          lem:rel_pic_sharp_groupoid}; the last two are in
    \cref{chap:Picard_LineBundlePullback} (A.1.b) and
    \cref{chap:Picard_RelPicFunctor} (A.1.c) respectively.
\end{itemize}
All \(\uses\) targets resolve either in this chapter, in
\cref{chap:Picard_LineBundlePullback}, or in
\cref{chap:Picard_RelPicFunctor}. No declaration cross-references a
chapter that does not yet exist on disk.

\section{Sheaf internal-hom of \(\mathcal{O}_X\)-modules (the \(\otimes\)-inverse
  infrastructure)}
\label{sec:tensorobj_dual_infra}

The PRIMARY group law (\cref{lem:tensorobj_isoclass_commgroup}) needs a
\(\otimes\)-inverse for every line bundle, supplied by
\cref{lem:tensorobj_inverse_invertible} as the dual
\(L^{-1} = \mathcal{H}om_{\mathcal{O}_X}(L, \mathcal{O}_X)\). The standard route
constructs that dual as a special case of the \emph{sheaf internal hom} of
\(\mathcal{O}_X\)-modules; over a fixed ring the internal hom of \(R\)-modules is the
linear-coyoneda module \(\mathtt{ihom}\,M\,N = (M \to N)\), with dual
\(\mathtt{ihom}\,M\,\mathbf{1} = M^\vee\), but for the \emph{varying} structure sheaf
no such
internal-hom / monoidal-closed structure exists at the
\(\mathtt{PresheafOfModules}\), \(\mathtt{SheafOfModules}\), or general-categorical
level. This section decomposes the construction of the missing object into
individually-formalizable sub-steps, following the slice presentation of the internal
hom; each block names the intended Lean declaration and its dependencies. The
specialisation that closes \cref{lem:tensorobj_inverse_invertible} is recorded in
\cref{rem:dual_discharges_inverse}; an alternative descent-theoretic route is noted in
\cref{rem:dual_via_stack}.

\emph{Status: deferred, off the group's critical path.} The blocks of this section
--- the dual object \(\mathtt{dual}\) and its restriction compatibility
\cref{lem:dual_restrict_iso}, the local-triviality of the dual
\cref{lem:dual_isLocallyTrivial}, the gluing engine
\texttt{sheafofmodules\_hom\_of\_local\_compat}, and the discharge remark
\cref{rem:dual_discharges_inverse} --- together constitute the \emph{deferred}
\(\mathtt{IsInvertible} \Leftrightarrow \mathtt{IsLocallyTrivial}\) bridge: on a
scheme an invertible \(\mathcal{O}_X\)-module is locally free of rank one, and
conversely (Stacks tag 0B8M), so the two carriers agree, and the locally-trivial
carrier's missing inverse (\cref{lem:tensorobj_inverse_invertible}) is recovered
precisely from the dual built here.
% SOURCE: [Stacks Project], Tag 0B8M (Modules on Ringed Spaces, \S Invertible
% modules), Lemma lemma-invertible-is-locally-free-rank-1
% (read from references/stacks-modules.tex, L4159-L4165).
% SOURCE QUOTE:
% "\begin{lemma}
% \label{lemma-invertible-is-locally-free-rank-1}
% Let $(X, \mathcal{O}_X)$ be a ringed space. Any locally free
% $\mathcal{O}_X$-module of rank $1$ is invertible.
% If all stalks $\mathcal{O}_{X, x}$ are local rings, then
% the converse holds as well (but in general this is not the case).
% \end{lemma}"
However, the relative Picard \emph{group law} is now carried on the
\(\otimes\)-invertibility predicate \(\mathtt{IsInvertible}\)
(\cref{sec:tensorobj_pic_carrier}, \texttt{pic\_commgroup}), where the inverse is
the membership witness and \emph{no} dual object is constructed. This section is
therefore needed only by a downstream consumer that specifically requires the
locally-trivial form of an inverse --- not by the group law. The blocks are retained
intact for that future use.

A structural point governs the whole construction. The naive rule
\(U \mapsto \Hom_{\mathcal{O}_X(U)}(M(U), N(U))\) is \emph{contravariant} in the
restriction maps of \(M\) (a restriction \(M(U) \to M(V)\) for \(V \subseteq U\)
induces precomposition \(\Hom(M(V), N(V)) \to \Hom(M(U), \cdot)\) in the wrong
direction), so it is \emph{not} a presheaf in the covariant sense and cannot be
sheafified like the tensor product \cref{def:scheme_modules_tensorobj}. The remedy is
the \emph{slice} formula \(U \mapsto (M|_U \to N|_U)\) --- the module of morphisms of
\emph{restricted} modules over the restricted ring --- whose presheaf functoriality
(restrict a global-over-\(U\) morphism to a global-over-\(V\) morphism, for
\(V \subseteq U\)) \emph{is} covariant. Constructing that slice presheaf, and
proving it satisfies the sheaf condition, is the substance of the build.

\begin{definition}

\leanok
[The global-ring module on morphisms of presheaves of modules (value module)]
  \label{def:presheaf_internal_hom_value}
  \lean{PresheafOfModules.InternalHom.homModule}
  \uses{def:scheme_modules_tensorobj}
  % SOURCE: [Stacks Project], "Modules on Ringed Spaces", \S Internal Hom
  % (\texttt{section-internal-hom}, tag area 01CM)
  % (read from references/stacks-modules.tex, L3508--L3514).
  % SOURCE QUOTE:
  % "abelian groups. In addition, given an element
  % $\varphi \in \Hom_{\mathcal{O}_X|_U}(\mathcal{F}|_U, \mathcal{G}|_U)$
  % and a section $f \in \mathcal{O}_X(U)$ then we can define
  % $f\varphi \in \Hom_{\mathcal{O}_X|_U}(\mathcal{F}|_U, \mathcal{G}|_U)$
  % by either precomposing with multiplication by $f$ on $\mathcal{F}|_U$
  % or postcomposing with multiplication by $f$ on $\mathcal{G}|_U$ (it gives
  % the same result). Hence we in fact get a sheaf of $\mathcal{O}_X$-modules."
  \textit{Source: [Stacks Project], ``Modules on Ringed Spaces'', \S Internal Hom
  (tag area 01CM): the section ring acts on a morphism of restricted modules by
  post-composition with multiplication by the section.}
  Let \(R\) be a presheaf of commutative rings on a base category \(\mathcal{C}\)
  that has a terminal object \(T\), and let \(M, N \in \mathtt{PresheafOfModules}\,R\).
  The abelian group of morphisms \(M \longrightarrow N\) of presheaves of \(R\)-modules
  carries a canonical \(R(T)\)-module structure --- the \emph{morphism module}
  \(\mathtt{homModule}\,M\,N\). A global scalar \(f \in R(T)\) acts on a morphism
  \(\varphi : M \to N\) by post-composition with multiplication by \(f\),
  \[
    f \cdot \varphi \;:=\; m_f^{N} \circ \varphi,
  \]
  where \(m_f^{N} : N \to N\) is the \emph{multiplication-by-\(f\)} endomorphism of
  \(N\): the endomorphism that, on each object \(U \in \mathcal{C}\), is multiplication
  by the restriction \(f|_U \in R(U)\) of \(f\) along the unique morphism \(U \to T\)
  to the terminal object. The module axioms hold because the assignment
  \(f \mapsto m_f^{N}\) is a ring homomorphism \(R(T) \to \mathrm{End}(N)\) into the
  endomorphism ring of \(N\) --- a product \(f \cdot g\) goes to the \emph{composite}
  \(m_f^{N} \circ m_g^{N}\) of the two multiplication endomorphisms, the order in which
  the factors compose being reversed by the fact that the scalar acts by
  \emph{post}-composition --- and composition of morphisms is biadditive in the
  preadditive category \(\mathtt{PresheafOfModules}\,R\). Distributivity and the
  associativity of the action then follow from bilinearity of composition together with
  the homomorphism property of \(f \mapsto m_f^{N}\). This is the genuinely new core of
  the construction: Mathlib supplies the \emph{fixed}-ring internal hom
  \(\mathtt{ihom}\,M\,N = (M \to N)\) over a single base ring, but nothing equipping the
  morphism group with the module structure for the \emph{varying} structure presheaf at
  the \(\mathtt{PresheafOfModules}\) level.
\end{definition}

\begin{definition}

\leanok
[Slice specialisation of the morphism module]
  \label{def:presheaf_internal_hom_slice_value}
  \lean{PresheafOfModules.InternalHom.internalHomObjModule}
  \uses{def:presheaf_internal_hom_value}
  Specialise \cref{def:presheaf_internal_hom_value} to a slice category. For an object
  (open) \(U\) of \(\mathcal{C} = \Opens X\), restricting along the forgetful functor
  \(\mathtt{Over.forget}\,U : \mathcal{C}/U \to \mathcal{C}\) --- the open-immersion
  restriction \(M \mapsto M|_U\) --- carries \(M, N\) to presheaves of modules
  \(M|_U, N|_U\) over the restricted structure ring on the slice \(\mathcal{C}/U\). The
  slice \(\mathcal{C}/U\) has the terminal object \(\mathtt{Over.mk}\,(\mathrm{id}_U)\),
  at which the restricted ring takes the value \(R(U)\); hence
  \cref{def:presheaf_internal_hom_value} makes the morphism group
  \(M|_U \longrightarrow N|_U\) an \(R(U)\)-module, the
  \(\mathtt{internalHomObjModule}\). This \(R(U)\)-module \((M|_U \to N|_U)\) is exactly
  the value \(\mathcal{H}om(M, N)(U)\) of \cref{def:presheaf_internal_hom}.
\end{definition}

\begin{definition}

\leanok
[Presheaf internal hom of \(\mathcal{O}_X\)-modules]
  \label{def:presheaf_internal_hom}
  \lean{PresheafOfModules.InternalHom.internalHom}
  \uses{def:scheme_modules_tensorobj, def:presheaf_internal_hom_value, def:presheaf_internal_hom_slice_value, lem:presheaf_internal_hom_restriction}
  % SOURCE: [Stacks Project], "Modules on Ringed Spaces", \S Internal Hom
  % (\texttt{section-internal-hom}, tag area 01CM)
  % (read from references/stacks-modules.tex, L3500--L3524).
  % SOURCE QUOTE:
  % "Let $(X, \mathcal{O}_X)$ be a ringed space.
  % Let $\mathcal{F}$, $\mathcal{G}$ be $\mathcal{O}_X$-modules.
  % Consider the rule
  % $$
  % U \longmapsto \Hom_{\mathcal{O}_X|_U}(\mathcal{F}|_U, \mathcal{G}|_U).
  % $$
  % It follows from the discussion in Sheaves, Section
  % \ref{sheaves-section-glueing-sheaves} that this is a sheaf of
  % abelian groups. In addition, given an element
  % $\varphi \in \Hom_{\mathcal{O}_X|_U}(\mathcal{F}|_U, \mathcal{G}|_U)$
  % and a section $f \in \mathcal{O}_X(U)$ then we can define
  % $f\varphi \in \Hom_{\mathcal{O}_X|_U}(\mathcal{F}|_U, \mathcal{G}|_U)$
  % by either precomposing with multiplication by $f$ on $\mathcal{F}|_U$
  % or postcomposing with multiplication by $f$ on $\mathcal{G}|_U$ (it gives
  % the same result). Hence we in fact get a sheaf of $\mathcal{O}_X$-modules.
  % We will denote this sheaf
  % $\SheafHom_{\mathcal{O}_X}(\mathcal{F}, \mathcal{G})$.
  % There is a canonical ``evaluation'' morphism
  % $$
  % \mathcal{F}
  % \otimes_{\mathcal{O}_X}
  % \SheafHom_{\mathcal{O}_X}(\mathcal{F}, \mathcal{G})
  % \longrightarrow
  % \mathcal{G}.
  % $$"
  \textit{Source: [Stacks Project], ``Modules on Ringed Spaces'', \S Internal Hom
  (tag area 01CM).}
  Let \(R\) be a presheaf of commutative rings on \(\Opens X\) and let
  \(M, N \in \mathtt{PresheafOfModules}\,R\). The \emph{presheaf internal hom}
  \(\mathcal{H}om(M, N)\) is the presheaf of \(R\)-modules whose value on an open
  \(U\) is the module of morphisms of restricted modules
  \[
    \mathcal{H}om(M, N)(U)
      \;:=\;
    \mathtt{ModuleCat.of}\;\bigl(R(U)\bigr)\;\bigl(M|_U \longrightarrow N|_U\bigr),
  \]
  the \(R(U)\)-module of \(R|_U\)-linear morphisms \(M|_U \to N|_U\) of presheaves of
  modules on \(U\), formed via the open-immersion restriction \(U \hookrightarrow X\).
  This value module is the slice specialisation
  \cref{def:presheaf_internal_hom_slice_value} of the morphism module
  \cref{def:presheaf_internal_hom_value}.

  Explicitly, \(\mathcal{H}om(M, N)\) is the presheaf of \(R\)-modules assembled from
  three pieces. \emph{(a) Value modules.} On each open \(U\) the value is the
  \(R(U)\)-module \((M|_U \to N|_U)\) of \cref{def:presheaf_internal_hom_value} /
  \cref{def:presheaf_internal_hom_slice_value}. \emph{(b) Restriction maps.} For an
  inclusion \(V \subseteq U\) the restriction map
  \(\mathcal{H}om(M, N)(U) \to \mathcal{H}om(M, N)(V)\) is the further-restriction map
  \cref{lem:presheaf_internal_hom_restriction}, \(\varphi \mapsto \varphi|_V\), which
  is additive and semilinear over the ring restriction \(R(g) : R(U) \to R(V)\)
  induced by \(g : V \to U\) --- the compatibility datum that the presheaf-of-modules
  assembly (\(\mathtt{PresheafOfModules.mk}\), via \(\mathtt{ofPresheaf}\)) consumes.
  \emph{(c) Functoriality.} Restriction along the identity \(U \to U\) is the identity
  map, and restriction along a composite \(V \xrightarrow{g} U \xrightarrow{h} W\) of
  inclusions \(V \subseteq U \subseteq W\) (the morphism \(g : V \to U\) followed by
  \(h : U \to W\) in \(\mathcal{C} = \Opens X\)) is the composite of the two
  restriction maps \(\mathcal{H}om(M,N)(W) \to \mathcal{H}om(M,N)(U) \to
  \mathcal{H}om(M,N)(V)\); both identity and composition are inherited from the
  functoriality of further restriction, since \(\mathtt{Over.map}\) is a functor. These three pieces
  make \(\mathcal{H}om(M, N)\) a genuine, covariant presheaf of \(R\)-modules; this
  assembled presheaf is the construction target \(\mathtt{PresheafOfModules.internalHom}\).

  The slice formula is forced by contravariance: the pointwise rule
  \(U \mapsto \Hom_{R(U)}(M(U), N(U))\) varies \emph{against} the restriction maps of
  \(M\) and so is not a covariant presheaf, whereas the module of morphisms of
  \emph{restricted} objects \(M|_U \to N|_U\) restricts covariantly. Thus, unlike the
  tensor product \cref{def:scheme_modules_tensorobj} --- which is covariant in the
  restriction maps and sheafifies directly --- the internal hom must be built through
  the slice presentation, and there is no presheaf-level dual to sheafify naively.
\end{definition}

\begin{lemma}

\leanok
[Restriction map of the presheaf internal hom]
  \label{lem:presheaf_internal_hom_restriction}
  \lean{PresheafOfModules.InternalHom.restrictionMap}
  \uses{def:presheaf_internal_hom_value, def:presheaf_internal_hom_slice_value}
  % SOURCE: [Stacks Project], "Modules on Ringed Spaces", \S Internal Hom
  % (\texttt{section-internal-hom}, tag area 01CM)
  % (read from references/stacks-modules.tex, L3502--L3508).
  % SOURCE QUOTE:
  % "Consider the rule
  % $$
  % U \longmapsto \Hom_{\mathcal{O}_X|_U}(\mathcal{F}|_U, \mathcal{G}|_U).
  % $$
  % It follows from the discussion in Sheaves, Section
  % \ref{sheaves-section-glueing-sheaves} that this is a sheaf of
  % abelian groups."
  \textit{Source: [Stacks Project], ``Modules on Ringed Spaces'', \S Internal Hom
  (tag area 01CM): the rule \(U \mapsto \Hom(\mathcal{F}|_U, \mathcal{G}|_U)\), with its
  further-restriction maps, is a (pre)sheaf of abelian groups.}
  Let \(g : V \to U\) be a morphism of \(\mathcal{C} = \Opens X\) (an inclusion
  \(V \subseteq U\) of opens), and let \(M, N \in \mathtt{PresheafOfModules}\,R\). There
  is a restriction map
  \[
    \mathcal{H}om(M, N)(U) \longrightarrow \mathcal{H}om(M, N)(V),
    \qquad \varphi \longmapsto \varphi|_V,
  \]
  the \(\mathtt{restrictionMap}\), sending a morphism of restricted modules
  \(M|_U \to N|_U\) to its further restriction \(M|_V \to N|_V\). Concretely, the
  functor \(\mathtt{Over.map}\,g : \mathcal{C}/V \to \mathcal{C}/U\) realises
  ``restrict from \(U\) to \(V\)'', and \(\varphi|_V\) is \(\varphi\) pulled back along
  it --- the further restriction of the same morphism of presheaves of modules. This
  map is additive, and it is \emph{semilinear} over the ring restriction
  \(R(g) : R(U) \to R(V)\): for \(f \in R(U)\),
  \[
    (f \cdot \varphi)\big|_V \;=\; R(g)(f) \cdot \bigl(\varphi|_V\bigr).
  \]
  This semilinearity is the compatibility datum that the presheaf assembly of
  \cref{def:presheaf_internal_hom} (through \(\mathtt{PresheafOfModules.ofPresheaf}\) /
  \(\mathtt{mk}\)) requires.
\end{lemma}

\begin{proof}
  \uses{def:presheaf_internal_hom_value}
  \leanok
  Further restriction along \(\mathtt{Over.map}\,g\) is functorial, so it carries a
  morphism \(M|_U \to N|_U\) to a morphism \(M|_V \to N|_V\) and respects identities and
  composites; in particular the induced map on morphism groups is additive. For the
  semilinearity, recall that the global-scalar action of
  \cref{def:presheaf_internal_hom_value} acts on \(\varphi\) by post-composition with
  the multiplication-by-\(f\) endomorphism \(m_f^{N}\); restricting that endomorphism
  from \(U\) to \(V\) is multiplication by \(f|_V = R(g)(f)\) on \(N|_V\). Since further
  restriction commutes with composition of morphisms, restricting \(f \cdot \varphi =
  m_f^{N} \circ \varphi\) yields \(m_{R(g)(f)}^{N|_V} \circ (\varphi|_V) = R(g)(f) \cdot
  (\varphi|_V)\), which is the asserted compatibility.
\end{proof}

\begin{definition}

\leanok
[Presheaf dual against the structure presheaf]
  \label{def:presheaf_dual}
  \lean{PresheafOfModules.dual}
  \uses{def:presheaf_internal_hom}
  With \(R\) and \(M\) as in \cref{def:presheaf_internal_hom}, the
  \emph{presheaf dual} of \(M\) is the internal hom into the unit (structure)
  presheaf,
  \[
    M^\vee \;:=\; \mathcal{H}om(M, R),
  \]
  i.e.\ \(M^\vee(U) = \mathtt{ModuleCat.of}\,(R(U))\,(M|_U \to R|_U)\), the
  \(R(U)\)-module of \(R|_U\)-linear functionals on \(M|_U\). This is the
  presheaf-of-modules analogue of the linear dual
  \(\mathtt{Module.Dual}\,R\,M = (M \to R)\) that serves as the \(\otimes\)-inverse of
  an invertible module over a fixed ring; here the fixed-ring dual is taken
  open-by-open and assembled by the slice presentation of
  \cref{def:presheaf_internal_hom}.
\end{definition}

\begin{lemma}

\leanok
[Evaluation/contraction of the internal hom]
  \label{lem:internal_hom_eval}
  \lean{PresheafOfModules.internalHomEval}
  % NOTE (iter-224): \texttt{PresheafOfModules.internalHomEval} is CLOSED axiom-clean
  % (\(\{\texttt{propext}, \texttt{Classical.choice}, \texttt{Quot.sound}\}\)). The full
  % morphism assembles the per-open contraction \texttt{internalHomEvalApp}
  % (the \(R(U)\)-bilinear \((M|_U) \otimes_{R(U)} (M|_U \to R|_U) \to R(U)\),
  % \(s \otimes \varphi \mapsto \varphi(s)\)) sectionwise; naturality follows from the
  % evaluate/restrict-commutation argument in the proof body. The iter-222/223 ``whnf
  % heartbeat-bomb'' obstacle was a stale Mathlib-version artifact and is gone.
  \uses{def:presheaf_dual, def:scheme_modules_tensorobj}
  % SOURCE: [Stacks Project], "Modules on Ringed Spaces", \S Internal Hom
  % (\texttt{section-internal-hom}, tag area 01CM)
  % (read from references/stacks-modules.tex, L3517--L3524).
  % SOURCE QUOTE:
  % "There is a canonical ``evaluation'' morphism
  % $$
  % \mathcal{F}
  % \otimes_{\mathcal{O}_X}
  % \SheafHom_{\mathcal{O}_X}(\mathcal{F}, \mathcal{G})
  % \longrightarrow
  % \mathcal{G}.
  % $$"
  \textit{Source: [Stacks Project], ``Modules on Ringed Spaces'', \S Internal Hom
  (tag area 01CM).}
  With \(R\) and \(M\) as above, there is a canonical \emph{evaluation} (contraction)
  morphism of presheaves of \(R\)-modules
  \[
    \mathrm{ev}_M \;\colon\; M \otimes_R M^\vee \;\longrightarrow\; R,
    \qquad s \otimes \phi \longmapsto \phi(s),
  \]
  the counit of the tensor--hom adjunction
  \(\mathcal{H}om(M \otimes_R -, R) \cong \mathcal{H}om(-, M^\vee)\) specialised at
  \(R\). On each open \(U\) it is the open-by-open contraction
  \((M|_U) \otimes_{R|_U} (M|_U \to R|_U) \to R|_U\), \(s \otimes \phi \mapsto
  \phi(s)\), and these are compatible with restriction, so they assemble to a morphism
  of presheaves of modules. The tensor product is the presheaf tensor underlying
  \cref{def:scheme_modules_tensorobj}.
\end{lemma}

\begin{proof}
  \uses{def:presheaf_dual, def:scheme_modules_tensorobj, lem:presheaf_internal_hom_restriction}
  \leanok
  On a fixed open \(U\), the value \(M^\vee(U) = (M|_U \to R|_U)\) is the module of
  \(R|_U\)-linear functionals on \(M|_U\), and the evaluation
  \(s \otimes \phi \mapsto \phi(s)\) is the standard contraction
  \((M|_U) \otimes_{R(U)} (M|_U \to R|_U) \to R(U)\) bilinear over \(R(U)\), hence a
  well-defined \(R(U)\)-linear map out of the tensor product. For \(V \subseteq U\)
  the square relating the \(U\)- and \(V\)-contractions commutes because evaluating
  then restricting equals restricting both arguments then evaluating
  (\(\phi(s)|_V = (\phi|_V)(s|_V)\)); hence the open-by-open contractions are natural
  in the restriction maps and define a morphism of presheaves of \(R\)-modules.
  Concretely, the morphism is defined by specifying, on each open \(U\), the per-open
  contraction \((M|_U) \otimes_{R(U)} (M|_U \to R|_U) \to R(U)\), \(s \otimes \phi \mapsto \phi(s)\).
  The single remaining datum is the naturality square, which reduces to the
  evaluate/restrict-commutation identity \(\phi(s)|_V = (\phi|_V)(s|_V)\) above: the
  restriction \(\phi|_V\) of a functional inside the dual \(M^\vee\) is its restriction
  to the open \(V\), and restriction is functorial, so evaluating the restricted
  functional on the restricted section agrees with restricting the value of the
  evaluation. This is the same restriction-functoriality argument already used to make
  the internal-hom restriction maps functorial in
  \cref{lem:presheaf_internal_hom_restriction}. This is the counit of the tensor--hom
  adjunction of \cref{def:presheaf_internal_hom} specialised to the unit presheaf.
\end{proof}

\begin{lemma}


\leanok
[The internal hom is a sheaf; the sheaf-level dual]
  \label{lem:internal_hom_isSheaf}
  \lean{AlgebraicGeometry.Scheme.Modules.dual}
  \uses{def:presheaf_dual, lem:internal_hom_eval}
  % SOURCE: [Stacks Project], "Modules on Ringed Spaces", \S Internal Hom
  % (\texttt{section-internal-hom}, tag area 01CM)
  % (read from references/stacks-modules.tex, L3502--L3514).
  % SOURCE QUOTE:
  % "Consider the rule
  % $$
  % U \longmapsto \Hom_{\mathcal{O}_X|_U}(\mathcal{F}|_U, \mathcal{G}|_U).
  % $$
  % It follows from the discussion in Sheaves, Section
  % \ref{sheaves-section-glueing-sheaves} that this is a sheaf of
  % abelian groups. ... Hence we in fact get a sheaf of $\mathcal{O}_X$-modules."
  \textit{Source: [Stacks Project], ``Modules on Ringed Spaces'', \S Internal Hom
  (tag area 01CM).}
  Let \(X\) be a scheme, \(R = X.\mathtt{ringCatSheaf}\), and let
  \(M, N \in \Scheme.\mathtt{Modules}\,X\) with \(N\) a sheaf of modules. Then the
  presheaf internal hom \(\mathcal{H}om(M, N)\) of \cref{def:presheaf_internal_hom}
  satisfies the sheaf condition: for an open cover \(\{U_i\}\) of \(U\), a family of
  morphisms \(M|_{U_i} \to N|_{U_i}\) agreeing on the overlaps
  \(U_i \cap U_j\) glues to a unique morphism \(M|_U \to N|_U\), because \(N\) is a
  sheaf and morphisms of sheaves of modules are determined and gluable section-wise.
  Hence \(\mathcal{H}om(M, N)\) descends to an object of \(\Scheme.\mathtt{Modules}\,X\).
  Specialising to \(N = \mathcal{O}_X = \mathtt{SheafOfModules.unit}\,R\) gives the
  \emph{sheaf-level dual}
  \[
    \mathtt{dual}\,M \;:=\; \mathcal{H}om_{\mathcal{O}_X}(M, \mathcal{O}_X)
    \;\in\; \Scheme.\mathtt{Modules}\,X,
  \]
  and the evaluation \cref{lem:internal_hom_eval} descends to a morphism of sheaves of
  modules \(M \otimes_X \mathtt{dual}\,M \to \mathcal{O}_X\).
\end{lemma}

\begin{proof}
  \uses{def:presheaf_dual, lem:internal_hom_eval}
  \leanok
  Fix an open \(U\) and an open cover \(\{U_i\}\) of \(U\). A morphism
  \(M|_U \to N|_U\) is a section of \(\mathcal{H}om(M, N)\) over \(U\); restricting it
  to the \(U_i\) gives a compatible family. Conversely, given morphisms
  \(\varphi_i : M|_{U_i} \to N|_{U_i}\) with \(\varphi_i|_{U_i \cap U_j} =
  \varphi_j|_{U_i \cap U_j}\), section-wise gluing in the sheaf \(N\) produces, for
  each open \(W \subseteq U\) and section \(s \in M(W)\), a well-defined section of
  \(N(W)\) by gluing the \(\varphi_i(s|_{W \cap U_i})\) (these agree on overlaps by
  hypothesis, and \(N\) is a sheaf so the glued section is unique). The resulting
  assignment is \(R(W)\)-linear and natural in \(W\), hence a morphism
  \(M|_U \to N|_U\) restricting to the \(\varphi_i\); uniqueness is the separatedness
  of \(N\). This is exactly the sheaf condition for \(\mathcal{H}om(M, N)\), so it is a
  sheaf of \(\mathcal{O}_X\)-modules. Taking \(N = \mathcal{O}_X\) yields
  \(\mathtt{dual}\,M\), and the evaluation morphism of \cref{lem:internal_hom_eval},
  being a morphism of presheaves into the sheaf \(\mathcal{O}_X\), is a morphism of
  the corresponding sheaves of modules.
\end{proof}

\begin{lemma}

\leanok
[Base change along a ring isomorphism commutes with the linear dual (H2$'$)]
  \label{lem:restrictscalars_ringiso_dualequiv}
  \lean{restrictScalarsRingIsoDualEquiv}
  \uses{lem:restrictscalars_ringiso_tensorequiv}
  Let \(e : R \simeq S\) be an isomorphism of commutative rings, and let \(M\) be an
  \(S\)-module. Restriction of scalars along \(e\) commutes with the formation of the
  linear dual: the canonical map
  \[
    \mathtt{restrictScalars}_e\bigl(M \to_S S\bigr)
      \;\xrightarrow{\ \sim\ }\;
    \bigl(\mathtt{restrictScalars}_e M \to_R R\bigr),
    \qquad \varphi \;\longmapsto\; e^{-1} \circ \varphi,
  \]
  is an \(R\)-linear equivalence, with inverse \(\psi \mapsto e \circ \psi\). On both
  sides the \(R\)-module structure is the \(S\)-module structure of \(M\) pulled back
  along \(e\), so that \(r \cdot m = e(r) \cdot m\): the left-hand side is the
  \(S\)-linear dual \((M \to_S S)\) regarded as an \(R\)-module by restriction, and the
  right-hand side is the \(R\)-linear dual of the restricted module
  \(\mathtt{restrictScalars}_e M\) into the restricted ground ring. This is the dual
  analogue of \cref{lem:restrictscalars_ringiso_tensorequiv}: it is to the internal hom
  / dual exactly what that lemma is to the tensor product. It is the \emph{H2$'$
  ingredient} of the restriction-compatibility isomorphism of the C-bridge
  \cref{lem:dual_isLocallyTrivial} --- the ring-iso/dual commutation that, composed with
  the open-immersion pushforward comparison H1, identifies the dual of a restriction
  with the restriction of the dual.
\end{lemma}

\begin{proof}

\leanok
  The two assignments \(\varphi \mapsto e^{-1} \circ \varphi\) and
  \(\psi \mapsto e \circ \psi\) are mutually inverse, since \(e \circ e^{-1} =
  \mathrm{id}_S\) and \(e^{-1} \circ e = \mathrm{id}_R\) as ring homomorphisms, so the
  two composites are the identity on functionals. Each assignment is \(R\)-linear:
  additivity is immediate, and for a scalar \(r \in R\) the defining relation
  \(r \cdot m = e(r) \cdot m\) of the pulled-back module structure makes each linearity
  check a single rewrite --- \(e^{-1}\) (resp.\ \(e\)) intertwines the \(R\)-action on
  the source with the \(R\)-action on the target precisely because \(e\) and \(e^{-1}\)
  are inverse ring homomorphisms identifying \(R\) with \(S\). The forward map is
  therefore an \(R\)-linear bijection, i.e.\ an \(R\)-linear equivalence; packaging the
  two one-sided maps with the inverse identities gives
  \(\mathtt{restrictScalarsRingIsoDualEquiv}\).
\end{proof}

\begin{definition}

\leanok
[Precomposition equivalence on dual sections]
  \label{def:presheaf_dual_precomp_equiv}
  \lean{PresheafOfModules.dualPrecompEquiv}
  \uses{def:presheaf_dual}
  Let \(R\) be a presheaf of commutative rings and let \(e : M \xrightarrow{\sim} M'\)
  be an isomorphism of presheaves of \(R\)-modules. On a fixed open \(U\), precomposition
  with \(e\) gives an \(R(U)\)-linear equivalence between the dual sections
  \[
    \mathtt{dualPrecompEquiv}\,e\,U \;\colon\;
      \bigl(M'|_U \to R|_U\bigr) \;\xrightarrow{\ \sim\ }\; \bigl(M|_U \to R|_U\bigr),
      \qquad \varphi \;\longmapsto\; \varphi \circ e|_U,
  \]
  with inverse \(\psi \mapsto \psi \circ (e^{-1})|_U\). This is the section-level core of
  the contravariant functoriality of the presheaf dual \(M \mapsto M^\vee\)
  (\cref{def:presheaf_dual}) in isomorphisms: precomposing a functional on \(M'\) with
  the comparison isomorphism produces a functional on \(M\), \(R(U)\)-linearly and
  invertibly.
\end{definition}

\begin{definition}

\leanok
[The presheaf dual carries an iso \(M \cong M'\) to an iso \(M'^\vee \cong M^\vee\)]
  \label{def:presheaf_dual_iso_of_iso}
  \lean{PresheafOfModules.dualIsoOfIso}
  \uses{def:presheaf_dual, def:presheaf_dual_precomp_equiv}
  An isomorphism \(e : M \xrightarrow{\sim} M'\) of presheaves of \(R\)-modules induces a
  canonical isomorphism of presheaf duals
  \[
    \mathtt{dualIsoOfIso}\,e \;\colon\; M'^\vee \;\xrightarrow{\ \sim\ }\; M^\vee
  \]
  --- note the reversal of direction, the dual being contravariant --- assembled
  open-by-open from the precomposition equivalences
  \(\mathtt{dualPrecompEquiv}\,e\,U\) of \cref{def:presheaf_dual_precomp_equiv}.
  Naturality in \(U\), i.e.\ compatibility of these sectionwise equivalences with
  restriction, is the routine check that precomposition by \(e\) commutes with
  restricting functionals.
\end{definition}

\begin{definition}


\leanok
[The sheaf-level dual carries an iso \(M \cong M'\) to an iso \(\mathtt{dual}\,M' \cong \mathtt{dual}\,M\)]
  \label{def:scheme_modules_dual_iso_of_iso}
  \lean{AlgebraicGeometry.Scheme.Modules.dualIsoOfIso}
  \uses{lem:internal_hom_isSheaf, def:presheaf_dual_iso_of_iso}
  Let \(X\) be a scheme and \(e : M \xrightarrow{\sim} M'\) an isomorphism in
  \(\Scheme.\mathtt{Modules}\,X\). Then there is a canonical isomorphism of sheaf-level
  duals
  \[
    \mathtt{dualIsoOfIso}\,e \;\colon\; \mathtt{dual}\,M' \;\xrightarrow{\ \sim\ }\;
      \mathtt{dual}\,M,
  \]
  obtained by sheafifying the presheaf isomorphism \(\mathtt{dualIsoOfIso}\,e\) of
  \cref{def:presheaf_dual_iso_of_iso} through the sheaf-level dual of
  \cref{lem:internal_hom_isSheaf}. It is the dual (contravariant) analogue of the
  functoriality of the tensor product in isomorphisms: where an isomorphism of a tensor
  factor induces an isomorphism of tensor products, an isomorphism \(M \cong M'\)
  induces an isomorphism of duals in the reverse direction.
\end{definition}

\begin{lemma}




\leanok
[The per-section dual transport \(\mathtt{sliceDualTransport}\)]
  \label{lem:slice_dual_transport}
  \lean{AlgebraicGeometry.Scheme.Modules.sliceDualTransport}
  \uses{lem:internal_hom_isSheaf, lem:restrictscalars_ringiso_dualequiv, def:dual_unit_ring_swap, lem:isiso_eps_restrictscalars_appiso, def:dual_unit_iso_gen, lem:image_preimage_of_le, def:dual_unit_ring_swap_hom, lem:isiso_eps_restrictscalars_appiso_hom, lem:slice_dual_transport_inv, def:dual_unit_ring_swap_inv, lem:dual_unit_ring_swap_comp_inv, lem:dual_unit_ring_swap_inv_comp, lem:slice_dual_transport_naturality_apply, lem:appiso_hom_naturality_apply}
  Let \(f : Y \hookrightarrow X\) be an open immersion of schemes and let
  \(M \in \Scheme.\mathtt{Modules}\,X\). Write \(\alpha\) for the open-immersion
  structure-presheaf natural transformation with components
  \(\alpha_U = (f.\mathtt{appIso}\,U).\mathrm{inv}\), a \emph{\(\mathtt{CommRingCat}\)}-level
  ring isomorphism, and \(\beta = \mathtt{whiskerRight}\,\alpha\,
  (\mathtt{forget}_2\,\mathtt{CommRingCat}\,\mathtt{RingCat})\) for its
  \(\mathtt{RingCat}\)-level shadow, with component
  \(\beta_V \colon \mathcal{O}_Y(V) \xrightarrow{\sim} \mathcal{O}_X(fV)\)
  (\(fV := f.\mathtt{opensFunctor}(V)\)). Then for every open \(V \subseteq Y\) there is an
  \(\mathcal{O}_Y(V)\)-linear isomorphism
  \[
    \bigl((\mathtt{pushforward}\,\beta).\mathtt{obj}\,(\mathtt{dual}\,M.\mathtt{val})\bigr)(V)
      \;\xrightarrow{\ \sim\ }\;
    \bigl(\mathtt{dual}\,((\mathtt{pushforward}\,\beta).\mathtt{obj}\,M.\mathtt{val})\bigr)(V),
  \]
  realised as a single \(\mathtt{LinearEquiv.toModuleIso}\) that packages \emph{both} leg (A)
  --- the slice-Hom base-change reindexing across \(f.\mathtt{opensFunctor}\) --- and leg (B)
  --- the unit codomain ring-iso swap. The naturality of the section family in \(W\) has
  \emph{two} distinct parts: the underlying base-morphism uniqueness in
  \((\mathtt{Over}\,V.\mathrm{unop})^{\mathrm{op}}\) is \(\mathtt{Subsingleton.elim}\) (a thin
  poset has at most one inclusion between objects), but the accompanying equation of
  \(\mathcal{O}_Y(V)\)-module maps is a genuine, separate obligation: it is the
  \(\varepsilon\)-naturality of \(\mathtt{restrictScalars}\) along the structure-ring iso
  \(\beta_W\), which \(\mathtt{Subsingleton.elim}\) does \emph{not} discharge.
\end{lemma}

\begin{proof}
  \uses{lem:internal_hom_isSheaf, lem:restrictscalars_ringiso_dualequiv, def:dual_unit_ring_swap, lem:isiso_eps_restrictscalars_appiso, def:dual_unit_iso_gen, lem:image_preimage_of_le, def:dual_unit_ring_swap_hom, lem:isiso_eps_restrictscalars_appiso_hom, lem:slice_dual_transport_inv, def:dual_unit_ring_swap_inv, lem:dual_unit_ring_swap_comp_inv, lem:dual_unit_ring_swap_inv_comp, lem:slice_dual_transport_inv_app_apply, lem:slice_dual_transport_naturality_apply, lem:appiso_hom_naturality_apply}
  \leanok
  The isomorphism is produced from an \(\mathcal{O}_Y(V)\)-linear equivalence between the
  left-hand module --- the \(\mathcal{O}_Y(V)\)-module inherited from
  \(\mathtt{pushforward}\,\beta\) --- and the right-hand \(\mathtt{internalHomObjModule}\),
  assembled via \(\mathtt{LinearEquiv.toModuleIso}\), with both module structures
  identified explicitly.

  \emph{Leg (A): the slice-Hom reindexing, built categorically.} A dual section
  \(\varphi\) of the left-hand value is a morphism \(\mathtt{restr}_{fV}\,M.\mathtt{val} \to
  \mathtt{restr}_{fV}\,\mathbf{1}_X\) over \(\mathtt{Over}\,(fV)\). Its image is the dual
  section over \(\mathtt{Over}\,V\) whose component at \(W \le V\) (with image
  \(fW := f.\mathtt{opensFunctor}(W)\)) is the categorical image
  \[
    (\mathtt{restrictScalars}\,\beta_W).\mathrm{map}\,
      \bigl(\varphi_{(\mathtt{Over.mk}\,(f.\mathtt{opensFunctor}(W \le V)))}\bigr),
  \]
  i.e.\ the component of \(\varphi\) at the \(f\)-image of \(W\), reindexed across
  \(f.\mathtt{opensFunctor}\) by restriction of scalars along \(\beta_W\), and then
  post-composed with the leg-(B) codomain swap below.

  \emph{Leg (B): the unit codomain ring-iso swap, via the lax-monoidal unit \(\varepsilon\).}
  The codomain of the reindexed component is the restriction of scalars of the unit section,
  which must be carried to the unit section over \(Y\). This is exactly the inverse of the
  lax-monoidal unit comparison of the restriction-of-scalars functor,
  \[
    \mathtt{codomainMap} \;:=\;
      \mathtt{inv}\bigl(\varepsilon(\mathtt{restrictScalars}\,g)\bigr),
      \qquad g := (f.\mathtt{appIso}\,W).\mathrm{inv}.\mathrm{hom},
  \]
  (\cref{def:dual_unit_ring_swap})
  where \(g\) is taken at the \emph{\(\mathtt{CommRingCat}\)} level. The unit \(\varepsilon\)
  is invertible because \(g\) is a ring isomorphism
  (\cref{lem:isiso_eps_restrictscalars_appiso}),
  so \(\varepsilon(\mathtt{restrictScalars}\,g)\) is an isomorphism. The presheaf-level
  identification underlying the unit handling is the dual-of-unit isomorphism for the unit
  section (\cref{def:dual_unit_iso_gen}).

  \emph{Inverse.} The inverse reindexing depends on a set-theoretic down-set fact: because
  \(f\) is an open immersion, \(fV \subseteq \mathrm{range}\,f\), so every open
  \(W'' \le fV\) already lies in the image of \(f.\mathtt{opensFunctor}\), namely
  \(W'' = f.\mathtt{opensFunctor}(f^{-1}W'')\). This is the open-image down-set coverage of an
  \(\mathtt{IsOpenImmersion}\): an open contained in the (open) image of the immersion is the
  functor-image of its own preimage. The project records it as the down-set identity
  \(\iota_!(\iota^{-1}V) = V\) for \(V \le W\) (the helper
  \(\mathtt{image\_preimage\_of\_le}\) (\cref{lem:image_preimage_of_le})),
  and it is exactly the lemma the inverse reindexing relies on to define the components on the
  whole down-set of \(fV\).

  With this in hand the inverse \(\mathtt{PresheafOfModules.Hom}\) mirrors the forward
  component formula, with two changes: it uses \((f.\mathtt{appIso}\,W'').\mathrm{hom}\) in
  place of its inverse, and its \(\varepsilon\)-codomain swap is
  \(\mathtt{dualUnitRingSwapHom} = \mathtt{inv}\bigl(\varepsilon(\mathtt{restrictScalars}\,
  (f.\mathtt{appIso}\,W'').\mathrm{hom})\bigr)\) (\cref{def:dual_unit_ring_swap_hom})
  --- that is, the \emph{inverse} of \(\varepsilon\) in the \(\mathrm{hom}\)-direction (where
  \(\varepsilon\) is invertible by
  \cref{lem:isiso_eps_restrictscalars_appiso_hom}), \emph{not}
  \(\varepsilon\) of the \(\mathrm{inv}\)-direction \(\beta\). (The reverse reindex turns an
  \(\mathcal{O}_X(fW'')\)-section map into an \(\mathcal{O}_Y(W'')\)-map, so it must restrict
  along the \(\mathrm{hom}\)-direction ring iso, whose codomain swap is therefore \(\mathtt{inv}\,
  \varepsilon\) of \emph{that} functor.) The whole family is reindexed by the inverse of the
  down-set bijection \(W'' \leftrightarrow f^{-1}W''\) above. The round-trip identities \(\mathtt{left\_inv}\)
  and \(\mathtt{right\_inv}\) then close component-wise by \(\mathtt{Iso.inv\_hom\_id}\) and
  \(\mathtt{Iso.hom\_inv\_id}\) of \(f.\mathtt{appIso}\) (the forward/inverse \(\beta\)-swaps
  cancel) together with the cancellation of the down-set bijection
  \(\iota_!(\iota^{-1}V) = V\) against its inverse.

  For additivity and \(\mathcal{O}_Y(V)\)-linearity, one must be careful about \emph{which}
  module structure carries the left-hand value. The additive and scalar structure on a dual
  section of the left-hand value
  \(\bigl((\mathtt{pushforward}\,\beta).\mathtt{obj}\,(\mathtt{dual}\,M.\mathtt{val})\bigr)(V)\)
  is the \(\mathcal{O}_Y(V)\)-module structure on the internal-hom object
  (\cref{def:presheaf_internal_hom_slice_value}) --- the pointwise module structure in which the
  sum and scalar multiple of dual sections are formed at the terminal object \(\mathrm{id}_{fV}\)
  of the slice over \(fV\). This agrees with, but is \emph{a priori} a distinct presentation of,
  the additive and scalar structure of the underlying presheaf morphism (the pointwise add and
  scalar action on morphisms of presheaves of modules). The verification of additivity and
  \(\mathcal{O}_Y(V)\)-linearity of the transport therefore proceeds in three distinct steps, the
  first of which is a required first move rather than an afterthought:
  \begin{enumerate}
    \item[(i)] \emph{Identify the two module presentations.} The internal-hom module structure's
      addition and scalar action are identified with the underlying pointwise add and scalar action
      on presheaf morphisms. This is a definitional identification of the two
      \(\mathcal{O}_Y(V)\)-module presentations of the same group of dual sections: the
      internal-hom operations, formed at the terminal object of the slice, are precisely the
      pointwise morphism-level operations of the presheaf morphism they package. Until this
      identification is made, the additivity and scalar-compatibility obligations of the transport
      cannot be read off the change-of-rings compatibility, because their sums and scalar multiples
      are taken in the \emph{internal-hom} structure, not the morphism structure.
    \item[(ii)] \emph{Match the X-side and Y-side scalars by the \(\beta\)-naturality ring
      identity.} The scalar acting on the X-side (the pushforward-restricted scalar \(s\) by
      which \(\mathtt{restrictScalars}\,\beta_W\) multiplies the reindexed component) and the
      scalar acting on the Y-side (the \(\mathcal{O}_Y(V)\)-scalar \(c\) of the internal-hom
      action) are \emph{a priori} elements of different rings and cannot be compared directly.
      The bridge is the ring identity
      \[
        s \;=\; (\beta.\mathtt{app}\,W').\mathrm{hom}\, c,
      \]
      obtained from \(\mathtt{InternalHom.termRingMap\_naturality}\) --- which records how the
      internal-hom term-ring map transports the scalar action across the slice --- together with
      the naturality of \(\beta\) on the thin poset \(\mathrm{Opens}\,Y\). Without this identity
      the two sides of the scalar-compatibility equation are not even of matching type, so it
      must be supplied before the change-of-rings compatibility of step (iii) can be invoked.
      The restricted scalar action is made explicit by
      \(\mathtt{ModuleCat.restrictScalars.smul\_def'}\), which exhibits the restricted smul
      as the original smul along \(\beta_W\).
    \item[(iii)] \emph{Apply the change-of-rings compatibility on the morphism level.} Once the
      operations are expressed pointwise on the underlying morphisms and the scalars matched by
      (ii), additivity and
      \(\mathcal{O}_Y(V)\)-linearity follow from the post-composition compatibility of the
      change-of-rings reindexing: the functor \(\mathtt{restrictScalars}\,\beta_W\) preserves sums
      and scalar multiples of morphisms, and the structure scalar is intertwined by the ring
      isomorphism \(\beta_W\). This last is the presheaf-level shadow of
      \cref{lem:restrictscalars_ringiso_dualequiv}.
  \end{enumerate}

  \emph{Naturality of the section family in \(W\).} This splits into two obligations of
  different character, and it is a mistake to discharge the whole square by
  \(\mathtt{Subsingleton.elim}\) alone. (a) The underlying \emph{base}-morphism naturality ---
  that the two routes around a square agree on the \emph{indexing} morphisms of
  \((\mathtt{Over}\,V.\mathrm{unop})^{\mathrm{op}}\) --- holds by \(\mathtt{Subsingleton.elim}\),
  since \(\mathrm{Opens}\,Y\) is a thin poset and there is at most one inclusion between any
  two objects; this discharges \emph{only} the uniqueness of the base morphisms. (b) The
  actual naturality obligation, however, is an equation between
  \(\mathcal{O}_Y(V)\)-\emph{module} maps, and it is \emph{not} settled by (a). Once the base
  morphisms are pinned by (a), the module-map equation reduces to the \(\varepsilon\)-naturality
  square of \(\mathtt{restrictScalars}\) --- the naturality of the lax-monoidal unit
  \(\varepsilon\) with respect to the restriction maps of \emph{both} the source and the target
  presheaf, along the structure-ring isomorphism \(\beta_W\) --- post-composed with the
  definition of the codomain swap \(\mathtt{dualUnitRingSwap} = \mathtt{inv}\,
  \varepsilon(\mathtt{restrictScalars}\,\beta_W)\). This \(\varepsilon\)-naturality step is a
  genuine, separate obligation; the thin-poset argument of (a) does \emph{not} discharge it.
  The required \(\varepsilon\)-commutativity-with-restriction-maps square is delivered by the
  extracted lemma \(\mathtt{sliceDualTransport\_naturality\_apply}\), which closes the
  module-map square \emph{pointwise} rather than through a single natural-transformation
  field. Once the base morphisms are pinned by (a), it reduces step (b) to two genuine
  ingredients: the ring-level square \(\mathtt{appIso\_hom\_naturality\_apply}\) --- that
  \((f.\mathtt{appIso}).\mathrm{hom}\) intertwines the \(X\)- and \(Y\)-side restriction maps,
  this being the \(\beta_W\)-naturality of the structure-ring isomorphism --- post-composed with
  the codomain swap \(\mathtt{dualUnitRingSwap} = \mathtt{inv}\,\varepsilon(\mathtt{restrictScalars}\,
  \beta_W)\) evaluated through its proven action \(\mathtt{dualUnitRingSwap\_apply}\). The
  thin-poset \(\mathtt{Subsingleton.elim}\) of (a) settles only the base morphisms.
\end{proof}

\subsection{Dual-inverse lane: the slice-transport helper declarations}
\label{subsec:dual_inverse_helpers}

These are the individual ingredients out of which \cref{lem:slice_dual_transport} and its
inverse leg are assembled: the leg-(B) unit-ring swaps, their invertibility and cancellation
identities, and the presheaf-level dual-of-unit identification. Each is a self-contained
internal construction.

\begin{lemma}\leanok
[Leg-(B) unit \(\varepsilon\) is invertible, \(\mathrm{inv}\)-direction]
  \label{lem:isiso_eps_restrictscalars_appiso}
  \lean{AlgebraicGeometry.Scheme.Modules.isIso_ε_restrictScalars_appIso}
  \uses{def:restrict_scalars_lax_epsilon}
  Let \(f : Y \hookrightarrow X\) be an open immersion and \(W' \subseteq Y\) an open. The
  lax-monoidal unit \(\varepsilon\) of the change-of-rings functor
  \(\mathtt{restrictScalars}\,\bigl((f.\mathtt{appIso}\,W').\mathrm{inv}.\mathrm{hom}\bigr)\),
  taken along the \(\mathtt{CommRingCat}\)-level structure-ring isomorphism, is an
  isomorphism.
\end{lemma}
\begin{proof}
  \uses{def:restrict_scalars_lax_epsilon}
  \leanok
  Proved directly in Lean --- the underlying ring map
  \((f.\mathtt{appIso}\,W').\mathrm{inv}.\mathrm{hom}\) is an isomorphism, hence bijective, so
  its restriction-of-scalars unit \(\varepsilon\) (\cref{def:restrict_scalars_lax_epsilon}) is
  invertible by the bijectivity criterion for \(\varepsilon\).
\end{proof}

\begin{lemma}\leanok
[Leg-(B) unit \(\varepsilon\) is invertible, \(\mathrm{hom}\)-direction]
  \label{lem:isiso_eps_restrictscalars_appiso_hom}
  \lean{AlgebraicGeometry.Scheme.Modules.isIso_ε_restrictScalars_appIso_hom}
  \uses{def:restrict_scalars_lax_epsilon}
  In the same setting, the lax-monoidal unit \(\varepsilon\) of
  \(\mathtt{restrictScalars}\,\bigl((f.\mathtt{appIso}\,W').\mathrm{hom}.\mathrm{hom}\bigr)\)
  --- the \(\mathrm{hom}\)-direction of the structure-ring iso --- is likewise an
  isomorphism. This is the variant powering the reverse (\(\mathtt{invFun}\)) leg, which
  reindexes a \(\mathcal{O}_Y\)-section map back to a \(\mathcal{O}_X\)-map.
\end{lemma}
\begin{proof}
  \uses{def:restrict_scalars_lax_epsilon}
  \leanok
  Proved directly in Lean --- identical to \cref{lem:isiso_eps_restrictscalars_appiso} with
  the \(\mathrm{hom}\)-direction (also bijective) ring map.
\end{proof}

\begin{lemma}


\leanok
[Unit \(\varepsilon\) is invertible, section-relabel direction]
  \label{lem:isiso_eps_restrictscalars_presheafmap}
  \lean{AlgebraicGeometry.Scheme.Modules.isIso_ε_restrictScalars_presheafMap}
  \uses{def:restrict_scalars_lax_epsilon}
  Let \(X\) be a scheme and let \(a = b\) be an equality of opens of \(X\) (taken in
  \((\mathrm{Opens}\,X)^{\mathrm{op}}\)), inducing the \(\mathtt{eqToHom}\) section-ring relabel
  \(X.\mathtt{presheaf}.\mathrm{map}\,(\mathtt{eqToHom}\,e) \colon \mathcal{O}_X(b) \to
  \mathcal{O}_X(a)\). The lax-monoidal unit \(\varepsilon\) of the change-of-rings functor
  \(\mathtt{restrictScalars}\,\bigl((X.\mathtt{presheaf}.\mathrm{map}\,(\mathtt{eqToHom}\,e)).\mathrm{hom}\bigr)\)
  along that relabel is an isomorphism. This is the \(\mathtt{eqToHom}\)-induced-ring-map
  variant of the \(\varepsilon\)-invertibility criterion, the section-relabel sibling of
  \cref{lem:isiso_eps_restrictscalars_appiso}.
\end{lemma}
\begin{proof}
  \uses{def:restrict_scalars_lax_epsilon}
  \leanok
  Proved directly in Lean --- the underlying ring map
  \((X.\mathtt{presheaf}.\mathrm{map}\,(\mathtt{eqToHom}\,e)).\mathrm{hom}\) is the image of an
  \(\mathtt{eqToHom}\) under a functor, hence an isomorphism and bijective, so its
  restriction-of-scalars unit \(\varepsilon\) (\cref{def:restrict_scalars_lax_epsilon}) is
  invertible by the bijectivity criterion for \(\varepsilon\).
\end{proof}

\begin{definition}\leanok
[Leg-(B) codomain unit ring-iso swap, \(\mathrm{inv}\)-direction]
  \label{def:dual_unit_ring_swap}
  \lean{AlgebraicGeometry.Scheme.Modules.dualUnitRingSwap}
  \uses{lem:isiso_eps_restrictscalars_appiso}
  The codomain unit ring-iso swap
  \[
    \mathtt{dualUnitRingSwap}\,f\,W'
      \;:=\; \mathtt{inv}\bigl(\varepsilon(\mathtt{restrictScalars}\,
        (f.\mathtt{appIso}\,W').\mathrm{inv}.\mathrm{hom})\bigr)
    \;\colon\;
    (\mathtt{restrictScalars}\,(f.\mathtt{appIso}\,W')^{-1})(\mathbf 1_X(fW'))
      \;\longrightarrow\; \mathbf 1_Y(W'),
  \]
  the inverse of the (invertible, \cref{lem:isiso_eps_restrictscalars_appiso}) lax-monoidal
  unit \(\varepsilon\), carrying the restriction-of-scalars of the \(X\)-unit section to the
  \(Y\)-unit section. It is the leg-(B) ingredient of \cref{lem:slice_dual_transport}.
\end{definition}
\begin{proof}
  \uses{lem:isiso_eps_restrictscalars_appiso}
  \leanok
  Proved directly in Lean --- the categorical inverse of an isomorphism.
\end{proof}

\begin{definition}\leanok
[Leg-(B) unit ring-iso swap, reverse (\(\varepsilon\) itself)]
  \label{def:dual_unit_ring_swap_inv}
  \lean{AlgebraicGeometry.Scheme.Modules.dualUnitRingSwapInv}
  \uses{lem:isiso_eps_restrictscalars_appiso}
  The reverse map
  \(\mathtt{dualUnitRingSwapInv}\,f\,W' := \varepsilon(\mathtt{restrictScalars}\,
  (f.\mathtt{appIso}\,W').\mathrm{inv}.\mathrm{hom})\) --- the lax-monoidal unit
  \(\varepsilon\) \emph{itself} (not its inverse) --- from \(\mathbf 1_Y(W')\) back to the
  restriction-of-scalars of the \(X\)-unit section. It is the two-sided inverse of
  \cref{def:dual_unit_ring_swap}.
\end{definition}
\begin{proof}
  \uses{lem:isiso_eps_restrictscalars_appiso}
  \leanok
  Proved directly in Lean --- the unit \(\varepsilon\), invertible by
  \cref{lem:isiso_eps_restrictscalars_appiso}.
\end{proof}

\begin{definition}\leanok
[Reverse leg codomain unit ring-iso swap, \(\mathrm{hom}\)-direction]
  \label{def:dual_unit_ring_swap_hom}
  \lean{AlgebraicGeometry.Scheme.Modules.dualUnitRingSwapHom}
  \uses{lem:isiso_eps_restrictscalars_appiso_hom}
  The \(\mathrm{hom}\)-direction codomain swap
  \[
    \mathtt{dualUnitRingSwapHom}\,f\,W'
      \;:=\; \mathtt{inv}\bigl(\varepsilon(\mathtt{restrictScalars}\,
        (f.\mathtt{appIso}\,W').\mathrm{hom}.\mathrm{hom})\bigr)
    \;\colon\;
    (\mathtt{restrictScalars}\,(f.\mathtt{appIso}\,W').\mathrm{hom})(\mathbf 1_Y(W'))
      \;\longrightarrow\; \mathbf 1_X(fW'),
  \]
  the inverse of the (invertible, \cref{lem:isiso_eps_restrictscalars_appiso_hom}) unit
  \(\varepsilon\) in the \(\mathrm{hom}\)-direction. It is the leg-(B) ingredient of the
  reverse transport \cref{lem:slice_dual_transport_inv}.
\end{definition}
\begin{proof}
  \uses{lem:isiso_eps_restrictscalars_appiso_hom}
  \leanok
  Proved directly in Lean --- the categorical inverse of an isomorphism.
\end{proof}

\begin{definition}


\leanok
[Unit-section relabel swap (the fourth leg of the reverse transport)]
  \label{def:unit_relabel_swap}
  \lean{AlgebraicGeometry.Scheme.Modules.unitRelabelSwap}
  \uses{lem:isiso_eps_restrictscalars_presheafmap}
  For an equality \(a = b\) of opens of \(X\) (in \((\mathrm{Opens}\,X)^{\mathrm{op}}\)), the
  unit-section relabel swap
  \[
    \mathtt{unitRelabelSwap}\,e
      \;:=\; \mathtt{inv}\bigl(\varepsilon(\mathtt{restrictScalars}\,
        (X.\mathtt{presheaf}.\mathrm{map}\,(\mathtt{eqToHom}\,e)).\mathrm{hom})\bigr)
    \;\colon\;
    \bigl(\mathtt{restrictScalars}\,(X.\mathtt{presheaf}.\mathrm{map}\,(\mathtt{eqToHom}\,e))\bigr)
      (\mathbf 1_X(b))
      \;\longrightarrow\; \mathbf 1_X(a),
  \]
  the inverse of the (invertible, \cref{lem:isiso_eps_restrictscalars_presheafmap}) unit
  \(\varepsilon\) of the section-ring relabel. It transports the codomain unit section
  \(\mathbf 1_X\) across the cross-fiber section-ring relabel \(\mathcal{O}_X(b) \cong
  \mathcal{O}_X(a)\), and is the fourth leg of the reverse transport
  \cref{lem:slice_dual_transport_inv} --- the unit-section mirror of
  \cref{def:dual_unit_ring_swap_hom}, needed because the source relabel
  \(M.\mathtt{val}(\mathtt{eqToHom})\) is semilinear and the in-fiber core must be conjugated
  back into the slice section's fiber.
\end{definition}
\begin{proof}
  \uses{lem:isiso_eps_restrictscalars_presheafmap}
  \leanok
  Proved directly in Lean --- the categorical inverse of an isomorphism.
\end{proof}

\begin{lemma}\leanok
[Round-trip of the swap pair, \(\mathrm{inv}\circ\mathrm{swap}\)]
  \label{lem:dual_unit_ring_swap_comp_inv}
  \lean{AlgebraicGeometry.Scheme.Modules.dualUnitRingSwap_comp_dualUnitRingSwapInv}
  \uses{def:dual_unit_ring_swap, def:dual_unit_ring_swap_inv}
  \(\mathtt{dualUnitRingSwap}\,f\,W' \circ \mathtt{dualUnitRingSwapInv}\,f\,W' =
  \mathrm{id}\): the swap followed by its reverse is the identity.
\end{lemma}
\begin{proof}
  \uses{def:dual_unit_ring_swap, def:dual_unit_ring_swap_inv}
  \leanok
  Proved directly in Lean --- the \(\mathtt{Iso.inv\_hom\_id}\) cancellation of
  \(\varepsilon\) with its inverse.
\end{proof}

\begin{lemma}\leanok
[Round-trip of the swap pair, \(\mathrm{swap}\circ\mathrm{inv}\)]
  \label{lem:dual_unit_ring_swap_inv_comp}
  \lean{AlgebraicGeometry.Scheme.Modules.dualUnitRingSwapInv_comp_dualUnitRingSwap}
  \uses{def:dual_unit_ring_swap, def:dual_unit_ring_swap_inv}
  \(\mathtt{dualUnitRingSwapInv}\,f\,W' \circ \mathtt{dualUnitRingSwap}\,f\,W' =
  \mathrm{id}\): the reverse followed by the swap is the identity.
\end{lemma}
\begin{proof}
  \uses{def:dual_unit_ring_swap, def:dual_unit_ring_swap_inv}
  \leanok
  Proved directly in Lean --- the \(\mathtt{Iso.hom\_inv\_id}\) cancellation of
  \(\varepsilon\) with its inverse.
\end{proof}

\begin{definition}\leanok
[Section equivalence for the dual of the unit]
  \label{def:unit_dual_section_equiv}
  \lean{PresheafOfModules.unitDualSectionEquiv}
  \uses{def:eval_lin, def:global_smul, def:presheaf_internal_hom_slice_value}
  For a presheaf of modules over a base ring presheaf \(R_0\) and an object \(X\), the
  \(R_0(X)\)-linear equivalence
  \[
    \bigl(\mathtt{restr}_X\,\mathbf 1 \to \mathtt{restr}_X\,\mathbf 1\bigr)
      \;\simeq_{R_0(X)}\; R_0(X)
  \]
  identifying endomorphisms of the (restricted) monoidal unit with the ground ring, via
  evaluation at the global section \(1\); the inverse is multiplication by a global scalar
  (\cref{def:global_smul}).
\end{definition}
\begin{proof}
  \uses{def:eval_lin, def:global_smul, def:presheaf_internal_hom_slice_value}
  \leanok
  Proved directly in Lean. The forward map is \(\mathtt{evalLin}\) at \(1\)
  (\cref{def:eval_lin}); the inverse is the global-scalar action \(\mathtt{globalSMul}\)
  (\cref{def:global_smul}). The substantive field is \(\mathtt{left\_inv}\): every
  endomorphism of the unit is multiplication by its value at \(1\), which follows from the
  \(\varphi\)-naturality of the section toward the terminal object of the slice. The module
  structure on the hom-space is the slice internal-hom structure
  (\cref{def:presheaf_internal_hom_slice_value}).
\end{proof}

\begin{definition}\leanok
[Presheaf dual of the unit is the unit (general base)]
  \label{def:dual_unit_iso_gen}
  \lean{PresheafOfModules.dualUnitIsoGen}
  \uses{def:unit_dual_section_equiv, def:presheaf_dual, lem:internal_hom_eval}
  The presheaf-level isomorphism
  \(\mathtt{dual}\,\mathbf 1 \cong \mathbf 1\) of presheaves of modules over a general
  single-universe base ring presheaf \(R_0\), assembled sectionwise from
  \cref{def:unit_dual_section_equiv} with the evaluation-at-\(1\) naturality.
\end{definition}
\begin{proof}
  \uses{def:unit_dual_section_equiv, def:presheaf_dual, lem:internal_hom_eval}
  \leanok
  Proved directly in Lean. The sectionwise components are the section equivalences
  \cref{def:unit_dual_section_equiv} packaged by \(\mathtt{isoMk}\); the naturality square is
  the evaluation-at-\(1\) naturality of the internal-hom evaluation
  (\cref{lem:internal_hom_eval}) specialised at \(M = \mathbf 1\).
\end{proof}

\begin{definition}\leanok
[Scheme-level presheaf dual of the unit is the unit]
  \label{def:presheaf_dual_unit_iso}
  \lean{AlgebraicGeometry.Scheme.Modules.presheafDualUnitIso}
  \uses{def:dual_unit_iso_gen}
  The scheme-level instance \(\mathtt{dual}\,\mathbf 1 \cong \mathbf 1\) of the
  presheaf-of-modules dual-of-unit isomorphism, obtained by specialising
  \cref{def:dual_unit_iso_gen} at \(R_0 = Y.\mathtt{presheaf}\). It is the presheaf core of
  the third leg \(\mathtt{dual\_unit\_iso}\) (\cref{lem:dual_unit_iso}).
\end{definition}
\begin{proof}
  \uses{def:dual_unit_iso_gen}
  \leanok
  Proved directly in Lean --- a direct instance of \cref{def:dual_unit_iso_gen}.
\end{proof}

\begin{lemma}\leanok
[Reverse slice dual transport (the \(\mathtt{invFun}\) of \cref{lem:slice_dual_transport})]
  \label{lem:slice_dual_transport_inv}
  \lean{AlgebraicGeometry.Scheme.Modules.sliceDualTransportInv}
  \uses{def:dual_unit_ring_swap_hom, lem:image_preimage_of_le, lem:isiso_eps_restrictscalars_appiso_hom, lem:restrictscalars_ringiso_dualequiv, def:presheaf_dual_unit_iso, def:unit_relabel_swap}
  \textit{Source: internal categorical construction; no external reference.}
  In the setting of \cref{lem:slice_dual_transport}, with \(f : Y \hookrightarrow X\) an open
  immersion, \(M \in \Scheme.\mathtt{Modules}\,X\), \(V \subseteq Y\) an open and \(\beta\)
  the open-immersion structure-ring morphism, there is a reverse transport: a dual section
  \(\psi : \mathtt{restr}_V\,((\mathtt{pushforward}\,\beta)\,M.\mathtt{val}) \to
  \mathtt{restr}_V\,\mathbf 1_Y\) over \(\mathtt{Over}\,V\) is sent to an \(X\)-slice dual
  section \(\mathtt{restr}_{fV}\,M.\mathtt{val} \to \mathtt{restr}_{fV}\,\mathbf 1_X\) over
  \(\mathtt{Over}\,(fV)\), where \(fV = f.\mathtt{opensFunctor}(V)\). This is the inverse leg
  of the forward equivalence \cref{lem:slice_dual_transport}.
\end{lemma}
\begin{proof}
  \uses{def:dual_unit_ring_swap_hom, lem:image_preimage_of_le, lem:isiso_eps_restrictscalars_appiso_hom, lem:restrictscalars_ringiso_dualequiv, def:presheaf_dual_unit_iso, lem:slice_dual_transport_inv_naturality_apply}
  \leanok
  We build the underlying \(\mathtt{PresheafOfModules.Hom}\) component-by-component over the
  \(X\)-slice \(\mathtt{Over}\,(fV)\), mirroring the forward component formula of
  \cref{lem:slice_dual_transport} with two changes.

  \emph{The component at \(W''\).} Fix \(W'' \le fV\) and set \(W' := W''.\mathrm{left}\) and
  \(P := f^{-1}(W')\) for its preimage. Because \(f\) is an open immersion and
  \(fV \subseteq \mathrm{range}\,f\), the open \(W' \le fV\) lies in the image of
  \(f.\mathtt{opensFunctor}\), so the down-set identity
  \(f.\mathtt{opensFunctor}(P) = W'\) holds --- but only \emph{propositionally}
  (\cref{lem:image_preimage_of_le}); this is the single set-theoretic fact the inverse
  reindexing depends on. The component of the reverse section at \(W''\) is the \(X\)-slice
  mirror of the forward component:
  \[
    M.\mathtt{val}(\mathtt{eqToHom})
      \;\circ\;
    (\mathtt{restrictScalars}\,(f.\mathtt{appIso}\,P).\mathrm{hom}).\mathrm{map}
      \bigl(\psi_{(\mathtt{Over.mk}\,(P \le V))}\bigr)
      \;\circ\;
    \mathtt{dualUnitRingSwapHom}\,f\,P
      \;\circ\;
    \mathtt{unitRelabelSwap}\,(\mathtt{eqToHom}\,he.\mathrm{symm}),
  \]
  conjugated on source and target by the \(\mathtt{eqToHom}\) transports arising from the
  propositional equality \(f.\mathtt{opensFunctor}(P) = W'\) (the mirror of
  \(\mathtt{homLocalSection}\)). Compared with the forward formula, this map uses
  \((f.\mathtt{appIso}\,P).\mathrm{hom}\) in place of its inverse --- because the reverse
  reindex turns a \(\mathcal{O}_Y(P)\)-section map into a \(\mathcal{O}_X(fP)\)-map --- and
  its codomain swap is \(\mathtt{dualUnitRingSwapHom}\) (\cref{def:dual_unit_ring_swap_hom}),
  the inverse of \(\varepsilon\) in the \(\mathrm{hom}\)-direction
  (\cref{lem:isiso_eps_restrictscalars_appiso_hom}). The fourth leg
  \(\mathtt{unitRelabelSwap}\,(\mathtt{eqToHom}\,he.\mathrm{symm})\) transports the codomain
  unit \(\mathbf 1_X\) across the cross-fiber section-ring relabel \(\mathcal{O}_X(W') \cong
  \mathcal{O}_X(fP)\): it is the \(\mathtt{inv}\,\varepsilon\) of that \(\mathtt{eqToHom}\)-induced
  relabel, the unit-section mirror of \(\mathtt{dualUnitRingSwapHom}\), needed because the
  source relabel \(M.\mathtt{val}(\mathtt{eqToHom})\) is semilinear and the in-fiber core must be
  conjugated back into the fiber \(\mathcal{O}_X(W')\) of the slice section.

  This reverse component is \emph{not} available for an arbitrary \(\beta\): collapsing the two
  stacked change-of-rings layers on \(M.\mathtt{val}(fP)\) into the identity --- the step
  \(\mathtt{restrictScalarsComp'App}\) that the in-fiber core depends on --- requires the
  compatibility identity
  \[
    h\beta : \forall P,\quad
      (\beta.\mathrm{app}\,(\mathtt{op}\,P)).\mathrm{hom}
        \;\circ\;
      (f.\mathtt{appIso}\,P).\mathrm{hom}.\mathrm{hom}
        \;=\;
      \mathrm{id},
  \]
  which the lemma takes as an explicit hypothesis. The equality is \emph{false} for a general
  structure morphism \(\beta\); at the unique call site \(\beta\) is the open-immersion
  structure-ring morphism \((f.\mathtt{appIso})^{-1}\), and the identity is discharged by
  \(\mathtt{Iso.hom\_inv\_id}\).

  \emph{Additivity and \(\mathcal{O}_Y(V)\)-linearity} follow exactly as in the forward leg
  (steps (i)--(iii) of \cref{lem:slice_dual_transport}): identify the two module
  presentations, match the X- and Y-side scalars by the \(\beta\)-naturality ring identity,
  and apply the change-of-rings compatibility, the module-level shadow of
  \cref{lem:restrictscalars_ringiso_dualequiv}.

  \emph{Naturality} of the component family across morphisms of \((\mathtt{Over}\,(fV))^{\mathrm{op}}\)
  reduces, after the thin-poset \(\mathtt{Subsingleton.elim}\) on the indexing morphisms, to
  the \(\varepsilon\)-naturality of \(\mathtt{restrictScalars}\) along the structure-ring iso,
  exactly as in the forward leg. The unit handling is the presheaf dual-of-unit
  identification \cref{def:presheaf_dual_unit_iso}. The whole family is reindexed by the
  inverse of the down-set bijection \(W'' \leftrightarrow f^{-1}(W'')\) supplied by
  \cref{lem:image_preimage_of_le}.
\end{proof}

\begin{lemma}\leanok
[Pointwise naturality square of the reverse slice transport]
  \label{lem:slice_dual_transport_inv_naturality_apply}
  \lean{AlgebraicGeometry.Scheme.Modules.sliceDualTransportInv_naturality_apply}
  \uses{def:dual_unit_ring_swap_hom, lem:image_preimage_of_le}
  \textit{Source: internal categorical construction; no external reference.}
  With notation as in \cref{lem:slice_dual_transport_inv}, the reverse-transport component
  family satisfies the pointwise naturality square: for a morphism \(g\) of the slice
  \((\mathtt{Over}\,(fV))^{\mathrm{op}}\) and a section \(z\), the two ways of relabelling
  and transporting \(z\) agree, with the two \(\mathtt{inv}\,\varepsilon\) legs
  (\(\mathtt{unitRelabelSwap}\), \(\mathtt{dualUnitRingSwapHom}\)) kept \emph{shallow} (the
  down-set facts passed explicitly). It is the clean-statement extraction whose \(\mathtt{exact}\)
  (by defeq) closes the \(\mathtt{naturality}\) field of \(\mathtt{sliceDualTransportInv}\).
\end{lemma}
\begin{proof}
  \uses{lem:slice_dual_transport_inv, def:dual_unit_ring_swap_hom, lem:image_preimage_of_le}
  \leanok
  After rewriting the two \(\mathtt{inv}\,\varepsilon\) legs into their shallow \(\mathtt{\_apply}\)
  forms (so no \(\mathtt{whnf}\) of the deep \(\mathtt{inv}\,\varepsilon\) composite occurs) the
  goal is a concrete element equation, closed by three ingredients: the \(M\)-side coherence
  (both legs are \(M.\mathtt{val}\)-restrictions of \(z\), fused by \(M.\mathtt{val}.\mathtt{map\_comp}\)
  and the thin-poset \(\mathtt{Subsingleton.elim}\)); \(\psi\)-naturality
  (\(\mathtt{PresheafOfModules.naturality\_apply}\)) at the slice morphism; and
  \(\mathtt{appIso}\)-naturality (\(\mathtt{Scheme.Hom.appIso\_inv\_naturality}\)), the
  \(X\)-presheaf round-trip closing by a \(\mathtt{CommRingCat}\) morphism identity.
\end{proof}

\begin{lemma}\leanok
[Clean pointwise form of the reverse-transport component]
  \label{lem:slice_dual_transport_inv_app_apply}
  \lean{AlgebraicGeometry.Scheme.Modules.sliceDualTransportInv_app_apply}
  \uses{def:dual_unit_ring_swap_hom}
  \textit{Source: internal categorical construction; no external reference.}
  The \(\mathtt{app}\) component of \(\mathtt{sliceDualTransportInv}\) evaluated at a section
  equals its explicit reverse-transport formula (the leg reductions are definitional). This
  is the \(\mathtt{rfl}\) bridge reused by the \(\mathtt{left\_inv}\) and \(\mathtt{right\_inv}\)
  round-trip proofs of \cref{lem:slice_dual_transport} to expose the inverse leg in clean form
  before the \(\mathtt{appIso}\) round-trip cancellation.
\end{lemma}
\begin{proof}
  \uses{lem:slice_dual_transport_inv}
  \leanok
  Definitional: every leg reduction in the component formula is a \(\mathtt{rfl}\)-identity, so
  the equation holds by \(\mathtt{rfl}\) once the down-set proofs are supplied.
\end{proof}

\begin{lemma}


\leanok
[Pointwise \(\mathrm{hom}\)-naturality of the structure ring iso]
  \label{lem:appiso_hom_naturality_apply}
  \lean{AlgebraicGeometry.Scheme.Modules.appIso_hom_naturality_apply}
  \textit{Source: internal categorical construction; no external reference.}
  Let \(f : Y \hookrightarrow X\) be an open immersion, \(i : U \to V\) an inclusion of opens of
  \(Y\), and \(w\) a section of \(\mathcal{O}_X\) over \(f.\mathtt{opensFunctor}(U)\). Then
  \((f.\mathtt{appIso}\,V).\mathrm{hom}\) intertwines the \(X\)- and \(Y\)-side restriction maps:
  \[
    (f.\mathtt{appIso}\,V).\mathrm{hom}\bigl(
      (X.\mathtt{presheaf}.\mathrm{map}\,(f.\mathtt{opensFunctor}^{\mathrm{op}}\,i))\,w\bigr)
      \;=\;
    (Y.\mathtt{presheaf}.\mathrm{map}\,i)\bigl((f.\mathtt{appIso}\,U).\mathrm{hom}\,w\bigr).
  \]
  This is the single ring-level square underlying the forward naturality and round-trip pastes
  of the slice transport (the \(\mathrm{hom}\)-direction \(\beta_W\)-naturality of the
  structure-ring isomorphism).
\end{lemma}
\begin{proof}


\leanok
  Proved directly in Lean. It is the \(\beta_W\)-naturality of the structure-ring isomorphism,
  obtained by rotating the forward naturality square of \((f.\mathtt{appIso}).\mathrm{inv}\)
  through the iso: \(f.\mathtt{appIso}\,V\) is an isomorphism, so its \(\mathrm{inv}\)-leg is
  bijective; applying that injectivity and cancelling the \(\mathrm{hom}\)/\(\mathrm{inv}\)
  round-trips on both sides reduces the claim to the \(\mathrm{inv}\)-direction naturality
  square, which holds.
\end{proof}

\begin{lemma}


\leanok
[Pointwise naturality square of the forward slice transport]
  \label{lem:slice_dual_transport_naturality_apply}
  \lean{AlgebraicGeometry.Scheme.Modules.sliceDualTransport_naturality_apply}
  \uses{def:dual_unit_ring_swap, lem:appiso_hom_naturality_apply}
  \textit{Source: internal categorical construction; no external reference.}
  With notation as in \cref{lem:slice_dual_transport}, the forward-transport component family
  satisfies the pointwise naturality square: for a morphism \(f_1\) of the slice
  \((\mathtt{Over}\,V)^{\mathrm{op}}\) and a section \(z\), the \(\mathtt{dualUnitRingSwap}\)-conjugated
  reindex of a dual section \(\varphi\) commutes with the slice restriction maps. It is
  the \(\mathtt{toFun}\)-side mirror of \cref{lem:slice_dual_transport_inv_naturality_apply},
  serving as the naturality certificate for \(\mathtt{sliceDualTransport}\).
\end{lemma}
\begin{proof}
  \uses{lem:slice_dual_transport, def:dual_unit_ring_swap, lem:appiso_hom_naturality_apply}
  \leanok
  After substituting the \(\mathtt{inv}\,\varepsilon\) leg via
  \(\mathtt{dualUnitRingSwap\_apply}\), the goal reduces to a concrete element equation closed by two
  ingredients: \(\varphi\)-naturality (\(\mathtt{PresheafOfModules.naturality\_apply}\)) at the
  \(f\)-image of \(f_1\) in the \(X\)-slice (the indexing triangle pinned by the thin-poset
  \(\mathtt{Subsingleton.elim}\)), and the ring-level square \cref{lem:appiso_hom_naturality_apply}
  --- that \((f.\mathtt{appIso}).\mathrm{hom}\) intertwines the \(X\)- and \(Y\)-side
  restriction maps.
\end{proof}

\begin{lemma}





\leanok
[Restriction commutes with the dual along an open immersion (the C-bridge)]
  \label{lem:dual_restrict_iso}
  \lean{AlgebraicGeometry.Scheme.Modules.dual_restrict_iso}
  \uses{lem:internal_hom_isSheaf, lem:restrictscalars_ringiso_dualequiv, lem:slice_dual_transport}
  % SOURCE: [Stacks Project], "Modules on Ringed Spaces", \S Internal Hom,
  % Lemma~\texttt{lemma-pullback-internal-hom} (tag area 01CM)
  % (read from references/stacks-modules.tex, L3710--L3721).
  % SOURCE QUOTE:
  % "Let $f : (X, \mathcal{O}_X) \to (Y, \mathcal{O}_Y)$ be a morphism
  % of ringed spaces. Let $\mathcal{F}$, $\mathcal{G}$ be $\mathcal{O}_Y$-modules.
  % If $\mathcal{F}$ is finitely presented and $f$ is flat,
  % then the canonical map
  % $$
  % f^*\SheafHom_{\mathcal{O}_Y}(\mathcal{F}, \mathcal{G})
  % \longrightarrow
  % \SheafHom_{\mathcal{O}_X}(f^*\mathcal{F}, f^*\mathcal{G})
  % $$
  % is an isomorphism."
  \textit{Source: [Stacks Project], ``Modules on Ringed Spaces'', \S Internal Hom,
  Lemma~\texttt{lemma-pullback-internal-hom} (tag area 01CM).}
  Let \(f : Y \hookrightarrow X\) be an open immersion of schemes, with image the open
  \(U \subseteq X\), and let \(M \in \Scheme.\mathtt{Modules}\,X\). Then there is a
  canonical isomorphism of \(\mathcal{O}_Y\)-modules
  \[
    (\mathtt{dual}\,M)\big|_f \;\xrightarrow{\ \sim\ }\; \mathtt{dual}\bigl(M|_f\bigr),
  \]
  natural in \(M\), between the restriction of the sheaf-level dual
  (\cref{lem:internal_hom_isSheaf}) and the dual of the restriction. (Only opens
  \(V \subseteq U\) intervene downstream, where the comparison is needed.)
\end{lemma}

\begin{proof}
  \uses{lem:internal_hom_isSheaf, lem:restrictscalars_ringiso_dualequiv, lem:slice_dual_transport}
  \leanok
  This is the classical internal-hom base-change isomorphism
  \(f^*\mathcal{H}\!om_{\mathcal{O}_X}(M, \mathcal{O}_X) \cong
  \mathcal{H}\!om_{\mathcal{O}_Y}(f^*M, f^*\mathcal{O}_X)\) of
  Lemma~\texttt{lemma-pullback-internal-hom}, specialised to the structure sheaf
  \(\mathcal{G} = \mathcal{O}_X\). For an open immersion the morphism \(f\) is flat and
  restriction \((-)|_f\) is the exact pullback \(f^*\), so the finite-presentation and
  flatness hypotheses of the general ringed-space statement are automatic; this is the
  \emph{favorable} (open-immersion) case of the base-change comparison, in which the map
  is an isomorphism for \emph{every} \(M\), not only the finitely-presented ones.

  We emphasise at the outset that this is a genuine \emph{natural isomorphism}, not a
  definitional equality, and --- crucially --- that the dual is \emph{not sectionwise}.
  Unlike the tensor product, whose value is computed fibrewise
  (\((M \otimes N)(V) = M(V) \otimes N(V)\)), the dual evaluates to a Hom over an entire
  down-set: for an open \(V \subseteq U\),
  \[
    (\mathtt{dual}\,M)(V)
      \;=\; \mathcal{H}om\bigl(M|_V,\ \mathcal{O}|_V\bigr)\big|_{\mathtt{Over}\,V},
  \]
  a module of morphisms of presheaves over the whole slice \(\mathtt{Over}\,V \subseteq
  \mathrm{Opens}\,Y\), \emph{not} a single fibre \(M(V) \to \mathcal{O}(V)\). Consequently
  the assertion that restriction along the open immersion \(f\) commutes with the dual does
  \emph{not} reduce to a single ground-ring reconciliation. The proof proceeds in two
  stages: a preliminary adjunction-uniqueness rewrite (stage~H1) that re-presents the
  pulled-back dual in pushforward form, followed by the resolution of the resulting
  \emph{residual} into \emph{two} legs, carried out objectwise in the open \(V\) (with
  \(V \subseteq U\)) and assembled over the slice with poset-thin naturality.

  \emph{Stage H1 --- adjunction-uniqueness rewrite.} The pulled-back dual is initially
  presented in the \(\mathtt{pullback}\,\varphi\)-form inherited from the abstract
  definition of the restriction functor, and no sectionwise comparison can begin against
  that opaque form. We first rewrite it into a \(\mathtt{pushforward}\,\beta\)-form: the
  restriction-along-an-open-immersion pullback and the pushforward along the structure-ring
  map \(\beta\) are both left adjoints to one and the same functor, so by uniqueness of left
  adjoints --- the adjunction \(\mathtt{pushforwardPushforwardAdj}\) composed with
  \(\mathtt{leftAdjointUniq}\) --- they are canonically isomorphic, and this isomorphism
  transports the pulled-back dual into pushforward form. After H1 the residual obligation is
  \emph{exactly}
  \[
    (\mathtt{pushforward}\,\beta).\mathtt{obj}\,(\mathtt{dual}\,M.\mathtt{val})
      \;\xrightarrow{\ \sim\ }\;
    \mathtt{dual}\bigl((\mathtt{pushforward}\,\beta).\mathtt{obj}\,M.\mathtt{val}\bigr),
  \]
  i.e.\ ``\(\mathtt{pushforward}\,\beta\) commutes with the dual''. It is this residual ---
  \emph{not} the original Step-4 isomorphism --- that legs (A) and (B) below resolve; legs
  (A) and (B) are the content remaining \emph{after} the H1 rewrite, not a standalone
  two-leg split of the whole step.

  \emph{Leg (A) --- slice-site Hom base-change (Beck--Chevalley).} The domain presheaf must
  first be transported across the open immersion. Because \(f\) is an open immersion, its
  direct-image functor on open sets is fully faithful with image
  the down-set \(\{W : W \le U\}\); this carries the restricted domain
  \(\mathtt{restr}_V(f_*M)\) over \(\mathtt{Over}\,V \subseteq \mathrm{Opens}\,Y\) to the
  restricted domain \(\mathtt{restr}_{fV}\,M\) over \(\mathtt{Over}\,(fV) \subseteq
  \mathrm{Opens}\,X\), and hence identifies the two Hom-modules computed over these matched
  slices. This is a genuine slice-site Hom base-change --- a Beck--Chevalley square on the
  internal hom across the fully-faithful open-immersion functor, restricted to the down-set
  of \(fV\) --- \emph{not} a fibrewise identity.

  \emph{Leg (B) --- ground-ring reconciliation.} The codomain of the internal hom is in
  both cases the structure sheaf. Once leg (A) has matched the domains, the unit codomain
  is reconciled along the open-immersion ring isomorphism
  \(\beta_V = (\mathtt{toRingCatSheafHom}\,f).\mathrm{hom}_V \colon
  \mathcal{O}_X(fV) \xrightarrow{\sim} \mathcal{O}_Y(V)\): restriction of scalars along the
  ring \emph{isomorphism} \(\beta_V\) is an equivalence, and its ModuleCat-level
  shadow on duals is exactly \cref{lem:restrictscalars_ringiso_dualequiv}, carrying the
  \(\mathcal{O}_X(fV)\)-linear dual onto the \(\mathcal{O}_Y(V)\)-linear dual,
  \(R(V)\)-linearly and invertibly. This is the dual analogue, at the module level, of the
  ring-isomorphism tensor equivalence (\texttt{restrictscalars\_ringiso\_strongmonoidal}).

  \emph{The combined leg-(A)\(\circ\)(B) atom \(\mathtt{sliceDualTransport}\).} Both legs are
  realised together by a single \(\mathcal{O}_Y(V)\)-linear isomorphism
  \(\mathtt{sliceDualTransport}\,f\,M\,V\) (\cref{lem:slice_dual_transport}), built
  \emph{by hand} --- it is \emph{not} obtained by consuming a sheaf-level slice equivalence.
  This is essential: the equivalence of \emph{sheaves of modules} over an open immersion
  compares the restriction functors and the units, but carries \emph{no} internal-hom or
  dual content; the assertion realized by leg (A) --- that the dual commutes with slice
  reindexing along the open-immersion direct-image functor \(f_{!}\) on opens
  (\(f.\mathtt{opensFunctor}\)) --- cannot be supplied by such an object without a
  monoidal-closed structure on sheaves of modules, which this development deliberately does
  \emph{not} carry (cf.\ \cref{rem:scheme_modules_monoidal_off_path}). The atom must
  therefore be constructed directly.

  Concretely (\cref{lem:slice_dual_transport}), the forward component of
  \(\mathtt{sliceDualTransport}\,f\,M\,V\) at \(W \le V\) is the categorical image
  \((\mathtt{restrictScalars}\,\beta_W).\mathrm{map}\,(\varphi_{\bullet})\) of the dual-section
  component of \(\varphi\) at the \(f\)-image of \(W\), reindexed across
  \(f.\mathtt{opensFunctor}\) by restriction of scalars along \(\beta_W\), post-composed with
  the leg-(B) codomain swap \(\mathtt{inv}(\varepsilon(\mathtt{restrictScalars}\,\beta_W))\).
  Leg (A) is built categorically via \((\mathtt{restrictScalars}\,\beta_W).\mathrm{map}\):
  a genuine change-of-rings functor acts here, so the reindexing must be applied as a
  functor map rather than pointwise. Leg (B)'s unit comparison \(\varepsilon\) is invertible
  because \(\beta_W\) is a ring isomorphism, as detailed in \cref{lem:slice_dual_transport}.
  Because \(\mathrm{Opens}\,Y\) is a thin poset, the
  naturality squares commute by \(\mathtt{Subsingleton.elim}\). The forward and inverse maps
  identify the slice-Hom value over \(\mathtt{Over}\,(fV)\) with the slice-Hom value over
  \(\mathtt{Over}\,V\), \(\mathcal{O}_Y(V)\)-linearly.

  The Step-4 residual is then \(\mathtt{PresheafOfModules.isoMk}\) applied \emph{directly} to
  the family \(V \mapsto \mathtt{sliceDualTransport}\,f\,M\,V\) --- which already packages
  \emph{both} legs (A) and (B) --- with the outer naturality square thin-poset-trivial. No
  separate leg-(B) step is interposed in \(\mathtt{isoMk}\): the per-section atom of
  \cref{lem:slice_dual_transport} carries the full leg-(A)\(\circ\)(B) content. The outer
  \(\mathtt{isoMk}\) naturality square follows from the per-section family naturality
  \(\mathtt{sliceDualTransport.naturality}\) of \cref{lem:slice_dual_transport}: the
  \(\varepsilon\)-naturality is established first at the level of individual sections, and the
  outer presheaf-morphism naturality of \(\mathtt{isoMk}\) then follows.

  The residual of stage~H1 (``\(\mathtt{pushforward}\,\beta\) commutes with the dual'') is
  the sectionwise composite of leg (A) followed by leg (B), packaged over the slice with
  poset-thin naturality; the full Step-4 isomorphism is then the H1 isomorphism followed by
  this composite. The slice-reindexing coherence is \emph{light}, because \(\mathrm{Opens}\)
  is a thin poset and the relevant squares commute by \(\mathtt{Subsingleton.elim}\),
  whereas H1 and legs (A) and (B) carry the substantive content.

  It is worth recording \emph{why} the strong-monoidal restriction
  (\texttt{presheaf\_pushforward\_adj\_substrate}) has no direct dual analogue here. For
  the tensor product the transport \emph{collapsed to leg (B) alone}: because \(\otimes\)
  is sectionwise, pushing a tensor forward reduced, section by section, to a single
  ground-ring tensor equivalence, with no domain slice left to re-index. The dual's
  non-sectionwise nature removes that collapse and forces the additional leg (A): the
  domain \(\mathtt{restr}_V(f_*M)\) genuinely lives over a varying slice and must be
  transported before the ground ring can be reconciled.

  We caution that the naive ``fixed-value-category slice equivalence'' of
  \texttt{open\_immersion\_slice\_sheaf\_equiv} is \emph{not} applicable to leg (A). That instrument
  operates in a fixed value category over a fixed base ring, whereas the residual of
  leg (A) lives at the \emph{presheaf} level, with a per-section ring \(\mathcal{O}_Y(V)\)
  that varies with \(V\) and a finer slicing than the sheaf-level equivalence resolves;
  leg (A) must therefore be built directly at the presheaf-of-modules / varying-ring level.

  Composing legs (A) and (B) gives, for each \(V \subseteq U\), an
  \(\mathcal{O}_Y(V)\)-linear isomorphism between the two values; assembling these into a
  presheaf-of-modules natural isomorphism and sheafifying through the sheaf-level dual of
  \cref{lem:internal_hom_isSheaf} resolves the H1 residual. Precomposing with the stage-H1
  rewrite then yields the stated isomorphism
  \((\mathtt{dual}\,M)|_f \cong \mathtt{dual}(M|_f)\).

  \emph{The inverse-uniqueness shortcut is not viable.} One might hope to derive the
  isomorphism from the already-established \cref{lem:tensorobj_restrict_iso} by uniqueness of
  monoidal inverses --- the dual being the \(\otimes\)-inverse of an invertible module and
  restriction commuting with \(\otimes\) --- sidestepping leg (A). The abstract statements of
  this idiom (\(\mathtt{hasRightDualOfEquivalence}\), \(\mathtt{rightDualIso}\)) require,
  however, four structures the present setting lacks: a registered \(\mathtt{MonoidalCategory}\)
  instance on \(\mathtt{PresheafOfModules}\) (the monoidal structure here is carried entirely
  by hand), a \emph{strong} (not merely lax/oplax) monoidal restriction functor, restriction
  being an \emph{equivalence} (false for a non-surjective open immersion, which is a
  localization), and an exact pairing \(\mathtt{ExactPairing}\,M\,(\mathtt{dual}\,M)\) (which
  exists only for dualizable \(M\), whereas the lemma is stated for general \(M\)). Even
  restricted to the invertible consumer this would require building a coevaluation and
  verifying the zig-zag identities --- strictly \emph{more} infrastructure than leg (A) --- so
  the construction commits to legs (A) and (B) above.
\end{proof}

\begin{lemma}\leanok
[Dual of the structure sheaf]
  \label{lem:dual_unit_iso}
  \lean{AlgebraicGeometry.Scheme.Modules.dual_unit_iso}
  \uses{lem:internal_hom_eval, lem:internal_hom_isSheaf, lem:tensorobj_unit_iso}
  Let \(Y\) be a scheme. The sheaf-level dual of the structure sheaf is canonically the
  structure sheaf itself,
  \[
    \mathtt{dual}\,\mathcal{O}_Y \;\xrightarrow{\ \sim\ }\; \mathcal{O}_Y,
  \]
  obtained by sheafifying the presheaf-level ``evaluation at \(1\)'' isomorphism
  \(\mathcal{H}om(\mathbf{1}, \mathbf{1}) \cong \mathbf{1}\): every endomorphism of the
  unit presheaf is multiplication by its value at the global section \(1\), so evaluation
  at \(1\) identifies the internal hom of the unit into itself with the unit.
\end{lemma}

\begin{proof}
  \uses{lem:internal_hom_eval, lem:internal_hom_isSheaf, def:presheaf_dual_unit_iso}
  \leanok
  At the presheaf level the dual of the unit \(\mathbf{1} = \mathcal{O}_Y\) is the internal
  hom \(\mathcal{H}om(\mathbf{1}, \mathbf{1})\), whose value on an open \(U\) is the module
  of \(\mathcal{O}(U)\)-linear endomorphisms of \(\mathbf{1}|_U\). Each such endomorphism is
  determined by its value at the global section \(1\) and equals multiplication by that
  value; hence \emph{evaluation at \(1\)} is itself the sectionwise isomorphism
  \(\mathcal{H}om(\mathbf{1}, \mathbf{1}) \cong \mathbf{1}\) of presheaves of modules, with
  inverse the scalar-action map \(r \mapsto (\,\cdot\,r)\). (No left unitor is needed: the
  identification is the direct evaluation-at-\(1\) / global-scalar-multiplication
  isomorphism, not the composite through \(\mathbf{1} \otimes \mathbf{1} \cong \mathbf{1}\).)
  Sheafifying through the sheaf-level dual of \cref{lem:internal_hom_isSheaf} yields
  \(\mathtt{dual}\,\mathcal{O}_Y \cong \mathcal{O}_Y\). The scheme-level alias of this
  presheaf evaluation-at-\(1\) isomorphism is
  \(\mathtt{presheafDualUnitIso}\) (\cref{def:presheaf_dual_unit_iso}).
\end{proof}

\begin{lemma}





\leanok
[The dual of a line bundle is a line bundle]
  \label{lem:dual_isLocallyTrivial}
  \lean{AlgebraicGeometry.Scheme.Modules.dual_isLocallyTrivial}
  \uses{lem:internal_hom_isSheaf, lem:dual_restrict_iso, lem:dual_unit_iso, def:scheme_modules_dual_iso_of_iso, lem:restrictscalars_ringiso_dualequiv}
  % SOURCE: [Stacks Project], Tag 01CR (Modules on Ringed Spaces, \S Invertible
  % modules), Lemma~lemma-constructions-invertible (item~2)
  % (read from references/stacks-modules.tex, L4200--L4213).
  % SOURCE QUOTE:
  % "\item If $\mathcal{L}$ is an invertible $\mathcal{O}_X$-module, then so is
  % $\SheafHom_{\mathcal{O}_X}(\mathcal{L}, \mathcal{O}_X)$ and the evaluation map
  % $\mathcal{L} \otimes_{\mathcal{O}_X}
  % \SheafHom_{\mathcal{O}_X}(\mathcal{L}, \mathcal{O}_X) \to \mathcal{O}_X$
  % is an isomorphism."
  \textit{Source: [Stacks Project], Tag~01CR,
  Lemma~\texttt{lemma-constructions-invertible} (item~2).}
  Let \(X\) be a scheme and \(L \in \Scheme.\mathtt{Modules}\,X\) with
  \(\mathtt{LineBundle.IsLocallyTrivial}\,L\). Then the sheaf-level dual
  \(\mathtt{dual}\,L\) of \cref{lem:internal_hom_isSheaf} is again locally trivial of
  rank one, i.e.\ \(\mathtt{LineBundle.IsLocallyTrivial}\,(\mathtt{dual}\,L)\).
\end{lemma}

\begin{proof}
  \uses{lem:internal_hom_isSheaf, lem:dual_restrict_iso, lem:dual_unit_iso, def:scheme_modules_dual_iso_of_iso, lem:restrictscalars_ringiso_dualequiv}
  \leanok
  Local triviality of \(L\) furnishes, around each point, an affine open \(U\) with a
  trivialisation \(L|_U \cong \mathcal{O}_U\); writing \(f : U \hookrightarrow X\) for the
  associated open immersion, it suffices to exhibit an isomorphism
  \((\mathtt{dual}\,L)|_U \cong \mathcal{O}_U\). This is the following three-step chain:
  \[
    (\mathtt{dual}\,L)\big|_f
      \;\xrightarrow{\ \sim\ }\; \mathtt{dual}\bigl(L|_f\bigr)
      \;\xrightarrow{\ \sim\ }\; \mathtt{dual}\,\mathcal{O}_U
      \;\xrightarrow{\ \sim\ }\; \mathcal{O}_U.
  \]
  \begin{enumerate}
    \item The first isomorphism is the restriction-compatibility isomorphism
      \cref{lem:dual_restrict_iso}: restriction along the open immersion \(f\) commutes
      with the formation of the dual. This is the C-bridge proper; its substantive content
      is the sectionwise ring-iso transport of the module structure
      (\cref{lem:restrictscalars_ringiso_dualequiv}), as detailed there.
    \item The second isomorphism is the dual applied to the trivialisation
      \(L|_U \cong \mathcal{O}_U\), via the contravariant functoriality of the sheaf-level
      dual in isomorphisms (\(\mathtt{dualIsoOfIso}\),
      \cref{def:scheme_modules_dual_iso_of_iso}): an isomorphism \(L|_U \cong \mathcal{O}_U\)
      induces \(\mathtt{dual}(L|_U) \cong \mathtt{dual}\,\mathcal{O}_U\).
    \item The third isomorphism trivialises the dual of the unit,
      \(\mathtt{dual}\,\mathcal{O}_U \cong \mathcal{O}_U\) (\cref{lem:dual_unit_iso}):
      the dual of the structure sheaf is the structure sheaf itself, the comparison being
      evaluation at \(1 \in \mathcal{O}_U\), under which
      \(\mathcal{H}om_{\mathcal{O}_U}(\mathcal{O}_U, \mathcal{O}_U) \cong \mathcal{O}_U\).
  \end{enumerate}
  Composing the three exhibits \((\mathtt{dual}\,L)|_U \cong \mathcal{O}_U\). Hence
  \(\mathtt{dual}\,L\) is locally trivial of rank one, as required by item~2 of Stacks
  Lemma~\texttt{lemma-constructions-invertible} (tag 01CR).
\end{proof}

\begin{lemma}


\leanok
[Local isomorphisms of \(\mathcal{O}_X\)-modules glue to a global isomorphism]
  \label{lem:isiso_of_isiso_restrict}
  \lean{AlgebraicGeometry.Scheme.Modules.isIso_of_isIso_restrict}
  Let \(X\) be a scheme and \(\varphi \colon M \to N\) a morphism in
  \(\Scheme.\mathtt{Modules}\,X = \mathtt{SheafOfModules}\,X.\mathtt{ringCatSheaf}\).
  Suppose that for every point \(x \in X\) there is an open neighbourhood \(U_x\) of
  \(x\) such that the restriction of \(\varphi\) to \(U_x\) --- the image
  \((\mathtt{restrictFunctor}\,(U_x).\mathtt{i}).\mathtt{map}\,\varphi\) under the
  restriction functor along the open immersion \((U_x).\mathtt{i}\colon U_x \to X\) ---
  is an isomorphism in \(\Scheme.\mathtt{Modules}\,U_x\). Then \(\varphi\) is an
  isomorphism.
\end{lemma}

\begin{proof}



\leanok
  Being an isomorphism for a morphism of sheaves of \(\mathcal{O}_X\)-modules can be
  checked on the underlying morphism of sheaves of abelian groups, and there it can be
  checked stalkwise. Fix a point \(x \in X\). Choose a preimage \(x' \in U_x\) of \(x\)
  under the open immersion \((U_x).\mathtt{i}\). By hypothesis the restriction
  \((\mathtt{restrictFunctor}\,(U_x).\mathtt{i}).\mathtt{map}\,\varphi\) is an
  isomorphism, hence so is its image under any functor --- in particular under the
  stalk-at-\(x'\) functor on the underlying ab-sheaves. The natural isomorphism
  ``restriction commutes with stalks'' (\(\mathtt{restrictStalkNatIso}\)) identifies
  this with the stalk of (the underlying ab-sheaf morphism of) \(\varphi\) at the image
  point \(x\); transporting along it shows that stalk is an isomorphism. As this holds
  at every \(x\), and a morphism of sheaves of abelian groups all of whose stalks are
  isomorphisms is itself an isomorphism (stalkwise iso-detection,
  \(\mathtt{isIso\_of\_stalkFunctor\_map\_iso}\) on \(\mathtt{TopCat.Sheaf}\)), the
  underlying ab-sheaf morphism of \(\varphi\) is an isomorphism. Finally the forgetful
  functor \(\Scheme.\mathtt{Modules.toPresheaf}\) reflects isomorphisms, so \(\varphi\)
  itself is an isomorphism.
\end{proof}

\begin{lemma}\leanok
[Down-set image identity \(\iota_!(\iota^{-1}V) = V\)]
  \label{lem:image_preimage_of_le}
  \lean{AlgebraicGeometry.Scheme.Modules.image_preimage_of_le}
  For opens \(V \le W\) of a scheme \(X\), the image under the open immersion \(W.\iota\) of
  the preimage of \(V\) is \(V\) again: \(W.\iota\,{}_!(W.\iota^{-1}(V)) = V\). This is the
  equality powering the \(\mathtt{eqToHom}\)-conjugations of \(\mathtt{homLocalSection}\)
  (\texttt{scheme\_modules\_hom\_local\_section}) and of the reverse slice transport
  (\cref{lem:slice_dual_transport_inv}).
\end{lemma}
\begin{proof}




\leanok
  Proved directly in Lean --- the image of the preimage is the intersection with the
  (open) range of the immersion, which equals \(V\) since \(V \le W\) lies in that range.
\end{proof}

\begin{definition}




\leanok
[Local section of the hom-sheaf from a compatible family]
  \label{def:scheme_modules_hom_local_section}
  \lean{AlgebraicGeometry.Scheme.Modules.homLocalSection}
  \uses{lem:image_preimage_of_le}
  For a family of open subsets \(\{U_i\}_{i \in \iota}\) of a scheme \(X\) and a compatible
  family of morphisms \(f_i : M|_{U_i} \to N|_{U_i}\) of \(\mathcal{O}_X\)-modules, the
  local section `homLocalSection` packages the section of the hom-sheaf over \(U_i\) whose
  value on a smaller open \(V \le U_i\) is obtained by conjugating the component of \(f_i\)
  along the canonical down-set identity \(\iota_!(\iota^{-1}V) = V\). The naturality in the
  slice is automatic on the thin poset of opens once the down-set identity is in place.
\end{definition}

\begin{definition}




\leanok
[Top section to global hom]
  \label{def:scheme_modules_top_section_to_hom}
  \lean{AlgebraicGeometry.Scheme.Modules.topSectionToHom}
  A section of the hom-sheaf over the terminal open \(\top\) determines a global
  morphism of presheaves \(F \to G\) by the standard sections-to-morphisms equivalence:
  `topSectionToHom` is the transport of that section through the canonical
  hom-sections equivalence, so that a compatible top section yields a global morphism.
\end{definition}

\begin{lemma}




\leanok
[Sectionwise value of `topSectionToHom`]
  \label{lem:scheme_modules_top_section_to_hom_app}
  \lean{AlgebraicGeometry.Scheme.Modules.topSectionToHom_app}
  \uses{def:scheme_modules_top_section_to_hom}
  Let \(s\) be a section over the terminal open. Then the component of the induced morphism
  at an open \(W\) is exactly the \(W\)-component of \(s\), evaluated at the unique map
  \(W \to \top\). In particular, the recovered morphism is sectionwise the same as the input
  section, just reindexed along the terminal inclusion.
\end{lemma}

\begin{definition}


\leanok
[Promote an \(\mathcal{O}_X\)-linear presheaf morphism to a module morphism]
  \label{def:scheme_modules_homMk}
  \lean{AlgebraicGeometry.Scheme.Modules.homMk}
  A presheaf morphism of the underlying abelian-group sheaves that is sectionwise
  \(\mathcal{O}_X\)-linear packages as a morphism of \(\mathcal{O}_X\)-modules. This
  is the promotion step that turns the glued \(Ab\)-valued morphism into an actual
  module morphism.
\end{definition}

\begin{lemma}




\leanok
[Compatible local morphisms glue to a global morphism]
  \label{lem:sheafofmodules_hom_of_local_compat}
  \lean{AlgebraicGeometry.Scheme.Modules.homOfLocalCompat}
  \uses{def:scheme_modules_homMk, def:scheme_modules_hom_local_section, def:scheme_modules_top_section_to_hom, lem:scheme_modules_top_section_to_hom_app, lem:image_preimage_of_le}
  Given a cover \(\{U_i\}\) of \(X\), a family of local morphisms \(f_i : M|_{U_i} \to N|_{U_i}\)
  that agree on overlaps glues to a unique global \(\mathcal{O}_X\)-linear morphism \(M \to N\).
  The proof glues the associated hom-sheaf sections using the sheaf property, converts the
  glued top section to a global morphism, and checks sectionwise \(\mathcal{O}_X\)-linearity on
  each open of the cover.
\end{lemma}

\begin{lemma}



\leanok
[The glued morphism restricts to its local datum]
  \label{lem:hom_of_local_compat_restrict}
  \lean{AlgebraicGeometry.Scheme.Modules.homOfLocalCompat_restrictFunctor_map}
  \uses{lem:sheafofmodules_hom_of_local_compat, lem:scheme_modules_top_section_to_hom_app}
  In the situation of \cref{lem:sheafofmodules_hom_of_local_compat} --- a cover \(\{U_i\}\)
  of \(X\) and a compatible family \(f_i : M|_{U_i} \to N|_{U_i}\) --- the restriction of the
  glued global morphism back to a cover member recovers its local datum:
  \[
    (\mathtt{restrictFunctor}\,U_i).\mathtt{map}\,(\mathtt{homOfLocalCompat}\,f) \;=\; f_i .
  \]
  This is the companion gluing-value identity needed to apply the B-connector
  \cref{lem:isiso_of_isiso_restrict}: when every \(f_i\) is an isomorphism, this identity
  transports that to the restriction of the glued morphism, so the global morphism is
  locally an isomorphism and hence an isomorphism.
\end{lemma}
\begin{proof}



\leanok
  By construction, \(\mathtt{homOfLocalCompat}\) produces its global morphism by gluing
  the local sections of the hom-sheaf associated to each \(f_i\) along the cover; the sheaf
  gluing axiom states that the resulting global section restricts on \(U_i\) to the section
  corresponding to \(f_i\). Converting the global section back to a morphism is sectionwise
  faithful (\cref{lem:scheme_modules_top_section_to_hom_app}), and the section-to-morphism
  construction preserves the underlying section, so the restriction of the global morphism to
  \(U_i\) is sectionwise equal to \(f_i\) --- hence equal as a morphism.
\end{proof}

\begin{remark}
[How the dual infrastructure discharges the \(\otimes\)-inverse]
  \label{rem:dual_discharges_inverse}
  \uses{lem:tensorobj_inverse_invertible, lem:dual_isLocallyTrivial, lem:sheafofmodules_hom_of_local_compat, lem:hom_of_local_compat_restrict, lem:isiso_of_isiso_restrict, lem:tensorobj_restrict_iso, lem:tensorobj_unit_iso, lem:tensorobjisoofiso_trans, lem:tensorobjisoofiso_refl, lem:dualisoofiso_trans, lem:dualisoofiso_refl, lem:trivialisation_restrict_compat, lem:tensorobj_unit_self_duality_collapse}
  Once the dual package \cref{lem:dual_isLocallyTrivial} and the A-bridge
  \cref{lem:sheafofmodules_hom_of_local_compat} are available,
  \cref{lem:tensorobj_inverse_invertible} is discharged by taking \(L^{-1} := \mathtt{dual}\,L\)
  (\cref{lem:internal_hom_isSheaf}), which is a line
  bundle by \cref{lem:dual_isLocallyTrivial}. The global contraction \(\varepsilon_L
  \colon L \otimes_X L^{-1} \to \mathcal{O}_X\) is \emph{not} obtained by sheafifying
  the presheaf-level evaluation; it is instead built by \emph{gluing the canonical local
  contraction morphisms}. On each trivialising open \(U\), the
  restriction-compatibility iso \cref{lem:tensorobj_restrict_iso} carries
  \((L \otimes_X L^{-1})|_U\) to \(\mathcal{O}_U \otimes_U \mathcal{O}_U\), and the
  canonical local contraction \(\mathcal{O}_U \otimes_U \mathcal{O}_U \to \mathcal{O}_U\)
  of the free rank-one module with its dual is the left unitor
  \cref{lem:tensorobj_unit_iso}, an isomorphism. These local contractions are canonical,
  hence agree on overlaps, and so glue --- by
  \cref{lem:sheafofmodules_hom_of_local_compat} --- to a single global morphism
  \(\varepsilon_L \colon L \otimes_X L^{-1} \to \mathcal{O}_X\). Since ``being an
  isomorphism'' is local on \(X\) and \(\varepsilon_L\) restricts to an isomorphism on
  each \(U\), it is a global isomorphism \(L \otimes_X L^{-1} \xrightarrow{\sim}
  \mathcal{O}_X\) by the B-connector \cref{lem:isiso_of_isiso_restrict}.

  This descent construction avoids the ``sheafify-the-evaluation'' route, which would
  require commuting the sheafification unit past tensor-stalk whiskering; the gluing strategy
  computes no tensor stalk anywhere.

  Two conditions make the gluing of \cref{lem:sheafofmodules_hom_of_local_compat} applicable.
  First, \emph{overlap compatibility}: on each overlap \(V \subseteq U_i \cap U_j\) the
  sectionwise restrictions of the two local contractions \(f_i\) and \(f_j\) to \(V\) must
  agree. They do because each \(f_x\) is the canonical evaluation
  \(L \otimes L^{-1} \to \mathcal{O}\), independent of the chosen trivialisation: the dual
  trivialisation \(e^N_x\) is built from the primal \(e^M_x\) via
  \cref{lem:tensorobj_restrict_iso}, \(\mathtt{tensorObjIsoOfIso}\), \(\mathtt{dualIsoOfIso}\),
  and \cref{lem:tensorobj_unit_iso}, so the transition isomorphism \(g_{ij}\) entering
  \(e^M_x\) and its inverse \(g_{ij}^{-1}\) entering \(e^N_x\) cancel through
  \(\mathtt{tensorObjIsoOfIso}\); this cancellation is a genuine computation of abelian-group
  homomorphisms. Concretely, the mechanism realising the cancellation is the
  \emph{functoriality of the iso-transport operators}: bifunctoriality of
  \(\mathtt{tensorObjIsoOfIso}\) (\cref{lem:tensorobjisoofiso_trans}, \cref{lem:tensorobjisoofiso_refl})
  and contravariant functoriality of \(\mathtt{dualIsoOfIso}\) (\cref{lem:dualisoofiso_trans},
  \cref{lem:dualisoofiso_refl}). Pushing both legs to a pure tensor \(a \otimes b\), these laws
  rewrite \(e^N_x\) so that the \(g_{ij}\) inside \(e^M_x\) meets \(g_{ij}^{-1} = \mathtt{dualIsoOfIso}\,
  (g_{ij})\) inside \(e^N_x\) as adjacent factors, where \(\mathtt{tensorObjIsoOfIso}\,g_{ij}\,
  (\mathtt{dualIsoOfIso}\,g_{ij})\) followed by \(\mathtt{tensorObj\_unit\_iso}\) reduces to
  \(\mathtt{tensorObj\_unit\_iso}\) alone (the unit self-duality evaluation collapse); the
  iso-level identities above thereby reduce sectionwise to the required cocycle equality.
  Second, the \emph{local-iso} step: by the gluing-value identity
  \cref{lem:hom_of_local_compat_restrict}, the restriction of \(\varepsilon_L\) to \(U_x\)
  equals the local isomorphism \(f_x\), so \cref{lem:isiso_of_isiso_restrict} upgrades
  \(\varepsilon_L\) to a global isomorphism.
\end{remark}

The two mechanism steps named in the overlap-compatibility argument of
\cref{rem:dual_discharges_inverse} are isolated below as formal target lemmas, so the
cocycle obligation (residual~A) of \cref{lem:tensorobj_inverse_invertible} has a named
formalization target for each. The iso-transport operators
\(\mathtt{tensorObjIsoOfIso}\) and \(\mathtt{dualIsoOfIso}\) and the unit comparison
\(\mathtt{tensorObj\_unit\_iso}\) are the existing helpers
(\cref{lem:tensorobjisoofiso_trans}, \cref{lem:dualisoofiso_trans},
\cref{lem:tensorobj_unit_iso}).

The five restriction-naturality squares constituting \cref{lem:trivialisation_restrict_compat}
are isolated here, \emph{in chain order}, as named, independently-provable sub-lemmas; the target
then telescopes them. Throughout, fix opens \(V \subseteq U\) of \(X\) and write
\(j : V \hookrightarrow U\) for the open immersion \(\Scheme.\mathtt{Hom.resLE}\,(\id_X)\,U\,V\),
the chart morphism with \(j \mathbin{;} \iota_U = \iota_V\); restriction along \(j\) is the
functor \((-)|_j = \mathtt{restrict}\,j\). The two a-priori-distinct opens
\(\iota_U^{-1}(V)\) and \(\iota_V^{-1}(V)\) are identified as \(V\) only through the
direct-image equality \(\iota_U\,''(\iota_U^{-1}V) = V = \iota_V\,''(\iota_V^{-1}V)\)
(\cref{lem:image_preimage_of_le}); write \(\varrho\) for the canonical equality-of-opens
reindexing isomorphism it induces, sitting on \emph{both} flanks of each constituent's square.

The four squares S2/S3/S4a/S4b are \emph{not} four independent chases: each is a
\emph{base-change naturality} square for one and the same comparison morphism --- the
sheaf-level pullback--tensor map \(\mathtt{pullbackTensorMap}\,f\)
(\cref{lem:pullback_tensor_map}) --- transported from the pullback world onto the
restriction world along the comparison \(\mathtt{restrictFunctorIsoPullback}\,f :
\mathtt{restrict}\,f \cong \mathtt{pullback}\,f\) (\cref{def:restrictfunctorisopullback})
and reindexed by \(\varrho\). For an open immersion \(f\) the comparison
\(\mathtt{pullbackTensorMap}\,f\) is invertible
(\cref{lem:pullback_tensor_map_isiso_open_immersion}, a public witness), so every
transported square is a genuine isomorphism square. Crucially there is \emph{no} monoidal
structure on the sheaf categories \(\Scheme.\mathtt{Modules}\) in play: the only monoidal
data used lives one level down, on presheaves of modules, where the pullback is oplax
monoidal with oplax tensorator \(\delta = \mathtt{pullbackTensorMap}\), and the whole
coherence is governed by uniqueness of left adjoints rather than by a sheaf-level
\(\mathtt{MonoidalCategory}\).

Two supporting bridges organise the squares. \emph{Bridge B1}
(\cref{lem:tensorobj_restrict_iso_eq_pullback_tensor_map}) identifies
\(\mathtt{tensorObj\_restrict\_iso}\,f\) with the transported comparison: it is the
\emph{direct} conjugate factorisation
\(\mathtt{tensorObj\_restrict\_iso}\,f = \mathtt{restrictFunctorIsoPullback}\,f \mathbin{;}
\mathtt{asIso}(\mathtt{pullbackTensorMap}\,f) \mathbin{;} (\text{per-leg reindex})\),
proved without any sheaf-level monoidal transport --- both sides share the
\(\mathtt{restrictFunctorIsoPullback}\,f \mathbin{;} \mathtt{sheafificationCompPullback}\,f\)
prefix, and after cancelling it the residue is the presheaf-level identity that the
sheafified oplax tensorator \(\delta(\mathtt{pullback}\,\varphi')\) equals the
left-adjoint-transported tensorator building the tail of \(\mathtt{tensorObj\_restrict\_iso}\),
discharged by the \(\delta\)-conjugation identity
(\cref{lem:pushforward_mu_appiso_collapse}) together with uniqueness of left adjoints.
\emph{Bridge B2} (\cref{lem:restrictfunctorisopullback_comp_compat}) is the reindex
coherence: \(\mathtt{restrictFunctorIsoPullback} : \mathtt{restrict}\,f \cong \mathtt{pullback}\,f\)
(the left-adjoint uniqueness isomorphism) is pseudonatural across \(j \mathbin{;} \iota_U\),
identifying \(\varrho = \mathtt{restrictCompReindex}\,j\) (\cref{def:restrictcompreindex}) with
\(\mathtt{pullbackComp}\) conjugated by \(\mathtt{restrictFunctorIsoPullback}\) on each factor; its
closed form is \(\mathtt{leftAdjointUniq\_trans}\) threaded by \(\mathtt{mateEquiv\_vcomp}\).
With B1 in hand each of S2/S3/S4a/S4b reduces to pure base-change bookkeeping: substitute
B1 to turn both the \(U\)- and \(V\)-built \(\mathtt{tensorObj\_restrict\_iso}\) into
\(\mathtt{asIso}(\mathtt{pullbackTensorMap})\) conjugates, whereupon the morphism equation
is the already-proven composition law \(\mathtt{pullbackTensorMap\_restrict}\)
(\cref{lem:pullback_tensor_map_basechange}) together with naturality
\(\mathtt{pullbackTensorMap\_natural}\) (\cref{lem:pullback_tensor_map_natural}),
transported back along \(\mathtt{restrictFunctorIsoPullback}\) and reindexed by \(\varrho\).
The remaining global-unit square S4c rests on B2 alone, as the unit composition law
\(\mathtt{pullbackObjUnitToUnit\_comp}\) (\cref{lem:pullbackObjUnitToUnit_comp}) re-wrapped by the
reindex coherence. The constituents are therefore landed in the dependency order
\[
  \mathrm{B2} \;\to\; \mathrm{B1} \;\to\; \mathrm{S4c} \;\to\; \mathrm{S2}
    \;\to\; \mathrm{S4b} \;\to\; \mathrm{S3} \;\to\; \mathrm{S4a}.
\]

\begin{definition}


\leanok
[Restriction-composite reindexing \(\varrho\)]
  \label{def:restrictcompreindex}
  \lean{AlgebraicGeometry.Scheme.Modules.restrictCompReindex}
  \uses{lem:image_preimage_of_le}
  Let \(j : V \hookrightarrow U\) be an open immersion of opens of \(X\) with
  \(j \mathbin{;} \iota_U = \iota_V\), and let \(A \in \Scheme.\mathtt{Modules}\,X\). The
  \emph{restriction-composite reindexing} is the isomorphism
  \[
    \varrho_A := \mathtt{restrictCompReindex}\,j\,A :
      A|_{\iota_V} \;\cong\; (A|_{\iota_U})\big|_j ,
  \]
  defined as the composite \((\mathtt{restrictFunctorCongr}\,h_{j\iota})^{-1} \mathbin{;}
  \mathtt{restrictFunctorComp}\,j\,\iota_U\): the (Mathlib) restriction-composition isomorphism
  \(\mathtt{restrictFunctorComp}\) precomposed with the congruence along
  \(j \mathbin{;} \iota_U = \iota_V\). It is the \(\varrho\) appearing on both flanks of every
  square below.
\end{definition}

\begin{definition}


\leanok
[Unit-restriction identification \(u_\iota\)]
  \label{def:unitrestrictiso}
  \lean{AlgebraicGeometry.Scheme.Modules.unitRestrictIso}
  \uses{lem:pullback_unit_iso}
  For an open immersion \(f : Y \hookrightarrow X\) the restriction of the global unit to \(Y\) is
  the structure sheaf of \(Y\): the \emph{unit-restriction identification} is
  \[
    u_\iota(f) := \bigl(\mathtt{restrictFunctorIsoPullback}\,f\bigr).\mathrm{app}\,\mathcal{O}_X
      \mathbin{;} \mathtt{pullbackUnitIso}\,f :
      \mathcal{O}_X|_f \;\cong\; \mathcal{O}_Y ,
  \]
  the comparison \(\mathtt{restrictFunctorIsoPullback}\) into the pullback world followed by the
  structure-sheaf comparison \(\mathtt{pullbackUnitIso}\) (\cref{lem:pullback_unit_iso}). It is the
  unit identification on the chart-scheme side of S4a/S4b and the final leg \(u_\iota^{-1}\) of S4c.
\end{definition}

\begin{lemma}


\leanok
[Bridge B2: reindex coherence of \(\mathtt{restrictFunctorIsoPullback}\)]
  \label{lem:restrictfunctorisopullback_comp_compat}
  \lean{AlgebraicGeometry.Scheme.Modules.restrictFunctorIsoPullback_comp_compat}
  \uses{def:restrictcompreindex, lem:chart_isiso, lem:image_preimage_of_le, lem:restrictfunctorisopullback_comp_compat_hom}
  Let \(j : V \hookrightarrow U\) with \(j \mathbin{;} \iota_U = \iota_V\). The natural isomorphism
  \(\mathtt{restrictFunctorIsoPullback} : \mathtt{restrict}\,f \cong \mathtt{pullback}\,f\) is
  pseudonatural across the composite \(j \mathbin{;} \iota_U\): the restriction-world reindexing
  \(\varrho = \mathtt{restrictCompReindex}\,j\) (\cref{def:restrictcompreindex}) corresponds, under
  \(\mathtt{restrictFunctorIsoPullback}\) applied on each factor, to the pullback-world reindexing
  \(\mathtt{pullbackComp}\,j\,\iota_U\), threaded by the equality-of-opens transport along
  \(j \mathbin{;} \iota_U = \iota_V\) (\cref{lem:image_preimage_of_le}).
\end{lemma}
\begin{proof}
  \uses{def:restrictcompreindex, lem:chart_isiso, lem:adjunction_leftadjointuniq_trans, lem:restrictfunctorisopullback_comp_compat_hom, lem:unit_leftadjointuniq_hom_app, lem:conjugate_equiv_restrictfunctorcomp_inv, lem:image_preimage_of_le}
  This is a general coherence of the pseudofunctor comparison
  \(\mathtt{restrictFunctorIsoPullback}\). Both \(\mathtt{restrictFunctorComp}\) and
  \(\mathtt{pullbackComp}\) are the canonical isomorphism comparing a composite base change with the
  composite of base changes, and \(\mathtt{restrictFunctorIsoPullback}\) intertwines them by
  naturality of the comparison in each factor. The equality-of-opens transport of
  \(j \mathbin{;} \iota_U = \iota_V\) matches on the two sides because both reindexings are defined by
  transport along the \emph{same} equality of opens. The isomorphism-existence shadow of this
  identity is already used in \cref{lem:chart_isiso}; the present statement promotes it to an
  equality of morphisms. In closed form this equality holds at the level of natural transformations:
  it is exactly the hom-level factorisation
  \(\bigl(\mathtt{restrictFunctorIsoPullback}\,\iota_V\bigr).\mathrm{hom}.\mathrm{app}\,A
  = c_1 \mathbin{;} c_2 \mathbin{;} c_3 \mathbin{;} c_4 \mathbin{;} c_5 \mathbin{;} c_6\) of
  \cref{lem:restrictfunctorisopullback_comp_compat_hom}, evaluated at \(A\). That factorisation is
  obtained not from a horizontal-composite mate identity but by conjugating both sides onto the
  common right adjoint \(\mathtt{pushforward}\,\iota_V\) and distributing the conjugation
  \emph{leg-by-leg} through \(\mathtt{conjugateEquiv\_comp}\) (\cref{lem:conjugateequiv_comp}): every
  intermediate functor is a left adjoint \(X.\mathtt{Modules} \to V.\mathtt{Modules}\) of the common
  \(\mathtt{pushforward}\,\iota_V\), each leg maps to a known pushforward-world morphism, and their
  reversed product telescopes to the identity. The \(\mathrm{app}\,A\) statement of B2 then follows
  by transporting that natural-transformation identity across the unit triangle of
  \(\mathtt{restrictFunctorIsoPullback}\,\iota_V\) (\cref{lem:unit_leftadjointuniq_hom_app}),
  together with the transitivity \(\mathtt{leftAdjointUniq\_trans}\)
  (\cref{lem:adjunction_leftadjointuniq_trans}) aligning the two transports of the comparison.
\end{proof}

% ---------------------------------------------------------------------------
% B2 hom-level content and its per-leg conjugate decomposition
% ---------------------------------------------------------------------------
\begin{lemma}\leanok
[Bridge B2, hom level: leg-by-leg factorisation of \(\mathtt{restrictFunctorIsoPullback}\)]
  \label{lem:restrictfunctorisopullback_comp_compat_hom}
  \lean{AlgebraicGeometry.Scheme.Modules.restrictFunctorIsoPullback_comp_compat_hom}
  \uses{def:restrictcompreindex, lem:conjugateequiv_comp, lem:conjugate_equiv_restrictfunctorcomp_inv,
    lem:conjugateequiv_pullbackcomp_hom, lem:conjugateequiv_restrictfunctorisopullback_whiskerright,
    lem:conjugateequiv_restrictfunctorisopullback_whiskerleft, lem:conjugateequiv_reindexcongr,
    lem:pushforwardcomp_reindex_telescope, lem:unit_leftadjointuniq_hom_app, lem:image_preimage_of_le,
    lem:conjugateequiv_restrictfunctorisopullback_hom}
  Let \(j : V \hookrightarrow U\) with \(j \mathbin{;} \iota_U = \iota_V\) and let
  \(A \in \Scheme.\mathtt{Modules}\,X\). Evaluated at \(A\), the single restriction-to-pullback
  comparison map factors as the composite of the six reindex legs
  \[
    \bigl(\mathtt{restrictFunctorIsoPullback}\,\iota_V\bigr).\mathrm{hom}.\mathrm{app}\,A
      \;=\; c_1 \mathbin{;} c_2 \mathbin{;} c_3 \mathbin{;} c_4 \mathbin{;} c_5 \mathbin{;} c_6,
  \]
  where
  \(c_1 = \bigl((\mathtt{restrictFunctorCongr}\,h_{j\iota})^{-1}\bigr).\mathrm{hom}.\mathrm{app}\,A\),
  \(c_2 = (\mathtt{restrictFunctorComp}\,j\,\iota_U).\mathrm{hom}.\mathrm{app}\,A\),
  \(c_3 = (\mathtt{restrict}\,j)\bigl((\mathtt{restrictFunctorIsoPullback}\,\iota_U).\mathrm{hom}.\mathrm{app}\,A\bigr)\),
  \(c_4 = (\mathtt{restrictFunctorIsoPullback}\,j).\mathrm{hom}.\mathrm{app}\,\bigl((\mathtt{pullback}\,\iota_U).\mathrm{obj}\,A\bigr)\),
  \(c_5 = (\mathtt{pullbackComp}\,j\,\iota_U).\mathrm{hom}.\mathrm{app}\,A\) and
  \(c_6 = (\mathtt{pullbackCongr}\,h_{j\iota}).\mathrm{hom}.\mathrm{app}\,A\). This is the genuine
  natural-transformation content to which Bridge B2
  (\cref{lem:restrictfunctorisopullback_comp_compat}) reduces.
\end{lemma}
\begin{proof}
  \uses{def:restrictcompreindex, lem:conjugateequiv_comp, lem:conjugate_equiv_restrictfunctorcomp_inv,
    lem:conjugateequiv_pullbackcomp_hom, lem:conjugateequiv_restrictfunctorisopullback_whiskerright,
    lem:conjugateequiv_restrictfunctorisopullback_whiskerleft, lem:conjugateequiv_reindexcongr,
    lem:pushforwardcomp_reindex_telescope, lem:image_preimage_of_le,
    lem:conjugateequiv_restrictfunctorisopullback_hom}
  It suffices to prove the underlying identity of natural transformations
  \(\mathtt{restrictFunctor}\,\iota_V \Rightarrow \mathtt{pullback}\,\iota_V\); the \(\mathrm{app}\,A\)
  statement follows by evaluating the whisker legs \(c_3, c_4\) componentwise. Every functor in the
  right-hand chain is a functor \(X.\mathtt{Modules} \to V.\mathtt{Modules}\) and a \emph{left}
  adjoint of the common right adjoint \(\mathtt{pushforward}\,\iota_V\); hence both sides are
  determined by their image under the conjugation bijection
  \(\mathtt{conjugateEquiv}\,(\mathtt{pullbackPushforwardAdjunction}\,\iota_V)\,
  (\mathtt{restrictAdjunction}\,\iota_V)\) onto natural transformations
  \(\mathtt{pushforward}\,\iota_V \Rightarrow \mathtt{pushforward}\,\iota_V\), and it is enough to
  compare the two images.

  The left-hand image is \(\mathbf{1}_{\mathtt{pushforward}\,\iota_V}\): the comparison
  \(\mathtt{restrictFunctorIsoPullback}\,\iota_V\) is itself the left-adjoint-uniqueness comparison,
  and the conjugate of such a hom is the identity. On the right-hand side the chain factors through
  the seven intermediate left adjoints of \(\mathtt{pushforward}\,\iota_V\),
  \(G_0 = \mathtt{restrictAdjunction}\,\iota_V\),
  \(G_1 = \mathtt{restrictAdjunction}\,(j \mathbin{;} \iota_U)\),
  \(G_2 = (\mathtt{restrictAdjunction}\,\iota_U).\mathrm{comp}\,(\mathtt{restrictAdjunction}\,j)\),
  \(G_3 = (\mathtt{pullbackPushforwardAdjunction}\,\iota_U).\mathrm{comp}\,(\mathtt{restrictAdjunction}\,j)\),
  \(G_4 = (\mathtt{pullbackPushforwardAdjunction}\,\iota_U).\mathrm{comp}\,(\mathtt{pullbackPushforwardAdjunction}\,j)\),
  \(G_5 = \mathtt{pullbackPushforwardAdjunction}\,(j \mathbin{;} \iota_U)\),
  \(G_6 = \mathtt{pullbackPushforwardAdjunction}\,\iota_V\), so the conjugation distributes
  \emph{leg-by-leg} through the composition law \(\mathtt{conjugateEquiv\_comp}\)
  (\cref{lem:conjugateequiv_comp}); no horizontal-composite mate identity is required. Each leg is a
  known pushforward-world map:
  \begin{itemize}
    \item \(c_2 \mapsto (\mathtt{pushforwardComp}\,j\,\iota_U).\mathrm{hom}\) by the keystone
      \cref{lem:conjugate_equiv_restrictfunctorcomp_inv};
    \item \(c_5 \mapsto (\mathtt{pushforwardComp}\,j\,\iota_U).\mathrm{inv}\) by
      \cref{lem:conjugateequiv_pullbackcomp_hom};
    \item \(c_3, c_4 \mapsto\) the pushforward-whiskered units of \(\iota_U\) and \(j\) by
      \cref{lem:conjugateequiv_restrictfunctorisopullback_whiskerright} and
      \cref{lem:conjugateequiv_restrictfunctorisopullback_whiskerleft};
    \item \(c_1, c_6 \mapsto\) the pushforward congruences along
      \(j \mathbin{;} \iota_U = \iota_V\) by \cref{lem:conjugateequiv_reindexcongr}.
  \end{itemize}
  The resulting reversed product of pushforward-world legs telescopes to
  \(\mathbf{1}_{\mathtt{pushforward}\,\iota_V}\) by \cref{lem:pushforwardcomp_reindex_telescope}: the
  two \(\mathtt{pushforwardComp}\) legs are mutually inverse, and the unit and congruence legs cancel
  against the reindexing \(j \mathbin{;} \iota_U = \iota_V\) (\cref{lem:image_preimage_of_le}). Both
  conjugate images equal \(\mathbf{1}_{\mathtt{pushforward}\,\iota_V}\), so the two natural
  transformations agree.
\end{proof}

The telescope above is assembled from the following atomic per-leg conjugate computations, each a
small fail-fast obligation shared with the Bridge B1 crux
(\cref{lem:h1inv_app_eq_pullbackval_restrict}).

\begin{lemma}\leanok
[\(c_5\): conjugate of the pullback-composition hom]
  \label{lem:conjugateequiv_pullbackcomp_hom}
  \lean{AlgebraicGeometry.Scheme.Modules.conjugateEquiv_pullbackComp_hom}
  \uses{lem:conjugate_equiv_pullbackcomp_inv}
  Let \(f : X \hookrightarrow Y\), \(g : Y \hookrightarrow Z\) be open immersions. Under the
  conjugation between the composite pullback adjunction
  \((\mathtt{pullbackPushforwardAdjunction}\,g).\mathrm{comp}\,(\mathtt{pullbackPushforwardAdjunction}\,f)\)
  and the pullback adjunction of the composite, the hom of the pullback-composition isomorphism
  corresponds to the \emph{inverse} of the pushforward-composition isomorphism:
  \[
    \mathtt{conjugateEquiv}\;\dots\;(\mathtt{pullbackComp}\,f\,g).\mathrm{hom}
      \;=\; (\mathtt{pushforwardComp}\,f\,g).\mathrm{inv}.
  \]
\end{lemma}
\begin{proof}
  \uses{lem:conjugate_equiv_pullbackcomp_inv}
  The conjugation bijection is an additive-inverse-respecting equivalence on the relevant hom-groups
  only formally; concretely, apply it to the identity
  \((\mathtt{pullbackComp}\,f\,g).\mathrm{hom} \mathbin{;} (\mathtt{pullbackComp}\,f\,g).\mathrm{inv}
  = \mathbf{1}\) and use that the conjugate of \((\mathtt{pullbackComp}\,f\,g).\mathrm{inv}\) is
  \((\mathtt{pushforwardComp}\,f\,g).\mathrm{hom}\) (\cref{lem:conjugate_equiv_pullbackcomp_inv}).
  Since \(\mathtt{conjugateEquiv}\) sends the hom of an isomorphism to a map inverse to the image of
  its inverse, the image of \((\mathtt{pullbackComp}\,f\,g).\mathrm{hom}\) is the inverse of
  \((\mathtt{pushforwardComp}\,f\,g).\mathrm{hom}\), namely \((\mathtt{pushforwardComp}\,f\,g).\mathrm{inv}\).
\end{proof}

\begin{lemma}\leanok
[LHS-collapse: the conjugate of the comparison hom is the identity]
  \label{lem:conjugateequiv_restrictfunctorisopullback_hom}
  \lean{AlgebraicGeometry.Scheme.Modules.conjugateEquiv_restrictFunctorIsoPullback_hom}
  \uses{def:restrictfunctorisopullback}
  For an open immersion \(f\), the conjugate of the comparison hom under
  \(\mathtt{conjugateEquiv}\,(\mathtt{pullbackPushforwardAdjunction}\,f)\,(\mathtt{restrictAdjunction}\,f)\)
  is the identity:
  \[
    \mathtt{conjugateEquiv}\,\bigl((\mathtt{restrictFunctorIsoPullback}\,f).\mathrm{hom}\bigr)
      \;=\; \mathbf{1}_{\mathtt{pushforward}\,f}.
  \]
  This is the left-hand collapse used by the B2 telescope and by both whisker legs.
\end{lemma}
\begin{proof}
  \uses{def:restrictfunctorisopullback}
  By definition \(\mathtt{restrictFunctorIsoPullback}\,f\) is the left-adjoint-uniqueness
  comparison isomorphism of \(\mathtt{restrictAdjunction}\,f\) against
  \(\mathtt{pullbackPushforwardAdjunction}\,f\). The conjugation bijection sends the hom of this
  uniqueness comparison to the identity natural transformation of the common right adjoint
  \(\mathtt{pushforward}\,f\), since conjugation transports the defining unit-triangle of the
  comparison to the identity.
\end{proof}

\begin{lemma}\leanok
[\(c_3\): conjugate of the \(\iota_U\)-comparison whiskered by \(\mathtt{restrict}\,j\)]
  \label{lem:conjugateequiv_restrictfunctorisopullback_whiskerright}
  \lean{AlgebraicGeometry.Scheme.Modules.conjugateEquiv_restrictFunctorIsoPullback_whiskerRight}
  \uses{def:restrictfunctorisopullback, lem:unit_leftadjointuniq_hom_app, lem:conjugateequiv_comp,
    lem:conjugateequiv_restrictfunctorisopullback_hom}
  For the chart \(j : V \hookrightarrow U\), the leg
  \(c_3 = \mathtt{whiskerRight}\,(\mathtt{restrictFunctorIsoPullback}\,\iota_U).\mathrm{hom}\,
  (\mathtt{restrict}\,j)\) maps, under the conjugation through the adjunctions \(G_2 \to G_3\), to the
  \(j\)-pushforward of the unit-triangle datum of
  \(\mathtt{restrictFunctorIsoPullback}\,\iota_U\) --- i.e. to the pushforward-world transport of the
  left-adjoint-uniqueness unit of the \(\iota_U\) comparison.
\end{lemma}
\begin{proof}
  \uses{def:restrictfunctorisopullback, lem:unit_leftadjointuniq_hom_app}
  Since \(\mathtt{restrictFunctorIsoPullback}\,\iota_U\) is the left-adjoint-uniqueness comparison
  of \(\mathtt{restrictAdjunction}\,\iota_U\) against
  \(\mathtt{pullbackPushforwardAdjunction}\,\iota_U\), its conjugate is computed by the defining unit
  triangle \(\mathtt{unit\_leftAdjointUniq\_hom\_app}\) (\cref{lem:unit_leftadjointuniq_hom_app}); the
  outer \(\mathtt{whiskerRight}\) by \(\mathtt{restrict}\,j\) transports this through the second factor
  of \(G_2, G_3\), yielding the displayed \(j\)-pushforward whisker of the unit.
\end{proof}

\begin{lemma}\leanok
[\(c_4\): conjugate of the \(j\)-comparison whiskered into \(\mathtt{pullback}\,\iota_U\)]
  \label{lem:conjugateequiv_restrictfunctorisopullback_whiskerleft}
  \lean{AlgebraicGeometry.Scheme.Modules.conjugateEquiv_restrictFunctorIsoPullback_whiskerLeft}
  \uses{def:restrictfunctorisopullback, lem:unit_leftadjointuniq_hom_app, lem:conjugateequiv_comp,
    lem:conjugateequiv_restrictfunctorisopullback_hom}
  For the chart \(j : V \hookrightarrow U\), the leg
  \(c_4 = \mathtt{whiskerLeft}\,(\mathtt{pullback}\,\iota_U)\,(\mathtt{restrictFunctorIsoPullback}\,j).\mathrm{hom}\)
  maps, under the conjugation through the adjunctions \(G_3 \to G_4\), to the unit-triangle datum of
  \(\mathtt{restrictFunctorIsoPullback}\,j\) precomposed with \(\mathtt{pullback}\,\iota_U\) --- the
  pushforward-world transport of the left-adjoint-uniqueness unit of the \(j\) comparison.
\end{lemma}
\begin{proof}
  \uses{def:restrictfunctorisopullback, lem:unit_leftadjointuniq_hom_app}
  As in \cref{lem:conjugateequiv_restrictfunctorisopullback_whiskerright}, the comparison
  \(\mathtt{restrictFunctorIsoPullback}\,j\) is a \(\mathtt{leftAdjointUniq}\) hom; its conjugate is
  given by the unit triangle \(\mathtt{unit\_leftAdjointUniq\_hom\_app}\)
  (\cref{lem:unit_leftadjointuniq_hom_app}). The \(\mathtt{whiskerLeft}\) by \(\mathtt{pullback}\,\iota_U\)
  threads this datum through the first factor of \(G_3, G_4\), giving the displayed map.
\end{proof}

\begin{lemma}\leanok
[\(c_1, c_6\): conjugate of the reindex congruences along \(j \mathbin{;} \iota_U = \iota_V\)]
  \label{lem:conjugateequiv_reindexcongr}
  \lean{AlgebraicGeometry.Scheme.Modules.conjugateEquiv_reindexCongr}
  \uses{lem:image_preimage_of_le}
  The two flanking congruence legs --- \(c_1\) the inverse hom of
  \(\mathtt{restrictFunctorCongr}\,h_{j\iota}\) (relating \(G_0\) and \(G_1\)) and \(c_6\) the hom of
  \(\mathtt{pullbackCongr}\,h_{j\iota}\) (relating \(G_5\) and \(G_6\)) --- both arise by transport
  along the single equality of opens \(j \mathbin{;} \iota_U = \iota_V\)
  (\cref{lem:image_preimage_of_le}). Under conjugation they map to the corresponding
  \(\mathtt{pushforward}\) congruences (\(\mathtt{eqToHom}\)-transports of
  \(\mathtt{pushforward}\) along that same equality), inverse to one another modulo the reindex.
\end{lemma}
\begin{proof}
  \uses{lem:image_preimage_of_le}
  Both congruences are \(\mathtt{eqToHom}\)/\(\mathtt{eqToIso}\) transports along
  \(j \mathbin{;} \iota_U = \iota_V\); conjugation commutes with such equality-of-functor transports,
  sending each to the \(\mathtt{pushforward}\) congruence along the same equality. As the two
  endpoints \(G_0, G_6\) of the chain are both \(\mathtt{pushforward}\,\iota_V\), the two congruence
  images are mutually inverse transports and cancel in the telescope.
\end{proof}

\begin{lemma}[Telescope of the reindexed pushforward legs]
  \label{lem:pushforwardcomp_reindex_telescope}
  % NOTE: (iter-051) exposition sub-lemma — no standalone \lean{} pin. This telescope content is
  % proved INLINE in the tail of `restrictFunctorIsoPullback_comp_compat_hom` (Iso.inv_hom_id_assoc
  % cancel + conjugateEquiv_reindexCongr), so the dangling \lean{} was removed; the block documents
  % the reindex-cancellation step of the B2 telescope.
  \uses{lem:conjugate_equiv_restrictfunctorcomp_inv, lem:conjugateequiv_pullbackcomp_hom, lem:image_preimage_of_le}
  With \(j \mathbin{;} \iota_U = \iota_V\), the reversed product of the six per-leg conjugate images
  --- the two \(\mathtt{pushforward}\) congruences along \(j \mathbin{;} \iota_U = \iota_V\), the two
  whiskered units, and the pair
  \((\mathtt{pushforwardComp}\,j\,\iota_U).\mathrm{hom} \mathbin{;}
  (\mathtt{pushforwardComp}\,j\,\iota_U).\mathrm{inv}\) --- equals the identity natural
  transformation \(\mathbf{1}_{\mathtt{pushforward}\,\iota_V}\).
\end{lemma}
\begin{proof}
  \uses{lem:conjugate_equiv_restrictfunctorcomp_inv, lem:conjugateequiv_pullbackcomp_hom, lem:image_preimage_of_le}
  The central pair \((\mathtt{pushforwardComp}\,j\,\iota_U).\mathrm{hom} \mathbin{;}
  (\mathtt{pushforwardComp}\,j\,\iota_U).\mathrm{inv}\) cancels to the identity. The two whiskered
  unit legs are then absorbed by the pushforward triangle identities, leaving the two
  \(\mathtt{pushforward}\) congruences along \(j \mathbin{;} \iota_U = \iota_V\)
  (\cref{lem:image_preimage_of_le}), which are mutually inverse transports and cancel. The product
  is therefore \(\mathbf{1}_{\mathtt{pushforward}\,\iota_V}\).
\end{proof}

% ---------------------------------------------------------------------------
% Mathlib dependency anchors for the reindex bridge B2
% ---------------------------------------------------------------------------
Bridge B2 (the reindex coherence of \(\mathtt{restrictFunctorIsoPullback}\)) rests on the
following Mathlib facts about uniqueness of left adjoints and the compositionality of
mates, recorded here as dependency anchors.

\begin{lemma}[Transitivity of the left-adjoint uniqueness isomorphism]
  \label{lem:adjunction_leftadjointuniq_trans}
  \lean{CategoryTheory.Adjunction.leftAdjointUniq_trans}
  \mathlibok
  \textit{Provided by Mathlib.}
  For adjunctions \(F \dashv G\), \(F' \dashv G\), \(F'' \dashv G\) sharing the right adjoint, the
  left-adjoint uniqueness comparisons compose:
  \[
    (\mathtt{leftAdjointUniq}\,\mathrm{adj}_1\,\mathrm{adj}_2)_{\mathrm{hom}}
      \mathbin{;} (\mathtt{leftAdjointUniq}\,\mathrm{adj}_2\,\mathrm{adj}_3)_{\mathrm{hom}}
      = (\mathtt{leftAdjointUniq}\,\mathrm{adj}_1\,\mathrm{adj}_3)_{\mathrm{hom}}.
  \]
\end{lemma}

\begin{lemma}[Mate of a horizontal composite is the vertical composite of mates]
  \label{lem:mateequiv_vcomp}
  \lean{CategoryTheory.mateEquiv_vcomp}
  \mathlibok
  \textit{Provided by Mathlib.}
  For adjunctions \(\mathrm{adj}_1, \mathrm{adj}_2, \mathrm{adj}_3\) and two-squares \(\alpha,
  \beta\), the mate of the horizontal composite equals the vertical composite of the mates:
  \(\mathtt{mateEquiv}\,\mathrm{adj}_1\,\mathrm{adj}_3\,(\alpha \cdot_h \beta)
    = (\mathtt{mateEquiv}\,\mathrm{adj}_1\,\mathrm{adj}_2\,\alpha)
      \cdot_v (\mathtt{mateEquiv}\,\mathrm{adj}_2\,\mathrm{adj}_3\,\beta)\).
\end{lemma}

\begin{lemma}[Conjugation distributes over composition]
  \label{lem:conjugateequiv_comp}
  \lean{CategoryTheory.conjugateEquiv_comp}
  \mathlibok
  \textit{Provided by Mathlib.}
  For adjunctions \(\mathrm{adj}_1 : L_1 \dashv R_1\), \(\mathrm{adj}_2 : L_2 \dashv R_2\),
  \(\mathrm{adj}_3 : L_3 \dashv R_3\) all between the same pair of categories, and natural
  transformations \(\alpha : L_2 \Rightarrow L_1\), \(\beta : L_3 \Rightarrow L_2\), the conjugation
  bijection takes a composite to the (reversed) composite of conjugates:
  \[
    \mathtt{conjugateEquiv}\,\mathrm{adj}_1\,\mathrm{adj}_2\,\alpha
      \mathbin{;} \mathtt{conjugateEquiv}\,\mathrm{adj}_2\,\mathrm{adj}_3\,\beta
      \;=\; \mathtt{conjugateEquiv}\,\mathrm{adj}_1\,\mathrm{adj}_3\,(\beta \mathbin{;} \alpha).
  \]
\end{lemma}

\begin{lemma}[Conjugate of the pullback-composition inverse]
  \label{lem:conjugate_equiv_pullbackcomp_inv}
  \lean{AlgebraicGeometry.Scheme.Modules.conjugateEquiv_pullbackComp_inv}
  \mathlibok
  \textit{Provided by Mathlib.}
  For open immersions \(f, g\) of schemes, under the conjugation between the composite pullback
  adjunction \((\mathtt{pullbackPushforwardAdjunction}\,g).\mathrm{comp}\,
  (\mathtt{pullbackPushforwardAdjunction}\,f)\) and the pullback adjunction of the composite, the
  \emph{inverse} of the pullback-composition isomorphism corresponds to the hom of the
  pushforward-composition isomorphism:
  \[
    \mathtt{conjugateEquiv}\;\dots\;(\mathtt{pullbackComp}\,f\,g).\mathrm{inv}
      \;=\; (\mathtt{pushforwardComp}\,f\,g).\mathrm{hom}.
  \]
  It holds on the nose because \(\mathtt{pullbackComp}\) is \emph{defined} as the left-adjoint
  composition isomorphism \(\mathtt{leftAdjointCompIso}\).
\end{lemma}

\begin{lemma}[Unit triangle of the left-adjoint uniqueness comparison]
  \label{lem:unit_leftadjointuniq_hom_app}
  \lean{CategoryTheory.Adjunction.unit_leftAdjointUniq_hom_app}
  \mathlibok
  \textit{Provided by Mathlib.}
  For adjunctions \(\mathrm{adj}_1 : F \dashv G\) and \(\mathrm{adj}_2 : F' \dashv G\) sharing the
  right adjoint, the left-adjoint uniqueness comparison satisfies the unit triangle
  \[
    \mathrm{adj}_1.\mathrm{unit}.\mathrm{app}\,x
      \mathbin{;} G\bigl((\mathtt{leftAdjointUniq}\,\mathrm{adj}_1\,\mathrm{adj}_2).\mathrm{hom}.\mathrm{app}\,x\bigr)
      \;=\; \mathrm{adj}_2.\mathrm{unit}.\mathrm{app}\,x.
  \]
\end{lemma}

\begin{lemma}[Restrict-side conjugate of the restriction-composition isomorphism]
  \label{lem:conjugate_equiv_restrictfunctorcomp_inv}
  \lean{AlgebraicGeometry.Scheme.Modules.conjugateEquiv_restrictFunctorComp_inv}
  \leanok
  \uses{def:restrictcompreindex}
  Let \(f : X \hookrightarrow Y\) and \(g : Y \hookrightarrow Z\) be open immersions. Under the
  conjugation equivalence between the composite restriction adjunction
  \((\mathtt{restrictAdjunction}\,g) \circ (\mathtt{restrictAdjunction}\,f)\) and the restriction
  adjunction of the composite \(\mathtt{restrictAdjunction}\,(f \mathbin{;} g)\), the hom of the
  restriction-composition isomorphism corresponds to the hom of the pushforward-composition
  isomorphism:
  \[
    \mathtt{conjugateEquiv}\;
      \bigl((\mathtt{restrictAdjunction}\,g).\mathrm{comp}\,(\mathtt{restrictAdjunction}\,f)\bigr)\;
      (\mathtt{restrictAdjunction}\,(f \mathbin{;} g))\;
      (\mathtt{restrictFunctorComp}\,f\,g).\mathrm{hom}
    = (\mathtt{pushforwardComp}\,f\,g).\mathrm{hom}.
  \]
  This is the restriction-world mirror of Mathlib's
  \(\mathtt{conjugateEquiv\_pullbackComp\_inv}\) (which is available for the pullback world only
  because \(\mathtt{pullbackComp}\) is \emph{definitionally equal} to \(\mathtt{leftAdjointCompIso}\),
  so the pullback identity is immediate from the definition). It is the foundational lemma on which
  both Bridge B2 (\cref{lem:restrictfunctorisopullback_comp_compat}) and the Bridge B1 crux
  \(\mathtt{H1inv\_app\_eq\_pullbackVal\_restrict}\) depend.
\end{lemma}
\begin{proof}
  \leanok
  Both \(\mathtt{restrictFunctorComp}\) and \(\mathtt{pushforwardComp}\) are sectionwise
  \(\mathtt{eqToHom}\)/identity maps: on each open the restriction-composition hom is
  \(\mathtt{presheaf.map}\) of an equality-of-opens \(\mathtt{eqToHom}\), and the
  pushforward-composition hom is the identity. The composite restriction adjunction's unit and
  counit are sectionwise the presheaf restriction maps \(\mathtt{presheaf.map}\,(\mathtt{homOfLE}
  \cdots)\) along the canonical inclusions of preimage opens. Evaluating the
  \(\mathtt{conjugateEquiv}\) app-formula therefore reduces the claim to a composite endomorphism of
  each section \(\Gamma(M', U)\) built entirely from \(\mathtt{presheaf.map}\) of morphisms of the
  thin poset \(\mathtt{Opens}^{\mathrm{op}}\). By functoriality these collapse to
  \(\mathtt{presheaf.map}\) of a single opens-morphism, which is forced to be the identity by
  thin-poset uniqueness (every parallel pair of opens-morphisms is equal), whence the composite is
  \(\mathtt{presheaf.map}\,\mathbf{1} = \mathbf{1}\). The computation is carried out at the level of
  \(\mathtt{presheaf.map}\) — composing opens-morphisms before applying the functor, so the module
  carrier is never forced — which keeps it free of the \(\mathtt{pushforward}\)-carrier
  definitional diamond.
\end{proof}

% ---------------------------------------------------------------------------
% Direct base-change anchors for Bridge B1
% ---------------------------------------------------------------------------
\begin{definition}[Restriction--pullback comparison]
  \label{def:restrictfunctorisopullback}
  \lean{AlgebraicGeometry.Scheme.Modules.restrictFunctorIsoPullback}
  \mathlibok
  \textit{Provided by Mathlib.}
  For an open immersion \(f : Y \hookrightarrow X\), presented by the ring-presheaf map
  \(\varphi\), the restriction functor \(\mathtt{restrict}\,f\) on \(\mathcal{O}\)-modules is
  naturally isomorphic to the abstract presheaf-of-modules pullback,
  \[
    \mathtt{restrictFunctorIsoPullback}\,f :
      \mathtt{restrict}\,f \;\cong\; \mathtt{pullback}\,f .
  \]
  It is the left-adjoint uniqueness comparison between the two left adjoints of the
  pushforward, and it is the comparison through which the pullback-world base-change data are
  transported onto the restriction world throughout this section.
\end{definition}

% ---------------------------------------------------------------------------
% Realized delta-conjugation machinery (shared by K1 and Bridge B1)
% ---------------------------------------------------------------------------
The strong-monoidality of pullback-along-an-open-immersion, used both by K1
(\cref{lem:pullback_tensor_map_isiso_open_immersion}) and by Bridge B1
(\cref{lem:tensorobj_restrict_iso_eq_pullback_tensor_map}), rests on three pieces of machinery: the
presheaf-level strong-monoidal upgrade \(\mathbf{H2}\) of restriction of scalars, an abstract
conjugation of one left adjoint's oplax tensorator into another's across uniqueness of left
adjoints, and the resulting "read \(\mathtt{IsIso}\,\delta\) off the conjugation" principle. They are
recorded here as load-bearing nodes.

\begin{lemma}\leanok
[H2: strong-monoidal restriction of scalars along a sectionwise iso]
  \label{lem:restrictscalars_monoidal_of_bijective}
  \lean{PresheafOfModules.restrictScalarsMonoidalOfBijective}
  \uses{lem:restrictscalars_laxmonoidal}
  Let \(\alpha : R \to S\) be a morphism of presheaves of \emph{commutative} rings whose every
  section \(\alpha_U\) is \emph{bijective}. Then the restriction-of-scalars functor
  \(\mathtt{restrictScalars}\,\alpha\) on presheaves of modules is \emph{strong} monoidal: its lax
  tensorator \(\mu\) and lax unit \(\varepsilon\) (\cref{lem:restrictscalars_laxmonoidal}) are
  assembled sectionwise from the \(\mathtt{ModuleCat}\)-level structure maps, which are isomorphisms
  for a bijective ground-ring map; hence \(\mu\) and \(\varepsilon\) are sectionwise isomorphisms,
  hence isomorphisms, and the lax structure upgrades to a strong (genuinely monoidal) one. This is
  the presheaf-level \(\mathbf{H2}\) ingredient: along an open immersion the structure map of the
  restriction functor is sectionwise the bijective structure-ring isomorphism, so restriction is
  strong monoidal and carries an invertible tensorator.
\end{lemma}

\begin{lemma}[Strong-monoidal tensorator iso]
  \label{lem:functor_monoidal_muiso}
  \lean{CategoryTheory.Functor.Monoidal.μIso}
  \mathlibok
  \textit{Provided by Mathlib.}
  A strong (genuinely) monoidal functor \(F\) carries an \emph{invertible} tensorator
  \(\mu_{\mathrm{Iso}}\,F\,A\,B : F A \otimes F B \cong F(A \otimes B)\); its inverse is the oplax
  tensorator \(\delta\,F\,A\,B\). This is the abstract source of every "\(\delta\) is an
  isomorphism" statement for a strong-monoidal functor.
\end{lemma}

\begin{lemma}\leanok
[Abstract \(\delta\)-conjugation across uniqueness of left adjoints]
  \label{lem:deltaconjofmucomparison}
  \lean{AlgebraicGeometry.Scheme.Modules.deltaConjOfMuComparison}
  \uses{lem:adjunction_leftadjointuniq_trans, lem:mateequiv_vcomp, lem:functor_monoidal_muiso}
  Let \(F_1, F_2\) be two left adjoints of a common right adjoint \(G\), with adjunctions
  \(\mathrm{adj}_1 : F_1 \dashv G\) and \(\mathrm{adj}_2 : F_2 \dashv G\); suppose \(F_1, F_2\) are
  oplax monoidal, \(G\) is lax monoidal, and \(\mathrm{adj}_2\) is a \emph{monoidal} adjunction.
  Write \(H_1 := \mathtt{leftAdjointUniq}\,\mathrm{adj}_1\,\mathrm{adj}_2 : F_1 \cong F_2\) for the
  canonical comparison of the two left adjoints. If the lax tensorator of \(G\) induced by
  \(\mathrm{adj}_1\) (its \(\mathtt{rightAdjointLaxMonoidal}\) structure) agrees with the ambient
  lax tensorator of \(G\) on the image of \(F_1\),
  \[
    \mu\bigl(\mathtt{rightAdjointLaxMonoidal}\,\mathrm{adj}_1\bigr)\,(F_1 A)\,(F_1 B)
      \;=\; \mu\,G\,(F_1 A)\,(F_1 B),
  \]
  then the oplax tensorator of \(F_1\) is the \(H_1\)-conjugate of that of \(F_2\):
  \[
    \delta\,F_1\,A\,B
      \;=\; H_1.\mathrm{hom}.\mathrm{app}\,(A \otimes B)
        \mathbin{;} \delta\,F_2\,A\,B
        \mathbin{;} \bigl(H_1.\mathrm{inv}.\mathrm{app}\,A \otimes_{\mathrm m}
          H_1.\mathrm{inv}.\mathrm{app}\,B\bigr).
  \]
\end{lemma}
\begin{proof}
  \uses{lem:adjunction_leftadjointuniq_trans, lem:mateequiv_vcomp}
  \leanok
  Transpose the claimed identity across the hom-equivalence of \(\mathrm{adj}_1\). The oplax
  tensorator \(\delta\,F_1\) transposes to the lax tensorator
  \(\mu(\mathtt{rightAdjointLaxMonoidal}\,\mathrm{adj}_1)\) by the adjoint-mate identity relating a
  left adjoint's oplax tensorator to its right adjoint's lax tensorator. On the conjugate side, the
  two legs \(H_1.\mathrm{hom}\) and \(H_1.\mathrm{inv}\) of the left-adjoint uniqueness comparison
  collapse against the units of \(\mathrm{adj}_1\) and \(\mathrm{adj}_2\) by the triangle identity
  characterising \(\mathtt{leftAdjointUniq}\) (so that
  \(\mathrm{adj}_2.\mathrm{unit} \mathbin{;} G(H_1.\mathrm{inv}) = \mathrm{adj}_1.\mathrm{unit}\)),
  while the lax tensorator of \(G\) coming from \(\mathrm{adj}_2\) is the ambient one because
  \(\mathrm{adj}_2\) is monoidal; mate compositionality
  (\cref{lem:mateequiv_vcomp}) and transitivity of the uniqueness comparison
  (\cref{lem:adjunction_leftadjointuniq_trans}) align the two transports. The residue is exactly the
  hypothesis comparing the two lax tensorators of \(G\) on the image of \(F_1\), which closes the
  identity.
\end{proof}

\begin{lemma}\leanok
[Reading \(\mathtt{IsIso}\,\delta\) off a conjugation]
  \label{lem:isiso_oplaxdelta_of_conj}
  \lean{AlgebraicGeometry.Scheme.Modules.isIso_oplaxδ_of_conj}
  \uses{lem:deltaconjofmucomparison, lem:functor_monoidal_muiso}
  Let \(F_1\) be a strong monoidal functor and \(F_2\) an oplax monoidal functor, and let
  \(e : F_1 \cong F_2\) be a natural isomorphism. Suppose the oplax tensorator of \(F_1\) is the
  \(e\)-conjugate of that of \(F_2\),
  \[
    \delta\,F_1\,A\,B
      \;=\; e.\mathrm{hom}.\mathrm{app}\,(A \otimes B)
        \mathbin{;} \delta\,F_2\,A\,B
        \mathbin{;} \bigl(e.\mathrm{inv}.\mathrm{app}\,A \otimes_{\mathrm m}
          e.\mathrm{inv}.\mathrm{app}\,B\bigr).
  \]
  Then \(\delta\,F_2\,A\,B\) is an isomorphism.
\end{lemma}
\begin{proof}
  \uses{lem:deltaconjofmucomparison, lem:functor_monoidal_muiso}
  \leanok
  Solve the conjugation identity for \(\delta\,F_2\): it equals
  \[
    \mathtt{inv}\bigl(e.\mathrm{hom}.\mathrm{app}\,(A \otimes B)\bigr)
      \mathbin{;} \delta\,F_1\,A\,B
      \mathbin{;} \mathtt{inv}\bigl(e.\mathrm{inv}.\mathrm{app}\,A \otimes_{\mathrm m}
        e.\mathrm{inv}.\mathrm{app}\,B\bigr).
  \]
  Each factor is an isomorphism: the two flanking legs are inverses of components of the natural
  isomorphism \(e\) (and the tensor of two isos is an iso), and the middle factor \(\delta\,F_1\) is
  the inverse of the strong-monoidal tensorator \(\mu_{\mathrm{Iso}}\,F_1\)
  (\cref{lem:functor_monoidal_muiso}) since \(F_1\) is strong monoidal. A composite of isomorphisms
  is an isomorphism.
\end{proof}

\begin{lemma}\leanok
[Bridge B1 crux: sheafification intertwines the presheaf and sheaf adjoint-uniqueness comparisons]
  \label{lem:h1inv_app_eq_pullbackval_restrict}
  \lean{AlgebraicGeometry.Scheme.Modules.H1inv_app_eq_pullbackVal_restrict}
  \uses{def:restrictfunctorisopullback, def:pullback_val_iso,
    lem:sheafification_comp_pullback_eq_leftadjointuniq, lem:leftadjointuniq_app_unit_eta,
    lem:adjunction_leftadjointuniq_trans, lem:conjugateequiv_comp,
    lem:conjugate_equiv_restrictfunctorcomp_inv, lem:conjugateequiv_pullbackcomp_hom,
    lem:conjugateequiv_restrictfunctorisopullback_whiskerright,
    lem:conjugateequiv_restrictfunctorisopullback_whiskerleft, lem:conjugateequiv_reindexcongr,
    lem:unit_leftadjointuniq_hom_app, lem:sheaf_pullback_unit_forget_eq}
  Let \(f : Y \hookrightarrow X\) be an open immersion presented by the presheaf-ring map
  \(\varphi'\), and let \(M \in \Scheme.\mathtt{Modules}\,X\). Write \(H_1\) for the presheaf-level
  left-adjoint-uniqueness comparison \(\mathtt{pushforward}\,\beta \cong
  \mathtt{pullback}\,\varphi'\) (the linchpin of Step~4 of \(\mathtt{tensorObj\_restrict\_iso}\),
  where \(\beta\) is the structure-ring map of the inverse-image presentation). Then \(H_1.\mathrm{inv}\),
  evaluated at the underlying presheaf \(M^{\mathrm p}\) of \(M\), factors --- after the sheafification
  unit \(\eta\) --- as the sheaf-level per-leg reindex
  \[
    H_1.\mathrm{inv}.\mathrm{app}\,M^{\mathrm p}
      \;=\; \eta_{(\mathtt{pullback}\,\varphi')\,M^{\mathrm p}}
        \mathbin{;} \mathtt{forget}\bigl((\mathtt{pullbackValIso}\,f\,M).\mathrm{hom}\bigr)
        \mathbin{;} \mathtt{forget}\bigl(((\mathtt{restrictFunctorIsoPullback}\,f).\mathrm{app}\,M).\mathrm{inv}\bigr),
  \]
  where \(\eta\) is the unit of the sheafification adjunction on \(Y\), \(\mathtt{pullbackValIso}\,f\)
  (\cref{def:pullback_val_iso}) is the sheaf-level pullback comparison, and the maps are transported
  through the forgetful functor \(\Scheme.\mathtt{Modules}\,Y \to \mathtt{PresheafOfModules}\). This is
  the per-leg sheafification bookkeeping behind the \(M\)- and \(N\)-legs of the residual
  \(\mathtt{tensorHom}\) in Bridge B1 (\cref{lem:tensorobj_restrict_iso_eq_pullback_tensor_map}).
\end{lemma}
\begin{proof}
  \uses{def:restrictfunctorisopullback, def:pullback_val_iso,
    lem:sheafification_comp_pullback_eq_leftadjointuniq, lem:leftadjointuniq_app_unit_eta,
    lem:adjunction_leftadjointuniq_trans, lem:conjugateequiv_comp,
    lem:conjugate_equiv_restrictfunctorcomp_inv, lem:conjugateequiv_pullbackcomp_hom,
    lem:conjugateequiv_restrictfunctorisopullback_whiskerright,
    lem:conjugateequiv_restrictfunctorisopullback_whiskerleft, lem:conjugateequiv_reindexcongr,
    lem:unit_leftadjointuniq_hom_app, lem:sheaf_pullback_unit_forget_eq}
  Both sides are maps \((\mathtt{pullback}\,\varphi').\mathrm{obj}\,M^{\mathrm p} \to
  (\mathtt{pushforward}\,\beta).\mathrm{obj}\,M^{\mathrm p}\), and \(H_1.\mathrm{inv}\) is itself a
  \(\mathtt{leftAdjointUniq}\) hom of the presheaf adjunction onto \(\mathtt{pushforward}\,\varphi'\);
  so by \(\mathtt{homEquiv}\)-injectivity it suffices to compare the two images under that adjunction's
  hom-equivalence. The left image is the defining unit datum
  (\(\mathtt{unit\_leftAdjointUniq\_hom\_app}\), \cref{lem:unit_leftadjointuniq_hom_app}); the right image
  is expanded by the unit formula. Rewriting \(\mathtt{pullbackValIso}\)
  (\cref{def:pullback_val_iso}) and \(\mathtt{restrictFunctorIsoPullback}\)
  (\cref{def:restrictfunctorisopullback}) by their \(\mathtt{leftAdjointUniq}\) presentations, and
  identifying the sheafify-then-pullback comparison with the left-adjoint-uniqueness comparison
  (\cref{lem:sheafification_comp_pullback_eq_leftadjointuniq}), turns the obligation into a pure
  \(\mathtt{leftAdjointUniq}\)-coherence: the sheafification unit leg is supplied by the composite-unit
  transport \(\mathtt{leftAdjointUniqUnitEta}\) (\cref{lem:leftadjointuniq_app_unit_eta}), and the
  remaining comparison reduces to the \emph{same} leg-by-leg conjugate telescope as the Bridge B2
  helper (\cref{lem:restrictfunctorisopullback_comp_compat_hom}): the conjugation distributes through
  \(\mathtt{conjugateEquiv\_comp}\) (\cref{lem:conjugateequiv_comp}) and transitivity
  \(\mathtt{leftAdjointUniq\_trans}\) (\cref{lem:adjunction_leftadjointuniq_trans}), the
  \(\mathtt{restrictFunctorComp}\) factor is discharged by the keystone
  (\cref{lem:conjugate_equiv_restrictfunctorcomp_inv}), and the whisker, congruence and
  pullback-composition legs are the shared atomic computations
  (\cref{lem:conjugateequiv_restrictfunctorisopullback_whiskerright},
  \cref{lem:conjugateequiv_restrictfunctorisopullback_whiskerleft},
  \cref{lem:conjugateequiv_reindexcongr}, \cref{lem:conjugateequiv_pullbackcomp_hom}).
  In the realized formalization this comparison is split as a four-part \(\mathtt{forget}\)-transport
  of the sheaf \(\mathtt{unit\_leftAdjointUniq}\) triangle, whose genuine sheafification-boundary
  content is isolated in Part~III (\cref{lem:sheaf_pullback_unit_forget_eq}).
\end{proof}

\begin{lemma}

\leanok
[Bridge B1 crux, Part III: the sheaf pullback unit, forget-transported, factors through sheafification]
  \label{lem:sheaf_pullback_unit_forget_eq}
  \lean{AlgebraicGeometry.Scheme.Modules.sheafPullbackUnit_forget_eq}
  \uses{def:pullback_val_iso, lem:sheafification_comp_pullback_eq_leftadjointuniq}
  Let \(f : Y \hookrightarrow X\) be an open immersion presented by the presheaf-ring map
  \(\varphi'\), and let \(M \in \Scheme.\mathtt{Modules}\,X\). The unit of the \emph{sheaf}-level
  pullback--pushforward adjunction \(\mathtt{pullbackPushforwardAdjunction}\,f\), transported through
  the forgetful functor \(\mathtt{forget} : \Scheme.\mathtt{Modules}\,Y \to \mathtt{PresheafOfModules}\),
  equals the \emph{presheaf}-level pullback--pushforward unit followed by the sheafification unit and
  the pullback comparison \(\mathtt{pullbackValIso}\) (\cref{def:pullback_val_iso}):
  \[
    \mathtt{forget}\bigl((\mathtt{pullbackPushforwardAdjunction}\,f).\mathrm{unit}.\mathrm{app}\,M\bigr)
      = (\mathtt{pullbackPushforwardAdjunction}\,\varphi').\mathrm{unit}.\mathrm{app}\,M^{\mathrm p}
        \mathbin{;} (\mathtt{pushforward}\,\varphi').\mathrm{map}\bigl(
          \eta_{(\mathtt{pullback}\,\varphi')\,M^{\mathrm p}}
          \mathbin{;} \mathtt{forget}((\mathtt{pullbackValIso}\,f\,M).\mathrm{hom})\bigr),
  \]
  where \(\eta\) is the unit of the sheafification adjunction on \(Y\) and \(M^{\mathrm p}\) is the
  underlying presheaf of \(M\). This is the lone sheafification-boundary residual of the Bridge~B1
  crux \(H_1\) (\cref{lem:h1inv_app_eq_pullbackval_restrict}); the other three parts of that crux are
  definitional \(\mathtt{forget}\)/\(\mathtt{pushforward}\) functoriality identities.
\end{lemma}
\begin{proof}
  \uses{def:pullback_val_iso, lem:sheafification_comp_pullback_eq_leftadjointuniq,
    lem:unit_leftadjointuniq_hom_app}
  As shown in the Lean development the identity is, after computing the opaque sheaf unit via the
  concrete \(\mathtt{PullbackConstruction.adjunction}\) (using Mathlib's
  \(\mathtt{pullbackIso}\,\varphi = \mathtt{leftAdjointUniq}(\mathtt{pullbackPushforwardAdjunction}\,\varphi)\,
  (\mathtt{PullbackConstruction.adjunction}\,\varphi)\)), equivalent to the single comparison identity
  \[
    (\mathtt{pullbackIso}\,\varphi).\mathrm{inv}.\mathrm{app}\,M
      = (\mathtt{pullbackValIso}\,f\,M).\mathrm{hom}
      = \mathtt{sheafificationCompPullback}.\mathrm{inv}.\mathrm{app}\,M^{\mathrm p}
        \mathbin{;} (\mathtt{pullback}\,\varphi).\mathrm{map}\,
          (\varepsilon_M), \qquad \varepsilon := (\mathtt{sheafAdj}_X).\mathrm{counit},
  \]
  i.e.\ the two \(\mathtt{leftAdjointUniq}\) comparison isos --- one over the right adjoint
  \(\mathtt{pushforward}\,\varphi\), the other (\(\mathtt{sheafificationCompPullback}\)) over the
  composite right adjoint \(\mathtt{pushforward}\,\varphi \circ \mathtt{forget}\) --- agree on \(M\).
  Transposing along the \emph{whole composite} \(\mathtt{homEquiv}\) is circular (it returns the parent
  goal), and \(\mathtt{leftAdjointUniq\_trans}\) does not apply (distinct right adjoints).

  The identity is the \(\mathtt{SheafOfModules}\) analogue of Mathlib's
  \(\mathtt{Functor.toSheafify\_pullbackSheafificationCompatibility}\)
  (\(\mathtt{Sites/CoverLifting.lean}\): \(362\)--\(382\)), which proves the literal twin --- a
  \(\mathtt{leftAdjointUniq}\) of two \emph{composite} adjunctions, pushed through the forgetful functor
  \(\mathtt{sheafToPresheaf}\) and matched against the sheafification unit \(\mathtt{toSheafify}\). Its
  proof skeleton ports directly:
  % NOTE: (iter-053, review) The lemma is now CLOSED in Lean, but the REALIZED proof does NOT follow
  % the steps 1--4 below. It uses forget-faithfulness (fullyFaithfulForget.map_injective) + an inner
  % presheaf-pullback transpose telescope + the INVERSE-leftAdjointUniq triangle (hAcancel via
  % leftAdjointUniq_inv_app + unit_leftAdjointUniq_hom_app) + sheafificationCompPullback_eq_leftAdjointUniq.
  % It never invokes Adjunction.comp_unit_app (step 3) and leaves hpin dead. The intro paragraph and the
  % reduction-to-hKEY ARE faithful; only steps 1--4 are stale. Planner: dispatch blueprint-writer to
  % rewrite this enumerate to the realized route. See task_results/lean-vs-blueprint-checker-iter053.md.
  \begin{enumerate}
    \item Transpose along the \emph{innermost} adjunction
      \(\mathtt{pullbackPushforwardAdjunction}\,\varphi'\) (the presheaf-level one whose
      \(\mathtt{homEquiv}\) the crux \(H_1\) already uses), via
      \(\bigl((\mathtt{pullbackPushforwardAdjunction}\,\varphi').\mathtt{homEquiv}\,\_\,\_\bigr).\mathtt{injective}\)
      --- supplying the two objects explicitly --- \emph{not} the composite. This is the step that
      breaks circularity.
    \item Feed the defining unit triangle of the composite comparison,
      \(\mathtt{hpin} := \mathtt{unit\_leftAdjointUniq\_hom\_app}\,A\,B\,M^{\mathrm p}\)
      (\cref{lem:unit_leftadjointuniq_hom_app}), where
      \(A = \mathtt{sheafAdj}_X \mathbin{.}\!\mathtt{comp}\,\mathtt{pullbackPPAdj}_{\mathrm{sheaf}}\) and
      \(B = \mathtt{pullbackPPAdj}_{\mathrm{pre}} \mathbin{.}\!\mathtt{comp}\,\mathtt{sheafAdj}_Y\)
      (the two composite adjunctions whose left-adjoint comparison \emph{is}
      \(\mathtt{sheafificationCompPullback}\), \cref{lem:sheafification_comp_pullback_eq_leftadjointuniq}).
    \item Expand both composite units with \(\mathtt{Adjunction.comp\_unit\_app}\) (applied twice),
      exposing the sheafification unit \(\eta\) as an explicit separate factor --- the factor that is
      invisible from the composite \(\mathtt{homEquiv}\) and the reason the naive route is circular.
    \item Transpose back (\(\mathtt{homEquiv\_unit}\)/\(\mathtt{homEquiv\_counit}\) together with unit
      naturality), and discharge the residual sheafification--pullback leg with the \(X\)-counit
      naturality square
      \(\mathtt{hnat}\): \(\varepsilon_M \mathbin{;} \eta_M
      = \eta_{(\mathtt{a}_X)\,M^{\mathrm p}} \mathbin{;} (\mathtt{pushforward}\,\varphi).\mathrm{map}
      ((\mathtt{pullback}\,\varphi).\mathrm{map}\,\varepsilon_M)\).
  \end{enumerate}
  Both non-circular inputs \(\mathtt{hpin}\) (step~2) and \(\mathtt{hnat}\) (step~4) are exactly the two
  inputs the Mathlib precedent uses, and are already typed and compiling in the file. The remaining
  three parts of the crux \(H_1\) (\cref{lem:h1inv_app_eq_pullbackval_restrict}) are definitional
  \(\mathtt{forget}\)/\(\mathtt{pushforward}\) functoriality identities.
\end{proof}

\begin{lemma}


\leanok
[Bridge B1: \(\mathtt{tensorObj\_restrict\_iso}\) is the conjugated base-change comparison]
  \label{lem:tensorobj_restrict_iso_eq_pullback_tensor_map}
  \lean{AlgebraicGeometry.Scheme.Modules.tensorObj_restrict_iso_eq_pullbackTensorMap}
  \uses{lem:tensorobj_restrict_iso, lem:pullback_tensor_map, lem:pullback_tensor_map_isiso_open_immersion, lem:isiso_pullbacktensormap_of_sheafifydelta, lem:pushforward_mu_appiso_collapse, def:restrictfunctorisopullback}
  Let \(f : Y \hookrightarrow X\) be an open immersion and \(M, N \in \Scheme.\mathtt{Modules}\,X\).
  The tensor-restriction comparison is the \(\mathtt{restrictFunctorIsoPullback}\)-conjugate of the
  base-change tensor map, i.e. the conjugate factorisation
  \[
    \mathtt{tensorObj\_restrict\_iso}\,f\,M\,N
      = \bigl(\mathtt{restrictFunctorIsoPullback}\,f\bigr).\mathrm{app}\,(M \otimes_X N)
        \mathbin{;} \mathtt{asIso}\bigl(\mathtt{pullbackTensorMap}\,f\,M\,N\bigr)
        \mathbin{;} \mathtt{tensorObjIsoOfIso}\,\rho_M^{-1}\,\rho_N^{-1},
  \]
  where \(\rho_M := (\mathtt{restrictFunctorIsoPullback}\,f).\mathrm{app}\,M\) is the per-leg
  comparison and \(\mathtt{pullbackTensorMap}\,f\) (\cref{lem:pullback_tensor_map}) is an
  isomorphism for the open immersion \(f\)
  (\cref{lem:pullback_tensor_map_isiso_open_immersion}).
\end{lemma}
\begin{proof}
  \uses{lem:tensorobj_restrict_iso, lem:pullback_tensor_map, lem:pullback_tensor_map_isiso_open_immersion, lem:isiso_pullbacktensormap_of_sheafifydelta, lem:pushforward_mu_appiso_collapse, lem:deltaconjofmucomparison, def:restrictfunctorisopullback, lem:adjunction_leftadjointuniq_trans, lem:h1inv_app_eq_pullbackval_restrict}
  The identity is proved \emph{directly}, with no monoidal structure on the sheaf category. Both
  sides are built from the \emph{same} presheaf-level comparison through the \emph{same} transport
  prefix. By construction the tensor-restriction comparison
  \(\mathtt{tensorObj\_restrict\_iso}\,f\) (\cref{lem:tensorobj_restrict_iso}) factors as the shared
  prefix \(\mathtt{restrictFunctorIsoPullback}\,f \mathbin{;} \mathtt{sheafificationCompPullback}\,f\)
  --- reducing the restriction to the abstract pullback and moving it inside sheafification ---
  followed by a tail assembled from the structure-sheaf reconciliation; the base-change tensor map
  \(\mathtt{pullbackTensorMap}\,f\) (\cref{lem:pullback_tensor_map}) is built from the identical
  prefix followed by the sheafified oplax tensorator
  \(\delta\bigl(\mathtt{PresheafOfModules.pullback}\,\varphi'\bigr)\)
  (\cref{lem:isiso_pullbacktensormap_of_sheafifydelta}).

  Cancelling the common prefix, the goal becomes a \emph{presheaf-level} identity: the sheafified
  oplax tensorator \(\delta\bigl(\mathtt{PresheafOfModules.pullback}\,\varphi'\bigr)\) equals the
  left-adjoint--transported strong tensorator
  \(\mu\text{-iso}\bigl(\mathtt{pushforward}_0\mathtt{OfCommRingCat} \mathbin{;} \mathtt{restrictScalars}\,
  \beta'\bigr)\) that forms the tail of \(\mathtt{tensorObj\_restrict\_iso}\,f\). This is exactly the
  \(\delta\)-conjugation identity \(\mathtt{pushforward\_mu\_appIso\_collapse}\)
  (\cref{lem:pushforward_mu_appiso_collapse}), which expresses the strong tensorator
  \(\delta\,G_\beta\) as the \(H_1\)-conjugate of the oplax
  \(\delta\bigl(\mathtt{pullback}\,\varphi'\bigr)\) --- the very abstract conjugation
  \(\mathtt{deltaConjOfMuComparison}\) (\cref{lem:deltaconjofmucomparison}) specialised to the
  pullback adjunction --- combined with uniqueness of left adjoints
  (\(\mathtt{leftAdjointUniq\_trans}\), \cref{lem:adjunction_leftadjointuniq_trans}, and naturality of
  \(\mathtt{leftAdjointUniq}\)) to match the two transports of the presheaf comparison. Here all
  monoidal data genuinely \emph{exists} at the presheaf level: the category of presheaves of
  modules carries a symmetric monoidal structure, and no monoidal structure on the sheaf
  categories \(\Scheme.\mathtt{Modules}\) is invoked.

  Finally, for the open immersion \(f\) the comparison \(\mathtt{pullbackTensorMap}\,f\) is
  invertible (\cref{lem:pullback_tensor_map_isiso_open_immersion}), so
  \(\mathtt{asIso}(\mathtt{pullbackTensorMap}\,f)\) is defined; the per-leg comparisons
  \(\rho_M = (\mathtt{restrictFunctorIsoPullback}\,f).\mathrm{app}\,M\)
  (\cref{def:restrictfunctorisopullback}) reassemble the displayed conjugate factorisation
  \(\mathtt{tensorObj\_restrict\_iso}\,f = \rho_{M \otimes N} \mathbin{;}
  \mathtt{asIso}(\mathtt{pullbackTensorMap}\,f) \mathbin{;}
  \mathtt{tensorObjIsoOfIso}\,\rho_M^{-1}\,\rho_N^{-1}\).
\end{proof}

\begin{lemma}


\leanok
[S2: tensor-restriction comparison commutes with further restriction]
  \label{lem:tensorobj_restrict_iso_restrict_compat}
  \lean{AlgebraicGeometry.Scheme.Modules.tensorObj_restrict_iso_restrict_compat}
  \uses{lem:tensorobj_restrict_iso, lem:image_preimage_of_le}
  \textit{Source: [Stacks Project], Modules on Ringed Spaces, Lemma
  \texttt{lemma-tensor-product-pullback}: pullback commutes with tensor product
  \emph{functorially}, the functoriality on which the naturality square below rests.}
  Let \(M, N \in \Scheme.\mathtt{Modules}\,X\) be arbitrary \(\mathcal{O}_X\)-modules and let
  \(j : V \hookrightarrow U\) be the chart morphism above. Write
  \(T(f) := \mathtt{tensorObj\_restrict\_iso}\,f : (M \otimes_X N)|_f \cong M|_f \otimes N|_f\)
  (\cref{lem:tensorobj_restrict_iso}). The tensor-restriction comparison is \emph{natural under
  further restriction along \(j\)}: the square
  \[
    \begin{array}{ccc}
      \bigl((M \otimes_X N)|_{\iota_U}\bigr)\big|_j
        & \xrightarrow{\;(\mathtt{restrict}\,j)\,T(\iota_U)\;}
        & \bigl(M|_{\iota_U} \otimes_U N|_{\iota_U}\bigr)\big|_j \\[2pt]
      \big\downarrow{\scriptstyle \varrho}
        & & \big\downarrow{\scriptstyle \varrho \otimes \varrho} \\[2pt]
      (M \otimes_X N)|_{\iota_V}
        & \xrightarrow{\;T(\iota_V)\;}
        & M|_{\iota_V} \otimes_V N|_{\iota_V}
    \end{array}
  \]
  commutes, where \(\varrho\) is the equality-of-opens reindexing of
  \cref{lem:image_preimage_of_le}. Equivalently, modulo \(\varrho\),
  \((\mathtt{restrict}\,j)\,T(\iota_U) = T(\iota_V)\).
\end{lemma}
\begin{proof}
  \uses{lem:tensorobj_restrict_iso_eq_pullback_tensor_map, lem:pullback_tensor_map_basechange, lem:pullback_tensor_map_natural, def:restrictcompreindex, lem:image_preimage_of_le}
  % SOURCE: [Stacks Project], Modules on Ringed Spaces, Lemma
  % \ref{lemma-tensor-product-pullback} (pullback of $\mathcal{O}$-modules commutes
  % with tensor product, functorially in the two arguments)
  % (read from references/stacks-modules.tex, L2392-2400).
  % SOURCE QUOTE:
  % "\begin{lemma}
  % \label{lemma-tensor-product-pullback}
  % Let $f : (X, \mathcal{O}_X) \to (Y, \mathcal{O}_Y)$ be
  % a morphism of ringed spaces. Let $\mathcal{F}$, $\mathcal{G}$
  % be $\mathcal{O}_Y$-modules. Then
  % $f^*(\mathcal{F} \otimes_{\mathcal{O}_Y} \mathcal{G})
  % = f^*\mathcal{F} \otimes_{\mathcal{O}_X} f^*\mathcal{G}$
  % functorially in $\mathcal{F}$, $\mathcal{G}$.
  % \end{lemma}"
  % SOURCE QUOTE PROOF:
  % "\begin{proof}
  % Omitted.
  % \end{proof}"
  This square is \emph{not} a bespoke leg-by-leg chase: it is base-change bookkeeping built on
  Bridge B1 (\cref{lem:tensorobj_restrict_iso_eq_pullback_tensor_map}). Substitute B1 to rewrite
  both \(T(\iota_U) = \mathtt{tensorObj\_restrict\_iso}\,\iota_U\) and
  \(T(\iota_V) = \mathtt{tensorObj\_restrict\_iso}\,\iota_V\) as
  \(\mathtt{restrictFunctorIsoPullback}\)-conjugates of the base-change tensor maps
  \(\mathtt{asIso}(\mathtt{pullbackTensorMap}\,\iota_U)\) and
  \(\mathtt{asIso}(\mathtt{pullbackTensorMap}\,\iota_V)\). After this substitution the comparison
  factors \(\mathtt{restrictFunctorIsoPullback}\) cancel against the reindexing \(\varrho\) by the
  reindex coherence (Bridge B2, used here only through
  \(\mathtt{restrictFunctorIsoPullback}\)'s pseudonaturality), and the displayed square reduces to
  the corresponding equation between the \emph{pullback}-world tensor maps. That equation is the
  already-proven composition law \(\mathtt{pullbackTensorMap\_restrict}\)
  (\cref{lem:pullback_tensor_map_basechange}): for the factorisation
  \(j \mathbin{;} \iota_U = \iota_V\) it identifies
  \((\mathtt{pullback}\,j)\bigl(\mathtt{pullbackTensorMap}\,\iota_U\bigr)\) with
  \(\mathtt{pullbackTensorMap}\,\iota_V\), conjugated by the pullback pseudofunctoriality
  isomorphism \(\mathtt{pullbackComp}\,j\,\iota_U\). Naturality of \(\mathtt{pullbackTensorMap}\) in
  the two module arguments (\cref{lem:pullback_tensor_map_natural}) carries the per-leg comparisons
  \(\rho_M, \rho_N\) of B1 through the square. Finally the reindexing
  \(\varrho = \mathtt{restrictCompReindex}\,j\) (\cref{def:restrictcompreindex}) threads the
  equality of opens \(j \mathbin{;} \iota_U = \iota_V\) (\cref{lem:image_preimage_of_le}) on both
  flanks --- its tensor image \(\varrho \otimes \varrho\) on the right --- absorbing the
  congruence \(\mathtt{restrictFunctorCongr}\) and yielding the displayed square. No sheaf-level
  monoidal structure intervenes.
\end{proof}

\begin{lemma}
[Cone B / L1: the dual base-change comparison is sectionwise an internal-hom evaluation]
  \label{lem:presheafdual_pullback_comparison_eval_apply}
  \lean{AlgebraicGeometry.Scheme.Modules.presheafDual_pullback_comparison_eval_apply}
  \uses{lem:dual_restrict_iso, lem:internal_hom_eval, def:presheaf_dual}
  \textit{Source: internal categorical construction; no external reference.}
  Let \(j : V \hookrightarrow U\) be an open immersion with \(j \mathbin{;} \iota_U = \iota_V\) and
  \(M \in \Scheme.\mathtt{Modules}\,X\). On each open section the presheaf-level dual base-change
  comparison \(\theta_M\) (Step~4 of \(\mathtt{dual\_restrict\_iso}\), \cref{lem:dual_restrict_iso})
  acts as the internal-hom \emph{evaluation} of \cref{lem:internal_hom_eval}: a pulled-back dual
  section \(\varphi\) --- an \(\mathcal{O}\)-linear functional on the restricted module --- is carried
  by \(\theta_M\) to the functional on the pulled-back module whose value at a section \(s\) is
  \(\varphi\) evaluated at \(s\), reindexed across the immersion's direct-image
  \(j.\mathtt{opensFunctor}\). In short, the \(\mathtt{app}\) of \(\theta_M\) is, sectionwise, an
  evaluation map of the internal hom (the immersion analogue of the eval-at-\(1\) presentation of
  \(\mathtt{presheafDualUnitIso}\) in \cref{lem:presheafdualunitiso_naturality}).
\end{lemma}
\begin{proof}
  \uses{lem:dual_restrict_iso, lem:internal_hom_eval, def:presheaf_dual}
  Unfolding \(\theta_M\) through its construction (\cref{lem:dual_restrict_iso}), the forward
  component at an open \(W \le V\) is the categorical image of the dual-section component of
  \(\varphi\) at the \(j\)-image of \(W\), i.e.\ the evaluation map of the internal hom of
  \cref{lem:internal_hom_eval} reindexed across \(j.\mathtt{opensFunctor}\). This is a direct
  read-off of the definition, sectionwise; no naturality is required for this step.
\end{proof}

\begin{lemma}
[Cone B / L2: contravariant dual transport is precomposition by the restriction map]
  \label{lem:dualisoofiso_restrict_precomp_apply}
  \lean{AlgebraicGeometry.Scheme.Modules.dualIsoOfIso_restrict_precomp_apply}
  \uses{lem:presheaf_dualisoofiso_trans, def:presheaf_dual_iso_of_iso, def:presheaf_dual_precomp_equiv}
  \textit{Source: internal categorical construction; no external reference.}
  With \(j : V \hookrightarrow U\) an open immersion as above, the contravariant dual transport
  \(\mathtt{dualIsoOfIso}\,(\mathtt{restrict}\,j)\) is, on each section, \emph{precomposition} of a
  dual section (functional) by the underlying section map of \((\mathtt{restrict}\,j)\): for a dual
  section \(\varphi\), \(\bigl(\mathtt{dualIsoOfIso}\,(\mathtt{restrict}\,j)\bigr)(\varphi)
  = (\mathtt{restrict}\,j) \mathbin{;} \varphi\) sectionwise. This is the specialisation of the
  precomposition description underlying \cref{lem:presheaf_dualisoofiso_trans}
  (\cref{def:presheaf_dual_iso_of_iso}, \cref{def:presheaf_dual_precomp_equiv}) from a unit
  automorphism to the (non-invertible) restriction map, with no invertibility of the transported
  morphism used.
\end{lemma}
\begin{proof}
  \uses{lem:presheaf_dualisoofiso_trans, def:presheaf_dual_iso_of_iso, def:presheaf_dual_precomp_equiv}
  By construction \(\mathtt{dualIsoOfIso}\,e\) is assembled sectionwise from the precomposition
  equivalence \(\mathtt{dualPrecompEquiv}\,e\) (\cref{def:presheaf_dual_precomp_equiv}), which sends
  a dual section \(\varphi\) to \(e.\mathrm{hom} \mathbin{;} \varphi\). The proof of
  \cref{lem:presheaf_dualisoofiso_trans} uses only this precomposition description and never the
  invertibility of \(e\); applying it with \(e = (\mathtt{restrict}\,j)\) gives the claim.
\end{proof}

\begin{lemma}
[Cone B / L3a: internal-hom evaluation commutes with the immersion restriction map (genuine base-change atom)]
  \label{lem:presheafdual_eval_restrict_commute_apply}
  \lean{AlgebraicGeometry.Scheme.Modules.presheafDual_eval_restrict_commute_apply}
  \uses{lem:internal_hom_eval, lem:image_preimage_of_le}
  \textit{Source: internal categorical construction; no external reference.}
  Fix a dual section \(\varphi\) and a section \(s\). Internal-hom evaluation
  (\cref{lem:internal_hom_eval}) commutes with the immersion's restriction map: evaluating then
  restricting equals restricting both arguments then evaluating,
  \[
    \varphi(s)\big|_V \;=\; \bigl(\varphi\big|_V\bigr)\bigl(s\big|_V\bigr),
  \]
  now against the \emph{non-invertible} restriction map \((\mathtt{restrict}\,j)\) in place of the
  multiplication by a unit scalar that governs the unit-automorphism case
  (\cref{lem:presheafdualunitiso_naturality}). This is the genuine internal-hom base-change content ---
  the dual-flank analogue of Bridge~B1 (\cref{lem:tensorobj_restrict_iso_eq_pullback_tensor_map}) ---
  and is the immersion analogue of the evaluate/restrict commutation core of
  \cref{lem:presheafdualunitiso_naturality}. The
  down-set fact \(j \mathbin{;} \iota_U = \iota_V\) (\cref{lem:image_preimage_of_le}) pins the open
  along which both sides are restricted.
\end{lemma}
\begin{proof}
  \uses{lem:internal_hom_eval, lem:image_preimage_of_le}
  This is exactly the evaluate/restrict-commutation identity already established in the proof of
  \cref{lem:internal_hom_eval}: the restriction \(\varphi|_V\) of a functional inside the dual is its
  restriction to the open \(V\), restriction is functorial, and evaluating the restricted functional
  on the restricted section agrees with restricting the value of the evaluation
  (\(\varphi(s)|_V = (\varphi|_V)(s|_V)\)). Because the restriction map is not invertible on sections,
  this is a genuine naturality of evaluation against \((\mathtt{restrict}\,j)\), not a scalar
  commutation; the open \(V\) is pinned by \(j \mathbin{;} \iota_U = \iota_V\)
  (\cref{lem:image_preimage_of_le}).
\end{proof}

\begin{lemma}
[Cone B / L3b: pointwise naturality square of the dual base-change comparison (sectionwise form)]
  \label{lem:presheafdual_pullback_restrict_natural_apply}
  \lean{AlgebraicGeometry.Scheme.Modules.presheafDual_pullback_restrict_natural_apply}
  \uses{lem:presheafdual_pullback_comparison_eval_apply, lem:dualisoofiso_restrict_precomp_apply, lem:presheafdual_eval_restrict_commute_apply, lem:image_preimage_of_le}
  \textit{Source: internal categorical construction; no external reference.}
  With \(j : V \hookrightarrow U\) as above, the pointwise naturality square of the dual base-change
  comparison holds: for a section and an element, conjugating \(\theta\) by the contravariant
  transport \(\mathtt{dualIsoOfIso}\,(\mathtt{restrict}\,j)\) on the source equals transporting along
  \((\mathtt{restrict}\,j)\) on the target. Concretely, evaluating the source composite at a dual
  section \(\varphi\) and a section \(s\), L1 (\cref{lem:presheafdual_pullback_comparison_eval_apply})
  rewrites \(\theta\) as internal-hom evaluation and L2
  (\cref{lem:dualisoofiso_restrict_precomp_apply}) rewrites the transport as precomposition by
  \((\mathtt{restrict}\,j)\), so both composites become evaluations of \(\varphi\), and L3a
  (\cref{lem:presheafdual_eval_restrict_commute_apply}) identifies them. This is the immersion
  analogue of the eval-core of \cref{lem:presheafdualunitiso_naturality};
  the down-set factorisation \(j \mathbin{;} \iota_U = \iota_V\) is passed explicitly
  as in the support lemma \cref{lem:slice_dual_transport_naturality_apply}.
\end{lemma}
\begin{proof}
  \uses{lem:presheafdual_pullback_comparison_eval_apply, lem:dualisoofiso_restrict_precomp_apply, lem:presheafdual_eval_restrict_commute_apply, lem:image_preimage_of_le}
  Substitute L1 (\cref{lem:presheafdual_pullback_comparison_eval_apply}) and L2
  (\cref{lem:dualisoofiso_restrict_precomp_apply}): the source composite computes \(\varphi(s|_V)\),
  the target composite \(\varphi(s)|_V\); they agree by L3a
  (\cref{lem:presheafdual_eval_restrict_commute_apply}). The factorisation
  \(j \mathbin{;} \iota_U = \iota_V\) (\cref{lem:image_preimage_of_le}) pins the indexing.
\end{proof}

\begin{lemma}
[Immersion-naturality of the presheaf dual base-change comparison (Cone B)]
  \label{lem:presheafdual_pullback_restrict_natural}
  \lean{AlgebraicGeometry.Scheme.Modules.presheafDual_pullback_restrict_natural}
  \uses{lem:dual_restrict_iso, def:scheme_modules_dual_iso_of_iso, lem:presheafdual_pullback_restrict_natural_apply}
  Let \(j : V \hookrightarrow U\) be an open immersion with \(j \mathbin{;} \iota_U = \iota_V\), and let
  \(M \in \Scheme.\mathtt{Modules}\,X\). The presheaf-level dual base-change comparison
  \[
    \theta_M : (\mathtt{pullback}\,\varphi).\mathrm{obj}\,(\mathtt{dual}\,M_{\mathrm{val}})
      \;\cong\; \mathtt{dual}\bigl((\mathtt{pullback}\,\varphi).\mathrm{obj}\,M_{\mathrm{val}}\bigr)
  \]
  --- the internal-hom comparison
  \(f^{*}\mathcal{H}om(M, \mathcal{O}) \to \mathcal{H}om(f^{*}M, f^{*}\mathcal{O})\) forming Step~4 of
  \(\mathtt{dual\_restrict\_iso}\) (\cref{lem:dual_restrict_iso}) --- is natural in the immersion
  \(j\): conjugating \(\theta\) by the contravariant transport \(\mathtt{dualIsoOfIso}\)
  (\cref{def:scheme_modules_dual_iso_of_iso}) along \(j\) on the source matches its transport along
  \(j\) on the target. This is the immersion analogue of \cref{lem:presheafdualunitiso_naturality},
  with the unit automorphism there replaced by the restriction map \((\mathtt{restrict}\,j)\).
\end{lemma}
\begin{proof}
  \uses{lem:presheafdual_pullback_restrict_natural_apply, def:scheme_modules_dual_iso_of_iso}
  Exactly as in the proven unit case \cref{lem:presheafdualunitiso_naturality}: reduce by
  \(\mathrm{Iso.ext}\), sectionwise \(\mathtt{hom\_ext}\), and element extensionality to the pointwise
  square, then close by \cref{lem:presheafdual_pullback_restrict_natural_apply}.
\end{proof}

\begin{lemma}


\leanok
[S3: dual-restriction comparison commutes with further restriction (clean core)]
  \label{lem:dual_restrict_iso_restrict_compat}
  \lean{AlgebraicGeometry.Scheme.Modules.dual_restrict_iso_restrict_compat}
  \uses{lem:dual_restrict_iso, def:scheme_modules_dual_iso_of_iso, def:restrictcompreindex, lem:image_preimage_of_le}
  Let \(M \in \Scheme.\mathtt{Modules}\,X\) and let \(j : V \hookrightarrow U\) be the chart
  morphism above. Write \(D(f) := \mathtt{dual\_restrict\_iso}\,f :
  (\mathtt{dual}\,M)|_f \cong \mathtt{dual}\bigl(M|_f\bigr)\) (\cref{lem:dual_restrict_iso}). The
  \emph{bare} dual-restriction comparison is natural under further restriction along \(j\): modulo
  the reindexing \(\varrho = \mathtt{restrictCompReindex}\,j\) (\cref{def:restrictcompreindex}),
  with its contravariant image \(\mathtt{dualIsoOfIso}\,\varrho\)
  (\cref{def:scheme_modules_dual_iso_of_iso}) on the dual flank,
  \[
    D(\iota_V)
      \;=\; \varrho_{\mathtt{dual}\,M}
        \mathbin{;} (\mathtt{restrict}\,j)\,D(\iota_U)
        \mathbin{;} \mathtt{dual\_restrict\_iso}\,j\,\bigl(M|_{\iota_U}\bigr)
        \mathbin{;} \mathtt{dualIsoOfIso}\,\varrho_M .
  \]
  This square proves \emph{only} the \(\mathtt{dual\_restrict\_iso}\) immersion-naturality; the
  trivialisation step's \((\mathtt{dualIsoOfIso}\,e^M)^{-1}\) transport and the datum refinement
  \(e^M \mapsto \mathtt{restrictIsoUnitOfLE}\,h_{VU}\,e^M\) are \emph{not} part of the claim --- they
  telescope in the parent (\cref{lem:trivialisation_restrict_compat}).
\end{lemma}
\begin{proof}
  \uses{lem:dual_restrict_iso, lem:presheafdual_pullback_restrict_natural, lem:dualisoofiso_trans, def:scheme_modules_dual_iso_of_iso, def:restrictcompreindex, lem:image_preimage_of_le}
  This is the dual-side analogue of S2, routed through \emph{Cone~B} --- internal-hom base-change
  naturality --- and \emph{not} through any sheaf-level monoidal structure. The duality functor
  \(\mathtt{dual} = \mathcal{H}om(-, \mathcal{O})\) is the internal hom into the unit, and
  \(\mathtt{dual\_restrict\_iso}\,f\) (\cref{lem:dual_restrict_iso}) is its restriction comparison ---
  the favorable open-immersion case of the Stacks lemma
  \(f^{*}\mathcal{H}om(M, \mathcal{O}) \cong \mathcal{H}om(f^{*}M, f^{*}\mathcal{O})\). No general
  ``monoidal/closed functor preserves duals'' principle applies here: a pullback along an open
  immersion is neither faithful as stated nor an equivalence, so the only available functor-to-dual
  principles do not fire and the square is genuinely bespoke.

  Its naturality under further restriction along \(j\) mirrors S2 but rests on the \emph{dual analogue
  of Bridge B1}: the immersion-naturality of the presheaf dual base-change comparison \(\theta\)
  (\cref{lem:presheafdual_pullback_restrict_natural}). Chart-chasing the constituent steps of
  \(\mathtt{dual\_restrict\_iso}\) against the factorisation \(j \mathbin{;} \iota_U = \iota_V\), the
  Step-4 internal-hom comparison is carried through by
  \cref{lem:presheafdual_pullback_restrict_natural}, with the contravariant reindexing
  \(\mathtt{dualIsoOfIso}\,\varrho\) (\cref{def:scheme_modules_dual_iso_of_iso},
  \cref{lem:dualisoofiso_trans}) on the dual flank in place of \(\varrho \otimes \varrho\), while the
  sheafification-seam closers (identical to S2) carry the remaining steps. The reindexing
  \(\varrho = \mathtt{restrictCompReindex}\,j\) (\cref{def:restrictcompreindex}) threads the equality
  of opens \(j \mathbin{;} \iota_U = \iota_V\) (\cref{lem:image_preimage_of_le}) on both flanks. The
  trivialisation step's \((\mathtt{dualIsoOfIso}\,e^M)^{-1}\) transport and the datum refinement
  \(e^M \mapsto \mathtt{restrictIsoUnitOfLE}\,h_{VU}\,e^M\) are \emph{not} part of this square; they
  telescope in the parent (\cref{lem:trivialisation_restrict_compat}).
\end{proof}

\begin{lemma}


\leanok
[S4a: dual-unit contraction commutes with further restriction]
  \label{lem:dual_unit_iso_restrict_compat}
  \lean{AlgebraicGeometry.Scheme.Modules.dual_unit_iso_restrict_compat}
  \uses{lem:dual_unit_iso, lem:presheafdualunitiso_naturality, lem:image_preimage_of_le}
  With \(j : V \hookrightarrow U\) as above, the unit-duality contraction
  \(\mathtt{dual\_unit\_iso} : \mathcal{H}om(\mathcal{O}_U, \mathcal{O}_U) \cong \mathcal{O}_U\)
  (\cref{lem:dual_unit_iso}) commutes with further restriction along \(j\): modulo the reindexing
  \(\varrho\),
  \[
    (\mathtt{restrict}\,j)\,(\mathtt{dual\_unit\_iso}\ \text{over}\ U)
      \;=\; \mathtt{dual\_unit\_iso}\ \text{over}\ V.
  \]
\end{lemma}
\begin{proof}
  \uses{lem:dual_unit_iso, lem:dual_restrict_iso_restrict_compat, lem:presheafdual_pullback_restrict_natural, lem:presheafdualunitiso_naturality, def:unitrestrictiso, def:restrictcompreindex, lem:image_preimage_of_le}
  This is the unit companion of S3, routed through the same Cone~B. The unit-duality contraction
  \(\mathtt{dual\_unit\_iso} : \mathcal{H}om(\mathcal{O}, \mathcal{O}) \cong \mathcal{O}\) is the
  unit-evaluation component of the internal-hom restriction comparison
  \(\mathtt{dual\_restrict\_iso}\) (\cref{lem:dual_restrict_iso_restrict_compat}). Its
  restriction-naturality square therefore decomposes into the clean core S3 --- the bare
  \(\mathtt{dual\_restrict\_iso}\) immersion-naturality
  (\cref{lem:dual_restrict_iso_restrict_compat}) --- together with the chart-scheme unit
  identification \(u_\iota = \mathtt{unitRestrictIso}\,j\) (\cref{def:unitrestrictiso}) that
  evaluates the comparison at the structure sheaf. The immersion-naturality input is the
  specialisation of \cref{lem:presheafdual_pullback_restrict_natural} to the unit, which is exactly
  the proven \cref{lem:presheafdualunitiso_naturality} promoted from a unit automorphism to the
  immersion \(j\). Composing along the factorisation \(j \mathbin{;} \iota_U = \iota_V\) and threading
  the reindexing \(\varrho = \mathtt{restrictCompReindex}\,j\) (\cref{def:restrictcompreindex},
  \cref{lem:image_preimage_of_le}) through the unit comparison gives the claim.
\end{proof}

\begin{lemma}

\leanok
[The unit contraction is the left unitor at the unit (Cone A coverage)]
  \label{lem:tensorobj_unit_iso_eq_left_unitor}
  \lean{AlgebraicGeometry.Scheme.Modules.tensorObj_unit_iso_eq_left_unitor}
  \uses{lem:tensorobj_unit_iso}
  The structure-sheaf unit contraction
  \(\mathtt{tensorObj\_unit\_iso} : \mathcal{O}_X \otimes_X \mathcal{O}_X \cong \mathcal{O}_X\)
  coincides with the left unitor \(\mathtt{tensorObj\_left\_unitor}\) (\cref{lem:tensorobj_unit_iso})
  evaluated at the module argument \(M = \mathcal{O}_X\).
\end{lemma}
\begin{proof}
  \uses{lem:tensorobj_unit_iso}
  Both sides unfold to the same sheafified presheaf left unitor followed by the
  sheafification-adjunction counit; the identity is definitional.
\end{proof}

\begin{lemma}

\leanok
[Naturality of the left unitor (Cone A coverage)]
  \label{lem:tensorobj_left_unitor_naturality}
  \lean{AlgebraicGeometry.Scheme.Modules.tensorObj_left_unitor_naturality}
  \uses{lem:tensorobj_unit_iso}
  For a morphism \(g : M \to N\) of \(\mathcal{O}_X\)-modules, the left unitor
  (\cref{lem:tensorobj_unit_iso}) is natural:
  \[
    (\mathbf 1_{\mathcal{O}_X} \otimes g) \mathbin{;} (\mathtt{tensorObj\_left\_unitor}\,N).\mathrm{hom}
      = (\mathtt{tensorObj\_left\_unitor}\,M).\mathrm{hom} \mathbin{;} g .
  \]
\end{lemma}
\begin{proof}
  \uses{lem:tensorobj_unit_iso}
  The left unitor is the sheafification of the presheaf-level left unitor composed with the
  sheafification counit; both are natural transformations, so the square commutes by naturality of
  each leg.
\end{proof}

\begin{lemma}

\leanok
[$\eta$ mate-identification: the unit iso is the sheafified presheaf unit comparison (Cone A)]
  \label{lem:pullback_unit_iso_eq_sheafify_eta}
  \lean{AlgebraicGeometry.Scheme.Modules.pullbackUnitIso_eq_sheafify_eta}
  \uses{lem:pullback_unit_iso, lem:presheaf_pullback_oplaxmonoidal, lem:pullback_tensor_iso_unit, def:pullback_val_iso}
  Let \(f : Y \to X\) be a morphism of schemes and \(\varphi'\) its associated change-of-rings map.
  The sheaf-level unit comparison \(\mathtt{pullbackUnitIso}\,f\) (\cref{lem:pullback_unit_iso}) is the
  sheafification transport of the presheaf oplax unit comparison
  \(\eta\,(\mathtt{pullback}\,\varphi')\) (\cref{lem:presheaf_pullback_oplaxmonoidal}):
  \[
    (\mathtt{pullbackUnitIso}\,f).\mathrm{hom}
      = (\mathtt{pullbackValIso}\,f\,\mathcal{O}_X).\mathrm{inv}
        \mathbin{;} a_Y.\mathrm{map}\bigl(\eta\,(\mathtt{pullback}\,\varphi')\bigr)
        \mathbin{;} \mathtt{sheafifyUnitIso}.\mathrm{hom},
  \]
  where \(\mathtt{pullbackValIso}\) (\cref{def:pullback_val_iso}) reconciles the sheaf-level pullback
  value with the sheafified presheaf pullback and \(\mathtt{sheafifyUnitIso}\) --- the unit-specific
  sheafification reconciliation isomorphism --- closes the unit (\(\mathbf 1\)) seam.
\end{lemma}
\begin{proof}
  \uses{lem:pullback_unit_iso, lem:pullback_tensor_iso_unit, def:pullback_val_iso}
  This is the unit square already isolated in the proof that the tensor comparison on the unit pair is
  an isomorphism (\cref{lem:pullback_tensor_iso_unit}): there the right-hand composite is shown to
  equal \(\mathtt{pullbackObjUnitToUnit}\,\varphi\), which by definition is
  \(\mathtt{pullbackUnitIso}\,f\) (\cref{lem:pullback_unit_iso}). Reading that square as a defining
  equation for the sheafified \(\eta\) yields the displayed identity.
\end{proof}

\begin{lemma}

\leanok
[$\delta$ mate-identification: the tensor comparison is the sheafified cotensorator (Cone A)]
  \label{lem:pullback_tensor_map_eq_sheafify_delta}
  \lean{AlgebraicGeometry.Scheme.Modules.pullbackTensorMap_eq_sheafify_delta}
  \uses{lem:pullback_tensor_map, lem:presheaf_pullback_oplaxmonoidal, def:pullback_val_iso, def:scheme_modules_tensorobj}
  Let \(f : Y \to X\) and \(M, N \in \Scheme.\mathtt{Modules}\,X\). The sheaf-level tensor comparison
  \(\mathtt{pullbackTensorMap}\,f\,M\,N\) (\cref{lem:pullback_tensor_map}) is the sheafification
  transport of the presheaf oplax cotensorator
  \(\delta\,(\mathtt{pullback}\,\varphi')\,M_{\mathrm{val}}\,N_{\mathrm{val}}\)
  (\cref{lem:presheaf_pullback_oplaxmonoidal}):
  \[
    \mathtt{pullbackTensorMap}\,f\,M\,N
      = (\mathtt{pullbackValIso}\,f\,(M \otimes_X N)).\mathrm{inv}
        \mathbin{;} a_Y.\mathrm{map}\bigl(\delta\,(\mathtt{pullback}\,\varphi')\,M_{\mathrm{val}}\,N_{\mathrm{val}}\bigr)
        \mathbin{;} \beta_{M,N},
  \]
  where \(\beta_{M,N}\) is the tensor-of-\(\mathtt{pullbackValIso}\) reconciliation
  (\cref{def:pullback_val_iso}, \cref{def:scheme_modules_tensorobj}) on the \(\delta\)-codomain,
  identifying the \(a_Y\)-image of the presheaf tensor with \(f^{*}M \otimes_Y f^{*}N\).
\end{lemma}
\begin{proof}
  \uses{lem:pullback_tensor_map, def:pullback_val_iso}
  This is the definitional unfolding of \(\mathtt{pullbackTensorMap}\) (\cref{lem:pullback_tensor_map}):
  the sheaf comparison was constructed precisely as the sheafified presheaf cotensorator, pre- and
  post-composed with the \(\mathtt{pullbackValIso}\) reconciliations on source and target. The stated
  equation is that construction read back.
\end{proof}

\begin{lemma}

\leanok
[Cone A sub-lemma 1: \texttt{pullbackValIso} naturality at the sheaf left unitor (RHS reconciliation)]
  \label{lem:pullback_val_iso_naturality_left_unitor}
  \lean{AlgebraicGeometry.Scheme.Modules.pullbackValIso_naturality_leftUnitor}
  \uses{def:pullback_val_iso, lem:tensorobj_unit_iso, lem:presheaf_pullback_oplaxmonoidal}
  Let \(f : Y \to X\) and \(M \in \Scheme.\mathtt{Modules}\,X\); write
  \(F = \mathtt{pullback}\,\varphi'\) for the presheaf-level change-of-rings pullback, \(a_Y\) for the
  sheafification functor on \(Y\), and \(\mathtt{pbv}_{(-)} = \mathtt{pullbackValIso}\,f\,(-)\)
  (\cref{def:pullback_val_iso}). The reconciliation isomorphism \(\mathtt{pullbackValIso}\,f\) is
  natural; instantiated against the sheaf-level left unitor
  \((\mathtt{tensorObj\_left\_unitor}\,M).\mathrm{hom} : \mathcal{O}_X \otimes M \to M\), its naturality
  square identifies the \(f^{*}\)-image of that unitor --- the right-hand side of the identity of
  \cref{lem:pullback_tensor_map_left_unitality} --- as the \(a_Y\)-sheafification of the
  \(F\)-image of the unitor's underlying presheaf morphism, conjugated by \(\mathtt{pbv}\) on source
  and target:
  \[
    (\mathtt{pbv}_{\mathcal{O}_X \otimes M}).\mathrm{hom}
      \mathbin{;} f^{*}\bigl((\mathtt{tensorObj\_left\_unitor}\,M).\mathrm{hom}\bigr)
    = a_Y.\mathrm{map}\bigl(F.\mathrm{map}\,(\mathtt{tensorObj\_left\_unitor}\,M).\mathrm{hom}_{\mathrm{val}}\bigr)
      \mathbin{;} (\mathtt{pbv}_{M}).\mathrm{hom},
  \]
  where \((-)_{\mathrm{val}}\) denotes the underlying presheaf morphism (the forget-image) of a sheaf
  morphism.
\end{lemma}
\begin{proof}
  \uses{def:pullback_val_iso, lem:tensorobj_unit_iso}
  Object-wise, \(\mathtt{pullbackValIso}\,f\) is the inverse \(\mathtt{sheafificationCompPullback}\)
  comparison post-composed with the \(f^{*}\)-image of the sheafification counit
  (\cref{def:pullback_val_iso}); both legs are natural transformations. Hence the displayed equation
  is precisely the naturality of \(\mathtt{pullbackValIso}\,f\) against the sheaf morphism
  \((\mathtt{tensorObj\_left\_unitor}\,M).\mathrm{hom} : \mathcal{O}_X \otimes M \to M\). Unfolding the
  sheaf-level unitor as the sheafified presheaf unitor \(a_Y.\mathrm{map}\,\lambda\) followed by the
  counit (\cref{lem:tensorobj_unit_iso}) and cancelling that counit against the counit leg inside
  \(\mathtt{pullbackValIso}\) reduces the square to the naturality of the
  \(\mathtt{sheafificationCompPullback}\) comparison, which is structural.
\end{proof}

\begin{lemma}

\leanok
[Cone A sub-lemma 2: the pullback-object left unitor is the sheafified presheaf unitor]
  \label{lem:tensorobj_left_unitor_pullback_eq_sheafify}
  \lean{AlgebraicGeometry.Scheme.Modules.tensorObj_left_unitor_pullback_eq_sheafify}
  \uses{lem:tensorobj_unit_iso, lem:pullback_tensor_map_eq_sheafify_delta, def:pullback_val_iso}
  With notation as above, the sheaf-level left unitor at the pullback object \(f^{*}M\), pre-composed
  with the tensor/unit reconciliation device \(\beta = \mathtt{sheafifyTensorUnitIso}\) and the
  \(a_Y\)-image of the \(\mathtt{pullbackValIso}\) legs (the same right-hand wrapper as in the
  \(\delta\) identification \cref{lem:pullback_tensor_map_eq_sheafify_delta}), is the sheafified presheaf left unitor
  of \(F.\mathrm{obj}\,M_{\mathrm{val}}\), conjugated by \(\mathtt{pullbackValIso}\) on the target:
  \[
    \beta_{F.\mathrm{obj}\,\mathbf 1,\,F.\mathrm{obj}\,M_{\mathrm{val}}}.\mathrm{hom}
      \mathbin{;} a_Y.\mathrm{map}\bigl((\mathtt{pbv}\,\mathcal{O}_X)\otimes(\mathtt{pbv}\,M)\bigr)
      \mathbin{;} (\mathtt{tensorObj\_left\_unitor}\,(f^{*}M)).\mathrm{hom}
    = a_Y.\mathrm{map}\bigl(\lambda_{F.\mathrm{obj}\,M_{\mathrm{val}}}\bigr)
      \mathbin{;} (\mathtt{pbv}_{M}).\mathrm{hom},
  \]
  where \(\mathbf 1 = \mathcal{O}_X.\mathrm{val}\) is the presheaf tensor unit.
\end{lemma}
\begin{proof}
  \uses{lem:tensorobj_unit_iso, def:pullback_val_iso}
  By definition \(\mathtt{tensorObj\_left\_unitor}\,(f^{*}M)\) is the sheafification of the presheaf
  left unitor at \((f^{*}M).\mathrm{val}\) followed by the sheafification counit
  (\cref{lem:tensorobj_unit_iso}). The \(\beta\)-and-\(\mathtt{pullbackValIso}\) prefix is exactly the
  reconciliation that rewrites the sheafified pullback tensor \(a_Y.\mathrm{map}(F\,\mathbf 1 \otimes
  F\,M_{\mathrm{val}})\) as the sheaf tensor \(f^{*}\mathcal{O}_X \otimes f^{*}M\)
  (the right wrapper of the \(\delta\) identification \cref{lem:pullback_tensor_map_eq_sheafify_delta}).
  Naturality of the presheaf left unitor against the \(\mathtt{pullbackValIso}\) legs,
  followed by naturality of the sheafification counit, slides the unitor past the reconciliation and
  yields the displayed identity. No new monoidal structure is used --- only naturality of \(\lambda\)
  and of the counit.
\end{proof}

\begin{lemma}

\leanok
[Cone A sub-lemma 3: the whiskered unit comparison is the sheafified \(\eta\)-whisker]
  \label{lem:pullback_unit_iso_whisker_eq_sheafify_eta_whisker}
  \lean{AlgebraicGeometry.Scheme.Modules.pullbackUnitIso_whisker_eq_sheafify_eta_whisker}
  \uses{lem:pullback_unit_iso_eq_sheafify_eta, lem:pullback_tensor_map_eq_sheafify_delta, lem:presheaf_pullback_oplaxmonoidal, def:pullback_val_iso}
  With notation as above, the unit comparison whiskered into the left tensor factor,
  \((\mathtt{tensorObjIsoOfIso}\,(\mathtt{pullbackUnitIso}\,f)\,\mathbf 1_{f^{*}M}).\mathrm{hom}\),
  post-composed with the \(\delta\)-identification right wrapper of
  \(\cref{lem:pullback_tensor_map_eq_sheafify_delta}\)
  (\(\beta \mathbin{;} a_Y.\mathrm{map}((\mathtt{pbv}\,\mathcal{O}_X)\otimes(\mathtt{pbv}\,M))\)),
  equals the sheafified presheaf unit whisker, conjugated by the unit reconciliation
  \(\mathtt{sheafifyUnitIso}\) on the left factor:
  \[
    (\mathtt{tensorObjIsoOfIso}\,(\mathtt{pullbackUnitIso}\,f)\,\mathbf 1).\mathrm{hom}
      \mathbin{;} \bigl(\delta\text{-identification right wrapper}\bigr)
    \;\sim\;
    a_Y.\mathrm{map}\bigl(\eta\,F \rhd F.\mathrm{obj}\,M_{\mathrm{val}}\bigr),
  \]
  modulo the \(\mathtt{pullbackValIso}\)/\(\mathtt{sheafifyUnitIso}\)/\(\beta\) reconciliations, where
  \(\eta\,F = \mathtt{Functor.OplaxMonoidal.}\eta\,(\mathtt{pullback}\,\varphi')\) is the presheaf
  oplax unit (\cref{lem:presheaf_pullback_oplaxmonoidal}).
\end{lemma}
\begin{proof}
  \uses{lem:pullback_unit_iso_whisker_eq_sheafify_eta_whisker_3a, lem:pullback_unit_iso_whisker_eq_sheafify_eta_whisker_3b}
  The proof is the composite of two atomic reconciliations. By
  \cref{lem:pullback_unit_iso_whisker_eq_sheafify_eta_whisker_3a} (the bridge-1 substitution and the
  bifunctorial expansion of \(\mathtt{tensorObjIsoOfIso}\)) the central leg of the whiskered unit
  comparison --- the sheafified presheaf unit comparison \(a_Y.\mathrm{map}(\eta\,F)\) right-whiskered
  into the left tensor factor by \(\mathbf 1_{f^{*}M}\) --- is identified, through the tensor
  reconciliation, with the bare sheafified presheaf unit whisker
  \(a_Y.\mathrm{map}(\eta\,F \rhd F.\mathrm{obj}\,M_{\mathrm{val}})\). By
  \cref{lem:pullback_unit_iso_whisker_eq_sheafify_eta_whisker_3b} the flanking
  \(\mathtt{pullbackValIso}\)/\(\mathtt{sheafifyUnitIso}\) legs introduced on the left factor by that
  substitution cancel, in inverse pairs, against the \(\delta\)-identification right wrapper
  \(\beta \mathbin{;} a_Y.\mathrm{map}((\mathtt{pbv}\,\mathcal{O}_X)\otimes(\mathtt{pbv}\,M))\).
  Composing the two leaves exactly \(a_Y.\mathrm{map}(\eta\,F \rhd F.\mathrm{obj}\,M_{\mathrm{val}})\),
  the displayed identity. This is the unit-side analogue of the whiskering for the unit pair in
  \cref{lem:pullback_tensor_iso_unit}, now retaining the whisker rather than discarding it as an
  isomorphism.
\end{proof}

\begin{lemma}
[Cone A sub-lemma 3a: the whiskered sheafified unit comparison is the sheafified \(\eta\)-whisker]
  \label{lem:pullback_unit_iso_whisker_eq_sheafify_eta_whisker_3a}
  \lean{AlgebraicGeometry.Scheme.Modules.pullbackUnitIso_whisker_eq_sheafify_eta_whisker_3a}
  \uses{lem:pullback_unit_iso_eq_sheafify_eta, lem:presheaf_pullback_oplaxmonoidal}
  With notation as above, substitute bridge~1
  (\cref{lem:pullback_unit_iso_eq_sheafify_eta}) for \((\mathtt{pullbackUnitIso}\,f).\mathrm{hom}
  = (\mathtt{pbv}\,\mathcal{O}_X).\mathrm{inv} \mathbin{;} a_Y.\mathrm{map}(\eta\,F) \mathbin{;}
  \mathtt{sheafifyUnitIso}.\mathrm{hom}\) and expand the left-factor whisker
  \((\mathtt{tensorObjIsoOfIso}\,(\mathtt{pullbackUnitIso}\,f)\,\mathbf 1_{f^{*}M}).\mathrm{hom}\) by
  the bifunctorial action of the sheaf tensor. Then the central whiskered leg --- the sheafified
  presheaf unit comparison \(a_Y.\mathrm{map}(\eta\,F)\) right-whiskered into the left tensor factor by
  \(\mathbf 1_{f^{*}M}\) --- is, through the tensor reconciliation
  \(a_Y.\mathrm{map}((\mathtt{pbv}\,\mathcal{O}_X)\otimes(\mathtt{pbv}\,M))\), the sheafified presheaf
  unit whisker:
  \[
    a_Y.\mathrm{map}(\eta\,F) \otimes \mathbf 1_{f^{*}M}
      \;\sim\;
    a_Y.\mathrm{map}\bigl(\eta\,F \rhd F.\mathrm{obj}\,M_{\mathrm{val}}\bigr),
  \]
  modulo the \(\mathtt{pullbackValIso}\) tensor reconciliation identifying the sheaf tensor
  \(f^{*}\mathcal{O}_X \otimes f^{*}M\) with the sheafification of the presheaf tensor
  \(F.\mathrm{obj}\,\mathbf 1 \otimes F.\mathrm{obj}\,M_{\mathrm{val}}\).
\end{lemma}
\begin{proof}
  \uses{lem:pullback_unit_iso_eq_sheafify_eta, lem:presheaf_pullback_oplaxmonoidal}
  The whisker \((\mathtt{tensorObjIsoOfIso}\,e\,\mathbf 1).\mathrm{hom}\) is the bifunctorial action of
  the sheaf tensor on \(e.\mathrm{hom}\) and the identity. With \(e = \mathtt{pullbackUnitIso}\,f\)
  factored by bridge~1 (\cref{lem:pullback_unit_iso_eq_sheafify_eta}) into the three legs
  \((\mathtt{pbv}\,\mathcal{O}_X).\mathrm{inv}\), \(a_Y.\mathrm{map}(\eta\,F)\),
  \(\mathtt{sheafifyUnitIso}.\mathrm{hom}\), functoriality of the tensor product splits the whisker
  into the corresponding three whiskered legs. The central leg \(a_Y.\mathrm{map}(\eta\,F)\)
  whiskered by \(f^{*}M\), transported across the tensor reconciliation that identifies the sheaf
  tensor of the two \(f^{*}\)-values with the sheafification of the presheaf tensor
  \(F.\mathrm{obj}\,\mathbf 1 \otimes F.\mathrm{obj}\,M_{\mathrm{val}}\), becomes
  \(a_Y.\mathrm{map}(\eta\,F \rhd F.\mathrm{obj}\,M_{\mathrm{val}})\) by functoriality of \(a_Y\) and
  of the tensor product (the presheaf whisker \(\eta\,F \rhd F.\mathrm{obj}\,M_{\mathrm{val}}\)
  sheafified).
\end{proof}

\begin{lemma}
[Cone A sub-lemma 3b: the left-factor reconciliation devices cancel against the right wrapper]
  \label{lem:pullback_unit_iso_whisker_eq_sheafify_eta_whisker_3b}
  \lean{AlgebraicGeometry.Scheme.Modules.pullbackUnitIso_whisker_eq_sheafify_eta_whisker_3b}
  \uses{lem:pullback_tensor_map_eq_sheafify_delta, def:pullback_val_iso}
  With notation as above, the reconciliation devices flanking the central \(\eta\)-whisker leg --- the
  \((\mathtt{pbv}\,\mathcal{O}_X).\mathrm{inv}\) and \(\mathtt{sheafifyUnitIso}.\mathrm{hom}\) legs
  whiskered into the left tensor factor (introduced by the bridge-1 substitution of sub-lemma~3a),
  together with the \(\delta\)-identification right wrapper
  \(\beta \mathbin{;} a_Y.\mathrm{map}((\mathtt{pbv}\,\mathcal{O}_X)\otimes(\mathtt{pbv}\,M))\) of
  \cref{lem:pullback_tensor_map_eq_sheafify_delta} --- cancel in inverse pairs, leaving exactly the
  central leg:
  \[
    \bigl(\text{left-factor }\mathtt{pbv}/\mathtt{sheafifyUnitIso}\text{ legs}\bigr)
      \mathbin{;} a_Y.\mathrm{map}\bigl(\eta\,F \rhd F.\mathrm{obj}\,M_{\mathrm{val}}\bigr)
      \mathbin{;} \bigl(\delta\text{-identification right wrapper}\bigr)
    = a_Y.\mathrm{map}\bigl(\eta\,F \rhd F.\mathrm{obj}\,M_{\mathrm{val}}\bigr).
  \]
\end{lemma}
\begin{proof}
  \uses{lem:pullback_tensor_map_eq_sheafify_delta, def:pullback_val_iso}
  Pure iso-algebra. The \((\mathtt{pbv}\,\mathcal{O}_X)\)-component of the right wrapper
  \(a_Y.\mathrm{map}((\mathtt{pbv}\,\mathcal{O}_X)\otimes(\mathtt{pbv}\,M))\) is the inverse of the
  \((\mathtt{pbv}\,\mathcal{O}_X).\mathrm{inv}\) leg whiskered into the left factor, and
  \(\mathtt{sheafifyUnitIso}.\mathrm{hom}\) is the inverse of the \(\beta\) reconciliation on the unit
  factor (\cref{def:pullback_val_iso}); the \((\mathtt{pbv}\,M)\)-component pairs with the trivial
  identity whisker on the right factor. Reassociating the composition exposes each cancelling pair;
  the flanking reconciliation devices cancel in inverse pairs by
  \(\mathtt{Iso.hom\_inv\_id}\)/\(\mathtt{inv\_hom\_id}\), collapsing to the central whisker.
\end{proof}

\begin{lemma}

\leanok
[Sheaf-level left unitality of the pullback tensor comparison (Cone A core)]
  \label{lem:pullback_tensor_map_left_unitality}
  \lean{AlgebraicGeometry.Scheme.Modules.pullbackTensorMap_left_unitality}
  \uses{lem:pullback_val_iso_naturality_left_unitor, lem:tensorobj_left_unitor_pullback_eq_sheafify, lem:pullback_unit_iso_whisker_eq_sheafify_eta_whisker, lem:pullback_tensor_map_eq_sheafify_delta, lem:presheaf_pullback_oplaxmonoidal}
  Let \(f : Y \to X\) and \(M \in \Scheme.\mathtt{Modules}\,X\). The sheaf-level pullback functor
  \(f^{*} = \Scheme.\mathtt{Modules.pullback}\,f\) satisfies left unitality: the composite of the
  tensor comparison on the unit pair, the unit comparison whiskered into the left factor, and the
  target left unitor equals the \(f^{*}\)-image of the source left unitor,
  \[
    \mathtt{pullbackTensorMap}\,f\,\mathcal{O}_X\,M
      \mathbin{;} \bigl((\mathtt{pullbackUnitIso}\,f) \otimes \mathbf 1_{f^{*}M}\bigr)
      \mathbin{;} (\mathtt{tensorObj\_left\_unitor}\,(f^{*}M)).\mathrm{hom}
    = f^{*}\bigl((\mathtt{tensorObj\_left\_unitor}\,M).\mathrm{hom}\bigr).
  \]
\end{lemma}
\begin{proof}
  \uses{lem:pullback_val_iso_naturality_left_unitor, lem:tensorobj_left_unitor_pullback_eq_sheafify, lem:pullback_unit_iso_whisker_eq_sheafify_eta_whisker, lem:pullback_tensor_map_eq_sheafify_delta, lem:presheaf_pullback_oplaxmonoidal}
  Sheafify the presheaf oplax left-unitality coherence
  \(\mathtt{Functor.OplaxMonoidal.left\_unitality\_hom}\,(\mathtt{pullback}\,\varphi')\,
  M_{\mathrm{val}}\) (\cref{lem:presheaf_pullback_oplaxmonoidal}; pentagon-free): applying
  \(a_Y.\mathrm{map}\) and splitting by functoriality gives a three-leg identity with right-hand side
  \(a_Y.\mathrm{map}(F.\mathrm{map}\,\lambda_{M_{\mathrm{val}}})\). Rewrite each leg by its
  reconciliation: the \(\delta\)-leg to \(\mathtt{pullbackTensorMap}\,f\,\mathcal{O}_X\,M\) by
  \(\cref{lem:pullback_tensor_map_eq_sheafify_delta}\) (which also pins the shared right wrapper), the
  \(\eta\)-whisker leg to the whiskered \(\mathtt{pullbackUnitIso}\,f\) by sub-lemma~3
  (\cref{lem:pullback_unit_iso_whisker_eq_sheafify_eta_whisker}), the \(\lambda\)-leg to
  \(\mathtt{tensorObj\_left\_unitor}\,(f^{*}M)\) by sub-lemma~2
  (\cref{lem:tensorobj_left_unitor_pullback_eq_sheafify}), and the right-hand side to the
  \(f^{*}\)-image of the sheaf left unitor via the assembly-form reconciliation --- the
  \(\mathtt{sheafificationCompPullback}\)-headed naturality of \(\mathtt{pullbackValIso}\,f\) ---
  covered by sub-lemma~1 (\cref{lem:pullback_val_iso_naturality_left_unitor}). The
  \(\eta\)-whisker leg itself splits (sub-lemma~3 = 3a + 3b) into the bridge-1 substitution with
  \(\eta\)-whisker expansion and the left-factor device cancellation against the
  \(\mathtt{sheafifyTensorUnitIso}\)/\(\mathtt{pbv}\) wrapper. The interior reconciliation pairs
  cancel, and the displayed identity follows.
\end{proof}

\begin{lemma}

\leanok
[S4b inner seam: pullback-world left-unitor coherence at the chart unit (Cone A)]
  \label{lem:tensorobj_unit_iso_restrict_compat_inner}
  \lean{AlgebraicGeometry.Scheme.Modules.tensorObj_unit_iso_restrict_compat_inner}
  \uses{lem:pullback_tensor_map_left_unitality, lem:tensorobj_restrict_iso_eq_pullback_tensor_map, def:restrictfunctorisopullback, def:unitrestrictiso, lem:pullback_unit_iso}
  With \(j : V \hookrightarrow U\) the chart morphism (\(j \mathbin{;} \iota_U = \iota_V\)), the
  \emph{inner} seam of S4b is the pullback-world image of the left-unitor coherence at the chart unit.
  Transporting the restriction \((\mathtt{restrict}\,j)\,(\mathtt{tensorObj\_unit\_iso})\) through the
  chart comparison \(\mathtt{restrictFunctorIsoPullback}\,j\) (\cref{def:restrictfunctorisopullback})
  and promoting the tensor leg by Bridge~B1
  (\cref{lem:tensorobj_restrict_iso_eq_pullback_tensor_map}), the square reduces to the sheaf-level
  left unitality \(\mathtt{pullbackTensorMap\_left\_unitality}\,j\)
  (\cref{lem:pullback_tensor_map_left_unitality}) evaluated at \(M = \mathcal{O}_U\), with the unit leg
  supplied by \(\mathtt{unitRestrictIso}\,j\) (\cref{def:unitrestrictiso}) and
  \(\mathtt{pullbackUnitIso}\,j\) (\cref{lem:pullback_unit_iso}).
\end{lemma}
\begin{proof}
  \uses{lem:pullback_tensor_map_left_unitality, lem:tensorobj_restrict_iso_eq_pullback_tensor_map, def:restrictfunctorisopullback, def:unitrestrictiso}
  After bridging the chart restriction functor to the abstract pullback via
  \(\mathtt{restrictFunctorIsoPullback}\,j\) and rewriting the tensor leg by Bridge~B1, both flanks of
  the inner square are expressed in the pullback world: the left flank is the pullback image of the
  left unitor, the right flank is the cotensorator-then-unit-then-unitor composite. Their equality is
  precisely \cref{lem:pullback_tensor_map_left_unitality} at \(M = \mathcal{O}_U\), with the unit
  whiskering supplied by \(\mathtt{unitRestrictIso}\,j\) (\cref{def:unitrestrictiso}); iso-algebra then
  closes the seam. No sheaf-level monoidal instance is required --- the coherence is the free
  presheaf-level oplax left unitality transported through sheafification.
\end{proof}

\begin{lemma}


\leanok
[S4b: unit left-unitor commutes with further restriction]
  \label{lem:tensorobj_unit_iso_restrict_compat}
  \lean{AlgebraicGeometry.Scheme.Modules.tensorObj_unit_iso_restrict_compat}
  \uses{lem:tensorobj_unit_iso, lem:image_preimage_of_le}
  With \(j : V \hookrightarrow U\) as above, the unit-tensor contraction
  \(\mathtt{tensorObj\_unit\_iso} : \mathcal{O}_U \otimes_U (-) \cong (-)\) (the left unitor at the
  unit, \cref{lem:tensorobj_unit_iso}) commutes with further restriction along \(j\): modulo the
  reindexing \(\varrho\),
  \[
    (\mathtt{restrict}\,j)\,(\mathtt{tensorObj\_unit\_iso}\ \text{over}\ U)
      \;=\; \mathtt{tensorObj\_unit\_iso}\ \text{over}\ V.
  \]
  Unlike the tensor-restriction comparison \(\mathtt{tensorObj\_restrict\_iso}\), the contraction
  \(\mathtt{tensorObj\_unit\_iso}\) is \emph{not} built from the abstract pullback: by definition it
  is the \emph{sheafification of the presheaf-level left unitor at the unit}, followed by the
  sheafification-adjunction counit (\cref{lem:tensorobj_unit_iso}). Bridge B1 does not apply to the
  contraction directly; instead the square reduces, via the chart comparison
  \(\mathtt{restrictFunctorIsoPullback}\), to a pullback-world left-unitality coherence (Cone~A,
  \cref{lem:pullback_tensor_map_left_unitality}), established below.
\end{lemma}
\begin{proof}
  \uses{lem:tensorobj_unit_iso, lem:tensorobj_unit_iso_eq_left_unitor, lem:tensorobj_left_unitor_naturality, lem:tensorobj_unit_iso_restrict_compat_inner, lem:presheaf_pullback_oplaxmonoidal, def:unitrestrictiso, def:restrictcompreindex, lem:image_preimage_of_le}
  This is the \emph{unit} companion of S2, routed through \emph{Cone~A}. Because the unit object is
  fixed --- there is no module argument --- the square has \emph{no} naturality-in-\(M\) component,
  which makes it strictly simpler than B1.

  \emph{Peeling the contraction.} Write
  \(\mathtt{tensorObj\_unit\_iso} = \mathtt{sheafification.mapIso}(\lambda_{\mathbf 1})
  \mathbin{;} \varepsilon\), where \(\lambda_{\mathbf 1}\) is the presheaf-level left unitor of
  \(\mathtt{PresheafOfModules}\) at the tensor unit \(\mathbf 1 = \mathcal{O}\) and \(\varepsilon\) is
  the sheafification-adjunction counit. By \cref{lem:tensorobj_unit_iso_eq_left_unitor} the
  contraction is the sheaf-level left unitor \(\mathtt{tensorObj\_left\_unitor}\)
  (\cref{lem:tensorobj_unit_iso}) evaluated at \(M = \mathcal{O}\); restriction along \(j\) commutes
  through the outer \(\mathtt{Functor.mapIso}\)/\(\mathtt{toSheafify}\)/counit legs by functoriality
  and naturality of the unitor (\cref{lem:tensorobj_left_unitor_naturality}), so the square reduces to
  a single \emph{inner} seam.

  \emph{Inner seam (Cone~A).} The inner seam is the pullback-world left-unitor coherence at the chart
  unit, \cref{lem:tensorobj_unit_iso_restrict_compat_inner}. Crucially, this is \emph{not} a missing
  monoidal instance: the presheaf change-of-rings pullback \(\mathtt{pullback}\,\varphi'\) is
  \emph{already} oplax monoidal (\cref{lem:presheaf_pullback_oplaxmonoidal}, via
  \(\mathtt{Adjunction.leftAdjointOplaxMonoidal}\)), so its left-unitality coherence
  \(\mathtt{Functor.OplaxMonoidal.left\_unitality\_hom}\) holds for free. The inner seam is the
  \emph{sheafification transport} of that coherence to the sheaf-level maps
  \(\mathtt{pullbackTensorMap}\) (sheafified \(\delta\)) and \(\mathtt{pullbackUnitIso}\) (sheafified
  \(\eta\)); it is discharged by the sheaf-level left unitality
  \cref{lem:pullback_tensor_map_left_unitality} at \(M = \mathcal{O}_U\), through the chart comparison
  \(\mathtt{restrictFunctorIsoPullback}\,j\) and the unit conjugator \(\mathtt{unitRestrictIso}\,j\)
  (\cref{def:unitrestrictiso}). The pentagon plays no role.

  \emph{Outer seam (sheafify and counit).} The remaining legs are formal. The sheafification functor
  acts on the inner unitor square by \(\mathtt{Functor.mapIso}\), and the sheafification unit
  \(\mathtt{toSheafify}\) is natural, so the sheafified squares commute whenever the inner one does.
  The counit \(\varepsilon\) of the sheafification adjunction is a natural transformation, hence its
  naturality square against the restriction-induced morphism is free. Composing the inner agreement
  with these two formal squares yields the equality of the two contractions.

  \emph{Reindexing.} Finally the reindexing \(\varrho = \mathtt{restrictCompReindex}\,j\)
  (\cref{def:restrictcompreindex}) threads the equality of opens
  \(j \mathbin{;} \iota_U = \iota_V\) (\cref{lem:image_preimage_of_le}) through both flanks, absorbing
  the \(\mathtt{restrictFunctorCongr}\) congruence exactly as in S2, and giving the displayed
  identity. No sheaf-level monoidal structure is involved: the structure used lives at the presheaf
  level and is transported through sheafification.
\end{proof}

\begin{lemma}


\leanok
[S4c: global-unit comparison \(u_\iota\) commutes with further restriction]
  \label{lem:trivialisation_uiota_restrict_compat}
  \lean{AlgebraicGeometry.Scheme.Modules.trivialisation_uIota_restrict_compat}
  \uses{lem:pullback_unit_iso, lem:image_preimage_of_le}
  With \(j : V \hookrightarrow U\) as above, the global-unit comparison
  \[
    u_\iota(f) := \bigl(\mathtt{restrictFunctorIsoPullback}\,f\bigr).\mathrm{app}\,
      (\mathtt{SheafOfModules.unit}\,X) \mathbin{;} \mathtt{pullbackUnitIso}\,f
      : \mathtt{restrict}\,(\mathcal{O}_X)\,f \cong \mathcal{O}_U
  \]
  (\cref{lem:pullback_unit_iso}), whose inverse \(u_\iota(f)^{-1}\) is the final leg of the
  trivialisation chain, commutes with further restriction along \(j\): modulo the reindexing
  \(\varrho\),
  \[
    (\mathtt{restrict}\,j)\,u_\iota(\iota_U)^{-1} \;=\; u_\iota(\iota_V)^{-1}.
  \]
\end{lemma}
\begin{proof}
  \uses{def:unitrestrictiso, lem:restrictfunctorisopullback_comp_compat, lem:pullbackObjUnitToUnit_comp, lem:pullback_unit_iso, lem:image_preimage_of_le}
  This is the cheapest constituent: it is the project's proven \emph{unit composition law}
  \(\mathtt{pullbackObjUnitToUnit\_comp}\) (\cref{lem:pullbackObjUnitToUnit_comp}) re-wrapped by
  Bridge B2 alone, with no deep content. By definition
  \(u_\iota(f) = \mathtt{unitRestrictIso}\,f = \mathtt{restrictFunctorIsoPullback}\,f \mathbin{;}
  \mathtt{pullbackUnitIso}\,f\) (\cref{def:unitrestrictiso}), and
  \(\mathtt{pullbackUnitIso} = \mathtt{pullbackObjUnitToUnitIso}\) (\cref{lem:pullback_unit_iso}).
  Strip the \(\mathtt{restrictFunctorIsoPullback}\) factor across \(j \mathbin{;} \iota_U = \iota_V\)
  by Bridge B2 (\cref{lem:restrictfunctorisopullback_comp_compat}); what remains is exactly the
  pullback-world identity \(\mathtt{pullbackObjUnitToUnit}(j \mathbin{;} \iota_U) =
  \mathtt{pullbackComp}^{-1} \mathbin{;} (\mathtt{pullback}\,j)(\mathtt{pullbackObjUnitToUnit}\,
  \iota_U) \mathbin{;} \mathtt{pullbackObjUnitToUnit}\,j\) of
  \cref{lem:pullbackObjUnitToUnit_comp}. Taking inverses and threading the \(\varrho\) reindexing of
  \cref{lem:image_preimage_of_le} gives the stated equation for \(u_\iota^{-1}\).
\end{proof}

\begin{lemma}




\leanok
[Trivialisation iso-chain compatible with further restriction (cocycle residual A, step 1)]
  \label{lem:trivialisation_restrict_compat}
  \lean{AlgebraicGeometry.Scheme.Modules.trivialisation_restrict_compat}
  \uses{lem:tensorobj_restrict_iso_restrict_compat, lem:dual_restrict_iso_restrict_compat, lem:dual_unit_iso_restrict_compat, lem:tensorobj_unit_iso_restrict_compat, lem:trivialisation_uiota_restrict_compat, lem:restrictfunctorisopullback_comp_compat, def:restrictcompreindex, lem:image_preimage_of_le}
  Let \(V \subseteq U\) be opens of \(X\) (an inclusion of trivialising opens, e.g.
  \(V \subseteq U_i \cap U_j\)). The trivialisation iso-chain that builds the local
  contraction --- the restriction-compatibility iso \cref{lem:tensorobj_restrict_iso}
  carrying \((L \otimes_X L^{-1})|_U\) to \(\mathcal{O}_U \otimes_U \mathcal{O}_U\),
  composed with the unit contraction \cref{lem:tensorobj_unit_iso} --- is \emph{natural in
  the open}: its further restriction along \(V \hookrightarrow U\) equals the same chain
  built directly over \(V\). Consequently the sectionwise overlap equation of two local
  contractions \(f_i, f_j\) restricted to \(V\) lifts to an equation of \emph{restricted
  sheaf morphisms} over \(V\), the form on which the iso-level functoriality laws
  (\cref{lem:tensorobjisoofiso_trans}, \cref{lem:dualisoofiso_trans}) and
  \cref{lem:tensorobj_unit_self_duality_collapse} can act.
\end{lemma}
\begin{proof}
  \uses{lem:tensorobj_restrict_iso_restrict_compat, lem:dual_restrict_iso_restrict_compat, lem:dual_unit_iso_restrict_compat, lem:tensorobj_unit_iso_restrict_compat, lem:trivialisation_uiota_restrict_compat, lem:restrictfunctorisopullback_comp_compat, lem:dualisoofiso_trans, lem:tensorobjisoofiso_trans, lem:image_preimage_of_le}

  \emph{The reindexing obstacle.} The chain is assembled over \(U\) and read off over the sub-open
  \(V\); the bookkeeping the one-line statement hides is that the preimage \(\iota_U^{-1}(V)\) (over
  which the \(U\)-built chain evaluates after restriction along \(V \le U\)) and the preimage
  \(\iota_V^{-1}(V)\) (over which the \(V\)-built chain lives) coincide as subsets only up to the
  canonical equality-of-opens identification \(\varrho\) of the direct-image identity
  \((\iota_U)_!\bigl(\iota_U^{-1}(V)\bigr) = V\) (\cref{lem:image_preimage_of_le}). In the chain
  \(\varrho\) appears as a transport isomorphism on \emph{both} flanks of each constituent;
  threading it through every square is the essential bookkeeping.

  \emph{The five squares.} The trivialisation chain carrying \((L \otimes_X L^{-1})|_U\) to
  \(\mathcal{O}_U\), composed with the global-unit comparison \(u_\iota(\iota_U)^{-1}\), is --- \emph{in
  order} --- the composite of five constituents, each now carrying its own restriction-naturality
  square: (1) \(\mathtt{tensorObj\_restrict\_iso}\,\iota_U\)
  (\cref{lem:tensorobj_restrict_iso_restrict_compat}, S2); (2) the bare dual-restriction comparison
  \(\mathtt{dual\_restrict\_iso}\,\iota_U\) (\cref{lem:dual_restrict_iso_restrict_compat}, S3, clean
  core); (3) \(\mathtt{dual\_unit\_iso}\)
  (\cref{lem:dual_unit_iso_restrict_compat}, S4a); (4) \(\mathtt{tensorObj\_unit\_iso}\)
  (\cref{lem:tensorobj_unit_iso_restrict_compat}, S4b); (5) \(u_\iota(\iota_U)^{-1}\)
  (\cref{lem:trivialisation_uiota_restrict_compat}, S4c). The \((\mathtt{dualIsoOfIso}\,e^M)^{-1}\)
  transport that the full trivialisation step wraps around constituent (2) is \emph{not} carried by
  S3 itself; it telescopes here. Constituent (2), with that transport, sits inside the bifunctor
  \(\mathtt{tensorObjIsoOfIso}\,e^M\,(\cdots)\); bifunctoriality (\cref{lem:tensorobjisoofiso_trans})
  splits it into the \(e^M\)-leg --- whose \(V\)-side refinement is the
  \(\mathtt{restrictIsoUnitOfLE}\,h_{VU}\,e^M = (\mathtt{restrict}\,j)\,e^M\) identification --- and
  the dual-chain leg, where contravariant functoriality of \(\mathtt{dualIsoOfIso}\)
  (\cref{lem:dualisoofiso_trans}) transports \((\mathtt{dualIsoOfIso}\,e^M)^{-1}\) through
  \(\mathtt{restrict}\,j\), leaving the bare dual-restriction naturality covered by S3 and S4a.

  \emph{The telescope.} Each sub-lemma asserts that applying \(\mathtt{restrict}\,j\) (the chart
  morphism \(j : V \hookrightarrow U\), \(j \mathbin{;} \iota_U = \iota_V\)) to the \(U\)-built
  constituent equals the \(V\)-built one, with \(\varrho\) on both flanks. Composing the five
  squares in order, the target \(\varrho\) of each is the source \(\varrho\) of the next, so the
  internal reindexings cancel telescopically, leaving only the two outer reindexings
  \(\mathtt{hobjU} = \mathtt{image\_preimage\_of\_le}\,U\,h_{VU}\) and
  \(\mathtt{hobjV} = \mathtt{image\_preimage\_of\_le}\,V\,\mathrm{le\_rfl}\) recorded in the
  statement. Evaluated sectionwise over the preimage open,
  this is the asserted equality.

  \emph{Consequence.} Restricting the sectionwise overlap equation of two local contractions
  \(f_i, f_j\) along \(V \hookrightarrow U_i\) and \(V \hookrightarrow U_j\) through this
  compatibility lifts it from an equation of sections to an equation of \emph{restricted sheaf
  morphisms} over \(V\). That is precisely the form on which the iso-level functoriality laws
  --- bifunctoriality of the tensor iso-of-iso functor (\cref{lem:tensorobjisoofiso_trans}),
  contravariant functoriality of the dual iso-of-iso functor (\cref{lem:dualisoofiso_trans}), and
  the unit self-duality collapse \cref{lem:tensorobj_unit_self_duality_collapse} --- act to
  discharge the cocycle equality.
\end{proof}

\begin{lemma}



\leanok
[Mutually-inverse unit endomorphisms cancel against the left unitor]
  \label{lem:tensorhom_inv_comp_leftunitor}
  \lean{AlgebraicGeometry.Scheme.Modules.tensorHom_inv_comp_leftUnitor}
  In any monoidal category \(\mathcal C\), if \(s, s' : \mathbf 1 \to \mathbf 1\) are endomorphisms
  of the monoidal unit with \(s \mathbin{;} s' = \mathbf 1\), then
  \[
    (s \otimes_{\mathrm m} s') \;\mathbin{;}\; \lambda_{\mathbf 1} = \lambda_{\mathbf 1},
  \]
  where \(\lambda_{\mathbf 1} : \mathbf 1 \otimes \mathbf 1 \to \mathbf 1\) is the left unitor at the
  unit.
\end{lemma}
\begin{proof}
  Factor the tensor product \(s \otimes_{\mathrm m} s'\) through the interchange law. Sliding the
  right factor past the left unitor by its naturality, and the left factor past the right unitor
  (which coincides with the left unitor at the unit), the composite becomes \(\lambda_{\mathbf 1}\)
  precomposed with \(s \mathbin{;} s' = \mathbf 1\); cancelling that identity leaves
  \(\lambda_{\mathbf 1}\).
\end{proof}

\begin{lemma}



\leanok
[Pairing mutually-inverse unit autos through $\mathtt{tensorObjIsoOfIso}$ contracts (cocycle residual A, step 2b)]
  \label{lem:tensorobjisoofiso_comp_unit_iso}
  \lean{AlgebraicGeometry.Scheme.Modules.tensorObjIsoOfIso_comp_unit_iso}
  \uses{lem:tensorhom_inv_comp_leftunitor, def:tensorobjisoofiso, lem:tensorobj_unit_iso}
  Let \(t, s\) be self-isomorphisms of the trivial line bundle \(\mathcal O_V\) with
  \(t.\mathrm{hom} \mathbin{;} s.\mathrm{hom} = \mathbf 1\). Then pairing them through the
  iso-transport bifunctor and contracting via the unit comparison is trivial:
  \[
    \mathtt{tensorObjIsoOfIso}\,t\,s \;\mathbin{;}\; \mathtt{tensorObj\_unit\_iso}
      = \mathtt{tensorObj\_unit\_iso}.
  \]
\end{lemma}
\begin{proof}
  \uses{lem:tensorhom_inv_comp_leftunitor, def:tensorobjisoofiso, lem:tensorobj_unit_iso}
  Both \(\mathtt{tensorObjIsoOfIso}\,t\,s\) and \(\mathtt{tensorObj\_unit\_iso}\) are sheafifications
  of presheaf-level constructions, so the claim is the image under the sheafification functor of a
  presheaf-level coherence. At the presheaf level the pairing is the monoidal tensor
  \((\,\mathrm{forget}\,t \otimes_{\mathrm m} \mathrm{forget}\,s\,)\) of the two underlying unit
  endomorphisms, whose hom-composite is the identity, and the contraction is the left unitor at the
  unit; \cref{lem:tensorhom_inv_comp_leftunitor} collapses the composite to the bare unitor.
  Functoriality of sheafification transports this to the stated sheaf-level identity.
\end{proof}

\begin{lemma}


\leanok
[Naturality of the presheaf dual-of-unit iso against a unit automorphism (B1 eval-core)]
  \label{lem:presheafdualunitiso_naturality}
  \lean{AlgebraicGeometry.Scheme.Modules.presheafDualUnitIso_naturality}
  \uses{def:presheaf_dual_unit_iso, def:dual_unit_iso_gen, lem:internal_hom_eval, lem:presheaf_dualisoofiso_trans}
  Let \(\hat s\) be a presheaf-level automorphism of the unit \(\mathbf 1\) (the underlying map of a
  sheaf-level unit automorphism \(s\)). The presheaf dual-of-unit isomorphism
  \(\mathtt{presheafDualUnitIso} : \mathtt{dual}\,\mathbf 1 \cong \mathbf 1\) is natural in
  \(\hat s\) for the contravariant transport on the source:
  \[
    \mathtt{dualIsoOfIso}\,\hat s \;\mathbin{;}\; \mathtt{presheafDualUnitIso}
      = \mathtt{presheafDualUnitIso} \;\mathbin{;}\; \hat s.
  \]
\end{lemma}
\begin{proof}
  \uses{def:presheaf_dual_unit_iso, def:dual_unit_iso_gen, lem:internal_hom_eval, lem:presheaf_dualisoofiso_trans}
  \leanok
  By construction (\cref{def:dual_unit_iso_gen}, \cref{def:presheaf_dual_unit_iso})
  \(\mathtt{presheafDualUnitIso}\) is the sectionwise evaluation-at-\(1\) map of the internal hom of
  the unit into itself. On sections, \(\mathtt{dualIsoOfIso}\,\hat s\) is precomposition by the
  underlying map of \(\hat s\) (\cref{lem:presheaf_dualisoofiso_trans}), so the left-hand composite
  evaluates a dual section \(\varphi\) at \(1\) after precomposing by \(\hat s\), i.e.\ it computes
  \(\varphi(\hat s(1))\). Since \(\hat s\) is a unit automorphism --- multiplication by a unit scalar
  --- it commutes with evaluation, so this equals \(\hat s\) applied to \(\varphi(1)\), the
  right-hand composite. Equality of the two evaluation-at-\(1\) maps is the naturality of the
  internal-hom evaluation (\cref{lem:internal_hom_eval}) specialised at the unit.
\end{proof}

\begin{lemma}



\leanok
[Dual of a unit automorphism is itself (cocycle residual A, step 2a / B1)]
  \label{lem:dualunitiso_dualisoofiso}
  \lean{AlgebraicGeometry.Scheme.Modules.dualUnitIso_dualIsoOfIso}
  \uses{lem:dual_unit_iso, def:scheme_modules_dual_iso_of_iso, lem:presheafdualunitiso_naturality, def:presheaf_dual_unit_iso}
  Let \(s\) be a self-isomorphism of the trivial line bundle \(\mathcal O_V\). Conjugating its
  contravariant dual-transport by the unit-self-duality iso recovers \(s\):
  \[
    \mathtt{dual\_unit\_iso}^{-1}
      \;\mathbin{;}\; \mathtt{dualIsoOfIso}\,s
      \;\mathbin{;}\; \mathtt{dual\_unit\_iso}
    \;=\; s .
  \]
  Equivalently, the sheaf-level dual of a unit automorphism, read through
  \(\mathtt{dual}\,\mathcal O_V \cong \mathcal O_V\), is the automorphism itself.
\end{lemma}
\begin{proof}
  \uses{lem:dual_unit_iso, def:scheme_modules_dual_iso_of_iso, lem:presheafdualunitiso_naturality, def:presheaf_dual_unit_iso}
  By pure iso-algebra the claim is equivalent to the single naturality square
  \[
    (\mathrm N):\qquad
    \mathtt{dualIsoOfIso}\,s \;\mathbin{;}\; \mathtt{dual\_unit\_iso}
      = \mathtt{dual\_unit\_iso} \;\mathbin{;}\; s,
  \]
  the naturality of \(\mathtt{dual\_unit\_iso} : \mathtt{dual}\,\mathcal O_V \cong \mathcal O_V\) with
  respect to the unit automorphism \(s\), acting contravariantly through \(\mathtt{dualIsoOfIso}\,s\)
  on the source. This is the genuine open obligation. It decomposes through the two layers of the
  construction \(\mathtt{dual\_unit\_iso} = \mathtt{sheafify}(\mathtt{presheafDualUnitIso})
  \mathbin{;} \mathtt{counit}\):
  \begin{itemize}
    \item the presheaf eval-core naturality
      \(\mathtt{dualIsoOfIso}\,\hat s \mathbin{;} \mathtt{presheafDualUnitIso}
        = \mathtt{presheafDualUnitIso} \mathbin{;} \hat s\)
      (\cref{lem:presheafdualunitiso_naturality}, with \(\hat s\) the underlying presheaf map of
      \(s\)), pushed through the sheafification functor; and
    \item the sheafification-counit naturality
      \(\mathtt{sheafify}(\hat s) \mathbin{;} \mathtt{counit} = \mathtt{counit} \mathbin{;} s\),
      the naturality of the sheafification counit at the unit automorphism.
  \end{itemize}
  Composing the two squares gives \((\mathrm N)\), and the iso-algebra reduction then yields the
  stated identity.
\end{proof}

\begin{lemma}




\leanok
[Unit self-duality evaluation collapse (cocycle residual A, step 2)]
  \label{lem:tensorobj_unit_self_duality_collapse}
  \lean{AlgebraicGeometry.Scheme.Modules.tensorObj_unit_self_duality_collapse}
  \uses{lem:tensorobjisoofiso_trans, lem:tensorobjisoofiso_refl, lem:dualisoofiso_trans, lem:dualisoofiso_refl, lem:tensorobj_unit_iso, lem:dual_unit_iso, lem:dualunitiso_dualisoofiso, lem:tensorobjisoofiso_comp_unit_iso}
  Let \(t\) be a self-isomorphism of the trivial line bundle \(\mathcal{O}_V\) on a
  trivialising open \(V\). Pairing \(t\) on the \(M\)-leg with the dual-transported transition
  on the \(N\)-leg --- which lands at the node \(\mathtt{dual}\,\mathcal{O}_V\) and so must be
  conjugated back to \(\mathcal{O}_V\) through the unit-self-duality iso
  \(\mathtt{dual\_unit\_iso} : \mathtt{dual}\,\mathcal{O}_V \cong \mathcal{O}_V\)
  (\cref{lem:dual_unit_iso}) --- and then contracting via the unit comparison cancels:
  \[
    \mathtt{tensorObjIsoOfIso}\,t\,
      \bigl(\mathtt{dual\_unit\_iso}^{-1}
            \;\mathbin{;}\;(\mathtt{dualIsoOfIso}\,t)^{-1}
            \;\mathbin{;}\;\mathtt{dual\_unit\_iso}\bigr)
    \;\mathbin{;}\;
    \mathtt{tensorObj\_unit\_iso}
    \;=\;
    \mathtt{tensorObj\_unit\_iso}.
  \]
  This is the \(g_{ij}\) (inside \(e^M_x\)) meeting \(g_{ij}^{-1} = \mathtt{dualIsoOfIso}\,
  g_{ij}\) (inside \(e^N_x\)) cancellation: under the canonical evaluation of the free
  rank-one module against its dual, the transition iso and its inverse-transpose annihilate,
  leaving the trivialisation-independent unit contraction. The \(\mathtt{dual\_unit\_iso}\)
  conjugation on the \(N\)-leg is what makes both factors live on \(\mathcal{O}_V\) so the
  bifunctor \(\mathtt{tensorObjIsoOfIso}\) can pair them; under the eval-at-\(1\) semantics of
  \(\mathtt{presheafDualUnitIso}\) the conjugated dual leg reduces to \(t^{-1}\), so the pair is
  the symmetric \(t \otimes t^{-1}\) that the multiplication kills.
\end{lemma}
\begin{proof}
  \uses{lem:tensorobjisoofiso_trans, lem:tensorobjisoofiso_refl, lem:dualisoofiso_trans, lem:dualisoofiso_refl, lem:tensorobj_unit_iso, lem:dual_unit_iso, lem:dualunitiso_dualisoofiso, lem:tensorobjisoofiso_comp_unit_iso}
  The collapse factors through the two cocycle helpers rather than a direct sectionwise computation.
  First, the \(N\)-leg --- the dual-conjugated transition
  \(\mathtt{dual\_unit\_iso}^{-1} \mathbin{;} (\mathtt{dualIsoOfIso}\,t)^{-1} \mathbin{;}
  \mathtt{dual\_unit\_iso}\) --- is exactly \(t^{-1}\): inverting the dual-of-a-unit identity
  \cref{lem:dualunitiso_dualisoofiso} (applied to \(t\), read off its inverse) turns the conjugated
  transition into \(t^{-1}\). The collapse thus reduces to pairing \(t\) with \(t^{-1}\) on the two
  legs. Second, since \(t.\mathrm{hom} \mathbin{;} (t^{-1}).\mathrm{hom} = \mathbf 1\), the pairing
  contracts against the unit comparison by \cref{lem:tensorobjisoofiso_comp_unit_iso}, leaving
  \(\mathtt{tensorObj\_unit\_iso}\) alone. This is the \(g_{ij}\)-meets-\(g_{ij}^{-1}\) cancellation,
  now realised as the algebraic identity ``the dual of a unit-automorphism is itself, and a
  unit-automorphism paired with its inverse contracts to the unit''.
\end{proof}

\begin{remark}
[Alternative route via object-level descent (fallback)]
  \label{rem:dual_via_stack}
  Should the slice-and-sheafify route of
  \cref{def:presheaf_internal_hom,lem:internal_hom_isSheaf} bottom out, the dual can
  instead be assembled by \emph{object-level descent}: glue the local trivial duals
  \(\mathcal{O}_{U_i}\) along the inverse-transpose transition isomorphisms
  \(g_{ij}^{-\mathsf{T}}\) of \(L\) into a global \(L^{-1}\) on \(X\). Abstractly this
  is effective descent for sheaves of modules on the small Zariski site --- the
  essential-surjectivity of the descent-data functor of an \(\mathtt{IsStack}\)
  pseudofunctor \(U \mapsto (\mathcal{O}_U\text{-modules})\). This is the more
  principled and reusable construction but the larger build (it requires constructing
  the restriction pseudofunctor and proving the stack/effective-descent property),
  and is recorded only as a fallback; the slice internal-hom route above is the
  intended one. It is the same morphism-/object-level descent family on which the
  associator re-route of \cref{lem:tensorobj_assoc_iso} also waits.
\end{remark}

\subsection{Internal-hom and restrict-scalars helper declarations}
\label{subsec:tensorobj_internalhom_helpers}

The following are the internal Lean helper declarations assembling the morphism
module \cref{def:presheaf_internal_hom_value}, the slice internal hom
\cref{def:presheaf_internal_hom}, its restriction maps
\cref{lem:presheaf_internal_hom_restriction}, the presheaf dual
\cref{def:presheaf_dual}, the evaluation morphism \cref{lem:internal_hom_eval}, the
lax-monoidal restrict-scalars structure \cref{lem:restrictscalars_laxmonoidal}, the
ring-iso base-change equivalence \cref{lem:restrictscalars_ringiso_tensorequiv} and
the pushforward adjunction \texttt{presheaf\_pushforward\_adj\_substrate}. Each is
proved directly in Lean and carries the dependency edges of the construction.

\begin{definition}\leanok
[Canonical ring map from a terminal object]
  \label{def:term_ring_map}
  \lean{PresheafOfModules.InternalHom.termRingMap}
  \uses{def:presheaf_internal_hom_value}
  For a presheaf of commutative rings \(R\) on a base category with terminal object
  \(T\) and an object \(Y\), the unique morphism \(Y \to T\) induces the ring map
  \(R(T) \to R(Y)\) along which a global scalar acts on the value at \(Y\).
\end{definition}
\begin{proof}
  \uses{def:presheaf_internal_hom_value}
  \leanok
  Proved directly in Lean, as \(R\) applied to the opposite of the terminal morphism.
\end{proof}

\begin{lemma}\leanok
[Naturality of the terminal ring map]
  \label{lem:term_ring_map_naturality}
  \lean{PresheafOfModules.InternalHom.termRingMap_naturality}
  \uses{def:term_ring_map}
  The restriction map \(R(g) : R(X) \to R(Y)\) along \(g : X \to Y\) carries
  \(\mathtt{termRingMap}_X\,f\) to \(\mathtt{termRingMap}_Y\,f\), so the images of a
  global scalar \(f \in R(T)\) are compatible with restriction.
\end{lemma}
\begin{proof}
  \uses{def:term_ring_map}
  \leanok
  Proved directly in Lean, from uniqueness of morphisms into the terminal object and
  functoriality of \(R\).
\end{proof}

\begin{definition}\leanok
[Multiplication by a global scalar as an endomorphism]
  \label{def:global_smul}
  \lean{PresheafOfModules.InternalHom.globalSMul}
  \uses{def:term_ring_map, lem:term_ring_map_naturality}
  For \(f \in R(T)\), the endomorphism \(m_f^{N} : N \to N\) of a presheaf of modules
  \(N\) that at each object \(Y\) is multiplication by \(\mathtt{termRingMap}_Y\,f \in
  R(Y)\); its naturality is the heart of the global-scalar action on morphism modules.
\end{definition}
\begin{proof}
  \uses{def:term_ring_map, lem:term_ring_map_naturality}
  \leanok
  Proved directly in Lean: the naturality square commutes by
  \(\mathtt{termRingMap\_naturality}\) and the semilinearity of the restriction maps of
  \(N\).
\end{proof}

\begin{definition}\leanok
[Evaluation functional of a dual section]
  \label{def:eval_lin}
  \lean{PresheafOfModules.evalLin}
  \uses{def:presheaf_dual}
  At an object \(X\), the \(R(X)\)-linear map \(M(X) \to R(X)\) obtained by evaluating
  a dual section \(\varphi\) at its terminal component.
\end{definition}
\begin{proof}
  \uses{def:presheaf_dual}
  \leanok
  Proved directly in Lean, as the underlying linear map of the terminal component of
  \(\varphi\).
\end{proof}

\begin{definition}\leanok
[Lax-monoidal unit of restrict-scalars]
  \label{def:restrict_scalars_lax_epsilon}
  \lean{PresheafOfModules.restrictScalarsLaxε}
  \uses{lem:restrictscalars_laxmonoidal}
  The lax-monoidal unit \(\varepsilon\) of \(\mathtt{restrictScalars}\,\alpha\),
  assembled sectionwise from the \(\mathbf{ModuleCat}\)-level lax units.
\end{definition}
\begin{proof}
  \uses{lem:restrictscalars_laxmonoidal}
  \leanok
  Proved directly in Lean, sectionwise from
  \(\mathtt{ModuleCat.restrictScalars}\)'s lax unit.
\end{proof}

\begin{definition}\leanok
[Lax-monoidal tensorator of restrict-scalars]
  \label{def:restrict_scalars_lax_mu}
  \lean{PresheafOfModules.restrictScalarsLaxμ}
  \uses{lem:restrictscalars_laxmonoidal}
  The lax-monoidal tensorator \(\mu\) of \(\mathtt{restrictScalars}\,\alpha\),
  assembled sectionwise from the \(\mathbf{ModuleCat}\)-level lax tensorators.
\end{definition}
\begin{proof}
  \uses{lem:restrictscalars_laxmonoidal}
  \leanok
  Proved directly in Lean, sectionwise from
  \(\mathtt{ModuleCat.restrictScalars}\)'s lax tensorator.
\end{proof}

\subsection{Vestigial / off-path helper declarations}
\label{subsec:tensorobj_vestigial_helpers}

The following are the off-critical-path Lean helper declarations supporting the
flat-whiskering localizer stability \texttt{flat\_whisker\_localizer}, the
flatness-free whiskering \cref{lem:whisker_of_W}, the \(R_x\)-linear stalk map
\cref{lem:stalk_linear_map}, the stalkwise localizer characterisation
\texttt{W\_implies\_stalkwise\_iso} and the shared slice-site equivalence
\texttt{open\_immersion\_slice\_sheaf\_equiv}. Each is proved directly in Lean.

\section{Substrate helper declarations (wired lean-aux blocks)}
\label{sec:tensorobj_substrate_helpers}

This section gives a blueprint entry to each of the project-local helper declarations
of \texttt{TensorObjSubstrate.lean} that the substrate construction and the
pullback--tensor monoidality machinery invoke internally. Each had been written only
as an inline mid-proof \(\mathtt{\backslash lean}\) reference of an existing block,
which leaves it unregistered as a graph node; here each receives its own \(\mathtt{\backslash label}\)
and its \(\mathtt{\backslash uses}\) edges are read off its Lean body, and the consumers
above are wired to depend on them. All but one are sorry-free; the single residual is
\cref{lem:sheafificationcomppullback_comp_tail}, whose informal proof is given in full.

\begin{definition}\leanok
[Substrate \(\otimes\) respects a pair of isomorphisms]
  \label{def:tensorobjisoofiso}
  \lean{AlgebraicGeometry.Scheme.Modules.tensorObjIsoOfIso}
  \uses{def:scheme_modules_tensorobj}
  For \(M \cong M'\) and \(N \cong N'\) in \(\Scheme.\mathtt{Modules}\,X\), sheafifying the
  presheaf-of-modules tensor of the underlying presheaf isomorphisms produces an
  isomorphism \(M \otimes_X N \cong M' \otimes_X N'\). This is the functor-of-isomorphisms
  form of the substrate tensor, used to build the middle-four interchange and the
  well-definedness of the multiplication on iso-classes.
\end{definition}
\begin{proof}
  \uses{def:scheme_modules_tensorobj}
  \leanok
  Proved directly in Lean: it is \(\mathtt{sheafification}.\mathrm{mapIso}\) applied to the
  monoidal \(\mathtt{tensorIso}\) of the two forgetful images of the given isomorphisms.
\end{proof}

% Functoriality of the iso-transport operators: mechanism by which transition isomorphisms
% cancel in the overlap-compatibility step of rem:dual_discharges_inverse.
\begin{lemma}\leanok
[\(\mathtt{tensorObjIsoOfIso}\) is bifunctorial in composition]
  \label{lem:tensorobjisoofiso_trans}
  \lean{AlgebraicGeometry.Scheme.Modules.tensorObjIsoOfIso_trans}
  \uses{def:tensorobjisoofiso}
  For composable isomorphisms \(M \xrightarrow{e_1} M' \xrightarrow{e_2} M''\) and
  \(N \xrightarrow{e_1'} N' \xrightarrow{e_2'} N''\) in \(\Scheme.\mathtt{Modules}\,X\),
  \[
    \mathtt{tensorObjIsoOfIso}\,(e_1 \ggg e_2)\,(e_1' \ggg e_2')
      = \mathtt{tensorObjIsoOfIso}\,e_1\,e_1' \,\ggg\, \mathtt{tensorObjIsoOfIso}\,e_2\,e_2'.
  \]
  Tensor-transport of a composite of isomorphisms is the composite of the transports.
\end{lemma}
\begin{proof}
  \uses{def:tensorobjisoofiso}
  Both sides are sheafifications of presheaf tensor isomorphisms; the identity reduces, after
  applying the forgetful functor, to the interchange
  \((e_2 \circ e_1) \otimes (e_2' \circ e_1') = (e_2 \otimes e_2') \circ (e_1 \otimes e_1')\)
  of the monoidal tensor of morphisms, together with functoriality of sheafification
  (\(\mathrm{mapIso}\) of a composite is the composite of \(\mathrm{mapIso}\)s).
\end{proof}

\begin{lemma}\leanok
[\(\mathtt{tensorObjIsoOfIso}\) preserves identities]
  \label{lem:tensorobjisoofiso_refl}
  \lean{AlgebraicGeometry.Scheme.Modules.tensorObjIsoOfIso_refl}
  \uses{def:tensorobjisoofiso}
  \(\mathtt{tensorObjIsoOfIso}\,(\mathrm{Iso.refl}\,M)\,(\mathrm{Iso.refl}\,N)
    = \mathrm{Iso.refl}\,(M \otimes_X N)\): the tensor-transport of a pair of identities is the
  identity.
\end{lemma}
\begin{proof}
  \uses{def:tensorobjisoofiso}
  The monoidal tensor of two identity morphisms is the identity, and sheafification sends the
  identity isomorphism to the identity isomorphism.
\end{proof}

% Generic monoidal interchange plumbing (Archon-original): used by B1 to distribute a
% per-leg composite across the tensor and split the functor image in one step.
\begin{lemma}\leanok
[Threefold tensor/composition interchange]
  \label{lem:tensorhom_comp3}
  \lean{AlgebraicGeometry.Scheme.Modules.tensorHom_comp3}
  In a monoidal category \(C\), for morphisms \(a_i : X_i \to Y_i \to Z_i \to W_i\)
  (\(i = 1,2\)) presented as three composable legs on each tensor factor,
  \[
    (a_1^{(1)} \ggg a_1^{(2)} \ggg a_1^{(3)}) \otimes (a_2^{(1)} \ggg a_2^{(2)} \ggg a_2^{(3)})
      = (a_1^{(1)} \otimes a_2^{(1)}) \ggg (a_1^{(2)} \otimes a_2^{(2)}) \ggg (a_1^{(3)} \otimes a_2^{(3)}).
  \]
  The tensor of two threefold composites is the threefold composite of the tensors.
\end{lemma}
\begin{proof}
  Two applications of the bifunctoriality of the monoidal tensor product on morphisms
  (\(\mathtt{tensorHom\_comp\_tensorHom}\)).
\end{proof}

\begin{lemma}\leanok
[Functor image of the threefold interchange]
  \label{lem:map_tensorhom_comp3}
  \lean{AlgebraicGeometry.Scheme.Modules.map_tensorHom_comp3}
  \uses{lem:tensorhom_comp3}
  For a functor \(F : C \to D\) out of a monoidal category, the \(F\)-image of the
  threefold tensor/composition interchange (\cref{lem:tensorhom_comp3}) splits as the
  threefold composite of the \(F\)-images of the per-leg tensors.
\end{lemma}
\begin{proof}
  \uses{lem:tensorhom_comp3}
  Apply \cref{lem:tensorhom_comp3} inside \(F\), then functoriality of \(F\) on composition
  (\(F(g \circ f) = F(g) \circ F(f)\)) twice.
\end{proof}

\begin{lemma}\leanok
[Presheaf dual-transport is contravariantly functorial in composition]
  \label{lem:presheaf_dualisoofiso_trans}
  \lean{AlgebraicGeometry.Scheme.Modules.presheaf_dualIsoOfIso_trans}
  \uses{def:presheaf_dual_iso_of_iso}
  For composable presheaf-of-modules isomorphisms \(L \xrightarrow{e_1} L' \xrightarrow{e_2} L''\),
  \[
    \mathtt{dualIsoOfIso}\,(e_1 \ggg e_2)
      = \mathtt{dualIsoOfIso}\,e_2 \,\ggg\, \mathtt{dualIsoOfIso}\,e_1,
  \]
  the order reversing because the dual is built by \emph{pre}composition (it is contravariant).
\end{lemma}
\begin{proof}
  \uses{def:presheaf_dual_iso_of_iso}
  \(\mathtt{dualIsoOfIso}\,e\) is precomposition of internal-hom sections by the underlying map of
  \(e\). Precomposition by a composite is the composite of precompositions in reversed order, so
  the claim is the contravariant functoriality of \(\mathtt{Hom}(-,\,\mathcal{O})\) applied
  sectionwise; established by \(\mathrm{Iso.ext}\) followed by sectionwise extensionality.
\end{proof}

\begin{lemma}\leanok
[Presheaf dual-transport preserves identities]
  \label{lem:presheaf_dualisoofiso_refl}
  \lean{AlgebraicGeometry.Scheme.Modules.presheaf_dualIsoOfIso_refl}
  \uses{def:presheaf_dual_iso_of_iso}
  \(\mathtt{dualIsoOfIso}\,(\mathrm{Iso.refl}\,L) = \mathrm{Iso.refl}\,(L^{\vee})\) at the presheaf
  level.
\end{lemma}
\begin{proof}
  \uses{def:presheaf_dual_iso_of_iso}
  Precomposition by the identity is the identity on internal-hom sections.
\end{proof}

\begin{lemma}\leanok
[Sheaf-level dual-transport is contravariantly functorial in composition]
  \label{lem:dualisoofiso_trans}
  \lean{AlgebraicGeometry.Scheme.Modules.dualIsoOfIso_trans}
  \uses{def:scheme_modules_dual_iso_of_iso, lem:presheaf_dualisoofiso_trans}
  For composable isomorphisms \(L \xrightarrow{e_1} L' \xrightarrow{e_2} L''\) in
  \(\Scheme.\mathtt{Modules}\,X\),
  \[
    \mathtt{dualIsoOfIso}\,(e_1 \ggg e_2)
      = \mathtt{dualIsoOfIso}\,e_2 \,\ggg\, \mathtt{dualIsoOfIso}\,e_1.
  \]
\end{lemma}
\begin{proof}
  \uses{def:scheme_modules_dual_iso_of_iso, lem:presheaf_dualisoofiso_trans}
  The sheaf-level dual-transport is the presheaf-level one
  (\cref{lem:presheaf_dualisoofiso_trans}) transported through the sheaf-equivalence and the
  forgetful functor; functoriality of \(\mathrm{mapIso}\) carries the presheaf identity to the
  sheaf identity.
\end{proof}

\begin{lemma}\leanok
[Sheaf-level dual-transport preserves identities]
  \label{lem:dualisoofiso_refl}
  \lean{AlgebraicGeometry.Scheme.Modules.dualIsoOfIso_refl}
  \uses{def:scheme_modules_dual_iso_of_iso, lem:presheaf_dualisoofiso_refl}
  \(\mathtt{dualIsoOfIso}\,(\mathrm{Iso.refl}\,L) = \mathrm{Iso.refl}\,(L^{\vee})\) at the sheaf
  level.
\end{lemma}
\begin{proof}
  \uses{def:scheme_modules_dual_iso_of_iso, lem:presheaf_dualisoofiso_refl}
  Immediate from the presheaf-level statement \cref{lem:presheaf_dualisoofiso_refl} and
  functoriality of \(\mathrm{mapIso}\) on the identity.
\end{proof}

\begin{definition}\leanok
[\(\mathcal{O}_X \otimes_X \mathcal{O}_X \cong \mathcal{O}_X\)]
  \label{def:tensorobj_unit_self_iso}
  \lean{AlgebraicGeometry.Scheme.Modules.tensorObj_unit_iso}
  \uses{def:scheme_modules_tensorobj}
  The substrate tensor of the structure sheaf with itself is the structure sheaf:
  \(\mathtt{tensorObj}\,\mathcal{O}_X\,\mathcal{O}_X \cong \mathcal{O}_X\), where
  \(\mathcal{O}_X = \mathtt{SheafOfModules.unit}\,X.\mathtt{ringCatSheaf}\). It is the unit
  isomorphism consumed by the invertibility constructions.
\end{definition}
\begin{proof}
  \uses{def:scheme_modules_tensorobj}
  \leanok
  Proved directly in Lean: the sheafification of the presheaf left unitor
  \(\lambda_{\mathbf{1}}\), composed with the sheafification-adjunction counit at the
  already-sheaf unit.
\end{proof}

\begin{definition}\leanok
[Middle-four interchange for \(\otimes_X\)]
  \label{def:tensorobj_middlefour}
  \lean{AlgebraicGeometry.Scheme.Modules.tensorObj_middleFour}
  \uses{lem:tensorobj_assoc_iso, lem:tensorobj_comm_iso, def:tensorobjisoofiso}
  For arbitrary \(A,B,C,D \in \Scheme.\mathtt{Modules}\,X\) there is an isomorphism
  \((A \otimes B) \otimes (C \otimes D) \cong (A \otimes C) \otimes (B \otimes D)\),
  reassociating the four factors. It is the bookkeeping isomorphism used to show
  \(\otimes\)-invertibility is closed under tensor product.
\end{definition}
\begin{proof}
  \uses{lem:tensorobj_assoc_iso, lem:tensorobj_comm_iso, def:tensorobjisoofiso}
  \leanok
  Proved directly in Lean: an assembly of the unconditional associator
  \cref{lem:tensorobj_assoc_iso} and the braiding \cref{lem:tensorobj_comm_iso} through
  \(\mathtt{tensorObjIsoOfIso}\); no coherence pentagon is consumed.
\end{proof}

\begin{lemma}\leanok
[Unit comparison is an iso from a finality hypothesis]
  \label{lem:isiso_pbu_of_final}
  \lean{AlgebraicGeometry.Scheme.Modules.isIso_pbu_of_final}
  For \(g : Y \to X\) with \((\mathtt{Opens.map}\,g.\mathtt{base}).\mathtt{Final}\), the
  Mathlib unit comparison \(\mathtt{pullbackObjUnitToUnit}\,g\) is an isomorphism.
\end{lemma}
\begin{proof}


\leanok
  Proved directly in Lean by \(\mathtt{inferInstance}\): the lone finality hypothesis is
  isolated at a clean site so the Mathlib \(\mathtt{OfFinal}\) instance synthesises
  without colliding with other in-scope \(\mathtt{Final}\)/\(\mathtt{IsIso}\) hypotheses.
\end{proof}

\begin{definition}\leanok
[Bundled iso form of the unit comparison]
  \label{def:pullback_obj_unit_to_unit_iso}
  \lean{AlgebraicGeometry.Scheme.Modules.pullbackObjUnitToUnitIso}
  \uses{lem:isiso_pbu_of_final}
  For \(g : Y \to X\) with \((\mathtt{Opens.map}\,g.\mathtt{base}).\mathtt{Final}\), the
  isomorphism \(g^*\mathcal{O}_X \cong \mathcal{O}_Y\) bundling the unit comparison
  \(\mathtt{pullbackObjUnitToUnit}\,g\) of \cref{lem:isiso_pbu_of_final}. Handing out the
  bundled \(\mathtt{Iso}\) keeps downstream reasoning at the \(\mathtt{Iso}\) level.
\end{definition}
\begin{proof}
  \uses{lem:isiso_pbu_of_final}
  \leanok
  Proved directly in Lean: \(\mathtt{asIso}\) of the comparison with the explicit witness
  \cref{lem:isiso_pbu_of_final}.
\end{proof}

\begin{definition}\leanok
[Sheafification--monoidality reconciliation brick]
  \label{def:sheafify_tensor_unit_iso}
  \lean{AlgebraicGeometry.Scheme.Modules.sheafifyTensorUnitIso}
  \uses{lem:isiso_sheafification_map_of_W, lem:whisker_of_W}
  For presheaves of modules \(P,Q\) on \(X\), sheafifying the presheaf tensor \(P \otimes Q\)
  agrees with sheafifying the tensor of the underlying presheaves of their sheafifications.
  This is the ``sheafification is monoidal'' reconciliation bridging a presheaf-level
  tensorator with the substrate \(\otimes_X\).
\end{definition}
\begin{proof}
  \uses{lem:isiso_sheafification_map_of_W, lem:whisker_of_W}
  \leanok
  Proved directly in Lean, exactly as in \cref{lem:tensorobj_assoc_iso}: whisker the
  sheafification unit \(\eta\) (a \(J.W\)-morphism) on each side and invert under
  sheafification via \cref{lem:isiso_sheafification_map_of_W} together with the
  flatness-free whiskering stability \cref{lem:whisker_of_W}.
\end{proof}

\begin{lemma}\leanok
[\(\mathtt{hom}\) of the reconciliation brick as a whisker composite]
  \label{lem:sheafify_tensor_unit_iso_hom_eq}
  \lean{AlgebraicGeometry.Scheme.Modules.sheafifyTensorUnitIso_hom_eq}
  \uses{def:sheafify_tensor_unit_iso}
  The \(\mathtt{hom}\) of \cref{def:sheafify_tensor_unit_iso} equals the sheafification of
  \(\eta_P \rhd Q\) followed by the sheafification of \((aP).\mathrm{val} \lhd \eta_Q\),
  stated on the common \(\circ\,\mathtt{forget}_2\) carrier so the two whisker factors can
  be merged.
\end{lemma}
\begin{proof}
  \uses{def:sheafify_tensor_unit_iso}
  \leanok
  Proved directly in Lean by \(\mathtt{rfl}\), stripping the carrier cast.
\end{proof}

\begin{lemma}\leanok
[\(\mathtt{hom}\) of the reconciliation brick as one \(\mathtt{tensorHom}\)]
  \label{lem:sheafify_tensor_unit_iso_hom_eq_prime}
  \lean{AlgebraicGeometry.Scheme.Modules.sheafifyTensorUnitIso_hom_eq'}
  \uses{lem:sheafify_tensor_unit_iso_hom_eq, def:sheafify_tensor_unit_iso}
  The \(\mathtt{hom}\) of \cref{def:sheafify_tensor_unit_iso} equals the sheafification of
  the single \(\mathtt{tensorHom}\) \(\eta_P \otimes \eta_Q\) of the two sheafification-unit
  components --- keeping every term in one monoidal instance so naturality reduces to plain
  bifunctoriality.
\end{lemma}
\begin{proof}
  \uses{lem:sheafify_tensor_unit_iso_hom_eq}
  \leanok
  Proved directly in Lean: merge the two whisker factors of
  \cref{lem:sheafify_tensor_unit_iso_hom_eq} by \(\mathtt{\leftarrow Functor.map\_comp}\)
  and identify with \(\mathtt{tensorHom\_def}\).
\end{proof}

\begin{definition}\leanok
[Sheafified presheaf-pullback \(\cong\) abstract pullback on a value]
  \label{def:pullback_val_iso}
  \lean{AlgebraicGeometry.Scheme.Modules.pullbackValIso}
  \uses{lem:sheafification_comp_pullback_eq_leftadjointuniq}
  For \(f : Y \to X\) and \(M \in \Scheme.\mathtt{Modules}\,X\), sheafifying the
  presheaf-level pullback of the underlying presheaf \(M.\mathrm{val}\) recovers the
  abstract \(\mathcal{O}\)-module pullback \((\mathtt{pullback}\,f).\mathrm{obj}\,M\). It
  reconciles the right-hand side of \(\mathtt{pullbackTensorMap}\) with the substrate
  tensor of the pullbacks.
\end{definition}
\begin{proof}
  \uses{lem:sheafification_comp_pullback_eq_leftadjointuniq}
  \leanok
  Proved directly in Lean: the inverse of
  \(\mathtt{sheafificationCompPullback}\) (\cref{lem:sheafification_comp_pullback_eq_leftadjointuniq})
  at \(M.\mathrm{val}\), composed with the pullback of the sheafification counit at the
  already-sheaf \(M\).
\end{proof}

\begin{lemma}\leanok
[The sheafified \(\otimes\) of the two \(\mathtt{pullbackValIso}\) is an iso]
  \label{lem:isiso_sheafify_tensorhom_pullbackvaliso}
  \lean{AlgebraicGeometry.Scheme.Modules.isIso_sheafify_tensorHom_pullbackValIso}
  \uses{def:pullback_val_iso, def:scheme_modules_tensorobj}
  The fourth factor of \(\mathtt{pullbackTensorMap}\) --- the sheafification of the
  \(\mathtt{tensorHom}\) of the two \(\mathtt{pullbackValIso}\) comparisons (one for \(M\),
  one for \(N\)) --- is an isomorphism.
\end{lemma}
\begin{proof}
  \uses{def:pullback_val_iso}
  \leanok
  Proved directly in Lean: it is the \(\mathtt{hom}\) of the sheafification \(\mathtt{mapIso}\)
  of the monoidal \(\mathtt{tensorIso}\) of the two \(\mathtt{pullbackValIso}\) isomorphisms.
\end{proof}

\begin{lemma}\leanok
[Composite ring-map reconciliation]
  \label{lem:toringcatsheafhom_comp_hom_reconcile}
  \lean{AlgebraicGeometry.Scheme.Modules.toRingCatSheafHom_comp_hom_reconcile}
  For composable \(h : Z \to Y\), \(f : Y \to X\), the structure ring-presheaf map of the
  composite \(h \circ f\) factors as \(\varphi'_f\) followed by the whiskered \(\varphi'_h\).
  This feeds the presheaf \(\mathtt{pullbackComp}\) into the oplax \(\mathtt{comp\_}\delta\)
  decomposition (Sq2 of \cref{lem:pullback_tensor_map_basechange}).
\end{lemma}
\begin{proof}


\leanok
  Proved directly in Lean by \(\mathtt{rfl}\): the two sides hold up to definitional
  equality at default transparency.
\end{proof}

\begin{lemma}\leanok
[Sq2b --- monoidality of the presheaf \(\mathtt{pullbackComp}\)]
  \label{lem:pullbackcomp_delta}
  \lean{AlgebraicGeometry.Scheme.Modules.pullbackComp_δ}
  \uses{lem:presheaf_pullback_oplaxmonoidal, lem:presheaf_pushforward_laxmonoidal, lem:pushforwardcomp_lax_mu}
  The oplax cotensorator \(\delta\) of the pullback for a composite ring map transports,
  through the pullback pseudofunctor coherence \(\mathtt{pullbackComp}\), into the
  \(\mathtt{Functor.OplaxMonoidal.comp}\) cotensorator of the composite pullback. This is the
  \(\eta \to \delta\) analogue, at the \(\mathtt{PresheafOfModules}\) level, of the unit
  coherence \cref{lem:pullbackObjUnitToUnit_comp}.
\end{lemma}
\begin{proof}
  \uses{lem:presheaf_pullback_oplaxmonoidal, lem:presheaf_pushforward_laxmonoidal}
  \leanok
  Proved directly in Lean by the adjunction-mate calculus: transpose under the
  pullback--pushforward adjunction, rewrite the oplax \(\delta\) as the mate of the lax
  \(\mu\) of the pushforward (\cref{lem:presheaf_pushforward_laxmonoidal}), and use the
  conjugate identity for \(\mathtt{pullbackComp}.\mathrm{inv}\) (here \(\mathtt{pushforwardComp}
  = \mathtt{Iso.refl}\)); the sole residual is the lax-\(\mu\) composition coherence
  \(\mathtt{pushforwardComp\_lax\_}\mu\), discharged by a sectionwise pure-tensor collapse.
\end{proof}

\begin{lemma}


\leanok
[Sheafified split of the composite-pullback tensorator (step (ii) sheafified)]
  \label{lem:sheafifymap_deltacomp_split}
  \lean{AlgebraicGeometry.Scheme.Modules.sheafifyMap_δcomp_split}
  \uses{lem:pullbackcomp_delta}
  For composable \(h : Z \to Y\), \(f : Y \to X\) and \(M, N \in \Scheme.\mathtt{Modules}\,X\),
  the sheafification \(a_Z.\mathrm{map}\) of the oplax tensorator \(\delta_{\mathrm{comp}}\) of
  the \emph{composite} presheaf pullback \(\mathtt{pullback}\,\varphi'_{f} \,\circ\,
  \mathtt{pullback}\,\varphi'_{h}\) splits as
  \[
    a_Z.\mathrm{map}\,\delta_{\mathrm{comp}}
      \;=\; a_Z.\mathrm{map}\bigl((\mathtt{pullback}\,\varphi'_{h}).\mathrm{map}\,
            \delta_{f}(M.\mathrm{val}, N.\mathrm{val})\bigr)
        \;\mathbin{;}\;
        a_Z.\mathrm{map}\,\delta_{h}\bigl(F^{\mathrm p}_{f}M.\mathrm{val},\,
            F^{\mathrm p}_{f}N.\mathrm{val}\bigr).
  \]
  This is the sheafified form of step~(ii) of the residual chain in the proof of
  \cref{lem:pullback_tensor_map_basechange}.
\end{lemma}
\begin{proof}
  \uses{lem:pullbackcomp_delta}
  \leanok
  The oplax-monoidal structure on a composite functor \emph{defines} its tensorator
  \(\delta_{\mathrm{comp}}\) to be precisely the composite
  \((\mathtt{pullback}\,\varphi'_{h}).\mathrm{map}\,\delta_{f} \mathbin{;} \delta_{h}\)
  (\(\mathtt{Functor.OplaxMonoidal.comp\_}\delta\)), so the underlying equality holds
  by definition; applying \(a_Z.\mathrm{map}\) and distributing it over the composite by
  functoriality yields the displayed split.
\end{proof}

\begin{lemma}\leanok
[Sheaf-level conjugate of \(\mathtt{pullbackComp}.\mathrm{inv}\)]
  \label{lem:sheaf_unit_comp_pushforward_pullbackcomp_inv}
  \lean{AlgebraicGeometry.Scheme.Modules.sheaf_unit_comp_pushforward_pullbackComp_inv}
  For composable \(h : Z \to Y\), \(f : Y \to X\) and \(Q \in \Scheme.\mathtt{Modules}\,X\),
  the unit of the composite-pullback adjunction post-composed with the pushforward of
  \(\mathtt{pullbackComp}.\mathrm{inv}\) equals the unit of the composite of the \(f\)- and
  \(h\)-adjunctions post-composed with \(\mathtt{pushforwardComp}.\mathrm{hom}\). It is the
  cheap, sheafification-free R0-peel piece of the Sq1 mate calculus.
\end{lemma}
\begin{proof}


\leanok
  Proved directly in Lean: the sheaf-level instance of \(\mathtt{unit\_conjugateEquiv}\)
  combined with \(\mathtt{conjugateEquiv\_pullbackComp\_inv}\) (the mate of
  \(\mathtt{pullbackComp}.\mathrm{inv}\) is \(\mathtt{pushforwardComp}.\mathrm{hom}\)).
\end{proof}

\begin{lemma}\leanok
[Sq1 --- composition coherence of \(\mathtt{sheafificationCompPullback}\)]
  \label{lem:sheafificationcomppullback_comp}
  \lean{AlgebraicGeometry.Scheme.Modules.sheafificationCompPullback_comp}
  \uses{lem:sheafification_comp_pullback_eq_leftadjointuniq, lem:sheaf_unit_comp_pushforward_pullbackcomp_inv, lem:sheafificationcomppullback_comp_tail}
  For composable \(h : Z \to Y\), \(f : Y \to X\) and a presheaf \(P\) over \(X\), the
  sheafification--pullback comparison of \(h \circ f\) factors through the comparisons of
  \(f\) and \(h\), conjugated by the sheaf-level pullback pseudofunctor iso
  \(\mathtt{pullbackComp}\,h\,f\) and the sheafified presheaf pullback pseudofunctor iso.
  This is the \(\mathtt{sheafificationCompPullback}\) twin of \cref{lem:pullbackObjUnitToUnit_comp};
  it is the S1-foundational coherence consumed by \cref{lem:pullback_tensor_map_basechange}.
\end{lemma}
\begin{proof}
  \uses{lem:sheafification_comp_pullback_eq_leftadjointuniq, lem:sheaf_unit_comp_pushforward_pullbackcomp_inv, lem:sheafificationcomppullback_comp_tail}
  \leanok
  Proved directly in Lean: both \(\mathtt{sheafificationCompPullback}\) isos are
  \(\mathtt{leftAdjointUniq}\) of composite adjunctions
  (\cref{lem:sheafification_comp_pullback_eq_leftadjointuniq}); transpose the identity under
  the composite adjunction for \(h \circ f\), evaluate the left-hand side by
  \(\mathtt{homEquiv\_leftAdjointUniq\_hom\_app}\), peel the leading
  \(\mathtt{pullbackComp}.\mathrm{inv}\) factor by
  \cref{lem:sheaf_unit_comp_pushforward_pullbackcomp_inv}, and discharge the residual unit
  identity by \cref{lem:sheafificationcomppullback_comp_tail}.
\end{proof}

\begin{lemma}\leanok
[Sq1 tail --- the R1/R5 unit collapse]
  \label{lem:sheafificationcomppullback_comp_tail}
  \lean{AlgebraicGeometry.Scheme.Modules.sheafificationCompPullback_comp_tail}
  \uses{lem:sheafification_comp_pullback_eq_leftadjointuniq, lem:leftadjointuniq_app_unit_eta_general, lem:pullbackObjUnitToUnit_comp, lem:restrictscalarsid_map, lem:forget_map_pushforward_map}
  \textit{Source: internal categorical construction; no external reference.}
  After the R0-peel of \cref{lem:sheafificationcomppullback_comp}, the residual obligation is
  the unit identity for the two-layer composite adjunction (presheaf pullback--pushforward
  composed with sheafification). Writing \(B_\varphi\) for that sheaf-level composite
  adjunction, \(\mathtt{PresheafPullback}_f\) for the presheaf-level pullback along \(f\),
  \(R_1 = h^*\bigl((\mathtt{sheafCompPb}\,f).\mathrm{app}\,P\bigr)\) and
  \(R_5 = (\mathtt{sheafCompPb}\,h).\mathrm{app}\,(\mathtt{PresheafPullback}_f\,P)\) for the two
  sub-comparison factors, the tail asserts that the composite unit
  \(B_{h \circ f}.\mathrm{unit}.\mathrm{app}\,P\) equals the sheafification-unit of \(X\)
  followed by the (forgetful, change-of-rings-along-identity) image of the merged factor
  \(\bigl((B_f \mathbin{.\mathrm{comp}} B_h).\mathrm{unit} \mathbin{;}
  \mathtt{pushforwardComp}.\mathrm{hom}\bigr) \mathbin{;} (\text{pushforward of } R_1 \mathbin{;}
  R_5 \mathbin{;} a_Z.\mathrm{map}\,\delta_{\mathrm{pre}})\).
\end{lemma}
\begin{proof}
  \uses{lem:sheafification_comp_pullback_eq_leftadjointuniq, lem:leftadjointuniq_app_unit_eta_general, lem:pullbackObjUnitToUnit_comp, lem:restrictscalarsid_map, lem:forget_map_pushforward_map}
  \leanok
  \textit{Source: internal categorical construction; no external reference.}

  This is the sheafification-laden analogue of the unit composition coherence
  \cref{lem:pullbackObjUnitToUnit_comp}, one sheafification layer up. There the adjunctions
  are single transposes and the coherence is sectionwise definitional; here \(B_f\), \(B_h\)
  are \emph{composite} adjunctions \((\text{presheaf pullback--pushforward})\mathbin{.\mathrm{comp}}
  (\text{sheafification})\), so the coherence is a genuine mate-calculus identity rather than a
  reflexivity.

  Both \(\mathtt{sheafificationCompPullback}\) isos are
  \(\mathtt{Adjunction.leftAdjointUniq}\) of these composite adjunctions
  (\cref{lem:sheafification_comp_pullback_eq_leftadjointuniq}). The strategy is to keep the
  left-hand composite unit \(B_{h \circ f}.\mathrm{unit}.\mathrm{app}\,P\) intact and
  \emph{assemble} the right-hand side into it, in the following ordered steps.

  \emph{(a) Strip the identity wrapper.} The tail goal carries an outer
  \(\mathtt{restrictScalars}(\mathbf{1})\) change-of-rings factor from the identity ring map;
  it is removed first (restriction of scalars along the identity is the identity functor),
  exposing the genuine comparison factors.

  \emph{(b) Distribute the right-hand sheaf composite under \(\mathtt{forget}\).} Pushing the
  forgetful functor \(\mathtt{forget}\) (sheaf \(\to\) presheaf) through the right-hand
  composite separates the two sub-comparison factors \(R_1\) (the \(f\)-layer) and \(R_5\)
  (the \(h\)-layer), confined to the right-hand side so that the left-hand composite unit
  \(B_{h \circ f}.\mathrm{unit}\) is left intact for the final reassembly. (Distributing
  without confinement would unfold the left-hand sheafification unit into its
  \(\mathtt{toSheafify} \mathbin{;} \mathtt{restrictHomEquiv}\) expansion, the wrong normal
  form.)

  \emph{(c) Recover the sub-comparison units.} Apply the \(P\)-general
  uniqueness-of-left-adjoints brick \cref{lem:leftadjointuniq_app_unit_eta_general}
  (\(\mathtt{homEquiv\_leftAdjointUniq\_hom\_app}\) in its \(P\)-general form) to rewrite
  \(R_1\) as the composite-adjunction unit \(B_f.\mathrm{unit}.\mathrm{app}\,P\) and \(R_5\)
  as \(B_h.\mathrm{unit}.\mathrm{app}\,(\mathtt{PresheafPullback}_f\,P)\). This is the
  load-bearing step: it identifies each \(\mathtt{sheafCompPb}\) factor with a unit of the
  appropriate composite adjunction.

  \emph{(d) Slide \(\mathtt{pushforwardComp}\) past the recovered units.} The sheaf-level
  comparison \(\mathtt{pushforwardComp}\,h\,f\) is moved across the two recovered units
  \(B_f.\mathrm{unit}\), \(B_h.\mathrm{unit}\) by its naturality.

  \emph{(e) Reassemble the composite unit.} Expanding
  \(B_{h \circ f}.\mathrm{unit}\) by the composite-adjunction unit formula
  \(\mathtt{Adjunction.comp\_unit\_app}\) and applying \(\mathtt{Adjunction.unit\_naturality}\)
  identifies the slid expression of step (d) with the left-hand composite unit, closing the
  tail.

  The single non-mechanical bridge is the presheaf\(\leftrightarrow\)sheaf pushforward
  compatibility \(\mathtt{forget} \circ \mathtt{pushforward}^{\mathrm{sheaf}} =
  \mathtt{pushforward}^{\mathrm{pre}} \circ \mathtt{forget}\) (a definitional identity, since
  the sheaf pushforward is built sectionwise from the presheaf pushforward and
  \(\mathtt{forget}\) is the underlying-presheaf projection), interacting with the recovered
  \(B_f\)- and \(B_h\)-units of step (c); once it is in hand, steps (a)--(e) are a mechanical
  paste. This identity reduces, exactly as the unit twin
  \cref{lem:pullbackObjUnitToUnit_comp} does, to \(\mathtt{Iso.inv\_hom\_id}\) /
  unit-coherence at the unit presheaf. \emph{Warning:} transposing the \emph{whole} tail back
  under the composite adjunction --- instead of recovering the two sub-units individually in
  step (c) --- is circular: it returns the very identity one set out to prove.
\end{proof}

\subsection{Change-of-rings, lax-\(\mu\) and \(\mathtt{extendScalars}\) helpers}
\label{subsec:tensorobj_changeofrings_helpers}

This subsection blueprints the change-of-base-ring cluster of
\texttt{TensorObjSubstrate.lean}: the extension-of-scalars functor and its adjunction,
the lax-monoidal structure maps \(\mu\) of \(\mathtt{restrictScalars}\) and
\(\mathtt{pushforward}\), and the monoidality of the pushforward composition cell
(\(\mathtt{pushforwardComp}\)) --- the change-of-rings coherence that is the sole residual
of the presheaf-level Sq2b lemma \cref{lem:pullbackcomp_delta}. Each helper is sorry-free;
its \(\mathtt{\backslash uses}\) edges are read off its Lean body.

\begin{lemma}\leanok
[\(\mathtt{restrictScalars}(\mathbf{1})\) on morphisms]
  \label{lem:restrictscalarsid_map}
  \lean{AlgebraicGeometry.Scheme.Modules.restrictScalarsId_map}
  For a morphism \(g\) of presheaves of modules over \(R\),
  \((\mathtt{restrictScalars}\,(\mathbf{1}\,R)).\mathrm{map}\,g = g\): restriction of scalars
  along the identity ring map acts as the identity on morphisms.
\end{lemma}
\begin{proof}


\leanok
  Proved directly in Lean by \(\mathtt{rfl}\): \(\mathtt{restrictScalars}\,(\mathbf{1})\) is
  definitionally the identity functor, stated as a propositional rewrite so the wrapper can
  be stripped syntactically without a whole-term \(\mathtt{whnf}\).
\end{proof}

\begin{lemma}\leanok
[Sectionwise value of the \(\mathtt{restrictScalars}\) tensorator \(\mu\)]
  \label{lem:restrictscalars_mu_app}
  \lean{AlgebraicGeometry.Scheme.Modules.restrictScalars_μ_app}
  \uses{def:restrict_scalars_lax_mu}
  The lax tensorator \(\mu\) of \(\mathtt{restrictScalars}\,\alpha\), evaluated at a section
  \(W\), is the \(\mathtt{ModuleCat}\) lax \(\mu\) of \(\mathtt{restrictScalars}\,(\alpha_W)\)
  on the sections. It exposes the ambient presheaf \(\mu\) in a form on which the
  pure-tensor rule matches syntactically.
\end{lemma}
\begin{proof}
  \uses{def:restrict_scalars_lax_mu}
  \leanok
  Proved directly in Lean by \(\mathtt{rfl}\): the presheaf tensorator
  \(\mathtt{restrictScalarsLax}\mu\) (\cref{def:restrict_scalars_lax_mu}) is defined
  sectionwise as the \(\mathtt{ModuleCat}\) tensorator.
\end{proof}

\begin{lemma}\leanok
[Pure-tensor value of the \(\mathtt{ModuleCat}\) \(\mathtt{restrictScalars}\) tensorator]
  \label{lem:forget2_restrictscalars_mu_hom_tmul}
  \lean{AlgebraicGeometry.Scheme.Modules.forget₂_restrictScalars_μ_hom_tmul}
  On a pure tensor \(m \otimes_t n\), the \(\mathtt{ModuleCat}\) lax tensorator \(\mu\) of
  \(\mathtt{restrictScalars}\,f\) (with \(\mathtt{forget}_2\)-carrier rings coming from
  presheaves of \emph{commutative} rings) is the identity: \(\mu(m \otimes_t n) =
  m \otimes_t n\).
\end{lemma}
\begin{proof}


\leanok
  Proved directly in Lean: it is \(\mathtt{ModuleCat.restrictScalars\_}\mu\mathtt{\_tmul}\),
  the Mathlib pure-tensor formula for the change-of-rings tensorator, with the
  \(\mathtt{CommRing}\) instances supplied by the \(\mathtt{CommRingCat}\) carriers.
\end{proof}

\begin{lemma}\leanok
[Pure-tensor value of the presheaf \(\mathtt{restrictScalars}\) tensorator]
  \label{lem:restrictscalars_mu_app_tmul}
  \lean{AlgebraicGeometry.Scheme.Modules.restrictScalars_μ_app_tmul}
  \uses{lem:restrictscalars_mu_app}
  On a pure tensor, \((\mu(\mathtt{restrictScalars}\,\alpha)\,M_1\,M_2)_W
  (m \otimes_t n) = m \otimes_t n\): the presheaf \(\mathtt{restrictScalars}\) tensorator is
  the identity on pure tensors.
\end{lemma}
\begin{proof}
  \uses{lem:restrictscalars_mu_app}
  \leanok
  Proved directly in Lean: expose the \(\mathtt{ModuleCat}\) tensorator by
  \cref{lem:restrictscalars_mu_app}, then apply the Mathlib pure-tensor rule
  \(\mathtt{ModuleCat.restrictScalars\_}\mu\mathtt{\_tmul}\).
\end{proof}

\begin{lemma}

\leanok
[Pure-tensor value of the presheaf \(\mathtt{restrictScalars}\) \emph{oplax} tensorator \(\delta\)]
  \label{lem:restrictscalars_delta_app_tmul}
  \lean{AlgebraicGeometry.Scheme.Modules.restrictScalars_δ_app_tmul}
  \uses{lem:restrictscalars_mu_app_tmul}
  The \(\delta\)-twin of \cref{lem:restrictscalars_mu_app_tmul}. For a sectionwise-bijective
  ring-functor map \(\alpha\), \(\mathtt{restrictScalars}\,\alpha\) is strong monoidal, so its
  comonoidal (oplax) structure map \(\delta\) is the two-sided inverse of the lax \(\mu\). On a pure
  tensor, \((\delta(\mathtt{restrictScalars}\,\alpha)\,M_1\,M_2)_W(m \otimes_t n) = m \otimes_t n\):
  \(\delta\) sends a pure tensor to the same underlying tensor, viewed through
  \(\mathtt{restrictScalars}\). \emph{Remark (R-module instance).} The input \(m \otimes_t n\) lives
  in the \(S\)-tensor section, while the output \(m \otimes_t n\) must be read in the codomain object
  \(\bigl((\mathtt{restrictScalars}\,\alpha).\mathrm{obj}\,M_1 \otimes
  (\mathtt{restrictScalars}\,\alpha).\mathrm{obj}\,M_2\bigr)_W\), whose \(R\)-structure exists only
  via \(\mathtt{restrictScalars}\); the formalizer must ascribe the right-hand side to that codomain.
\end{lemma}
\begin{proof}
  \uses{lem:restrictscalars_mu_app_tmul}
  \leanok
  Proved directly in Lean: since \(\delta = \mu^{-1}\) (\(\mathtt{Functor.Monoidal.}\mu\mathtt{\_}\delta\)),
  rewrite the goal through the inverse and discharge the resulting \(\mu\)-value on the pure tensor by
  \cref{lem:restrictscalars_mu_app_tmul}.
\end{proof}

\begin{lemma}\leanok
[Pure-tensor value of the \(\mathtt{pushforward}\)-mapped \(\mathtt{restrictScalars}\) tensorator]
  \label{lem:pushforward_map_restrictscalars_mu_app_tmul}
  \lean{AlgebraicGeometry.Scheme.Modules.pushforward_map_restrictScalars_μ_app_tmul}
  \uses{lem:restrictscalars_mu_app_tmul}
  The outer leg of \cref{lem:pushforwardcomp_lax_mu}: applied to a pure tensor,
  \(((\mathtt{pushforward}\,\varphi).\mathrm{map}\,(\mu(\mathtt{restrictScalars}\,\psi)
  \,N_1\,N_2))_W\) is the identity.
\end{lemma}
\begin{proof}
  \uses{lem:restrictscalars_mu_app_tmul}
  \leanok
  Proved directly in Lean: reindex the section map through \(\mathtt{pushforward\_map\_app\_apply}\)
  (\(\mathtt{pushforward}\,\varphi\) is \(\mathtt{pushforward}_0 \cdot \mathtt{restrictScalars}\,\varphi\),
  so the section map at \(W\) is the \(\mu\) at \(F^{\mathrm{op}}.\mathrm{obj}\,W\)), then collapse
  by \cref{lem:restrictscalars_mu_app_tmul}.
\end{proof}

\begin{lemma}\leanok
[Reduction of the \(\mathtt{pushforward}\) tensorator to the \(\mathtt{restrictScalars}\) \(\mu\)]
  \label{lem:pushforward_mu_eq}
  \lean{AlgebraicGeometry.Scheme.Modules.pushforward_μ_eq}
  \uses{lem:presheaf_pushforward_laxmonoidal, def:restrict_scalars_lax_mu}
  The lax tensorator \(\mu\) of a single \(\mathtt{pushforward}\,\varphi\) equals the lax
  tensorator of the change-of-rings \(\mathtt{restrictScalars}\,\varphi'\) on the
  (strongly-monoidal, \(\mu_{\mathrm{Iso}} = \mathbf{1}\)) reindexed objects
  \(\mathtt{pushforward}_0^{\mathrm{Comm}}\,F\,R_0\). This unfolds the opaque pushforward
  lax-monoidal \(\mu\) into the directly computable \(\mathtt{restrictScalars}\) \(\mu\).
\end{lemma}
\begin{proof}
  \uses{lem:presheaf_pushforward_laxmonoidal, def:restrict_scalars_lax_mu}
  \leanok
  Proved directly in Lean by \(\mathtt{rfl}\): the pushforward lax-monoidal structure
  (\cref{lem:presheaf_pushforward_laxmonoidal}) is the \(\mathtt{Functor.LaxMonoidal.comp}\)
  of the (identity) \(\mathtt{pushforward}_0\) tensorator with the
  \(\mathtt{restrictScalars}\) tensorator (\cref{def:restrict_scalars_lax_mu}), which is
  definitionally the right-hand side.
\end{proof}

\begin{lemma}\leanok
[Sq2b residual --- monoidality of the \(\mathtt{pushforward}\) composition cell]
  \label{lem:pushforwardcomp_lax_mu}
  \lean{AlgebraicGeometry.Scheme.Modules.pushforwardComp_lax_μ}
  \uses{lem:pushforward_mu_eq, lem:restrictscalars_mu_app, lem:forget2_restrictscalars_mu_hom_tmul, lem:pushforward_map_restrictscalars_mu_app_tmul}
  The change-of-rings coherence ``\(\mathtt{pushforwardComp}\) is monoidal'': the lax
  tensorator \(\mu\) of the composite pushforward
  \(\mathtt{pushforward}\,\psi \cdot \mathtt{pushforward}\,\varphi\) (built by
  \(\mathtt{Functor.LaxMonoidal.comp}\)) agrees with the lax tensorator of the single
  pushforward \(\mathtt{pushforward}\,(\varphi \,;\, F^{\mathrm{op}} \!\lhd\! \psi)\). It is
  the sole genuine residual of the presheaf-level Sq2b lemma \cref{lem:pullbackcomp_delta}.
\end{lemma}
\begin{proof}
  \uses{lem:pushforward_mu_eq, lem:restrictscalars_mu_app, lem:forget2_restrictscalars_mu_hom_tmul, lem:pushforward_map_restrictscalars_mu_app_tmul}
  \leanok
  Proved directly in Lean by a sectionwise pure-tensor collapse: \(\mathtt{hom\_ext}\) then
  \(\mathtt{tensor\_ext}\), unfold the composite \(\mu\) by \(\mathtt{Functor.LaxMonoidal.comp\_}\mu\),
  lighten each \(\mu(\mathtt{pushforward}\,\_)\) to a \(\mathtt{restrictScalars}\) \(\mu\) by
  \cref{lem:pushforward_mu_eq}, expose the \(\mathtt{ModuleCat}\) \(\mu\) by
  \cref{lem:restrictscalars_mu_app}, and collapse the inner leg by
  \cref{lem:forget2_restrictscalars_mu_hom_tmul} and the outer
  \((\mathtt{pushforward}\,\varphi).\mathrm{map}\) leg by
  \cref{lem:pushforward_map_restrictscalars_mu_app_tmul}; the two collapsed pure tensors are
  then definitionally equal. (The equality is genuinely not \(\mathtt{rfl}\) at the presheaf
  level --- the \(\mathtt{restrictScalars}\) \(\mu\) on a pure tensor is real base-change
  content --- so the atomic-object helpers are applied as explicit arguments and matched by
  \(\mathtt{erw}\) to avoid a \(\mathtt{whnf}\)-explosion.)
\end{proof}

\begin{lemma}\leanok
[Forgetful image of a sheaf-pushforward map]
  \label{lem:forget_map_pushforward_map}
  \lean{AlgebraicGeometry.Scheme.Modules.forget_map_pushforward_map}
  The forgetful functor \(\mathtt{SheafOfModules.forget}\) intertwines the sheaf-level
  \(\mathtt{SheafOfModules.pushforward}\,\varphi\) with the presheaf-level
  \(\mathtt{PresheafOfModules.pushforward}\,\varphi.\mathrm{hom}\): for a morphism \(g\) of
  sheaves of modules, \(\mathtt{forget}.\mathrm{map}\,((\mathtt{pushforward}\,\varphi)
  .\mathrm{map}\,g) = (\mathtt{PresheafOfModules.pushforward}\,\varphi.\mathrm{hom})
  .\mathrm{map}\,(\mathtt{forget}.\mathrm{map}\,g)\). This is the STEP-1 binding
  reconciliation of the Sq1 tail.
\end{lemma}
\begin{proof}


\leanok
  Proved directly in Lean by \(\mathtt{rfl}\): the sheaf pushforward is built sectionwise
  from the presheaf pushforward and \(\mathtt{forget}\) is the underlying-presheaf
  projection, so the two sides are definitionally equal.
\end{proof}

\subsection{The Picard quotient: setoid, multiplication and inverse}
\label{subsec:tensorobj_pic_quotient_helpers}

This subsection blueprints the three quotient-level helpers underlying the
invertibility-carrier Picard group \cref{def:pic_carrier} and its abelian-group law
\texttt{pic\_commgroup}: the isomorphism setoid, the tensor-induced multiplication, and
the dual-induced inverse on iso-classes.

\begin{definition}\leanok
[The isomorphism setoid of invertible modules]
  \label{def:pic_setoid}
  \lean{AlgebraicGeometry.Scheme.Modules.picSetoid}
  \uses{def:scheme_modules_isinvertible}
  The setoid on \(\otimes\)-invertible \(\mathcal{O}_X\)-modules with \(M \sim M'\) iff there
  is an isomorphism \(M \cong M'\). The quotient of this setoid is the carrier of the Picard
  group \(\Pic X\).
\end{definition}
\begin{proof}
  \uses{def:scheme_modules_isinvertible}
  \leanok
  Proved directly in Lean: reflexivity, symmetry and transitivity of the
  \(\mathtt{Nonempty}\,(M \cong M')\) relation are witnessed by \(\mathtt{Iso.refl}\),
  \(\mathtt{Iso.symm}\) and \(\mathtt{Iso.trans}\).
\end{proof}

\begin{definition}\leanok
[Multiplication on \(\Pic X\) induced by \(\otimes_X\)]
  \label{def:pic_mul}
  \lean{AlgebraicGeometry.Scheme.Modules.picMul}
  \uses{def:pic_setoid, def:pic_carrier, def:tensorobjisoofiso, lem:isinvertible_tensor}
  The multiplication \([M] \cdot [M'] := [M \otimes_X M']\) on \(\Pic X\), well-defined on
  iso-classes.
\end{definition}
\begin{proof}
  \uses{def:pic_setoid, def:pic_carrier, def:tensorobjisoofiso, lem:isinvertible_tensor}
  \leanok
  Proved directly in Lean by \(\mathtt{Quotient.lift}_2\): the value lands in \(\Pic X\)
  because the tensor of invertibles is invertible (\cref{lem:isinvertible_tensor}), and it
  respects the setoid because \(\mathtt{tensorObjIsoOfIso}\) (\cref{def:tensorobjisoofiso})
  turns a pair of isomorphisms into an isomorphism of tensors.
\end{proof}

\begin{definition}\leanok
[Inverse on \(\Pic X\) induced by the tensor-inverse]
  \label{def:pic_inv}
  \lean{AlgebraicGeometry.Scheme.Modules.picInv}
  \uses{def:pic_setoid, def:pic_carrier, lem:isinvertible_inverse_welldef, lem:tensorobj_comm_iso}
  The inverse class \([M] \mapsto [N]\) on \(\Pic X\), where \(N\) is a chosen tensor-inverse
  of \(M\).
\end{definition}
\begin{proof}
  \uses{def:pic_setoid, def:pic_carrier, lem:isinvertible_inverse_welldef, lem:tensorobj_comm_iso}
  \leanok
  Proved directly in Lean by \(\mathtt{Quotient.lift}\): a chosen inverse witness
  (\(\mathtt{Classical.choose}\) of the invertibility) is itself invertible, with the
  braiding (\cref{lem:tensorobj_comm_iso}) supplying the membership isomorphism; the choice
  is independent of the representative because any two tensor-inverses of \(M\) are
  isomorphic (\cref{lem:isinvertible_inverse_welldef}).
\end{proof}
